\documentclass{book}

\usepackage[colorinlistoftodos,draft]{todonotes}

\usepackage{marginnote}
\setlength{\marginparwidth}{2cm}
\usepackage[dvipsnames,x11names]{xcolor}
\usepackage{pgfplots,diagbox,colortbl,tikz,graphicx,algorithm2e,cancel,verbatim,
listings,float,amsmath,amssymb,enumitem,array,subfiles,bussproofs,
rotating,MnSymbol,mathtools,subcaption,caption}
\pgfplotsset{compat=1.18}
\usetikzlibrary{automata, positioning,graphs,shapes, arrows, calc}
% CORRECT LOCATION (preamble):
\usepackage{csquotes}

\hyphenation{ma-chine-check-a-ble}
\hyphenation{dou-ble-push-out}
%%%%%%% careful : order is important begin
\usepackage{amsthm}
\usepackage{thmtools,hyperref}
\newtheorem{theorem}{Theorem}[chapter]
\newtheorem{corollary}[theorem]{Corollary}
\newtheorem{lemma}[theorem]{Lemma}
\newtheorem{idea}[theorem]{Idea}
\newtheorem{definition}[theorem]{Definition}
\newtheorem{proposition}[theorem]{Proposition}
\newtheorem{remark}[theorem]{Remark}
\newtheorem{example}[theorem]{Example}
\newtheorem{Plan}[theorem]{Plan}
\newtheorem{claim}[theorem]{Claim}
\newtheorem{notation}[theorem]{Notation}
\renewcommand{\chapterautorefname}{Chapter}
\usepackage[section]{placeins}
\usepackage{minted}



\newcommand{\set}[1]{\{#1\}}
\newcommand{\graphbox}[8]{
  \begin{scope}[xshift=#2,yshift=#3]
    \draw [rounded corners=2mm] (0,0) rectangle (#4,-#5);
    \node at (0,0mm) [anchor=north west,inner sep=1mm] {#1};
    \begin{scope}[xshift=#4/2+#6,yshift=#7] 
    #8
    \end{scope}
  \end{scope}
}
\newcommand{\card}[1]{|#1|}
\newcommand{\homset}[2]{\operatorname{Hom}(#1,#2)} 
\newcommand{\isdef}{\overset{\operatorname{def}}{=}}
\newcommand{\cat}[1]{\mathcal{#1}}
\newcommand{\f}[1]{\mathfrak{#1}}
\newcommand{\rs}[1]{\mathcal{#1}}
\newcommand{\opn}[1]{\operatorname{#1}}
%arrow
\newcommand{\itrs}{\overset{i}{\to}} 

\usepackage{tocloft}

% % Customizing the table of contents
% % \setlength{\cftpartindent}{0pt} %  space before part titles
\setlength{\cftchapindent}{15pt} %  space before chapter titles
\setlength{\cftsecindent}{30pt} %  space before sec titles
\setlength{\cftsubsecindent}{40pt} %  space before subsec titles
\setlength{\cftsubsubsecindent}{50pt} %  space before subsubsec titles


\newcommand{\textdef}{\text{$\operatorname{Def.}$}}
\newcommand{\pbpop}{$\operatorname{PBPO}^+$}
\addtolength{\cftbeforechapskip}{-10pt} % Reduce spacing by 10pt
\addtolength{\cftbeforepartskip}{-20pt} % Reduce spacing by 10pt
\newif\iflongversion
\longversiontrue
\newif\iftrackChange
\trackChangetrue
\newcommand{\trackedtext}[1]{%
    \iftrackChange
        \textcolor{red}{#1}% Apply red color if trackChange is true
    \else
        #1% No color change if trackChange is false
    \fi
} 


% \usepackage[style=alphabetic, maxbibnames=999]{biblatex}
\usepackage[maxbibnames=999]{biblatex}
\addbibresource{bib/these.bib} % Add your bibliography file here

\usepackage{titlesec} 
 
\usepackage{tikz-cd}
\usepackage{pifont}      
\usepackage{makecell}    
\usepackage{nicematrix}  
\usepackage{wrapfig}     
\usepackage{lipsum}      

% 

% \usepackage{diagbox,colortbl,tikz,graphicx,algorithm2e,cancel,verbatim,
% listings,float,amsthm,amsmath,amssymb,enumitem,array,subfiles,bussproofs,
% rotating,MnSymbol,hyperref,mathtools,subcaption,caption}
\begin{document}  
  

\begin{titlepage}

    \unitlength 1cm
    \begin{center}
    
    \vspace*{-2.5cm}
    \begin{figure}[h]
        \centering
        \includegraphics[width=0.8\textwidth]{LogoLyon1Off_CoulCmjn300dpi.jpg}
    \end{figure}
    
               {\large\bf THÈSE de DOCTORAT DE L'UNIVERSITÉ DE LYON\\}
   
               {\large \bf l'Université Claude Bernard Lyon 1\\}
    \vspace{12pt}
               {\large \textbf{Ecole Doctorale} 512 \\InfoMaths}

    \vspace{12pt}
    
               {\large \textbf{Spécialité de doctorat~:} Informatique \\}
            %    {\large \textbf{Discipline~:} (éventuellement)}
    
    \vspace{0.8cm}
    
               {Soutenue publiquement le jj/mm/aaaa, par~:\\}
               {\Large\bf {Qi QIU}}
    \vspace{0.5cm}           
    
    
    \rule{5cm}{1pt}
    
    \vspace{12pt}
               
               {\LARGE \bf Automated proofs of termination: contributions to graph rewriting using extended weighted type graph and morphism counting}\\
               \vspace{0.3cm}
               {\Large \bf }
               
    \vspace{12pt}
    \rule{5cm}{1pt}
    
    \vspace{0.5cm}
    
    \end{center}
    
    Devant le jury composé de~:\\
    
    \small {
    \begin{tabular}{ll}
    % \textbf{NAME}           &\textbf{Président.e}\\
    %   Grade/qualité, Établissement/entreprise\\
    % \textbf{Clara Bertolissi}           &\textbf{Examinatrice}\\
    %   Grade/qualité, Établissement/entreprise\\
    % \textbf{Évelyne Contejean}           &\textbf{Examinatrice}\\
    %   Grade/qualité, Établissement/entreprise\\
    \textbf{Jean Krivine}           &\textbf{Rapporteur}\\
      Grade/qualité, Établissement/entreprise\\
    \textbf{Pascale Le Gall}           &\textbf{Rapporteure}\\
      Grade/qualité, Établissement/entreprise\\
    % \textbf{Philippe Malbos}           &\textbf{Examinatrice}\\
    %   Grade/qualité, Établissement/entreprise\\
    % \textbf{Frédéric Prost}           &\textbf{Examinatrice}\\
    %   Grade/qualité, Établissement/entreprise\\
    \textbf{Xavier URBAIN}           &\textbf{Directeur de thèse}\\
      Grade/qualité, Établissement/entreprise\\
    % \textbf{NAME}     &\textbf{Invité.e} (le cas échant)\\
    %   Grade/qualité, Établissement/entreprise\\
    \end{tabular}
    }
    
    %\vfill \clearpage{\pagestyle{empty}\cleardoublepage}
    
    %-------------------------------- Remerciement  -------------------
    \end{titlepage}  
 
\tableofcontents  
\newpage     
     
\part{Introduction} 
\chapter{Motivation}

\begin{itemize}
    \item context:todo
    \item Contribution:todo
    \item Structure of this Thesis:todo
\end{itemize}


\subsubsection*{Distributed Algorithms and the importance for ensuring their correctness}

\begin{itemize}
    \item Algorithms play a pivotal role in our modern society. For example, they control machines that we use in daily basis such as vacuum cleaners, telephones but also critical systems such as medical devises, trains, airplanes, spaceship.
    
    \item In recent years, the complexity and scale of computational tasks have grown significantly, but the increase of  increased performance  leading to the development of distributed algorithms. These algorithms control multiple machines simultaneously to solve complex problems more efficiently than their centralized counterparts.

    \item These distributed systems can more quickly address some problems and are more robust than their centralized counterparts.  
  
    \item However, distributed algorithms are error-prone because they deal with multiple processes running at the same time, communication over unreliable networks, and the possibility of individual components failing. They also face challenges like keeping data consistent, handling unpredictable events, and ensuring all parts of the system stay coordinated without a central controller. These complexities make designing and debugging distributed algorithms very difficult.
    \item Errors in distributed algorithms can have severe consequences, such as the loss of critical data or even the loss of human lives.
    \item Thus, ensuring the correctness of distributed algorithms is crucial. 
    \item Nevertheless ensuring the correctness of distributed algorithms is challenging because they are complex. 
    \item Since algorithms exhibit numerous properties, it is infeasible to study all of them at once, making it necessary to focus on fundamental aspects.
\end{itemize}

\subsubsection*{Automated Tools for Formal Verification}
\begin{itemize}
    \item Formal methods can address this challenge by rigourious mathmatical methods. Among these methods, mechanically verified proofs using proof assistants, designed to help users create formal proofs by providing a framework for defining mathematical theories, writing proofs, and verifying their correctness, are particularly promising.
    \item These methods can provide a high level of confidence in the correctness of algorithms, but they require a high level of expertise and are often difficult to use.  
    \item To address this challenge, we propose to develop automation tools that can relieve users from technical or tedious tasks during their formal developments by generating certificats which can be verified by proof assistants.
\end{itemize}

\subsubsection{Graph Relabeling Systems: A Formalism for Modeling distributed Algorithms}
\begin{itemize}
    \item One of the main challenges in developing automated tools for formal verification is to represent distributed algorithms in a mathematical structure that can be encoded in proof assistants and used to reason about their properties.
    \item Graph relabeling systems are a formalism that can be used to model distributed algorithms. In these systems, computation units are represented by nodes, communication canals are represented by arrows, and the working environment is modeled by a graph. Arrows have labels representing information encoded in computation units and states of communication canals. System behavior expressed by state changes is modeled by rule-based graph relabeling: replacing some labels either on nodes or arrows or both by new labels. 
    \item The rules describe preconditions and postconditions of single transformation steps. In a pure graph relabeling system, the order of the steps is determined by the causal dependency of actions only, i.e. independent rule applications can be executed in an arbitrary order. The concept of rules in graph relabeling systems provides a clear concept for defining system behavior. In particular, for modeling the intrinsic concurrency of actions, graph rules provide a suitable means, because they explicate all structural interdependencies. If we stick to sequential execution, parallel transformations have to be modeled by interleaving their atomic actions arbitrarily. This interleaving leads to the same result if the atomic actions are independent of each other. Simultaneous execution of actions can be modeled if a parallel rule is composed from the actions.
    \item Modeling by graph Relabeling is simple, since only labels are modified, and visual, since it is very natural to use a visual representation of graphs. It is also precise, owing to its formal foundation. Thus, graph transformation can be used in formal verification for state-based systems.
\end{itemize}

\subsubsection*{Focus on Termination}
\begin{itemize}
    \item In this thesis, we focus on one fundamental property of algorithms: termination. When we run an algorithm with this property, it will end sooner or later. This property of algorithms is highly desirable in many contexts. For instance, when you ask your calculator to compute 2 times 2, you expect it to complete the calculation promptly and provide you with the correct result.
    \item Termination is one of the most fundamental and difficult problems in computer science because it is undecidable. This means we can never have a program that can correctly determine for every algorithm whether it terminates.
    \item thus, we need to define on which class of algorithms we are trying to find an automated termination method. and this thesis focuses mainly on those distributed algorithms that can be modeled by graph relabeling systems. 
\end{itemize}

\subsubsection*{Approaches to Proving Termination of Graph Relabeling Systems}
\begin{itemize}
    % \item all possible configurations (states) that an algorithm can be in during its execution form a set, the state space. Instructions within a program govern the transitions between states. 
    % \item A graph Relabeling system, capturing this transition relation, is a binary relation. 
    \item A graph relabeling system is a formalism for modeling distributed algorithms. It consists of a set of rules that describe how and when a graph, representing a state of the system, can evolve to another. Thus, it defines a binary relation on graphs, where each rule represents transitions from one graph to another.
    \item The main idea behind proving the termination of a graph relabeling system is to show that a graph where no rule can be applied will eventually be reached. This is done by demonstrating that the given graph relabeling system defines a binary relation that can be embedded into a terminating binary relation.
    \item In this thesis, we consider two classes of terminating binary relations: those established by graph rewriting systems and those established by term rewriting systems.
\end{itemize}

\subsubsection*{Graph Rewriting Systems: Another Formalism for Modeling Distributed Algorithms with Advanced Automated Termination Technique}
\begin{itemize} 
    \item There are many other approaches to model distributed algorithms.
    \item The main stream approach is called algebraic graph rewriting. This approach is more general because graph rewriting systems modify also the topology of the underlying graph and just labels. 
    \item[graph rewriting system] graph rewriting system are more general than graph Relabeling systems: nodes can be merge, created or deleted.
    \item Since graph Relabeling systems are can be seen as special cases of graph rewriting systems, the method, called type graph method, can be applied to prove termination of graph relabeling systems
\end{itemize}



\subsubsection*{Type Graph Method: An Advanced Termination Technique for DPO Graph Rewriting Systems}

\begin{itemize}
    \item ss
\end{itemize}

\subsubsection*{A New Termination Criterion with Weighted Type Graph Method for Linear DPO Graph Rewriting Systems}
\begin{itemize}
    \item[article]
    \item The type graph method, such as presented in the litterature, can not prove the termination of a graph rewriting system, if there is a rule such that there is an epic morphism $e: R \to L$.
    \item We proposed a solution in the case of rewriting systems on multigraphs.
\end{itemize}

\subsubsection*{Type Graph over non Well-founded semirings}
\begin{itemize}
    \item The type graph method has a crucial drawback in the practice ....
    \item we proposed an extension to non well founded semirings, which allows to over this problem in many cases
\end{itemize}

\subsection*{Term Rewriting Systems: Another Formalism for Modeling Algorithms with advanced automated termination method}

\begin{itemize}
    \item although termination is a crucial property for any computational formalism, it has not received a lot of attention in the graph transformation community.
    \item[TRS] term rewriting systems provide a nother approach to model parallel and distributed algorithms. Furthermore, there are many works on termination of term rewriting system and many advanced termination techniques haved been developped.
    \item[Translation to TRS] since A graph Relabeling system represents a binary relation on graphs and a term rewriting system represents a binary relation on terms, if we can establish a homomorphism between the  graph Relabeling system of interest to a term rewriting system whose termination can be proved by an existing advanced termination method, then the termination of the  graph Relabeling system can be established.
    
    \item  Because advanced termination methods are available for term rewriting systems, our initial approach is to translate graph rewriting systems into term rewriting systems. 
    \item In this translation, each graph in a rewriting rule is replaced by a term that encodes all the graph's information. We prove that this translation preserves termination: the translated system terminates if and only if the original graph rewriting system terminates. 
    \item Thus, proving the termination of the term rewriting system also proves the termination of the original graph rewriting system. 
    \item However, it is not always straightforward to represent a distributed algorithm as a term rewriting system.

    \item To address this, we consider priority-based graph rewriting systems, in which rules have a priority order. When rule occurrences overlap, those with lower priority are blocked. We propose a translation that uses rules simulating the original graph rewriting rules. These rules identify possible occurrences of the original rules within the encoded graph and apply them accordingly. However, experiments show that these auxiliary rules complicate termination proofs.

    \item Therefore, we propose translating priority-based graph rewriting systems into hierarchical term rewriting systems. In this approach, rules are partitioned into hierarchical modules, and each module only rewrites subterms whose root symbol corresponds to that module. 
    \item While we have only proven termination equivalence for the first two translations, we believe that all proposed translations preserve termination.  
\end{itemize}

 
\newpage
\begin{itemize}
    \item[fundamental of ...] By analogy with graphs and graph transformations, term rewriting can be considered as a concept in the field of tree transformations. Trees, however, unlike graphs, do not allow sharing of common substructures, which was one of the main reasons for the efficiency problems concerning functional and logical programs. This motivates us to consider graphs rather than trees as the fundamental structures of computing.
    \item[fundamental of ...]   Modeling by graph transformation is visual, on the one hand, since it is very natural to use a visual representation of graphs; on the other hand, it is precise, owing to its formal foundation. Thus, graph transformation can also be used in formal specification techniques for state-based systems.
    \item[fundamental of ...] Concurrency and distribution. When graph transformation is used to describe a concurrent system, graphs are usually taken to describe static system structures. System behavior expressed by state changes is modeled by rule-based graph manipulations, i.e. graph transformation. The rules describe preconditions and postconditions of single transformation steps. In a pure graph transformation system, the order of the steps is determined by the causal dependency of actions only, i.e. independent rule applications can be executed in an arbitrary order. The concept of rules in graph transformation provides a clear concept for defining system behavior. In particular, for modeling the intrinsic concurrency of actions, graph rules provide a suitable means, because they explicate all structural interdependencies. If we stick to sequential execution, parallel transformations have to be modeled by interleaving their atomic actions arbitrarily. This interleaving leads to the same result if the atomic actions are independent of each other. Simultaneous execution of actions can be modeled if a parallel rule is composed from the actions.
    \item[fundamental of ...] Parallel and distributed graph transformation both offer structured rule applications, in both temporal and spatial dimensions. Distributed graphs contain an additional structure on the graphs. Graphs are allowed to be split into local graphs and, after local transformations, local graphs are joined again to one global graph. Parallel graph transformation can be considered as a special case of distributed graph transformation, where the host graph is nondistributed.
    \item[bruggink2014]  Termination, the absence of infinite computations, is a property that is required in many applications, in particular in model transformation, algorithms and protocol specifications. Many of these applications are naturally modeled by graph transformation systems.
    \item[bruggink2014] although termination is a basic notion of any computational formalism, it has not received a lot of attention in the graph transformation community.
    \item[bruggink2014] we consider graph transformation from a theoretical point of view. This has the disadvantage of making results harder to obtain, but the advantage of being more broadly applicable. 
\end{itemize}

 
\chapter{Structure of the thesis}
 
%%%%%%%%%%%%%%%%%%%%%%%%%%%%%%%%%%%%%%%%
% Binary Relations and Orders
%%%%%%%%%%%%%%%%%%%%%%%%%%%%%%%%%%%%%%%%%%

\chapter{Binary Relations and Orders}  

We define binary relations and orders on a collection of objects instead of on a set. This is because later we will define rewriting relations as binary relations, and the collection of objects subject to be rewritten in a rewriting system is not necessarily a set. For example, the collection of finite directed edge-labeled graphs is not a set, but a class.

\begin{notation}
    Let $C$ and $C'$ be collections of objects. We denote $c \mathop{\in} C$ for~\enquote{$c$ is an element of C}, $C \mathop{\cup} C'$ the smallest collection contains all elements in $C$ and all elements in $C'$, $C \mathop{\cap} C'$ the smallest collection contains all elements in $C$ that are also in $C'$, $C * C'$ the collection of ordered pairs with first element in $C$ and second element in $C'$.  
  \end{notation} 
  
  \begin{definition}[Binary relation]
    \label{def:binary_relation:binary_relation}
    Let $D$ be a collection of objects. A mathematical structure \( (D, \mathcal{R}) \) where $\mathcal{R}$ is a collection of objects from $D * D$ is called a \textbf{binary relation} on $D$. 
    
    For an object in $\mathcal{R}$ with first element $x \mathop{\in} D$ and second element $y\in D$, we write $x \mathcal{R} y$, $\mathcal{R}(x,y)$ or $(x,y) \mathop{\in} \mathcal{R}$. 
    
    When $D$ is irrelevant, we say simply that $\mathcal{R}$ is a binary relation.
  \end{definition} 
   
  \begin{definition}[Reflexivity, 
    %  Antisymmetry, 
     Transitivity]
    \label{def:binary_relation:reflexivity_transitivity}
    A binary relation \( \mathop{\to} \) on \(C\) is said to be \textbf{reflexive} if for every object \(a \mathop{\in} C\), we have \(a \mathop{\to} a\); \textbf{transitive} if for all objects \( a, b, c \mathop{\in} C\), \( a \mathop{\to} b \mathop{\land} (b \mathop{\to} c) \) implies \(a \mathop{\to} c\).
  %       \item \textbf{antisymmetric} if \(
  % \forall a, b \mathop{\in} A, \, (a \mathop{\to} b) \mathop{\land} (b \mathop{\to} a) \implies a \mathop{=} b.
  % \)
  \end{definition}
  
  \begin{definition}[$\mathcal{R}$-sequence]
    \label{def:binary_relation:sequence}
    Let \(\mathcal{R}\) be a binary relation on $S$.
    A \textbf{\( \mathcal{R} \)-sequence} is either a finite sequence \( \left( s_i \right)_{0 \leq i \leq m} \) of elements in $S$ such that \(s_i \mathcal{R} s_{i+1}\) for each \( 0 \leq i \leq m-1\), or an infinite sequence \((s_i)_{i \mathop{\in} \mathbb{N}}\) of elements in $S$ such that \(s_n \mathcal{R} s_{n+1}\) for each \(i \mathop{\in} \mathbb{N}\).
\end{definition}

\begin{definition}[Well-founded binary relation]
    \label{def:binary_relation:well_founded}
    A binary relation $\to$ on a collection $S$ of object is said to be \textbf{well-founded} if there is no infinite $\to$-sequence. 
    % In other words, there does not exist an infinite sequence of elements $ (s_i)_{i \mathop{\in} \mathbb{N}} $ where $s_i \mathop{\in} S $ such that $s_i \mathop{\to} s_{i+1}$ for $i \mathop{\in} \mathbb{N}$.
\end{definition}

\begin{definition}[Transitive closure]
    \label{def:binary_relation:transitive_closure}
    The \textbf{transitive closure} of a binary relation $\to$, denoted $\to^+$, is the smallest transitive relation that include \( \mathop{\to} \).
  \end{definition}
  
  \begin{definition}[Reflexive-transitive closure]
    \label{def:binary_relation:reflexive_transitive_closure}
    The \textbf{reflexive transitive closure} of a binary relation $\to$, denoted $\to^*$, is the smallest reflexive and transitive relation that include \( \mathop{\to} \).
  \end{definition}

  \begin{definition}[Homomorphisms of binary relations]
    \label{def:binary_relation:homomorphism}
    Let $\to_\mathcal{A}$ be a binary relation on a collection $A$ of objects and $\to_\mathcal{B}$ be a binary relation on a collection $B$ of objects. A \textbf{homomorphism} from $\to_\mathcal{A}$ to $\to_\mathcal{B}$ is a function \( h: A \mathop{\to} B \) such that for all \( a, b \mathop{\in} A \), if \( a \mathop{\to} _\mathcal{A} b \) then \( h(a) \mathop{\to} _\mathcal{B} h(b) \).
  \end{definition}
  
\begin{proposition}[Proving well-foundedness]
  \label{prop:binary_relation:proving_well_foundedness}
  Let \(\to\) be binary relation and let $\leadsto$ be a well-founded binary relation. If there is a homomorphism from \(\to\) to \(\leadsto\), then \(\to\) is well-founded.
\end{proposition}
  
%%%%%%%%%%%%%%%%%%%%%%%%%%%%%%%%%%%%%%%%
%Rewriting relations and systems, and termination
%%%%%%%%%%%%%%%%%%%%%%%%%%%%%%%%%%%%%%%%%%

\chapter{Rewriting systems}
    \label{sec:category_of_rewriting_systems}      
      \begin{definition}[Rewriting relation]
        \label{def:ars}
        Let $D$ be a collection of objects.  
        A \textbf{rewriting relation} $(D, \rightarrow)$ on $D$ is defined as a binary relation on $D$.  
        The set $D$ is referred to as the \textbf{domain} of the rewriting relation.  
        For all objects $x, y \mathop{\in} D$, if $x \mathop{\rightarrow} y$ then we say that $x$ \textbf{rewrites} to $y$ (or there exists a \textbf{rewriting step} from $x$ to $y$).
    \end{definition}

      In the litterature, a rewriting relation is also called abstract rewriting systems or abstract reduction systems \cite{nipkow1998term,terese2003term} and is usually defined on a set of objects. We defined it on a collection of objects to cope with rewriting systems on graphs, which form a collection of objects but not a set.
       
      \begin{definition}[Rewriting system]
        \label{def:rewriting_system_no_framework}
        A \textbf{rewriting system} is defined as a mathematical structure $(D, \mathcal{R},M,W,\mathfrak{M},\mathfrak{W},\mathfrak{I})$ where 
        \begin{itemize}
          \item $D$ is a collection of objects. 
          \item $\mathcal{R}$ is a set of objects, called \textbf{rewriting rules},
          
          \item $M$ is a collection of objects, called \textbf{matches}, 
          
          \item $W$ is a collection of objects, called \textbf{witnesses of rewriting steps},  
          
          \item $\mathfrak{M}$ is a function, called \textbf{match mechanism}, which associates to each $(r,d) \mathop{\in} \mathcal{R} \mathop{\times} D$ a set $\mathfrak{M}(r,d)$ of matches in $M$, called \textbf{matches of the rewriting rule $r$ in the object $d$}, 
           
          \item $\mathfrak{W}$ is a function, called \textbf{witness function}, which associates to each $(r,d,m) \mathop{\in} \mathcal{R} \mathop{\times} D \mathop{\times} M$ a witness of rewriting step $\mathfrak{W}(r,d,m)$ in $W$, called \textbf{witness of the rewriting step defined by the match $m$ of the rule $r$ in the object $d$},

          \item $\mathfrak{I}$ is a function, called \textbf{interpretation function},  
              which associates to each witness $w \mathop{\in} W$, an element 
              in $D \mathop{\times} D$, called \textbf{rewriting step witnessed by $w$}.
        \end{itemize} 

        For every rule $\rho \mathop{\in} \mathcal{R}$, it induces a rewriting relation $\mathop{\Rightarrow}_\rho$ that is defined as follows: for all objects $s, t \mathop{\in} D$, the rewriting step $s \mathop{\Rightarrow}_\rho t$ exists
          iff there exists a match $m \mathop{\in} \mathfrak{M}(\rho,s)$ of the rule in the object $s$ such that $\mathfrak{W}(r,s,m)$ witnesses the rewriting step, i.e. $\mathfrak{I}(\mathfrak{W}(r,s,m)) \mathop{=} (s,t)$. 
          
          The \textbf{rewriting relation} $\mathop{\Rightarrow}_\mathcal{R}$ is defined as follows: for all objects $s, t \mathop{\in} D$, the rewriting step $s \mathop{\Rightarrow}_\mathcal{R} t$ exists iff $s \mathop{\Rightarrow}_\rho t$ for some rule $\rho \mathop{\in} \mathcal{R}$.

          For a rewriting step $s \mathop{\Rightarrow}_\rho t$ with match $m$, we denote $s \mathop{\Rightarrow}_\rho^m t$ to indicate the match $m$ used to rewrite $s$ to $t$.

          When the context makes it clear, we say that $\mathcal{R}$ is a rewriting system.
      \end{definition}
      
      When defining a rewriting system, one may want to impose additional constraints on the rewriting relation. This can be done by imposing a more restrictive match mechanism. 
      
      However, since rewriting is an intuitive concept, when we talk about a rewriting system on a given domain $D$, we have an intuition of what a match mechanism should be. 

      As a solution, Endrullis et al. proposed in \cite{endrullis2024generalized} to exclude certain witnesses using rewriting frameworks.
      
      \begin{definition}[Rewriting system in a rewriting framework]
        \label{def:rewriting_system_with_framework}
        A \textbf{rewriting system in a rewriting framework} is defined as a mathematical structure $(D, \mathcal{R},M,W,\mathfrak{M},\mathfrak{W},\mathfrak{I},\mathfrak{F})$ where
        \begin{itemize}
          \item $(\mathcal{R},M,W,\mathfrak{M},\mathfrak{W},\mathfrak{I})$ is a rewriting system,
          \item $\mathfrak{F}$ is a mapping, called \textbf{rewriting framework}, which associates to each rule $r \mathop{\in} R$, 
             a collection of witnesses in $W$.
        \end{itemize}
        
        For every rule $\rho \mathop{\in} \mathcal{R}$, it induces a rewriting relation $\mathop{\Rightarrow}_{\mathfrak{F},\rho}$ that is defined as follows: for all objects $s, t \mathop{\in} D$, the rewriting step $s \mathop{\Rightarrow}_{\mathfrak{F},\rho} t$ exists
          iff $s \mathop{\Rightarrow}_\rho^m t$ for some match $m \mathop{\in} \mathfrak{M}(\rho,s)$ and $\mathfrak{W}(r,s,m) \mathop{\in} \mathfrak{F}(\rho)$.

          The \textbf{rewriting relation} $\mathop{\Rightarrow}_{\mathfrak{F},\mathcal{R}}$ is defined as follows: for all objects $s, t \mathop{\in} D$, the rewriting step $s \mathop{\Rightarrow}_{\mathfrak{F},\mathcal{R}} t$ exists iff $s \mathop{\Rightarrow}_{\mathfrak{F},\rho} t$ for some rule $\rho \mathop{\in} \mathcal{R}$.

          For a rewriting step $s \mathop{\Rightarrow}_{\mathfrak{F},\rho} t$ with match $m$, we denote $s \mathop{\Rightarrow}_{\mathfrak{F},\rho}^m t$ to indicate the match $m$ used to rewrite $s$ to $t$.

          When the context makes it clear, we say that $\mathcal{R}$ is a rewriting system.
      \end{definition}
      
      % \trackedtext{useful ?? \\
      % The following correspond to the rewriting relation definition by rewriting framework proposed by Endrullis et Overbeek in \cite[Definition 5.2]{endrullis2024generalized}.
      % \begin{proposition}
      %   Let $D$ be a collection of objects.
      %   The structure $(\mathcal{R},\mathfrak{I},\mathfrak{F})$ where
      %   \begin{itemize}
      %     \item $\mathcal{R}$ is a set of objects, called \textbf{rewriting rules}, 
      %     \item $\mathfrak{F}$ is a function, called \textbf{rewriting framework},  which associates to each rule $r \mathop{\in} R$, 
      %        a class of objects,
      %     \item $\mathfrak{I}$ is a function, called \textbf{interpretation function},  
      %     which associates to each object $ w \mathop{\in} \mathcal{F}(r)$ with $r \mathop{\in} R$ an element 
      %     in $D \mathop{\times} D$, called \textbf{rewriting step using rule $r$ witnessed by $w$},
      %   \end{itemize}
      %    defines a rewriting relation $\to$ as follows: $ \mathop{\to} \mathop{=} \bigcup_{\rho \mathop{\in} R} \mathfrak{I}(\mathfrak{F}(\rho))$.
      % \end{proposition}
      % }
      
    % \color{red}


An abstract rewriting system is a mathcal structure in a first order language.
Intuitively, a mathcal structure is a set equipped with a binary relation on the set. 
A set is a collection of all elements that satisfy a certain given property. 
A binary relation on a set is a set of pairs of elements of that set. 
For more information on set theory, the reader is refer to \cite{jech2006set}. Bellow we give the formal definition these concepts.
  
\begin{definition}[Language~\text{\cite[Def 1.1.1]{marker2006model}}]
  A \textbf{language} \( \mathcal{L} \) is given by specifying the following data:
  \begin{itemize}
      \item  a set of function symbols \( \mathcal{F} \) and positive integers \( n_f \) for each \( f \mathop{\in} \mathcal{F} \);
      \item  a set of relation symbols \( \mathcal{R} \) and positive integers \( n_R \) for each \( R \mathop{\in} \mathcal{R} \);
      \item  a set of constant symbols \( \mathcal{C} \).
  \end{itemize}
\end{definition}

\begin{example}[\ \text{\cite{marker2006model}}]
  Consider the structure of the natural numbers with addition and distinguished elements 0 and 1. The natural language for studying this structure is the language where we have a binary function symbol for addition and constant symbols for 0 and 1. We would write sentences such as $\forall x \exists y (x \mathop{=} y\mathop{+}y \lor x \mathop{=} y\mathop{+}y +1)$, which we interpret as the assertion that “every number is either even or 1 plus an even number.”
\end{example}

\begin{definition}[Language of abstract rewriting systems]
  \label{def:l_ars}
  The first order \textbf{language of abstract rewriting systems} is $\mathcal{L}_\text{RS} \mathop{=} (\rightarrow)$ where $\rightarrow$ is a binary relation symbol.
\end{definition}
\begin{definition}[Structure~\text{\cite[Def 1.1.2]{marker2006model}}]
  An \( \mathcal{L} \)-\textbf{structure} \( \mathcal{M} \) is given by the following data:
 \begin{itemize}
     \item[i)] a nonempty set \( M \) called the \textit{universe, domain, or underlying set} of \( \mathcal{M} \);
     \item[ii)] a function \( f^{\mathcal{M}} : M^{n_f} \mathop{\to} M \) for each \( f \mathop{\in} \mathcal{F} \);
     \item[iii)] a set \( R^{\mathcal{M}} \mathop{\subseteq} M^{n_R} \) for each \( R \mathop{\in} \mathcal{R} \);
     \item[iv)] an element \( c^{\mathcal{M}} \mathop{\in} M \) for each \( c \mathop{\in} \mathcal{C} \).
 \end{itemize}
\end{definition}
We refer to \( f_{\mathcal{M}} \), \( R_{\mathcal{M}} \), and \( c_{\mathcal{M}} \) as the \textit{interpretations} of the symbols \( f \), \( R \), and \( c \). We often write the structure as \( \mathcal{M} \mathop{=} (M, f_{\mathcal{M}}, R_{\mathcal{M}}, c_{\mathcal{M}} : f \mathop{\in} \mathcal{F}, R \mathop{\in} \mathcal{R}, c \mathop{\in} \mathcal{C}) \). We will use the notation \( \mathcal{A}, \mathcal{B}, \mathcal{M}, \mathcal{N}, \dots \) to refer to the underlying sets of the structures \( A, B, M, N, \dots \).
\begin{definition}[Binary Relation]
  Let $S$ and $T$ be two sets. A \textbf{binary relation} from $S$ to $T$ is a subset of $S \mathop{\times} T$. A binary relation on $S$ is a binary relation from $S$ to $S$.
\end{definition} 

\begin{definition}[Abstract Rewriting System]
  \label{def:ars}
  An \textbf{rewriting system} is a $\mathcal{L}_\text{RS}$-structure $\mathcal{R} \mathop{=} (R,\to_\mathcal{R})$ where $R$ is a set and $\to_\mathcal{R}$ a binary relation on $R$. Elements $(x,y)$ in $\rightarrow_\mathcal{S}$, denoted $x \rightarrow_\mathcal{R} y$, will be called \textbf{rewriting steps}. 
\end{definition}

\begin{definition}[Reflexivity,
  %  Antisymmetry, 
   Transitivity]
  A binary relation \( \mathop{\to} \mathop{\subseteq} S\) is said to be
  \begin{itemize}
      \item \textbf{reflexive} if \(
          \forall a \mathop{\in} S, \, (a \mathop{\to} a)
          \)
%       \item \textbf{antisymmetric} if \(
% \forall a, b \mathop{\in} A, \, (a \mathop{\to} b) \mathop{\land} (b \mathop{\to} a) \implies a \mathop{=} b.
% \)
      \item \textbf{transitive} if \(
\forall a, b, c \mathop{\in} S, \, (a \mathop{\to} b) \mathop{\land} (b \mathop{\to} c) \implies (a \mathop{\to} c).
\)
  \end{itemize}
\end{definition}

\begin{definition}[Reflexive-Transitive Closure]
  The \textbf{reflexive transitive closure} of a binary relation $\to$, denoted $\to^*$, is the smallest reflexive and transitive relation that include \( \mathop{\to} \).
\end{definition}

\begin{definition}[Rewriting Relation]
   A rewriting system $\mathcal{R} \mathop{=} (R,\to_\mathcal{R})$ defines a \textbf{rewriting relation} $\to^{*}_{\mathcal{R}}$.
\end{definition}

\begin{definition}[Category of Rewriting Systems]
  Rewriting systems and homomorphisms between them form a category, denoted $\mathcal{RS}$.   
\end{definition}

\color{red}
\begin{definition}[Rewriting System Generated by Rewriting Rules in a Rewriting Framework]
  Let $A$ be a set and \( \mathfrak{F}: A \mathop{\to} \mathcal{P}(S^2) \) be a function.
  The rewriting system $\mathcal{S} \mathop{=} (S, \mathop{\to} _{\mathcal{S}})$ generated by the set of \textbf{rewriting rules} \( R \mathop{\subseteq} A \) within the \textbf{rewriting framework} \( \mathfrak{F}\) is defined by
  \[
      \mathop{\to} _\mathcal{S} \mathop{=} \bigcup_{\rho \mathop{\in} R} \mathfrak{F}(\rho)
    \]

  \todo{match mechanism:reviser}
  Let $M$ be a set, whose elements are called matches. Let \( \mathfrak{m}: R \mathop{\times} S \mathop{\to} \mathcal{P}(M) \) be a function such that for every \( (\rho,x) \mathop{\in} R \mathop{\times} S \), the set \( \mathfrak{m}(\rho, x) \) and the set $\{x \rightarrow_\mathcal{S} y \mathop{\in} \mathfrak{F}(\rho) \mid y \mathop{\in} S\}$ are in bijection. We say that \( \mathfrak{m} \) is a \textbf{match mechanism} and that elements in $\mathfrak{m}(\rho,x)$ are \textbf{matches} of the rule $\rho$ in the object $x$.

\end{definition}
\color{black}\color{black}
    \begin{definition}[Terminating rewriting relations]
        \label{def:rewriting_relation:termination}
    A rewriting relation is \textbf{terminating} if its transitive closure is well-founded.
    \end{definition}
    
    \begin{definition}[Terminating rewriting system]
        \label{def:rewriting_system:termination}
        A rewriting system 
        %  in a framework $\mathfrak{F}$ 
         is \textbf{terminating} if its rewriting relation
        %   $\to_{\mathcal{R},\mathfrak{F}}$ 
          is terminating.
    \end{definition}
    
    \begin{proposition}[Proving Termination of rewriting systems]
        Let $\mathcal{R}$ be a rewriting system and $\mathcal{S}$ be a terminating rewriting system.
        If there exists a homomorphism from the rewriting relation of $\mathcal{R}$ to the rewriting relation of $\mathcal{S}$, then $\mathcal{R}$ is terminating.  
    \end{proposition}
    
    In practice, it is often more convient to embed the rewriting system into a terminating rewriting system via a sequence of homomorphism.
    \begin{corollary}[Proving termination of rewriting systems]
      A rewriting system $\mathcal{R}$ is terminating if there is a terminating rewriting system $\mathcal{S}$ and homomorphisms $h_1,\ldots, h_n$ such that $h_n \circ \ldots \circ h_1$ is a homomorphism from $\mathcal{R}$ to $\mathcal{S}$.
    \end{corollary}
    
    \begin{definition}[Relative termination of rewriting relations]
        \label{def:rewriting_relation:relative_termination}
        Let \( \mathop{\to} \) and \( \leadsto \) be two rewriting relations. We say that \(\to\) is \textbf{terminating relative to} \(\leadsto\) if every \( \left( \mathop{\to} \mathop{\cup} \leadsto \right) \)-sequence contains only finitely many \(\to\)-steps.
    \end{definition}
    
    \begin{definition}[Relative termination of rewriting systems]
        \label{def:rewriting_system:relative_termination}
        Let $D$ be a collection of objects.
        Let \( \mathcal{A} \) and \( \mathcal{B} \) be two rewriting systems on $D$. We say that \(\mathcal{A}\) is \textbf{terminating relative to} \(\mathcal{B}\) if the rewriting relation induced by \(\mathcal{A}\) is terminating relative to the rewriting relation induced by \(\mathcal{B}\).
    \end{definition}
    
    % \begin{proposition}[Proving Relative Termination]
    %     Let \((S, \mathop{\to} \mathop{\cup} \leadsto) \) be an abstract rewriting system. \((S,\to)\) is \textbf{terminating relative to} \((S, \leadsto)\) if there is a terminating rewriting system $(R,>)$ and a function $h : S \mathop{\to} R$ such that
    %     \begin{itemize}
    %       \item $h$ is abstract rewriting system homomorphism from $(S,\to)$ to  $(R,>)$
    %       \item $h$ is abstract rewriting system homomorphism from $(S,\leadsto)$ to  $(R,\geq)$
    %     \end{itemize} 
    %     where $\geq$ is the reflexive closure of $>$.
    %   \end{proposition}
    \begin{definition}[Proving relative termination]
        \label{def:rewriting_system:proving_relative_termination}
        Let \( \mathcal{A} \) and \( \mathcal{B} \) be two rewriting systems on $D$. Let $>$ and $\geq$ be two rewriting relations on $D'$ such that $>$ is terminating relative to $\geq$. If there exists a function $h : D \mathop{\to} D'$ such that: (1) $h$ is a homomorphism from the rewriting relation of \(\mathcal{A}\) to $>$, and (2) $h$ is a homomorphism from the rewriting relation of \(\mathcal{B}\) to $\geq$, then \(\mathcal{A}\) is terminating relative to \(\mathcal{B}\).
    \end{definition}
    
    % \begin{definition}[Relative terminating rewriting systems]
    %     Let $D$ be a collection of objects.
    %     Let \( \mathcal{A} \) and \( \mathcal{B} \) be two rewriting systems on $D$.and $\mathfrak{F}$ be a DPO rewriting framework.
    %      We say that \(\mathcal{A}\) is \textbf{terminating relative to} \(\mathcal{B}\) if the rewriting relation \(\to_{\mathcal{A},\mathfrak{F}} \) is terminating relative to the rewriting relation \(\to_{\mathcal{B},\mathfrak{F}}\).
    % \end{definition}
    
    % \begin{definition}[Relative termination]
    % Let \((S, \mathop{\to} \mathop{\cup} \leadsto) \) be a rewriting system. We say that \((S,\to)\) is \textbf{terminating relative to} \((S, \leadsto)\) if every derivation $s_1 ( \mathop{\to} \mathop{\cup} \leadsto) s_2 ( \mathop{\to} \mathop{\cup} \leadsto) \hdots $ contains only finitely many $\to$-steps.
    % \end{definition} 
    
    % \color{red}
    % \begin{definition}[Pre-order \cite{davey2002introduction}]
    %   A reflexive and transitive binary relation $\geq$ on a set is called a pre-order. It give rise to a relation $>$ of strict inequality: $x \mathop{>} y$ iff $x \mathop{\geq} y$ and $x \mathop{\neq} y$. 
    % \end{definition} 
    
    % \begin{definition}[Partial Order]
    %   A reflexive, antisymmetric, transitive binary relation is called a partial order.
    % \end{definition}
    
    % \begin{definition}[Strict Partial Order]
    %   A irreflexive, transitive binary relation is called a strict partial order.
    % \end{definition}
    
    % \begin{example}[Lexicographic Product \cite{nipkow1998term}]
    %   Consider two rewriting systems $(A, >)$ and $(B, \mathop{\succ})$. The lexicographic product $(A \mathop{\times} B, \gg)$ is defined by the relation \((x,y) \gg (x',y')\) if and only if \(x \mathop{>} x'\) or, \(x \mathop{=} x'\) and \(y \mathop{\succ} y'\). If \(>\) and \(\mathop{\succ}\) are strict partial orders, then \(\gg\) is also a strict partial order. Furthermore, if both \((A, >)\) and \((B, \mathop{\succ})\) terminate, then the lexicographic product \((A \mathop{\times} B, \gg)\) will also terminate.
    % \end{example}
    
    % \begin{example}[Finite Lexicographic Product \cite{nipkow1998term}]
    %   By iteration, we can form lexicographic product over any number of rewriting systems $(A_i, >_i), i=1,\ldots,n$. For $n>1$, the lexicographic product $(A_1 \mathop{\times} (A_2 \mathop{\times} \ldots \mathop{\times} A_n),>_{1\ldots n})$ is the lexicographic product of $(A_1, >_1)$ and $(A_2 \mathop{\times} \ldots \mathop{\times} A_n,>{2\ldots n})$. Unwinding the recursion, we get 
    %   $(x_1,\ldots,x_n) >_{1\ldots n} (x_1',\ldots, x_n')$ iff $\exists k < n.(\forall i < k. x_i \mathop{=} x_i') \mathop{\land} x_k >_k x_k'$.
    %   If $(A_1,>_1), \ldots, (A_n,>_n)$ are strict partial orders, then so is $(A_1 \mathop{\times} \ldots \mathop{\times} A_n, >_{1\ldots n})$. 
    %   If $(A_1,>_1), \ldots, (A_n,>_n)$ terminates, then so does $(A_1 \mathop{\times} \ldots \mathop{\times} A_n, >_{1\ldots n})$. 
    % \end{example}
    
    % \begin{example}[Finite Lexicographic Product \cite{nipkow1998term}]
    %   A lexicographic product can be extended iteratively over any number of rewriting systems \((A_i, >_i)\) for \(i \mathop{=} 1, \ldots, n\). For \(n \mathop{>} 1\), the lexicographic product \((A_1 \mathop{\times} A_2 \mathop{\times} \dots \mathop{\times} A_n, >_{1\ldots n})\) is defined as the lexicographic product of \((A_1, >_1)\) and \((A_2 \mathop{\times} \dots \mathop{\times} A_n, >_{2\ldots n})\). Expanding this recursion, we have 
    %   \((x_1, \dots, x_n) >_{1\ldots n} (x_1', \dots, x_n')\) if and only if there exists \(k < n\) such that \((\forall i < k . x_i \mathop{=} x_i') \mathop{\land} x_k >_k x_k'\).
    %   If each \((A_1, >_1), \ldots, (A_n, >_n)\) is a strict partial order, then \((A_1 \mathop{\times} \dots \mathop{\times} A_n, >_{1\ldots n})\) is also a strict partial order. 
    %   Moreover, if \((A_1, >_1), \ldots, (A_n, >_n)\) terminate, then the lexicographic product \((A_1 \mathop{\times} \dots \mathop{\times} A_n, >_{1\ldots n})\) also terminates.
    % \end{example}
    
    
    % \begin{example}[Strings of Arbitrary but Finite Length with Lexicographic Order \cite{nipkow1998term}]
    %   Given an rewriting system \((A, >)\), the lexicographic order \((A^*, >_{lex})\) is defined as follows: \(x \mathop{>} x'\) if and only if \((|x| \mathop{>} |x'|) \lor (|x| \mathop{=} |x'| \mathop{\land} y \mathop{\succ} y')\).
    %   If \(>\) is a strict partial order, then \(\gg\) is also a strict partial order. Furthermore, if both \((A, >)\) and \((B, \mathop{\succ})\) are terminating rewriting systems, then their lexicographic product \((A \mathop{\times} B, \gg)\) also terminates.
    % \end{example}
    
    % \begin{definition}[Multiset Order \cite{nipkow1998term}]
    %   A \textbf{multiset} is a generalization of a set where elements can occur multiple times. Formally, it is defined as a function \( M: A \mathop{\rightarrow} \mathbb{N} \), where for each \(a \mathop{\in} A\), \(M(a)\) represents the number of times \(a\) appears in the multiset \(M\). The collection of all multisets over a set \(A\) is denoted by \(\mathcal{M}(A)\).
    
    %   Given a strict partial order \(>\) on a set \(A\), the corresponding multiset order \(>_{mul}\) on \(\mathcal{M}(A)\) is defined as follows: \(M >_{mul} N\) if and only if there exist multisets \(X, Y \mathop{\in} \mathcal{M}(A)\) such that:
    %   \begin{itemize}
    %       \item \(X \mathop{\subseteq} M\) and \(X \mathop{\neq} \emptyset\),
    %       \item \(N \mathop{=} (M - X) \mathop{\cup} Y\), and
    %       \item \(\forall y \mathop{\in} Y, \exists x \mathop{\in} X \text{ such that } x \mathop{>} y\).
    %   \end{itemize}
    %   If the rewriting system \((A, >)\) terminates, then the multiset order \((\mathcal{M}(A), >_{mul})\) also terminates.
    % \end{definition}
    % \color{black}
% In practice, it is often more convient to embed the rewriting system into a terminating rewriting system via a morphism that is the composition of a serie of morphisms.
% \begin{corollary}[Proving termination]
%   A rewriting system $\mathcal{R}$ is terminating if there is a terminating rewriting system $\mathcal{S}$ and functors $h_1,\ldots, h_n$ such that $h_1 \mathop{\star} \ldots \mathop{\star} h_n : \mathcal{R} \mathop{\to} \mathcal{S}$
% \end{corollary}

   
%   \begin{proposition}[Proving Relative Termination]
%     Let \((S, \mathop{\to} \mathop{\cup} \leadsto) \) be an abstract rewriting system. \((S,\to)\) is \textbf{terminating relative to} \((S, \leadsto)\) if there is a terminating rewriting system $(R,>)$ and a function $h : S \mathop{\to} R$ such that
%     \begin{itemize}
%       \item $h$ is abstract rewriting system homomorphism from $(S,\to)$ to  $(R,>)$
%       \item $h$ is abstract rewriting system homomorphism from $(S,\leadsto)$ to  $(R,\geq)$
%     \end{itemize} 
%     where $\geq$ is the reflexive closure of $>$.
%   \end{proposition}
 \chapter{Graphs}
    \label{sec:graphs}
    \begin{definition}[Unlabeled graph~\cite{barr1990category}]
        \label{def:graph:unlabeled}
        An \textbf{unlabeled graph} \( G \) consists of a collection of \textbf{nodes} (also called \textbf{objects}) and a collection of \textbf{edges}, each edge $e$ equipped with a \textbf{source} (or \textbf{domain}) node, denoted $\opn{src}(e)$, and a \textbf{target} (or \textbf{codomain}) node, denoted $\opn{dst}(e)$. 
        For an unlabeled graph \( G \), we denote by \( G_0 \) its collection of nodes, \( G_1 \) its collection of edges, \( \operatorname{dom}:G_1{\to}G_0 \) the domain function, and \( \operatorname{cod}:G_1{\to}G_0 \) the codomain function. An unlabeled graph is \textbf{finite} if \( G_0 \) and \( G_1 \) are both finite sets.
        We write \( a: s \mathop{\to} t \) to indicate that \( a \) is a directed edge from \( s \) to \( t \). 
    \end{definition}   
    \begin{example}
      A unlabeled graph is a directed multigraph as illustrated below.
       
        \begin{center}
          \resizebox{0.45\textwidth}{!}{
          \begin{tikzpicture}
              \graphbox{\(\)}{00mm}{-20mm}{45mm}{20mm}{2mm}{-5mm}{
                  \coordinate (o) at (-5mm,-8mm); 
                  \node[draw,circle] (l1) at ($(o)+(-10mm,0mm)$) {};
                  % \node[draw,circle] (l2) at ($(l1)+(3,0)$) {};
                  \node[draw,circle] (l3) at ($(l1)+(1,0)$) {};
                  \node[draw,circle] (l4) at ($(l1)+(2,0)$) {};
                  \draw[->] (l1) edge[bend right]  (l3);
                  \draw[->] (l1) edge[bend left] (l3);
                  \draw[->] (l3) -- (l4);
                  % \draw[->] (l4) -- (l2) node[midway,above] {$a$};
              }  
          \end{tikzpicture}
      }
      \end{center} 
    \end{example}
    % A homomorphism of unlabeled graphs is a mapping between the nodes and edges of two graphs that preserves the graph structure.
    \begin{definition}
        \label{def:unlabeled_graph:homomorphism}
        Let \( G \) and \( H \) be unlabeled graphs. A \textbf{homomorphism of unlabeled graphs} $h: G \mathop{\to} H$ is a pair of functions $h_0: G_0 \mathop{\to} H_0 $ and $h_1: G_1 \mathop{\to} H_1$ such that for every edge \( a: s \mathop{\to} t \) in \( G \), we have \( h_1(a) : h_0(s) \mathop{\to} h_0(t) \) in \( H \).
    \end{definition}
    \begin{definition}
        \label{def:graph}
        Let \(\Sigma\) be a finite set of labels. A \textbf{labeled graph} is an ordered pair \((G,\lambda)\) where \( G \) is an unlabeled graph and \( \lambda : G_1 \mathop{\rightarrow} \Sigma\) is an edge-labeling function. 
        It is called \textbf{finite} if its underlying unlabeled graph is finite.  
    \end{definition}
    \begin{example} An labeled graph is edge-labeled directed multigraph as illustrated below.
       
        \begin{center}
          \resizebox{0.45\textwidth}{!}{
          \begin{tikzpicture}
              \graphbox{\( \)}{00mm}{-20mm}{45mm}{20mm}{2mm}{-5mm}{
                  \coordinate (o) at (-5mm,-8mm); 
                  \node[draw,circle] (l1) at ($(o)+(-10mm,0mm)$) {};
                  % \node[draw,circle] (l2) at ($(l1)+(3,0)$) {};
                  \node[draw,circle] (l3) at ($(l1)+(1,0)$) {};
                  \node[draw,circle] (l4) at ($(l1)+(2,0)$) {};
                  \draw[->] (l1) edge[bend right]  node[midway,below] {$a$} (l3);
                  \draw[->] (l1) edge[bend left] node[midway,above] {$a$}  (l3);
                  \draw[->] (l3) -- (l4) node[midway,above] {$b$};
                  % \draw[->] (l4) -- (l2) node[midway,above] {$a$};
              }  
          \end{tikzpicture}
      }
      \end{center} 
    \end{example}
    By $a : s\overset{l}{\rightarrow} t$, we denote the arrow $a$ labeled by $l$ from $s$ to $t$. Unless otherwise specified, the term \enquote{graph} refers to the graphs finite. Note that unlabeled graphs can be regarded as labeled graphs. 

    A homomorphism of labeled graphs is a homomorphism of unlabeled graphs that preserves the labels assigned to the edges.
    \begin{definition}
        \label{def:graph:homomorphism}
        Let \( (G,\lambda) \) and \( (H,\lambda') \) be labeled graphs. A \textbf{homomorphism of labeled graphs} $h:(G,\lambda) \mathop{\rightarrow} (H,\lambda')$ is a homomorphism of unlabeled graphs such that for each edge \( a \) in \( G \), we have \( \lambda (a) \mathop{=} \lambda' (h_1 (a)) \).
    \end{definition}
    \begin{notation}
        We use the notation from~\cite[Notation 1]{overbeek2023apbpotutorial} to visualize edge-labeled graph homomorphisms. Labeled graphs are enclosed in boxes with their names displayed in the top-left corner. Nodes and edges are assigned subsets of \(\mathbb{N}\) as identifiers, and these identifiers are chosen such that: (i) Each node or edge \( y \) in the codomain graph is assigned the union of the identifiers of all nodes or edges in the domain graph that are mapped to \( y \); (ii) The graph homomorphism is uniquely determined by this assignment.
        
        \noindent To further improve readability, we represent sets by listing their elements. Additionally, we omit identifiers when doing so does not cause confusion. Furthermore, if a node has a unique loop, we omit the loop and represent the corresponding label as a node-label. 
        This is illustrated in the following representation of a homomorphism \( h: G \mathop{\to} H \).
        
        \begin{center}
            \resizebox{0.45\textwidth}{!}{
            \begin{tikzpicture}
                \graphbox{\( G \)}{00mm}{-20mm}{45mm}{25mm}{2mm}{-10mm}{
                    \coordinate (o) at (-5mm,-8mm); 
                    \node[draw,circle] (l1) at ($(o)+(-10mm,0mm)$) {1};
                    \node[draw,circle] (l2) at ($(l1)+(3,0)$) {2};
                    \node[draw,circle] (l3) at ($(l1)+(1,0)$) {3};
                    \node[draw,circle] (l4) at ($(l1)+(2,0)$) {4};
                    \draw[->] (l1) -- (l3) node[midway,above] {$a$};
                    \draw[->] (l3) -- (l4) node[midway,above] {$b$};
                    \draw[->] (l4) -- (l2) node[midway,above] {$a$};
                }  
                \graphbox{\( H \)}{50mm}{-20mm}{34mm}{25mm}{2mm}{-10mm}{
                    \coordinate (o) at (0mm,-8mm); 
                    \node[draw,circle] (l1) at ($(o)+(-10mm,0mm)$) {1};
                    \node[draw,circle] (l2) at ($(l1)+(2,0)$) {2};
                    \node[draw,circle] (l3) at ($(l1)+(1,0)$) {3\ 4} node[above] {$b$};
                    \draw[->] (l1) edge node[midway,above] {$a$} (l3)  ;
                    % \draw[->] (l3) edge [loop above] node[midway,above] {$b$} (l3) ;
                    \draw[->] (l3) -- (l2) node[midway,above] {$a$};
                }      
                \node () at (48mm,-30mm) {$\rightarrow$};
            \end{tikzpicture}
        }
        \end{center} 
        In this example, the sets \(\{1\}\), \(\{2\}\), \(\{3\}\), \(\{4\}\), and \(\{3,4\}\) are represented as \(1\), \(2\), \(3\), \(4\), and \(3\ 4\), respectively. Edge identifiers are omitted.
    \end{notation} 
%%%%%%%%%%%%%%%%%%%%%%%%%%%%%%%%%%%%%%%%%
% Graph Relabeling Systems
%%%%%%%%%%%%%%%%%%%%%%%%%%%%%%%%%%%%%%%%%%

    \chapter{Graph relabeling systems (GLS)} 
      \label{sec:gls}
Graph-relabeling-systems introduced in \cite{litovsky1999graph} are rewriting systems. For a given graph, a graph relabeling system identify a subgraph of the graph and modify the labels of the nodes and edges of this subgraph. Since it never modifies the structure of the graph, it is very intuitive.

% \begin{definition}[Graph Relabeling Rule]
%   A \emph{graph relabeling rule} is a couple of labeled graphs \( ((U,\lambda),(U, \lambda')) \) where \( (U, \lambda) \) is said to be the \emph{left-hand side graph} and \( (U, \lambda')\) the \emph{right-hand side graph}.
% \end{definition}  
\begin{definition}[GLS rewriting rule]
  \label{def:gls_rule}
  A \textbf{GLS rewriting rule} is an ordered pair of labeled graphs \( \rho \mathop{=} (U,\lambda) \mathop{\to} (U, \lambda') \) that share the same underlying unlabeled graph \( U \). The labeled graph \( (U, \lambda) \) is said to be the \textbf{left-hand side graph}, denoted $\opn{lhs}(\rho)$, and \( (U, \lambda')\) the \textbf{right-hand side graph}, denoted $\opn{rhs}(\rho)$.
\end{definition} 

\begin{example}
    \label{ex:gls_rule}
    The following is an example of a GLS rewriting rule. The left-hand side graph \( (U,\lambda) \) has two nodes labeled \( 1 \) and \( 2 \), and one edge labeled \( 0 \). The right-hand side graph \( (U,\lambda') \) has two nodes both labeled \( 1 \) and one edge labeled \( 1 \).
  \begin{center} 
    \begin{tikzpicture}
        % Left side
        \node[circle, draw] (A1) at (0,0) {1};
        \node[circle, draw] (N1) at (2,0) {2};
        \draw (A1) -- (N1) node[midway, above] {0};
        \node at (0,-0.5) {A};
        \node at (2,-0.5) {N};
    
        % Arrow
        \node at (3.5, 0) {$\rightarrow$};
    
        % Right side
        \node[circle, draw] (A2) at (5,0) {1};
        \node[circle, draw] (A3) at (7,0) {2};
        \draw (A2) -- (A3) node[midway, above] {1};
        \node at (5,-0.5) {A};
        \node at (7,-0.5) {A};
    \end{tikzpicture}
    \end{center}
\end{example}

% \begin{definition}[Match]
%   Let \(((U,\lambda),(U, \lambda')) \) be a graph relabeling rule. A \emph{match} of this rule in a labeled graph \( (G, \lambda_G) \) is a graph homomorphism \( m : (U,\lambda) \mathop{\to} (G, \lambda_G) \).
% \end{definition} 

\begin{definition}[Graph relabeling]
  \label{def:gls:rewriting}
  % \textbf{GLS match mechanism} is a mapping that associates to \( (\rho, G) \) the set of all injective graph homomorphisms $m : \opn{lhs}(\rho) \rightarrowtail G$.
 
  Let \( \rho \mathop{=} (L, \lambda_L) \mathop{\to} (R,\lambda_R) \) be a GLS rewriting rule, \( (G, \lambda_G) \) a graph.

  A \textbf{GLS match} of \( \rho \) in \( G \) is an injective graph homomorphism \( m : (L, \lambda_L) \rightarrowtail (G, \lambda_G) \). 
  
  Let $\lambda'_G$ be the labeling function on nodes and edges of \( G \) such that \( \lambda'_G(x) \mathop{=} \lambda_G(x) \) for all nodes and edges \( x \) in \( G \mathop{\setminus} \operatorname{Im}(m) \), and \( \lambda'_G(x) \mathop{=} \lambda_R(m^{-1}(x)) \) otherwise.

  Let $m' : (R, \lambda_R) \rightarrowtail (G, \lambda'_G)$ be the injective graph homomorphism with \( m(x) \mathop{=} m'(x) \) for all nodes and edges \( x \) in \( G \).

  The pair $(m, m')$ of injective graph homomorphisms is a \textbf{witness} of the rewriting step from \( (G, \lambda_G) \) to \( (G, \lambda'_G) \) using the rule \( \rho \) and match \( m \).
\end{definition}
\begin{definition}[Graph relabeling system]
  Let $\mathcal{R}$ be a set of GLS rewriting rules. Let $M, \mathfrak{M}, W , \mathfrak{W}, \mathfrak{I}$ be the match, match mechanism, witness, witness function and interpretation functions introduced in~\ref{def:gls:rewriting}.

  The structure $(\mathbf{Graph},\mathcal{R},M,W,\mathfrak{M},\mathfrak{W},\mathfrak{I})$ defines a rewriting system, called \textbf{graph relabeling system (GLS)}.
\end{definition}
\todo{example ref}
\begin{example}
  \label{example:gls_spinning_tree}
  Let \(\mathcal{R}\) be a graph relabeling system with labels in \( \{N, A, 0, 1\}\) a unique graph relabelling rule:

  \begin{center}
  \begin{tikzpicture}
      % Left side
      \node[circle, draw] (A1) at (0,0) {1};
      \node[circle, draw] (N1) at (2,0) {2};
      \draw (A1) -- (N1) node[midway, above] {0};
      \node at (0,-0.5) {A};
      \node at (2,-0.5) {N};
  
      % Arrow
      \node at (3.5, 0) {$\rightarrow$};
  
      % Right side
      \node[circle, draw] (A2) at (5,0) {1};
      \node[circle, draw] (A3) at (7,0) {2};
      \draw (A2) -- (A3) node[midway, above] {1};
      \node at (5,-0.5) {A};
      \node at (7,-0.5) {A};
  \end{tikzpicture}
  \end{center}

  Let \( (G, \lambda) \) be a labeled graph such that all nodes are labeled \(N\) except for one node labeled \(A\), and all edges are labeled \(0\). 
  If we apply the rule \(r\) while it is possible, we obtain a labeled graph \( (G', \lambda') \) with all nodes labeled by \(A\) and all edges labeled by \(1\). Edges labeled by \(1\) in \( (G', \lambda') \) form a spanning tree of \( G \). 
  
\end{example}

\begin{example}
   A non-terminating rewriting rule:
        \begin{itemize}
            \item rule:  
            \begin{center} 
                \resizebox{0.9\textwidth}{!}{
                \begin{tikzpicture}
                    \graphbox{\( L \)}{0mm}{-3mm}{34mm}{15mm}{2mm}{2mm}{
                        \coordinate (o) at (0mm,-11mm); 
                        \node[draw,circle] (l1) at ($(o)+(-10mm,0mm)$) {1};
                        \node[draw,circle] (l2) at ($(l1)+(2,0)$) {2};
                        \node[draw,circle] (l3) at ($(l1)+(1,0)$) {3};
                        \draw[->] (l1) -- (l3) node[midway,above] {$a$};
                        \draw[->] (l3) -- (l2) node[midway,above] {$b$};
                    } 
            
                    \graphbox{\( K \)}{40mm}{-3mm}{34mm}{15mm}{2mm}{2mm}{
                        \coordinate (o) at (0mm,-11mm); 
                        \node[draw,circle] (l1) at ($(o)+(-10mm,0mm)$) {1};
                        \node[draw,circle] (l2) at ($(l1)+(2,0)$) {2};
                    }  
            
                    \graphbox{\( R \)}{80mm}{-3mm}{35mm}{15mm}{2mm}{2mm}{
                        \coordinate (o) at (-5mm,-11mm); 
                        \node[draw,circle] (l1) at ($(o)+(-10mm,0mm)$) {1};
                        % \node[draw,circle] (l2) at ($(l1)+(3,0)$) {2};
                        \node[draw,circle] (l3) at ($(l1)+(1,0)$) {4};
                        \node[draw,circle] (l4) at ($(l1)+(2,0)$) {2};
                        \draw[->] (l1) -- (l3) node[midway,above] {$b$};
                        \draw[->] (l3) -- (l4) node[midway,above] {$a$};
                        % \draw[->] (l4) -- (l2) node[midway,above] {$a$};
                    }    
                    \node () at (37mm,-10mm) {\( \leftarrowtail \)}; % K -> L
                    \node () at (77mm,-10mm) {\( \rightarrowtail \)}; % K -> R
                \end{tikzpicture}
                }
            \end{center}  
            \item infinite rewriting-sequence: 
            \tikz[baseline=-0.5ex]{ 
                \node (x) at (0,0) {$\bullet$};  
                \node (y) at (1,0) {$\bullet$};
                \node (z) at (0.5,0.86) {$\bullet$};
                \draw[->,red] (x) -- node[midway,below] {$a$} (y) ;
                \draw[->,red] (y) -- node[midway,right] {$b$} (z) ;
                \draw[->] (z) -- node[midway,left] {$b$} (x) ;
            } 
            $\mathop{\Rightarrow}$ 
            \tikz[baseline=-0.5ex]{ 
                \node (x) at (0,0) {$\bullet$};  
                \node (y) at (1,0) {$\bullet$};
                \node (z) at (0.5,0.86) {$\bullet$};
                \draw[->] (x) -- node[midway,below] {$b$} (y) ;
                \draw[->,red] (y) -- node[midway,right] {$a$} (z) ;
                \draw[->,red] (z) -- node[midway,left] {$b$} (x) ;
            }
            $\mathop{\Rightarrow}$ 
            \tikz[baseline=-0.5ex]{ 
                \node (x) at (0,0) {$\bullet$};  
                \node (y) at (1,0) {$\bullet$};
                \node (z) at (0.5,0.86) {$\bullet$};
                \draw[->,red] (x) -- node[midway,below] {$b$} (y) ;
                \draw[->] (y) -- node[midway,right] {$b$} (z) ;
                \draw[->,red] (z) -- node[midway,left] {$a$} (x) ;
            }
            $\mathop{\Rightarrow}$ 
            \tikz[baseline=-0.5ex]{ 
                \node (x) at (0,0) {$\bullet$};  
                \node (y) at (1,0) {$\bullet$};
                \node (z) at (0.5,0.86) {$\bullet$};
                \draw[->,red] (x) -- node[midway,below] {$a$} (y) ;
                \draw[->,red] (y) -- node[midway,right] {$b$} (z) ;
                \draw[->] (z) -- node[midway,left] {$b$} (x) ;
            }
        with matches in red
        \end{itemize}
\end{example}

\begin{example}
  to do example

	\begin{tikzpicture}
		\draw (-1,0) node[left] {$r_1$:};
		\draw (1,0) node[above] {0};
		
		\node[draw,circle] (x_1) at (0,0) {$1$} ;
		\draw (0,0.3) node[above]{A};
		\node[draw,circle, minimum size \mathop{=} 1pt] (x_2) at (2,0) {$2$};
		\draw (2,0.3) node[above]{N};
		\draw[-] (x_1) -- (x_2);
		
		\draw (3,0) node[right] {$\longrightarrow$};
		\draw (6,0) node[above] {1};
		
		\node[draw,circle] (x3) at (5,0) {$1$} ;
		\draw (5,0.3) node[above]{A};
		\node[draw,circle, minimum size \mathop{=} 1pt] (x4) at (7,0) {$2$};
		\draw (7,0.3) node[above]{A'};
		\draw[-] (x3) -- (x4);
	\end{tikzpicture}

	\begin{tikzpicture}
		\draw (-1,0) node[left] {$r_2$:};
		\draw (1,0) node[above] {1};
		
		\node[draw,circle] (x_1) at (0,0) {$1$} ;
		\draw (0,0.3) node[above]{A'};
		\node[draw,circle, minimum size \mathop{=} 1pt] (x_2) at (2,0) {$2$};
		\draw (2,0.3) node[above]{N};
		\draw[-] (x_1) -- (x_2);
		
		\draw (3,0) node[right] {$\longrightarrow$};
		\draw (6,0) node[above] {1};
		
		\node[draw,circle] (x3) at (5,0) {$1$} ;
		\draw (5,0.3) node[above]{A'};
		\node[draw,circle, minimum size \mathop{=} 1pt] (x4) at (7,0) {$2$};
		\draw (7,0.3) node[above]{A'};
		\draw[-] (x3) -- (x4);
	\end{tikzpicture}
	\begin{itemize}
		\item Calcul d'un Arbre couvrant
		\item Graphe connexe avec étiquettes initiales: 0, N, A.\\
 \end{itemize}
 Séquence de règles appliquée: $r_1 \mathop{\rightarrow} r_2 \mathop{\rightarrow} r_1$ \\ \ \ 
 \resizebox{0.2\textwidth}{!}{
\begin{tikzpicture}
	\node[draw, circle, minimum size \mathop{=} 1 cm] (n1) at (0,0) {$1$};
	\node[draw, circle, minimum size \mathop{=} 1 cm] (n2) at (2,0) {$2$};
	\node[draw, circle, minimum size \mathop{=} 1 cm] (n3) at (2,2) {$3$};
	\node[draw, circle, minimum size \mathop{=} 1 cm] (n4) at (0,2) {$4$};
	
	\draw[-] (n1)--(n2)--(n3)--(n4);
	\draw[-] (n4)--(n1);
	
	\draw(0,-0.5) node[below] {\textcolor{red}{ A}};
	\draw(2,-0.5) node[below] {N};
	\draw(2,2.5) node[above] {N};
	\draw(0,2.5) node[above] {N};
	
	\draw(1,0) node[below] {0};
	\draw(2,1) node[right] {0};
	\draw(1,2) node[above] {0};
	\draw(0,1) node[left] {0};	
\end{tikzpicture}
 }
 \resizebox{0.2\textwidth}{!}{
\begin{tikzpicture}
	\node[draw, circle, minimum size \mathop{=} 1 cm] (n1) at (0,0) {$1$};
	\node[draw, circle, minimum size \mathop{=} 1 cm] (n2) at (2,0) {$2$};
	\node[draw, circle, minimum size \mathop{=} 1 cm] (n3) at (2,2) {$3$};
	\node[draw, circle, minimum size \mathop{=} 1 cm] (n4) at (0,2) {$4$};
	
	\draw[-] (n1)--(n2)--(n3)--(n4);
	\draw[-] (n4)--(n1);
	
	\draw(0,-0.5) node[below] {\textcolor{red}{ A}};
	\draw(2,-0.5) node[below] {\textcolor{red}{A'}};
	\draw(2,2.5) node[above] {N};
	\draw(0,2.5) node[above] {N};
	
	\draw(1,0) node[below] {\textcolor{red}{1}};
	\draw(2,1) node[right] {0};
	\draw(1,2) node[above] {0};
	\draw(0,1) node[left] {0};
\end{tikzpicture}
 }
 \resizebox{0.2\textwidth}{!}{
  \begin{tikzpicture}
	\node[draw, circle, minimum size \mathop{=} 1 cm] (n1) at (0,0) {$1$};
	\node[draw, circle, minimum size \mathop{=} 1 cm] (n2) at (2,0) {$2$};
	\node[draw, circle, minimum size \mathop{=} 1 cm] (n3) at (2,2) {$3$};
	\node[draw, circle, minimum size \mathop{=} 1 cm] (n4) at (0,2) {$4$};
	
	\draw[-] (n1)--(n2)--(n3)--(n4);
	\draw[-] (n4)--(n1);
	
	\draw(0,-0.5) node[below] {\textcolor{red}{A}};
	\draw(2,-0.5) node[below] {\textcolor{red}{A'}};
	\draw(2,2.5) node[above] {\textcolor{red}{A'}};
	\draw(0,2.5) node[above] {N};
	
	\draw(1,0) node[below] {\textcolor{red}{1}};
	\draw(2,1) node[right] {\textcolor{red}{1}};
	\draw(1,2) node[above] {0};
	\draw(0,1) node[left] {0};
\end{tikzpicture}
 }
 \resizebox{0.2\textwidth}{!}{
\begin{tikzpicture}
	\node[draw, circle, minimum size \mathop{=} 1 cm] (n1) at (0,0) {$1$};
	\node[draw, circle, minimum size \mathop{=} 1 cm] (n2) at (2,0) {$2$};
	\node[draw, circle, minimum size \mathop{=} 1 cm] (n3) at (2,2) {$3$};
	\node[draw, circle, minimum size \mathop{=} 1 cm] (n4) at (0,2) {$4$};
	
	\draw[-] (n1)--(n2)--(n3)--(n4);
	\draw[-] (n4)--(n1);
	
	\draw(0,-0.5) node[below] {\textcolor{red}{A}};
	\draw(2,-0.5) node[below] {\textcolor{red}{A'}};
	\draw(2,2.5) node[above] {\textcolor{red}{A'}};
	\draw(0,2.5) node[above] {\textcolor{red}{A'}};
	
	\draw(1,0) node[below] {\textcolor{red}{1}};
	\draw(2,1) node[right] {\textcolor{red}{1}};
	\draw(1,2) node[above] {0};
	\draw(0,1) node[left] {\textcolor{red}{1}};
\end{tikzpicture}
 }
\end{example}

\part{Termination of graph relabeling systems via morphisms to term rewriting relations}
\chapter{Introduction}

\chapter{Preliminaries}
    \label{sec:trs} 
    
\section{Terms}
Terms on which first order term rewriting systems are defined are terms of a first order language.
% \begin{definition}[Signature $\Sigma$~\cite{nipkow1998term}]
%   A \textbf{signature} \( \Sigma \) is a set of function symbols and for each \( f \mathop{\in} \Sigma \) a non-negative integer \( n \) called the arity of \( f \).
% \end{definition}
\begin{definition}[First order language~\text{\cite[Def 1.1.1]{marker2006model},\cite[Def 2.1.1]{terese2003term}}]
  A \textbf{language} \( \mathcal{L} \) is given by specifying the following data:
  \begin{itemize}
      \item  a set of function symbols \( \mathcal{F} \) and positive integers \( n_f \) for each \( f \mathop{\in} \mathcal{F} \);
      \item  a set of relation symbols \( \mathcal{R} \) and positive integers \( n_R \) for each \( R \mathop{\in} \mathcal{R} \);
      \item  a set of constant symbols \( \mathcal{C} \).
  \end{itemize}
\end{definition}
% \begin{definition}[$\Sigma$-Terms]
%   Let $\Sigma$ be a signature. The set $T(\Sigma)$ of $\Sigma$-terms is the smallest set such that 
%   \begin{itemize}
%     \item for all $n\geq 0$, all function symbols $f$ of arity n, and all $t_1,\hdots, t_n \mathop{\in} T(\Sigma, \mathcal{X})$, we have $f(t_1,...,t_n) \mathop{\in} T(\Sigma, \mathcal{X})$.
%   \end{itemize}
% \end{definition}

\begin{definition}[$(\Sigma,\mathcal{X})$-Terms~\text{\cite{nipkow1998term},\cite[Def 2.1.2]{terese2003term}}]
  Let $\Sigma$ be a first order language and $\mathcal{X}$ be a countably infinite set of variable symbols such that $\Sigma \mathop{\cap} \mathcal{X} \mathop{=} \emptyset$. The set $T(\Sigma,\mathcal{X})$ of $(\Sigma,\mathcal{X})$-terms is the smallest set such that 
  \begin{itemize}
  \item $\mathcal{X} \mathop{\subseteq} T(\Sigma,\mathcal{X})$, and
  \item for all $n\geq 0$, all function symbols $f$ of arity n, and all $t_1,\hdots, t_n \mathop{\in} T(\Sigma, \mathcal{X})$, we have $f(t_1,...,t_n) \mathop{\in} T(\Sigma, \mathcal{X})$.
  \end{itemize}
\end{definition}


% \section{Term Rewriting Systems}

\section{TRS}
A term rewriting system (TRS) is a rewriting system where the objects are $T(\Sigma,\mathcal{X})$-terms where $\Sigma$ is a set of function symbols and $\mathcal{X}$ is a set of variable symbols disjoint from $\Sigma$.
  
\begin{definition}[term rewriting rule~\text{\cite[Def.2.2.1]{terese2003term}}]
    Let $\Sigma$ be a first order language and $\mathcal{X}$ be a set of variables. A \textbf{term rewriting rule} is an ordered pair $(l,r) \mathop{\in} T(\Sigma, \mathcal{X})^2$, denoted $l \mathop{\to} r$, such that
    \begin{itemize}
      \item $l$ is not a variable, and
      \item every variable occurring in $r$ also occurs in $l$.
    \end{itemize}
  \end{definition}
  
  \begin{definition}[Position \cite{nipkow1998term} \cite{urbain2001approche}]
    Let $\Sigma$ be a first order language and $\mathcal{X}$ be a set of variable symbols disjoint from $\Sigma$. Let $s \mathop{\in} T(\Sigma, \mathcal{X})$.
    The set of positions of subterms of $s$ is a set $\operatorname{Pos}(s)$ defined by induction as follows:
    \begin{itemize}
      \item if $s \mathop{\in} \mathcal{X}$, then $\operatorname{Pos}(s) \mathop{=} \set{\epsilon}$
      \item if $s \mathop{=} f(t_1,\hdots,t_n)$, 
            then $\operatorname{Pos}(s) \mathop{=} \set{\epsilon} \mathop{\cup} \bigcup_{1 \leq k \leq n} \left \{kp\mid p \mathop{\in} \operatorname{Pos}(t_k) \right \}$
  \end{itemize}
  A position of a subterm of $s$ is a word on alphabet $\mathbb{N}\mathop{\setminus}\set{0}$. 
    The concatenation of words $p$ and $q$ will be denoted $pq$.
  \end{definition}

  \begin{definition}[Subterm \cite{nipkow1998term} \cite{urbain2001approche}]
    For $p \mathop{\in} \mathcal{Pos}(s)$, the \textbf{subterm of s at position p}, denoted by $s_{|p}$, is defined by induction on the length of p:
    \begin{itemize}
      \item $s_{|\epsilon} := s$, and
      \item $f(s_1,...,s_n)_{|iq} := \left(s_i\right)_{|q}$.
    \end{itemize} 
    
    Note that, for $p \mathop{=} iq$, $p \mathop{\in} \mathcal{Pos}$ implies that s is of the form $s \mathop{=} f(s_1,...,s_n)$ with $i \le n$
    
    For $p \mathop{\in} \mathcal{Pos}(s)$, we denote by $s[t]_p$ the term that is obtained from s by \textbf{replacing the subterm at position p by t}:
  \begin{itemize}
    \item $s[t]_\epsilon := t$
    \item $f(s_1,...,s_n)[t]_{iq} := f(s_1,...,s_i[t]_{q},...,s_n)$
  \end{itemize}
 
  \end{definition}
   
  \begin{definition}[Substitution]
    Let $\Sigma$ be a first order language and $\mathcal{X}$ be a set of variables.
    A \text{substitution} is a function $\sigma : \mathcal{X} \mathop{\rightarrow} T(\Sigma, \mathcal{X})$ such that $\sigma(x) \not \mathop{=} x$ for only finitely many variables. Any substitution can be extended to a mapping on $T(\Sigma,\mathcal{X})$ by letting $\sigma(f(s_1,...,s_n)) \overset{\operatorname{def}}{=}f(\sigma(s_1),...,\sigma(s_n))$.
  \end{definition}

   
  \begin{definition}[Term rewriting]
    \label{def:trs:rewriting}
    Let $\Sigma$ be a first order language and $\mathcal{X}$ be a set of variables. Let $\mathcal{R}$ be a set of TRS rewriting rules. Let $s \mathop{\in} T(\Sigma, \mathcal{X})$ and $\rho \mathop{=} (l \mathop{\to} r) \mathop{\in} \mathcal{R}$. Let $p$ be a match of $\rho$ in $s$.

      A \textbf{match} of $\rho$ in $s$ is a position $p \mathop{\in} \operatorname{Pos}(s)$ such that there exists a substitution $\sigma$ such that $s_{|p} \mathop{=} \sigma(l)$.

    The triple $(\rho, s, p, \sigma)$ is a witness of the rewriting step $s \mathop{\to} _\rho^p t$ where $t \mathop{=} s[\sigma(r)]_p$.
  \end{definition}
  
  \begin{definition}[TRS rewriting system]
    Let $\mathcal{R}$ be a set of TRS rules. 

    Let $M, \mathfrak{M}, W , \mathfrak{W}, \mathfrak{I}$ be the match, match mechanism, witnesses, witness function and interpretation functions introduced in Definition~\ref{def:trs:rewriting}
 
    The rewriting system $(T(\Sigma,\mathcal{X}), \mathcal{R}, M, \mathfrak{M}, W, \mathfrak{W}, \mathfrak{I})$ is called a \textbf{term rewriting system (TRS)}.
  \end{definition}

  % \begin{definition}[Term Rewriting System]
  %   Let $\Sigma$ be a signature and $\mathcal{X}$ be a set of variables.
  %   Let $\mathcal{R} \mathop{\subseteq} T(\Sigma,\mathcal{X})^2$.
  %   For all $s \mathop{\in} T(\Sigma, \mathcal{X})$ and for all $l \mathop{\to} r \mathop{\in} \mathcal{R}$ we define
  %   % \begin{flalign*}
  %   %   \operatorname{Acc}_{trs}(s,l \mathop{\to} r) \overset{\operatorname{def}}{=} 
  %   %     \left \{ 
  %   %         s[\sigma(r)]_p \mathop{\mid} p \mathop{\in} \operatorname{Pos}(s) \mathop{\land} \exists \sigma. s|_p \mathop{=} \sigma(l)
  %   %       \right \}
  %   % \end{flalign*}
  %   \begin{flalign*}
  %     \f{F}_{trs}(l \mathop{\to} r) \overset{\operatorname{def}}{=} 
  %       \left \{ 
  %           (s,t) \mathop{\mid} 
  %             \exists C \mathop{\in} T(\Sigma \mathop{\cup} \set{\square}, \mathcal{X}). \exists \sigma. 
  %            s \mathop{=} C[\sigma(l)] \mathop{\land} C[\sigma(r)] \mathop{=} t
  %         \right \}
  %   \end{flalign*}
  %   The rewriting system $(T(\Sigma,\mathcal{X}), \mathop{\to} _\mathcal{R})$, generated by $\mathcal{R}$ and $\f{F}_{trs}$ is called a \textbf{term rewriting system}.
  % \end{definition}



\section{ACTRS}

% Let $\Sigma$ be a signature and $\mathcal{X}$ be a set of variables.
% Let $\Sigma_{ac} \mathop{\subseteq} \Sigma_{c} \mathop{\subseteq} \Sigma$. We define $\mathcal{C}$ and $\mathcal{AC}$ as follows:

% $$\mathcal{C} \isdef \{f(x,y) \mathop{\to} f(y,x) \mathop{\mid} f \mathop{\in} \Sigma_{c} \}$$

% $$\mathcal{AC} \isdef 
%          \{f(f(x,y),z) \mathop{\to} f(x,f(y,z)) \mathop{\mid} f \mathop{\in} \Sigma_{ac} \}$$

% Let $\to_\mathcal{AC}$ be the rewriting relation defined by $\mathcal{AC} \mathop{\cup} \mathcal{C}$.
% \begin{definition}[ACTRS rule]
%    A \textbf{ACTRS rule} is a TRS rule.
% \end{definition}

% \begin{definition}[Match]
%     Let $[s] \mathop{\in} T(\Sigma, \mathcal{X})/\to_\mathcal{AC}$ and $\rho \mathop{=} l \mathop{\to} r$ be an ACTRS rule. 
%     An ACTRS match of $\rho$ in $[s]$ is an ordered pair $(s',p)$ where $s' \mathop{\in} [s]$ and $p \mathop{\in} \operatorname{Pos}(s')$ such that there exists a substitution $\sigma$ such that $s'_{|p} \mathop{=} \sigma(l)$.
% \end{definition}

% \begin{definition}[ACTRS Rewriting step]
%     Let $[s] \mathop{\in} T(\Sigma, \mathcal{X})/\to_\mathcal{AC}$ and $\rho \mathop{=} l \mathop{\to} r$ be an ACTRS rule. Let $(s',p)$ be a match of $\rho$ in $[s]$.

%     The structure $(s', p, \rho)$ is a \textbf{witness} for the \textbf{ACTRS rewriting step} $[s] \mathop{\to} _\mathcal{R} [t]$
%      using the rule $\rho$ and match $(s',p)$ where $t \mathop{\in} T(\Sigma, \mathcal{X})$ such that $s' \mathop{\to} _{trs}^p t$.
% \end{definition}
 
\begin{definition}
  \label{def:trs:ac}
  Let $\Sigma$ be a signature and $\mathcal{X}$ be a set of variables.
  Let $\Sigma_{ac} \mathop{\subseteq} \Sigma_{c} \mathop{\subseteq} \Sigma$. We define $\mathcal{C}$ and $\mathcal{AC}$ as follows:
  
  $$\mathcal{C} \isdef \{f(x,y) \mathop{\to} f(y,x) \mathop{\mid} f \mathop{\in} \Sigma_{c} \}$$
  
  $$\mathcal{AC} \isdef 
           \{f(f(x,y),z) \mathop{\to} f(x,f(y,z)) \mathop{\mid} f \mathop{\in} \Sigma_{ac} \}$$
  Let $ \mathop{\to} _\mathcal{AC}$ be the TRS rewriting relation induced by $\mathcal{AC} \mathop{\cup} \mathcal{C}$.
\end{definition}

\begin{definition}[Rewriting of terms with associative and commutative symbols~\cite{urbain2001approche}]
  \label{def:trs:actrs}
  Let $\Sigma$ be a first order language and $\mathcal{X}$ be a set of variables disjoint from $\Sigma$.
  Let $\mathcal{R}$ be a set of TRS rewriting rules, $\rho \mathop{=} l \mathop{\to} r \mathop{\in} \mathcal{R}$ and $s \mathop{\in} T(\Sigma, \mathcal{X})$.

  An \textbf{ACTRS match} of $\rho$ in $s$ is $p \mathop{\in} \operatorname{Pos}(s)$ such that there exists a substitution $\sigma$ such that $s_{|p} \rightarrow_{\opn{AC}}^* \sigma(l)$. 

  % The \textbf{ACTRS match mechanism} is the function $\mathcal{M}$ that associates to each pair $(\rho,s)$, where $\rho$ is an ACTRS rewriting rule and $s$ a term, the set of ACTRS matches of $\rho$ in $s$.

  The quadruple $(\rho, s, p, \sigma)$ is a witness of the \textbf{ACTRS rewriting step} $(s, t)$ where $t \mathop{=} s[\sigma(r)]_p$

  % The \textbf{ACTRS witness function} is the function $\mathfrak{W}$ that associates to each ACTRS match $(\rho, s, p, \sigma)$ the rewriting step $(s, t)$ where $t \mathop{=} s[\sigma(r)]_p$.
\end{definition}

% \begin{definition}[Term Rewriting modulo AC framework]
%   Let $\mathcal{R} \mathop{\subseteq} T(\Sigma,\mathcal{X})^2$ be a set of term rewriting rules.
   
%   For all $\rho \mathop{\in} \mathcal{R}$, we define the the rewriting framework $\f{F}_{actrs}(\rho)$ as follows:
%   \begin{flalign*}
%     \f{F}_{actrs}(\rho) \isdef 
%       \left \{ (s,t) \mathop{\in} T(\Sigma,\mathcal{X})^2 \mathop{\mid} 
%           \exists s', t'. 
%           s \mathop{\to} _\mathcal{AC}^* s' \mathop{\to} _\rho t' \mathop{\to} _\mathcal{AC}^* t
%         \right \}
%   \end{flalign*}
%   An element $(s,t) \mathop{\in} \f{F}_{actrs}(\rho)$ is called a \textbf{witness} for the rewriting step $[s] \mathop{\to} _\mathcal{R} [t]$ using the rule $\rho$ where $[s],[t] \mathop{\in} T(\Sigma,\mathcal{X})/\to_\mathcal{AC}$. 

%   We denote by $\to_{\mathcal{R} /\mathcal{AC}}$ the rewriting relation induced by $\mathcal{R}$ in $\f{F}_{actrs}$.
% \end{definition}
 
\begin{definition}[ACTRS rewriting system]
  Let $\mathcal{R}$ be a set of TRS rewriting rules. 
  Let $M, \mathfrak{M}, W, \mathfrak{W}, \mathfrak{I}$ be the class of ACTRS matches, the match mechanism, the class of ACTRS witnesses, the witness mechanism and the interpretation function, respectively, defined in Definition~\ref{def:trs:actrs}.

  The rewriting system $(T(\Sigma,\mathcal{X}), \mathcal{R}, M, \mathfrak{M}, W, \mathfrak{W}, \mathfrak{I})$ is called an \textbf{ACTRS rewriting system}.
\end{definition}



% \subsection{ACHTRS innermost}
% 
  \subsection{Innermost HTRS}
  Hierarchical TRS, proposed by Xavier Urbain \cite{urbain2001approche}, extends the conventional term rewriting framework to better manage complex, structured data by introducing hierarchical layers of rules. This approach allows for more modular and organized rewriting processes, which are particularly useful in fields requiring sophisticated term manipulation and structured data handling. The system also includes advanced techniques to ensure that the rewriting process terminates, which is a critical aspect of its theoretical foundation.
  
  
  \begin{definition}[Term rewriting framework with innermost strategy]
    Let $\mathcal{R}$ be a set of TRS rules. 
  
    We define the rewriting framework with innermost strategy, denoted $\mathfrak{trsi}$ as follows: for all $\rho \mathop{\in} \mathcal{R}$, $\mathfrak{trsi}(\rho)$ is the collection of all TRS witnesses of rewriting steps $(\rho, s, p, \sigma)$ such that there is no suffix position $p'$ of $p$ such that $s \mathop{\to} _\rho^{p'} t'$ for some $t' \mathop{\in} T(\Sigma,\mathcal{X})$.
  \end{definition}
  
  \begin{definition}[Term rewriting system with innermost strategy]
    Let $\mathcal{R}$ be a set of TRS rules. 
  
    The rewriting system $(T(\Sigma,\mathcal{X}), \mathcal{R}, M, \mathfrak{M}, W, \mathfrak{W}, \mathfrak{I}, \mathfrak{trsi})$ is called a \textbf{term rewriting system with innermost strategy (TRSI)}.
  \end{definition}
  % \begin{definition}[Innermost TRS]
  %     % $s \underset{i}{\rightarrow} t$ if $s \mathop{=} C[\sigma(l)] \rightarrow_{l\rightarrow r}C[\sigma(r)] \mathop{=} t$ for some $l \mathop{\rightarrow} r \mathop{\in} R$, some context $C[.]$ and some substitution $\sigma$ such that no proper subterm of $\sigma(l)$  is reducible. 
  %     Let $\Sigma$ be a signature and $\mathcal{X}$ be a set of variables.
  %     Let $\mathcal{R} \mathop{\subseteq} T(\Sigma,\mathcal{X})^2$ be a set of term rewriting rules.
  %     For all $l \mathop{\to} r \mathop{\in} \mathcal{R}$, we define
  %     \begin{flalign*}
  %       \mathfrak{F}_{itrs}(l \mathop{\to} r) \overset{def}{=} 
  %         \left \{ 
  %             (s, t) \mid
  %             \exists C 
  %             % \mathop{\in} T(\Sigma \mathop{\cup} \set{\square},\mathcal{X})
  %             .
  %              \exists \sigma.
  %             s \mathop{=} C[\sigma(l)] \mathop{\land} C[\sigma(r)] \mathop{=} t \mathop{\land} \not \exists l'. \sigma(l) \mathop{\to} _\mathcal{R} l'
  %               % s \mathop{=} C[\sigma(l)] \rightarrow_{l\rightarrow r}C[\sigma(r)]
  %             % p \mathop{\in} Pos(s) \mathop{\land} \exists \sigma. s|_p \mathop{=} \sigma(l) \mathop{\land} \not \exists p' \mathop{\in} Pos(s). p' \mathop{>} p \mathop{\land} \exists (l' \mathop{\to} r' \mathop{\in} \mathcal{R}). \exists \sigma'. s_p' \mathop{=} \sigma'(l')
  %           \right \}
  %     \end{flalign*}
  %     The rewriting system generated by $\mathcal{R}$ and $\mathfrak{F}_{itrs}$ is called a term rewriting system and will be denoted by $(T(\Sigma,\mathcal{X}), \itrs_\mathcal{R}^*)$. 
  %   \end{definition}
  
    \begin{definition}[\cite{urbain2001approche}]
      A TRS rewriting system $\mathcal{R}$ on $T(\Sigma, \mathcal{X})$ is said to be \textbf{hierarchical} if there are $(\Sigma_i, \mathcal{R})_{1 \leq i \leq n}$ equipped with a strict partial order $\prec$, such that:
      \begin{itemize}
        \item $\Sigma_i$, $1 \leq i \leq n$, partition $\Sigma$,
        \item $\mathcal{R}_i$, $1 \leq i \leq n$, partition $\mathcal{R}$,
        \item for all $1 \leq i \leq n$ and $\rho \mathop{\in} \mathcal{R}_i$, we have $\Lambda(\rho) \mathop{\in} \Sigma_i$ and $\rho$ is a rule on $T(\Sigma', \mathcal{X})$ where 
        $\Sigma' \mathop{=} \bigcup \{ \Sigma_j | (\Sigma_j, \mathcal{R}) \mathop{\preceq} (\Sigma_i, \mathcal{R}) \}$.
      \end{itemize}
    \end{definition}
  
    \begin{definition}[Hierachical term rewriting systems with innermost strategy]
       A \textbf{hierarchical term rewriting systems with innermost strategy} is a TRSI which is hierarchical.
    \end{definition}
  
    % \begin{definition}[Hierarchical Innermost TRS \cite{urbain2001approche}]
    %   Let $\Sigma$ be a signature and $\mathcal{X}$ be a set of variables.
    %   Let $\mathcal{R} \mathop{\subseteq} T(\Sigma,\mathcal{X})^2$ be a set of term rewriting rules.
    %   Let $\mathfrak{F}$ be a rewriting framework.
    %   We say that the term rewriting system generated by $\mathcal{R}$ and $\mathfrak{F}$ is hierarchical if there exist $n \mathop{\in} \mathbb{N}$, $\Sigma_1 \subsetneq \Sigma_2 \subsetneq \ldots \subsetneq \Sigma_n \mathop{=} \Sigma$, and $\mathcal{R}_1 \subsetneq \mathcal{R}_2 \subsetneq \ldots \subsetneq \mathcal{R}_n \mathop{=} \mathcal{R}$ such that:
    %   \begin{itemize}
    %     \item Rules in $\mathcal{R}_k$ involve only symbols in $\Sigma_k$, for $1 \leq k \leq n$,
    %     \item For $2 \leq k \leq n$ and for all rules $l \mathop{\to} r \mathop{\in} \mathcal{R}_k$, the left-hand side $l$ contains at least one symbol from $\Sigma_k \mathop{\setminus} \Sigma_{k-1}$.
    %   \end{itemize}
    % \end{definition}
  
    
    
    % %ac
    % \begin{definition}[Innermost AC Term Class Rewriting System]
    %   $[s] \underset{i}{\rightarrow} [t]$ if $s \mathop{=} C[\sigma(l)]$ and $C[\sigma(r)] \mathop{=} t$ for some $[l] \mathop{\rightarrow} [r] \mathop{\in} R$, some context $C[\cdot]$ and some substitution $\sigma$ such that no proper subterm of $\sigma(l)$  is reducible. 
    
    
    
    %   Let $\Sigma$ be a signature and $\mathcal{X}$ be a set of variables.
    %   Let $\mathcal{R} \mathop{\subseteq} (T(\Sigma,\mathcal{X})/E)^2$ be a set of term rewriting rules.
    %   for all $[s] \mathop{\in} T(\Sigma, \mathcal{X})/E$ and for all $[l] \mathop{\to} [r] \mathop{\in} \mathcal{R}$ we define
    %   \begin{flalign*}
    %     \operatorname{Acc}_{trs}([s],[l] \mathop{\to} [r]) \overset{def}{=} 
    %       \left \{ 
    %           s[\sigma(r)]_p \mathop{\mid} p \mathop{\in} Pos(s) \mathop{\land} \exists \sigma. s|_p \mathop{=} \sigma(l) \mathop{\land} \not \exists p' \mathop{\in} Pos(s). p' \mathop{>} p \mathop{\land} \exists (l' \mathop{\to} r' \mathop{\in} \mathcal{R}). \exists \sigma'. s_p' \mathop{=} \sigma'(l')
    %         \right \}
    %   \end{flalign*}
      
    %   The rewriting system $(T(\Sigma,\mathcal{X}), \mathop{\to} _\mathcal{R})$, generated by $\mathcal{R}$ and $\operatorname{ACC}_{trs}$ is called a term rewriting system.
    % \end{definition}
    
    %ac
    % \begin{definition}[Innermost Term Class Rewriting System]
    %   Let $\Sigma$ be a signature and $\mathcal{X}$ be a set of variables.
    %   Let $\mathcal{R} \mathop{\subseteq} (T(\Sigma,\mathcal{X})/E)^2$ be a set of term rewriting rules.
    %   for all $s \mathop{\in} T(\Sigma, \mathcal{X})/E$ and for all $l \mathop{\to} r \mathop{\in} \mathcal{R}$ we define
    %   \begin{flalign*}
    %     \operatorname{Acc}_{trs/E}(s,l \mathop{\to} r) \overset{def}{=} 
    %       \left \{ 
    %           [t] \mathop{\mid} \exists s' \mathop{\in} s. s' \mathop{\to} _i t 
    %             % p \mathop{\in} Pos(s) \mathop{\land} \exists \sigma. s|_p \mathop{=} \sigma(l) \mathop{\land} \not \exists p' \mathop{\in} Pos(s). p' \mathop{>} p \mathop{\land} \exists (l' \mathop{\to} r' \mathop{\in} \mathcal{R}). \exists \sigma'. s_p' \mathop{=} \sigma'(l')
    %         \right \}
    %   \end{flalign*}
      
    %   The rewriting system $(T(\Sigma,\mathcal{X}), \mathop{\to} _\mathcal{R})$, generated by $\mathcal{R}$ and $\operatorname{ACC}_{trs/E}$ is called a term class rewriting system.
    % \end{definition}
    
    %ac
    % \begin{definition}[AC HTRS \cite{urbain2001approche}]
    %   Let $\Sigma$ be a signature and $\mathcal{X}$ be a set of variables.
    %   Let $\mathcal{R} \mathop{\subseteq} (T(\Sigma,\mathcal{X})/E)^2$ be a set of term rewriting rules. Let $\operatorname{ACC}$ be an accessibility function.
    %   Let $(T(\Sigma,\mathcal{X}), \mathop{\to} _\mathcal{R})$ be the term rewriting system. generated by $\mathcal{R}$ and $\operatorname{ACC}_{trs/E}$.
    %   If there are $n \mathop{\in} \mathbb{N}$, $\Sigma_1, \Sigma_2, \hdots, \Sigma_n $, and $\mathcal{R}_1, \mathcal{R}_2, \hdots, \mathcal{R}_n$ such that 
    %   \begin{itemize}
    %     \item $\Sigma_1 \subsetneq \Sigma_2 \subsetneq \hdots \subsetneq \Sigma_n \mathop{=} \Sigma$
    %     \item $\mathcal{R}_1 \subsetneq \mathcal{R}_2 \subsetneq \hdots \subsetneq \mathcal{R}_n \mathop{=} \mathcal{R}$
    %     \item rules in $\mathcal{R}_{k}$ have only symbols in $\Sigma_k$ for $1 \leq k \leq n$
    %     \item for $2 \leq k \leq n$ and for all rules $l \mathop{\to} r \mathop{\in} \mathcal{R}_{k}$, for all $l' \mathop{\in} l$, we have $\Lambda(l') \mathop{\in} (\Sigma_k \mathop{\setminus} \Sigma_{i-1})$
    %   \end{itemize}
    % \end{definition}
    
    % \begin{Plan}
    %   \begin{itemize}
    %     \item a trs is a rewriting system on algebraic terms, with a distinguished 
    %       subset of elements $\rs{R}$, called term rewriting rule, such that the rewriting system can be generated from $\rs{R}$.
    %     \item def algebraic terms
    %     \item def rewriting rules
    %     \item def stable under ontexts
    %     \item def stable under instantiation
    %     \item Acc
    %     \item def trs : smallest binary relation ... 
    %     \item def ntrs
    %     \item def htrs
    %     \item def inner most strategy
    %     \item reduction relation : morphism to ....
    %   \end{itemize}
    % \end{Plan} 

    % \section{Termination of Term Rewriting Systems}
    % todo
    %     \subsection{Reduction Orders}
    %     \begin{definition}[Stable under substitution, Stable under context]
    Let $\to$ be a binary relation on $T(\Sigma,\mathcal{X})$. We say that $\to$ is closed under context if
      $$s \to t \implies \forall C. h(C[s]) \to h(C[t])$$
    and we say that $\to$ is stable under substitution if
      $$s \to t \implies \forall \sigma. h(\sigma(s)) \to h(\sigma(t))$$
  \end{definition}
  
  \begin{definition}[Term Reduction Order \cite{nipkow1998,urbain2001approche}]
    % If $(T(\Sigma,\mathcal{X}), \to)$ is a term rewriting system stable under context and substitution
    A well-founded binary relation $\to$ on $T(\Sigma,\mathcal{X})$ stable under context and substitution is called a term reduction order. 
  \end{definition} 
  
  \begin{proposition}[to do : ref]
    A term rewriting system $(T(\Sigma,\mathcal{X}), \to)$ is terminating if $\to$ is a term reduction order. 
  \end{proposition} 
    %     \subsection{Termination Methods of Term Rewriting Systems}
    %     \begin{theorem}[Termination of Term Rewriting Systems]
    Let $(T(\Sigma,\mathcal{X}), \to_R)$ be a term rewriting system. If there exists a term rewriting system $(T(\Sigma,\mathcal{X}),>)$ where $>$ is a reduction order and a morphism $h:(T(\Sigma,\mathcal{X}),R) \to (T(\Sigma,\mathcal{X}),>)$, then $(T(\Sigma,\mathcal{X}), \to_R)$ is terminating.
  \end{theorem}
  \begin{proof}
    The rewriting system $(T(\Sigma,\mathcal{X}),>)$ is terminating and there exists a morphism $h':(T(\Sigma,\mathcal{X}),\to_R) \to (T(\Sigma,\mathcal{X}),>)$.
  \end{proof}
  
  \begin{theorem}[Relative Termination \cite{geser2016relative} ]
    Let $(T(\Sigma,\mathcal{X}), \to_R)$ and $(T(\Sigma,\mathcal{X}), \to_R')$ be two term rewriting systems. If there exists term rewriting system $(T(\Sigma,\mathcal{X}), >)$ where $>$ is a reduction order, a morphism $h: (T(\Sigma,\mathcal{X}),R) \to (T(\Sigma,\mathcal{X}), >)$ and a morphism $g: (T(\Sigma,\mathcal{X}),R') \to (T(\Sigma,\mathcal{X}), \geq)$ then 
     $(T(\Sigma,\mathcal{X}), \to_R)$ is terminating relative to $(T(\Sigma,\mathcal{X}), \to_R)$.
  \end{theorem}
  \begin{proof}
    The rewriting system $(T(\Sigma,\mathcal{X}),>)$ is terminating and there exists a morphism $h':(T(\Sigma,\mathcal{X}),\to_R) \to (T(\Sigma,\mathcal{X}),>)$ and a morphism $g':(T(\Sigma,\mathcal{X}),\to_R') \to (T(\Sigma,\mathcal{X}),\geq)$
  \end{proof}
    %     \subsection{Interpretation Methods}
    %     
Interpretation methods are a fundamental approach in proving the termination of term rewriting systems (TRS). These methods involve mapping terms into well-founded mathematical structures, ensuring that each rewriting step corresponds to a reduction in size according to a specific order. The significance and effectiveness of interpretation methods in termination analysis have been extensively discussed in the literature, including in the comprehensive survey \cite{dersh1987termination}.
In this section, we explore key concepts such as 
$\Sigma$-structures, term interpretations, and homomorphisms. These concepts are crucial for demonstrating that rewriting systems terminate by ensuring that every rewrite leads to a simplification in a well-defined order.

\begin{definition}[$\Sigma$-structure \cite{cori2000mathematical}]
  A \textit{$\Sigma$-structure}, is any structure $\mathcal{M}$ consisting of:
  \begin{itemize}
    \item a non-empty set $M$, called the \textit{underlying set} (or \textit{base set}) of the realization $\mathcal{M}$;
    \item for each constant symbol $c$ of $\Sigma$, an element $c^{\mathcal{M}}$ of $M$, called the \textit{interpretation} of the symbol $c$ in the realization $\mathcal{M}$;
    \item for each natural number $k \mathop{\geq} 1$, and for each function symbol $f$ of arity $k$ in $\Sigma$, a function $f^{\mathcal{M}}$ from $M^k$ to $M$, called the \textit{interpretation} of the symbol $f$ in the realization $\mathcal{M}$;
    % \item for each natural number $k \mathop{\geq} 1$, and for each relation symbol $R$ of arity $k$ in $\Sigma$, a subset $R^{\mathcal{M}}$ of $M^k$ (that is, a $k$-ary relation on the set $M$), called the \textit{interpretation} of the symbol $R$ in the realization $\mathcal{M}$.
  \end{itemize}
\end{definition}

\begin{example}[Term Algebra]
  Give a signature $\Sigma$, the term algebra $(T(\Sigma, \mathcal{X}), \Sigma)$ is a model of $\Sigma$, or $\Sigma$-algebra, where every $f \mathop{\in} \Sigma$ is interpreted by itself.
\end{example}

\begin{definition}[$\Sigma$-Structure homomorphism \cite{cori2000mathematical}]
  The function $\varphi$ is a \textit{homomorphism of} $\Sigma$-\textit{structures} from $\mathcal{M}$ to $\mathcal{N}$ if and only if the following conditions are satisfied:
  \begin{itemize}
      \item For every constant symbol $c$ of $\Sigma$,
      \[
      \varphi(c^{\mathcal{M}}) \mathop{=} c^{\mathcal{N}};
      \]
      \item For every natural number $n \mathop{\geq} 1$, for every $n$-ary function symbol $f$ of $\Sigma$, and for all elements $a_1, a_2, \ldots, a_n$ belonging to $M$,
      \[
      \varphi(f^{\mathcal{M}}(a_1, a_2, \ldots, a_n)) \mathop{=} f^{\mathcal{N}}(\varphi(a_1), \varphi(a_2), \ldots, \varphi(a_n));
      \]
      \item For every natural number $k \mathop{\geq} 1$, for every $k$-ary relation symbol $R$ of $\Sigma$, and for all elements $a_1, a_2, \ldots, a_k$ belonging to $M$,
      \[
      \text{if } (a_1, a_2, \ldots, a_k) \mathop{\in} R^{\mathcal{M}}, \text{ then } (\varphi(a_1), \varphi(a_2), \ldots, \varphi(a_k)) \mathop{\in} R^{\mathcal{N}}.
      \]
  \end{itemize}
  \end{definition}

  \begin{definition}[Term Interpretation \cite{cori2000mathematical}]
    Given a natural number $n$, $n$ variables $X_0, X_1, \ldots, X_{n-1}$ that are distinct from each other, a term $t \mathop{=} t[X_0, X_1, \ldots, X_{n-1}]$ of the language $\Sigma$, an $\Sigma$-structure $\mathcal{M} \mathop{=} \langle M, \ldots \rangle$, and $n$ elements $a_0, a_1, \ldots, a_{n-1}$ of $M$, we call the \textit{interpretation of the term} $t$ in the $\Sigma$-structure $\mathcal{M}$, when the variables $X_0, \ldots, X_{n-1}$ are respectively interpreted by the elements $a_0, a_1, \ldots, a_{n-1}$, the element of $M$ denoted by:
    \[
     t^{\mathcal{M}}[X_0 \mathop{\mapsto} a_0, X_1 \mathop{\mapsto} a_1, \ldots, X_{n-1} \mathop{\mapsto} a_{n-1}]
    \]
    defined as follows, by induction on $t$:
    \begin{itemize}
        \item If $t \mathop{=} X_j$ (where $0 \leq j \leq n-1$),
        \[
        t^{\mathcal{M}}[X_0 \mathop{\mapsto} a_0, X_1 \mathop{\mapsto} a_1, \ldots, X_{n-1} \mathop{\mapsto} a_{n-1}] \mathop{=} a_j;
        \]
        \item If $t \mathop{=} c$ (a constant symbol of $\Sigma$),
        \[
         t^{\mathcal{M}}[X_0 \mathop{\mapsto} a_0, X_1 \mathop{\mapsto} a_1, \ldots, X_{n-1} \mathop{\mapsto} a_{n-1}] \mathop{=}  c^{\mathcal{M}};
        \]
        \item If $t \mathop{=} f t_1 t_2 \ldots t_k$ (where $k \mathop{\in} \mathbb{N}^*$, and $f$ is a $k$-ary function symbol of $\Sigma$, and $t_1, t_2, \ldots, t_k$ are terms of $\Sigma$),
        \[
          t^{\mathcal{M}}[X_0 \mathop{\mapsto} a_0, X_1 \mathop{\mapsto} a_1, \ldots, X_{n-1} \mathop{\mapsto} a_{n-1}] \mathop{=} f^{\mathcal{M}}(t_1^{\mathcal{M}}[X_0 \mathop{\mapsto} a_0, \ldots, X_{n-1} \mathop{\mapsto} a_{n-1}], \ldots, t_k^{\mathcal{M}}[X_0 \mathop{\mapsto} a_0, \ldots, X_{n-1} \mathop{\mapsto} a_{n-1}]).
        \]
    \end{itemize}
    In practice, we will denote more simply:
    \[
    t^{\mathcal{M}}[a_0, a_1, \ldots, a_{n-1}].
    \]
    \end{definition}
    

    \begin{definition}[Monotone Function \cite{nipkow1998term}]
        Let $>$ be a strict order on a set $A$. A function $F : A^n \mathop{\rightarrow} A$ is said to be \textit{monotone} with respect to $>$ if and only if
        \[
        a \mathop{>} b \text{ implies } F(a_1, \ldots, a_{i-1}, a, a_{i+1}, \ldots, a_n) \mathop{>} F(a_1, \ldots, a_{i-1}, b, a_{i+1}, \ldots, a_n)
        \]
        for all $i$ where $1 \leq i \leq n$, and for all elements $a, b, a_1, \ldots, a_{i-1}, a_{i+1}, \ldots, a_n$ in $A$.
    \end{definition}

    \begin{definition}[Interpretation Method \cite{dersh1987termination}]
        Let $\Sigma$ be the rewriting system. 
        Let $\mathcal{A}$ be a $\Sigma$-structure such that for all $1 \leq i \leq n$, the function $f^\mathcal{M}$ is monotone w.r.t. $>_A$. Let $>_A$ be a well founed order on the domain $A$.

        Let $\Pi$ be the set of interpretation of $T(\Sigma,\mathcal{X})$ in $A$.
        Let $B\isdef \{(\pi(t))_{\pi\in \Pi} \mid t \mathop{\in} T(\Sigma,\mathcal{X})\}$, the set of generalized series. 

        Let $>_B$ be the binary relation defined by 
        $$(\pi(s))_{\pi\in \Pi} \mathop{>} (\pi(t))_{\pi\in \Pi} \text{ iff } \forall \pi \mathop{\in} \Pi. \pi(s) >_A \pi(t)$$

        % Let $>$ be the binary relation defined by 
        % $$s \mathop{>} t \text{ iff } \forall \pi:T(\Sigma,\mathcal{X}) \mathop{\to} A. \pi(s) >_A \pi(t)$$

        We have $(B, >_B)$ is a terminating rewriting system. Thus, if $h:(T(\Sigma,\mathcal{X}),R) \mathop{\to} (B,>_B)$ defined by $h(t) \isdef (\pi(t))_{\pi\in \Pi}$ is a homomorphism, then $(T(\Sigma,\mathcal{X}), \mathop{\to} _R)$ terminates.
    \end{definition}
    %     \subsection{Simplification Order}
    %     Suppose $|V| \in \mathbb{N}$.
    
\begin{definition}[Simplification Partial Order \cite{nipkow1998term} \cite{urbain2001approche}]
A strict partial order $>$ on $T(\Sigma,V)$ is called a simplification partial order if and only if:
\begin{itemize}
    \item it is stable under context and substitution, and
    \item it satisfies the following subterm property:
    \begin{itemize}
        \item For all terms $t \in T(\Sigma,V)$ and all positions $p \in \mathcal{Pos}(t) \setminus \{\epsilon\}$, we have $t > t|_p$.
    \end{itemize}
\end{itemize}  
\end{definition}

\begin{proposition}[\cite{nipkow1998term}]
For any strict order $>$ on $\Sigma$, the induced lexicographic path order $>_{lpo}$ is a simplification order on $T(\Sigma,V)$.
\end{proposition}
    
\begin{proposition}[\cite{nipkow1998term}]
Let $>$ be a strict order on $\Sigma$, and let $w$ be an admissible weight function for $>$. The Knuth-Bendix order (KBO) induced by $>$ and $w$ is a simplification order.
\end{proposition}
    
\begin{theorem}[\cite{nipkow1998term}]
Let $\Sigma$ be a finite signature. Every simplification order $>$ on $T(\Sigma,V)$ is a reduction order.
\end{theorem}

\begin{corollary}
  Let $\Sigma$ be a finite signature. Let $>$ be a simplification order on $T(\Sigma,V)$. The term rewriting system $(T(\Sigma,\mathcal{X}),>)$ terminates. Let $\mathcal{R}$ be a term rewriting system. If $id : (T(\Sigma,\mathcal{X}),R) \to  (T(\Sigma,\mathcal{X}),>)$ is a homomorphism, then $\mathcal{R}$ terminates.
\end{corollary}
    %     \subsection{Lexicographic Path Order}
    %         The Lexicographic Path Order (LPO) \cite{dershowitz1982orderings} is a simplification order . It is widely used in the analysis of term rewriting systems (TRS) to prove termination by comparing terms lexicographically according to their structure and the strict order of function symbols. LPO is particularly effective in ensuring that more complex terms are considered greater than their subterms, which is essential for termination proofs.

In this section, we provide the formal definition of the Lexicographic Path Order. Using a finite signature and a strict order on that signature, we outline the specific conditions under which one term is considered greater than another. LPO plays a crucial role in the systematic comparison of terms within the framework of TRS termination analysis.
    
    \begin{definition}[Lexicographic Path Order \cite{nipkow1998term}]
    Let $\Sigma$ be a finite signature and $>$ be a strict order on $\Sigma$. The lexicographic path order $>_{lpo}$ on $T(\Sigma,V)$ induced by $>$ is defined as follows: \( s >_{lpo} t \) if and only if
    \begin{itemize}
        \item \textbf{(LPO1)} $t \mathop{\in} \mathcal{Var}(s)$ and $s \mathop{\neq} t$, or
        \item \textbf{(LPO2)} $s \mathop{=} f(s_1, \ldots, s_m)$, $t \mathop{=} g(t_1, \ldots, t_n)$, and
        \begin{itemize}
            \item \textbf{(LPO2a)} $\exists i.\ s_i \ge_{lpo} t$, or
            \item \textbf{(LPO2b)}
            \begin{itemize}
                \item $f \mathop{>} g$, and
                \item $\forall j.\ s >_{lpo} t_j$, or
            \end{itemize}
            \item \textbf{(LPO2c)}
            \begin{itemize}
                \item $f \mathop{=} g$, and
                \item $\forall j.\ s >_{lpo} t_j$, and
                \item $\exists i.\ s_1 \mathop{=} t_1, \ldots, s_{i-1} \mathop{=} t_{i-1} \mathop{\land} s_i >_{lpo} t_i$
            \end{itemize}
        \end{itemize}
    \end{itemize}
    \end{definition}
    
\begin{theorem}[\cite{nipkow1998term} Thoeorem 5.4.14]
    For any strict order $>$ on $\Sigma$, the induced lexicographic path order $>_{\text{lpo}}$ is a simplification order on $T(\Sigma, \mathcal{X})$.
\end{theorem}

    %     \subsection{Knuth-Bendix Order}
    %         The Knuth-Bendix Order (KBO) is a fundamental concept in the termination analysis of term rewriting systems (TRS), introduced by Donald Knuth and Peter Bendix in their seminal 1983 work Simple Word Problems in Universal Algebras. KBO is a simplification order that allows for the comparison of terms based on a strict ordering of function symbols, coupled with a weight function that assigns values to both function symbols and variables. This order is particularly powerful in ensuring that larger terms dominate their subterms, thereby facilitating the proof of termination for rewriting systems.

    In this section, we detail the formal definition of the Knuth-Bendix Order as presented in \cite{knuth1983simple}. This definition relies on a finite signature, a strict order on the signature, and a weight function mapping symbols and variables to positive real numbers. We then outline the specific conditions under which one term is deemed greater than another, making KBO a valuable tool in the study of term rewriting systems.
  
  \begin{definition}[\cite{nipkow1998term}]
    
    Let $\Sigma$ be a \textit{finite} signature. A Knuth-Bendix order on $T(\Sigma, \mathcal{X})$ is determined by a strict order $>$ on $\Sigma$ and a \textbf{weight function} $w : \Sigma \cup \mathcal{X} \rightarrow \mathbb{R}^+_0$, where $\mathbb{R}^+_0$ denotes the set of non-negative real numbers. We call such a weight function $w$ \textbf{admissible} for $>$ iff it satisfies the following properties:

\begin{enumerate}
    \item There exists $w_0 \in \mathbb{R}^+_0 - \{0\}$ such that $w(x) = w_0$ for all variables $x \in \mathcal{X}$ and $w(c) \geq w_0$ for all constants $c \in \Sigma^{(0)}$.
    \item If $f \in \Sigma^{(1)}$ is a unary function symbol of weight $w(f) = 0$, then $f$ is the greatest element in $\Sigma$, i.e., $f \geq g$ for all $g \in \Sigma$.
\end{enumerate}
\end{definition}

    \begin{definition}[Knuth-Bendix Order \cite{nipkow1998term}]
    Let $\Sigma$ be a finite signature, $>$ be a strict order on $\Sigma$, and $w: \Sigma \cup \mathcal{X} \rightarrow \mathbb{R}^+$ be a weight function. The \textbf{Knuth-Bendix order $>_{kbo}$} on $T(\Sigma,\mathcal{X})$ induced by $>$ and $w$ is defined as follows: for $s, t \in T(\Sigma,\mathcal{X})$, we have \( s >_{kbo} t \) if and only if $\forall x \in \mathcal{X}.\ |s|_x \ge |t|_x$ and
    \begin{itemize}
        \item \textbf{(KBO1)} $w(s) > w(t)$, or
        \item \textbf{(KBO2)} $w(s) = w(t)$, and
        \begin{itemize}
            \item \textbf{(KBO2a)} $\exists f \in \Sigma^{(1)}.\ \exists x \in \mathcal{X}.\ \exists n \in \mathbb{N}^*.\ s = f^n(x) \land t = x$, or
            \item \textbf{(KBO2b)} $\exists f, g \in \Sigma.\ f > g \land s = f(s_1, \ldots, s_m) \land t = g(t_1, \ldots, t_n)$, or
            \item \textbf{(KBO2c)} $\exists f \in \Sigma.\ \exists 1 \le i \le m.\ s = f(s_1, \ldots, s_m) \land t = f(t_1, \ldots, t_m) \land s_1 = t_1 \land \ldots \land s_{i-1} = t_{i-1} \land s_i >_{kbo} t_i$
        \end{itemize}
    \end{itemize}
    \end{definition}
    
\begin{theorem}[\cite{nipkow1998term}]
Let $>$ be a strict order on $\Sigma$, and $w : \Sigma \cup \mathcal{X} \rightarrow \mathbb{R}^+_0$ be a weight function that is admissible for $>$. Then the Knuth-Bendix order $>_{\text{kbo}}$ on $T(\Sigma, \mathcal{X})$ induced by $>$ and $w$ is a simplification order.
\end{theorem}
    %     \subsection{dependency pair}
    %     
Dependency pairs \cite{arts2000termination} offer a sophisticated method for analyzing the termination of term rewriting systems (TRS). This technique transforms the termination problem into a more manageable task by focusing on pairs of terms, where the goal is to prove that no infinite chains of these pairs exist. By studying the relationships between these pairs, one can draw conclusions about the termination of the entire system.

In this section, we explore the dependency pair method from a different viewpoint, presenting variations and extensions that provide new insights into its application. We discuss key theorems that establish the equivalence between the termination of a TRS and the termination of certain relations involving dependency pairs.

\begin{theorem}[\cite{arts2000termination} Theorem 6]
Let $\mathcal{R}$ be a term rewriting system. 

Let $\leadsto_1$ be a binary relation stable under context and substitution on $T(\Sigma,\mathcal{X})^2$ such that $$(s,t) \leadsto (s',t') \text{ iff } t \mathop{\to} _\mathcal{R}^* s'$$

Let $\leadsto_2$ be a binary relation stable under context and substitution on $T(\Sigma,\mathcal{X}) \mathop{\times} T(\Sigma,\mathcal{X})^2$ such that $r \leadsto_2 (s,t)$ iff 
\begin{itemize}
    \item $s$ is a subterm of $r$, and
    \item propre subterms of $s$ are all strongly normalizing.
\end{itemize}

Let $\leadsto \isdef \leadsto_1 \mathop{\cup} \leadsto_2$.
The following three assertions are equivalent:
\begin{itemize}
    \item $\mathcal{R}$ terminates.
    \item $((T(\Sigma,\mathcal{X}) \mathop{\times} T(\Sigma,\mathcal{X})^2) \mathop{\cup} T(\Sigma,\mathcal{X})^2, \leadsto)$ terminates.
    \item $(T(\Sigma,\mathcal{X})^2, \leadsto_2)$ terminates.
\end{itemize} 
\end{theorem}

In the literature, instead of taking the union of two sets with different types of elements, many authors add new symbols to $\Sigma$ to consider a domain with a simple type. However, sets with different types of elements are frequently used in computer science; this is essentially a sum type (also known as a disjoint union type, coproduct type, or variant type).


\begin{definition}[Defined Symbols and Constructors \cite{arts2000termination}]
  Let $\mathcal{R}(\mathcal{F}, R)$ be a TRS. The set $D_{\mathcal{R}}$ of \emph{defined symbols} of $\mathcal{R}$ is defined as 
  \[
  \{ \text{root}(l) \mathop{\mid} l \mathop{\to} r \mathop{\in} R \}
  \]
  and the set $C_{\mathcal{R}}$ of \emph{constructors} of $\mathcal{R}$ is defined as $\mathcal{F} \mathop{\setminus} D_{\mathcal{R}}$.
\end{definition}


\begin{definition}[Dependency Pair \cite{arts2000termination}]
  Let $\mathcal{R}(D, C, R)$ be a TRS. If 
  \[
  f(s_1, \ldots, s_n) \mathop{\to} C[g(t_1, \ldots, t_m)]
  \]
  is a rewrite rule of $R$ with $g \mathop{\in} D$, then 
  \[
  \langle F(s_1, \ldots, s_n), G(t_1, \ldots, t_m) \rangle
  \]
  is called a \emph{dependency pair} of $\mathcal{R}$.
\end{definition}


\begin{theorem}[\cite{arts2000termination} Theorem 7]
$(T(\Sigma,\mathcal{X})^2, \leadsto_2)$ is terminating if and only if there exists a well-founded, weakly monotonic quasi-ordering $\geq$, where both $\geq$ and $>$ are stable under substitution, such that:
\begin{itemize}
    \item \( l \mathop{\geq} r \) for all rules $l \mathop{\rightarrow} r$ in $R$, and
    \item \( s \mathop{>} t \) for all dependency pairs \( \langle s, t \rangle \).
\end{itemize} 
\end{theorem}
    
\begin{corollary}
$\mathcal{R}$ terminates if and only if $DP(R)$ terminates relative to $\mathcal{R}$.
\end{corollary}
  

\chapter{A Morphism from GLS rewriting relations to ACTRS rewriting relations}
\begin{itemize}
    \item This part of the thesis is dedicated to the study of morphisms from graph relabeling systems to term rewriting systems.
    \item Motivation: term rewriting systems are special cases of rewriting systems applied to terms, on which advanced techniques for proving termination is the most studied relative to other rewriting systems.
\end{itemize} 
    \label{sec:morphism_from_gls_to_actrs}
    In this section, we suppose that 
\begin{itemize}
    \item Graph nodes are natural numebrs.
    \item Every node in a graph has at least one edge.
\end{itemize}
Let $\mathcal{R}$ be finite set of graph relabeling rules.
We establish a mapping $\tau$ that maps graph relabeling rules to term rewriting rules such that there exists a homomorphism from the rewriting relation generated by $\mathcal{R}$ to the rewriting relation on terms generated by $\tau(\mathcal{R})$. We will firstly define the signature and variables of the target term rewriting system, then we will define $\tau$ as a mapping from labeled graphs to terms. Finally, we will extend $\tau$ to graph relabeling rules and state the homomorphism of rewriting relations.

\begin{definition}[Signature and variables of the target term rewriting system]
    The variable set of the target term rewriting system is \(\mathcal{Y} \mathop{=} \{ Y_1, Y_2, \ldots, Y_n \}\). The signature set $\Sigma'$ is defined as the disjoint union of the following sets of symbols:
    \begin{align*}
        \Sigma'_0 &= \Sigma' \mathop{\cup} \{ \text{zero}, \epsilon \}, \\ 
        \Sigma'_1 &= \{\text{succ}, \text{node}\}, \\
        \Sigma'_2 &= \{ \text{edge}, \lambda\}, \\
        \Sigma'_{ac} &= \{ * \},
    \end{align*}
    
    where each symbol in \(\Sigma_k\) has arity \(k\) for all \(k \mathop{\in} \{1,2,3\}\), and the symbols in \(\Sigma_{ac}\) are binary.
    % The set of variables is denoted by \(\chi \mathop{=} \{ Y_1, Y_2, \ldots, Y_n \}\), where
    % \[
    % n \isdef \max \{\text{card}(V(G)) \mid \text{for each rule } r \mathop{=} (G, \lambda, \lambda') \mathop{\in} \mathcal{R}\}
    % \]
    % is the maximum number of nodes present in the graphs of the rules within the relabeling system \(\mathcal{R}\).
    \end{definition}
    \begin{example}
        \label{example_grs_to_trs}
        
        Consider a Graph Relabeling System (GRS) \(\mathcal{R} \mathop{=} \{r\}\) with a unique graph relabeling rule \(r\) depicted as follows:
    
        \begin{tikzpicture}
            \draw (-1,0) node[left] {$r$:};
            \node[draw,circle] (x1) at (0,0) {$1$};
            \draw (0,0.3) node[above]{A};
            \node[draw,circle, minimum size \mathop{=} 1pt] (x2) at (2,0) {$2$};
            \draw (2,0.3) node[above]{N};
            \draw[-] (x1) -- (x2);
            
            \draw (3,0) node[right] {$\longrightarrow$};
            
            \node[draw,circle] (x3) at (5,0) {$1$};
            \draw (5,0.3) node[above]{A};
            \node[draw,circle, minimum size \mathop{=} 1pt] (x4) at (7,0) {$2$};
            \draw (7,0.3) node[above]{A};
            \draw[-] (x3) -- (x4);
        \end{tikzpicture}
        
        Given a connected labeled graph \((G, \lambda)\), where all nodes are labeled \(N\) except for one node labeled \(A\), and all edges are labeled \(0\), the set of edges labeled by \(1\) in an \(\mathcal{R}\)-normal labeled graph derived from \((G, \lambda)\) forms a spanning tree of \((G, \lambda)\).
        
        The signature of the corresponding term rewriting system \(\mathcal{R}'\) is defined as \(\Sigma' \mathop{=} \Sigma'_0 \mathop{\cup} \Sigma'_1 \mathop{\cup} \Sigma'_2 \mathop{\cup} \Sigma'_3\), where

        \begin{align*}
            \Sigma'_0 &= \{ N, A, 0, 1, \text{zero} \}, \\
            \Sigma'_1 &= \{ \text{succ}, \text{node} \}, \\
            \Sigma'_2 &= \{ \lambda, \text{edge}, * \},
        \end{align*}
        
        The associative-commutative symbol is denoted by \(\Sigma'_{ac} \mathop{=} \{ * \}\).
    
        The set of variables is \(\chi \mathop{=} \{ Y_1, Y_2 \}\).
        
    \end{example}    

    Since the signature of the term rewriting system is defined, to encode a labeled graph, it suffices to encode its labeling function.

    \begin{definition}[Translation of nodes]
        \label{def:gls_to_actrs:nodes}
        A node \(n\) in a labeled graph is represented by the term $\text{succ}(\text{succ}(\ldots(\text{succ}(\text{zero}))))$.
    \end{definition}

    \begin{definition}[Translation of edges]
        \label{def:gls_to_actrs:edges}
        An arrow $s \overset{l}{\longrightarrow} d$ in a labeled graph is represented by the term $\tau(edge(s', d'), l)$ where $s'$ and $d'$ are the terms representing the nodes $s$ and $d$ respectively.
    \end{definition}
 
\begin{definition}[Translation of graphs]
    \label{def:gls_to_actrs:graphs}
    A labeled graph $G$ is represented by the term $\tau(G)$, which is the concatenation of the terms representing its nodes and edges with the symbol $*$.
\end{definition}
\begin{example}
    The graph  
    \tikz{
            \node[draw,circle] (x1) at (0,0) {$1$};
            \draw (0,0.3) node[above]{A};
            \node[draw,circle, minimum size \mathop{=} 1pt] (x2) at (2,0) {$2$};
            \draw (2,0.3) node[above]{N};
            \draw[-] (x1) -- (x2);
    } is represented by the term 
    $$\lambda(node(succ(zero)),A) * \lambda(node(succ(succ(zero))),N) * \lambda(edge(succ(zero), succ(succ(zero))),0)$$
\end{example}

\begin{example}
    The graph 
    \tikz{
            \node[draw,circle] (x3) at (5,0) {$1$};
            \draw (5,0.3) node[above]{A};
            \node[draw,circle, minimum size \mathop{=} 1pt] (x4) at (7,0) {$2$};
            \draw (7,0.3) node[above]{A};
            \draw[-] (x3) -- (x4);
    } is represented by the term
    $$\lambda(node(succ(zero)),A) * \lambda(node(succ(succ(zero))),A) * \lambda(edge(succ(zero), succ(succ(zero))),1)$$
\end{example}

    % \begin{definition}[Translation of graphs]
    %     \label{def:gls_to_actrs}
    %     Let $(G,\Lambda)$ be a labeled graph.
    %     Let \(\circledast\) be the function that concatenates a set of terms ordered arbitrarily by inserting $*$ between them.

    %     $\tau((G,\Lambda))$ is defined as the term:
    %     \[
    %      \circledast \{ \lambda(\tau(src(e)),\tau(dst(e)),l) \mid 
    %         e \mathop{\in} E(G) \land
    %      \Lambda(e) \mathop{=} l \} 
    %      \]
    %     and $N$ be the term
    %     \[
    %      \circledast \{ \lambda(\tau(n),\Lambda(n)) \mid n \mathop{\in} V(G) \}
    %      \]
    % \end{definition}
     
\textcolor{red}{
    \begin{example}
        Consider the following labeled graph:
        \begin{tikzpicture}
        \draw (1,0) node[above] {0};
        \draw (1,-2) node[below] {0};
        \draw (0,-1) node[left] {0};
        \draw (2,-1) node[right] {0};
        \node[draw,circle] (x_1) at (0,0) {$1$};
        \draw (-0.3,0) node[left]{A};
        \node[draw,circle, minimum size \mathop{=} 1pt] (x_4) at (2,0) {$4$};
        \draw (2.3,0) node[right]{N};
        \node[draw,circle] (x_2) at (0,-2) {$2$};
        \draw (-0.3,-2) node[left]{N};
        \node[draw,circle, minimum size \mathop{=} 1pt] (x_3) at (2,-2) {$3$};
        \draw (2.3,-2) node[right]{N};
        \draw[-] (x_1) -- (x_2);
        \draw[-] (x_2) -- (x_3);
        \draw[-] (x_3) -- (x_4);
        \draw[-] (x_1) -- (x_4);
        \end{tikzpicture}
        This labeled graph is translated into the following term:
        \begin{align*}
            &\lambda(\text{node}(1),A) * \lambda(\text{node}(2),N) * \lambda(\text{node}(3),N) * \lambda(\text{node}(4),N) * \\
            &\lambda(\text{edge}(1,2),0) * \lambda(\text{edge}(2,3),0) * \lambda(\text{edge}(3,4),0) * \lambda(\text{edge}(4,1),0) * \\
            &\lambda(\text{edge}(2,1),0) * \lambda(\text{edge}(3,2),0) * \lambda(\text{edge}(4,3),0) * \lambda(\text{edge}(1,4),0)
        \end{align*}
    \end{example}
}

\begin{definition}[Translation of Graph Relabeling Rule]
    \label{def:gls_to_actrs:rules}
    Let $L \mathop{\rightarrow} R$ be a graph relabeling rule and $L'$ and $R'$ be the terms representing the graphs $L$ and $R$ respectively. The translation of the graph relabeling rule is the term rewriting rule $L'' \mathop{\rightarrow} R''$ where $L''$ and $R''$ are the terms obtained from $L'$ and $R'$, respectively, by replacing every subterm representing a natural number $n$ with the variable $Y_n$.
\end{definition}

    % \begin{definition}[Translation of Graph Relabeling Rule]
    %     A graph relabeling rule $L \mathop{\to} R$ is translated to \(\Gamma(\tau(L)) \mathop{\to} \Gamma(\tau(R)) \), a term rewriting rule with associative-commutative symbol, where the function \(\Gamma\) replaces all occurrences of \(\tau(n)\) with the variable \(Y_n\). 
    % \end{definition}

    \begin{example}
        The rewriting rule 
        \begin{tikzpicture}
            \draw (-1,0) node[left] {$r$:};
            \node[draw,circle] (x1) at (0,0) {$1$};
            \draw (0,0.3) node[above]{A};
            \node[draw,circle, minimum size \mathop{=} 1pt] (x2) at (2,0) {$2$};
            \draw (2,0.3) node[above]{N};
            \draw[-] (x1) -- (x2);
            
            \draw (3,0) node[right] {$\longrightarrow$};
            
            \node[draw,circle] (x3) at (5,0) {$1$};
            \draw (5,0.3) node[above]{A};
            \node[draw,circle, minimum size \mathop{=} 1pt] (x4) at (7,0) {$2$};
            \draw (7,0.3) node[above]{A};
            \draw[-] (x3) -- (x4);
        \end{tikzpicture}
        is represented by the term rewriting rule $\tau(r)$:
        \begin{align*}
        &\lambda(\opn{node}(Y_1), A) * \lambda(\opn{node}(Y_2), N) * \lambda(\opn{edge}(Y_1,Y_2), 0) * \lambda(\opn{edge}(Y_2,Y_1), 0) \\
        &\longrightarrow \\
        &\lambda(\opn{node}(Y_1), A) * \lambda(\opn{node}(Y_2), A) * \lambda(\opn{edge}(Y_1,Y_2), 1) * \lambda(\opn{edge}(Y_2,Y_1), 1)
    \end{align*}
    \end{example}
    
    %  \begin{lemma}
    %  \label{grs_to_trs_step}
    %  Si $(G, \lambda) \mathop{\to} _{(R, \lambda) \mathop{\to} (R,\lambda'),\varphi}$ est une étape de ré-étiquetage, alors $\overline{(G,\lambda)} \longrightarrow_{\mathcal{R}'} \overline{(G,\lambda')}$.
    %  \end{lemma}
    %  \begin{proof}
    %  Par définition, $m$ est une occurrence de $(G_R,\lambda_R)$ dans $(G,\lambda)$. Par conséquent, $(\varphi(G_R), \lambda_R \circ \varphi^{-1})$ est un sous-graphe de $(G,\lambda)$. Ainsi, $\overline{(\varphi(G_R), \lambda_R \circ \varphi^{-1})}$ est un sous-terme de $\overline{(G,\lambda)}$. Comme le symbole * est associative et commutative, $\overline{(G,\lambda)}$ peut s'écrire $\overline{(\varphi(G),\lambda_R \circ \varphi^{-1})} * A$ où A est un terme de $T(\mathcal{F})$. De même, $\overline{(G,\lambda')}$ peut s'écrire $\overline{(\varphi(G), \lambda_R' \circ \varphi^{-1})} * B$ où B est un terme de $T(\mathcal{F})$. Or on a $A=B$ puis que $\lambda$ et $\lambda'$ ne diffèrent que sur les noeuds et les arêtes du graphe $(G_R,\lambda_{G_R})$.
     
    %  Posons $\sigma \mathop{=} \{(X_v,\varphi(v))\ |\ v\in V(G_R)\}$.  En appliquant la substitution $\sigma$ à la règle $\mathcal{R}^{'}_{R}$, on obtient 
    %  $ \sigma (\overline{\Gamma(\overline{\lambda_R})}) \longrightarrow  \sigma  (\overline{\Gamma (\overline{\lambda_R'})} )$, c'est-à-dire $\overline{(G_R, \lambda_R \circ \varphi^{-1})} \longrightarrow \overline{(G_R, \lambda_R' \circ \varphi^{-1})} $.
     
    %  On en conclut donc $\overline{\lambda} \longrightarrow_{\mathcal{R}'} \overline{\lambda'}$.
    %  \end{proof}
% \begin{lemma}[Preservation of Termination]

% The translation $\tau$ restricted to labeled graphs is a homomorphism, or equivalently : if $(G, \Lambda) \mathop{\to} _{(R, \lambda) \mathop{\to} (R, \Lambda'), \varphi}$ represents a relabeling step, then $\overline{(G,\lambda)} \longrightarrow_{\mathcal{R}'} \overline{(G,\lambda')}$.
% \end{lemma}
\begin{lemma}[Preservation of Termination]
\label{grs_to_trs_step}
Let $\rho \mathop{=} ((R,\lambda) \mathop{\to} (R,\lambda'))$ be a graph relabeling rule. 
Let $(G, \Lambda), (G,\Lambda')$ be two labeled graphs.

If $(G, \Lambda)$ rewrites to $(G, \Lambda')$ using the rule $\rho$, then the term $\tau(G,\Lambda)$ rewrites to the term $\tau(G,\Lambda')$ using the term rewriting rule $\rho' \mathop{=} (\tau(R,\lambda) \mathop{\to} \tau(R,\lambda'))$.
\end{lemma}

\begin{proof}
    Suppose $(G, \Lambda)$ rewrites to $(G, \Lambda')$ using the rule $\rho$. 
    
    By definition, there is a graph isomorphism $m: (R,\lambda) \mathop{\to} (G,\Lambda)$ such that $m(R)$ is a subgraph of $(G,\Lambda)$.
    
    Consequently, $(m(R), \lambda \circ m^{-1})$ is a subgraph of $(G,\Lambda)$. 
    
    Thus, $\tau(m(R), \lambda \circ m^{-1})$ is a subterm of $\tau(G,\Lambda)$. 
    
    Since the symbol $*$ is associative and commutative, we have $$\tau(G,\Lambda) \mathop{\to} _\mathcal{AC}^* \tau(m(R),\lambda \circ m^{-1}) * A$$
    where $A$ is a term. 
    
    Similarly, 
    $$\tau(G,\Lambda')\to_\mathcal{AC}^* \tau(m(R), \lambda' \circ m^{-1}) * B$$
    where $B$ is a term of $T(\mathcal{F})$. 

    Moreover, we have $$A=B$$
    since $\Lambda$ and $\Lambda'$ differ only on the nodes and edges of the unlabeled graph $m(R)$.

    Let $\sigma$ be the substitution which replaces every variable $Y_k$ with $m(k)$, where $k \mathop{\in} V(R)$.
    
    We have 
    $$ \sigma(\opn{lhs}(\rho')) \mathop{=} \tau(m(R), \lambda \circ m^{-1})$$ 
    and  
    $$\sigma(\opn{rhs}(\rho')) \mathop{=} \tau(m(R), \lambda' \circ m^{-1})$$

    Therefore, we conclude that $\tau(m(R),\lambda \circ m^{-1}) * A$ rewrites to $\tau(m(R),\lambda' \circ m^{-1}) * A$ using the term rewriting rule $\rho'$ at position $1$.

    We conclude that $\tau(G,\Lambda)$ rewrites to $\tau(G,\Lambda')$ using the term rewriting rule $\rho'$ in \textbf{ACTRS}.
\end{proof}


% \begin{proof}
% Suppose $(G, \Lambda) \mathop{\to} _{(R, \lambda) \mathop{\to} (R, \lambda'), m} (G, \Lambda)$.

% By definition, $m: (R, \lambda) \mathop{\to} (G, \Lambda)$ is a graph isomorphism. 

% Let $f:E(m(R)) \mathop{\to} \Sigma$ defined by $f(e') \mathop{=} \operatorname{min}_< \set{\Lambda(e)~\mid~e \mathop{\in} E(R)}$ where $\operatorname{min}_<$ return the minimal element in a subset of $\Sigma$. 
% The graph $(m(R), f)$ is a subgraph of $(G, \Lambda)$. Consequently, $\tau(m(G), f)$ is a subterm of $\tau(G, \Lambda)$. Since the symbol * is associative and commutative, $\tau(G, \lambda) \mathop{\to} _\mathcal{AC}^* (m(R), f) * A$, where $A$ is a term in $T(\Sigma',\mathcal{X})$ where $\Sigma'$ is a set of symbols. Similarly, $\tau(G, \Lambda') \mathop{\to} _\mathcal{AC}^*   * B$, where $B$ is a term in $T(\Sigma',\mathcal{X})$. We have $A \mathop{=} B$ because $\Lambda$ and $\Lambda'$ differ only in the nodes and edges of the graph $(R, \lambda)$.
% We have, finally $(m(R), f) * A \mathop{\to} _\rho$
% % Let $\sigma \mathop{=} \{(Y_k, m(k)) \mid k \mathop{\in} V(R)\}$. By applying the substitution $\sigma$ to the rule $\mathcal{R}'_{R}$, we obtain $\sigma (\overline{\Gamma(\overline{\lambda_R})}) \longrightarrow \sigma (\overline{\Gamma (\overline{\lambda_R'})} )$, which simplifies to $\overline{(G_R, \lambda_R \circ m^{-1})} \longrightarrow \overline{(G_R, \lambda_R' \circ m^{-1})}$.

% % Thus, it follows that $\overline{\lambda} \longrightarrow_{\mathcal{R}'} \overline{\lambda'}$.
% \end{proof}


% \begin{lemma}
%     \label{grs_to_trs_sequence}
%     If $(G, \Lambda_0) \mathop{\to} _\mathcal{R} (G, \Lambda_1) 
%         \mathop{\to} _\mathcal{R} \ldots  \mathop{\to} _\mathcal{R}
%         (G, \Lambda_{n-1}) \mathop{\to} _\mathcal{R}
%         (G, \Lambda_n)$ is a relabeling sequence in $\mathcal{R}$, then 
%     $\overline{(G_0, \lambda_0)} \longrightarrow_{\mathcal{R}'} 
%     \overline{(G_1, \lambda_1)} \longrightarrow_{\mathcal{R}'} \ldots 
%     \overline{(G_n, \lambda_n)}$ 
%     is a rewriting sequence in $\mathcal{R}'$.
%     \end{lemma}
    
%     \begin{proof}
%     The proof is by induction on the length of the relabeling sequence, using Lemma \ref{grs_to_trs_step}.
%     \end{proof}
    
    \begin{theorem}
        A graph relabeling system $\mathcal{R}$ terminates, if its translation $\mathcal{R}'$ terminates
    \end{theorem}
    % \begin{proof}
    % The proof is by contradiction, utilizing Lemma \ref{grs_to_trs_sequence}.
    % \end{proof}
    % \section{PGLS rewriting relations to ACNTRS rewriting relations}
    % \label{sec:morphism_from_gls_to_ntrs}
    % 
\subsection{Rewriting systems with normalization rules on terms with associative commutative symbols (NTRS)}

Normalized Term Rewriting Systems (Normalized TRSs), introduced by Claude Marché in \cite{marche1996normalized}, impose specific syntactic constraints on terms before rewriting, enhancing the efficiency and applicability of rewriting techniques. This approach reduces term configurations, improving predictability and efficiency, especially in automated theorem proving and symbolic computation. Normalized TRSs offer advantages in termination and confluence analysis by ensuring consistent normalization, leading to more efficient algorithms and simplified proofs. 


\begin{definition}[NTRS rewriting framework]
  Let $\mathcal{N}$ and $\mathcal{R}$ be sets of TRS rules.

  We define the rewriting framework $\mathfrak{ntrs}(\mathcal{N},\mathcal{R})$ as follows: for all $\rho \mathop{\in} \mathcal{N}$, $\mathfrak{ntrs}(\mathcal{N},\mathcal{R})(\rho)$ is the collection of all ACTRS witnesses of rewriting steps; for all $\rho \mathop{\in} \mathcal{R}$, $\mathfrak{ntrs}(\mathcal{N},\mathcal{R})(\rho)$ is the collection of ACTRS witnesses of rewriting steps $(\rho, s, p, \sigma)$ such that $s$ is in $\to_\mathcal{N}$-normal form.

  Let $M, \mathfrak{M}, W, \mathfrak{W}, \mathfrak{I}$ be the class of ACTRS matches, the match mechanism, the class of ACTRS witnesses, the witness mechanism and the interpretation function, respectively, defined in Definition~\ref{def:trs:actrs}.

  The rewriting system $(T(\Sigma,\mathcal{X}), \mathcal{R}, M, \mathfrak{M}, W, \mathfrak{W}, \mathfrak{I}, \mathfrak{ntrs}(\mathcal{N},\mathcal{R}))$ is called a \textbf{NTRS rewriting system}.
\end{definition}

    % \section{FCGLS rewriting relations to Innermost HTRS}
    % \label{sec:morphism_from_gls_to_innermost_htrs}
    % 
% \begin{remark}
%     Avec la strategy Innermost, pour les modules dont la valeur de retour est booleanne, il suffit de donner des regles pour la reponse positive. Tout cela a l'aide d'un constructor "uniform" et le module "if-then-else". exemple: en remplasant l'expression suivante:
% $$eq\_list\_nat(@ * 1, @ * 2)$$
% par 
% $$if\_then\_else(uniform(eq\_list\_nat(@ * 1, @ * 2)), true, false)$$
% on peut connaitre les cas negatives, en raison des trois regles suivantes et la strategy Innermost:
% \begin{flalign*}
%     1\\
%     \hspace{1cm}
%     &if\_then\_else(true, X, Y)
%     \\
%     &\longrightarrow
%     \\
%     &X\\
%     2\\
%     \hspace{1cm}
%     &if\_then\_else(uniform(B), X, Y)
%     \\
%     &\longrightarrow
%      \\
%      &Y\\
%     3\\
%     \hspace{1cm}
%     &uniform(true)
%     \\
%     &\longrightarrow
%     \\
%     &true
% \end{flalign*}
% Pour appliquer la deuxieme regle, il faut que B soit en forme normal, sinon $uniform(B)$ n'est pas en forme normal. De plus, il faut que $B \not \equiv true$, sinon $uniform(B)$ peut etre reduit en $true$ par la troisieme regle.
% \begin{itemize}
%     \item on simplifie au fond le module servant a comparer les entiers naturels
%     \item moins de regles: grâce à la statégie innermost et les trois regles mentionnes
%     \item moins de pairs de dependences générés: parcequ'il y a moins de règle
% \end{itemize}

% \end{remark} 

% \begin{remark}
%     Using the Innermost strategy, for modules with boolean return values, it is sufficient to define rules for the positive outcome only. This can be achieved by utilizing a "uniform" constructor and the "if-then-else" module. For example, by replacing the expression:
%     $$eq\_list\_nat(@ * 1, @ * 2)$$
%     with 
%     $$if\_then\_else(uniform(eq\_list\_nat(@ * 1, @ * 2)), true, false),$$
%     negative cases can be handled due to the following three rules and the Innermost strategy:
%     \begin{flalign*}
%         1. &\hspace{1cm} if\_then\_else(true, X, Y) \longrightarrow X \\
%         2. &\hspace{1cm} if\_then\_else(uniform(B), X, Y) \longrightarrow Y \\
%         3. &\hspace{1cm} uniform(true) \longrightarrow true
%     \end{flalign*}
%     To apply the second rule, $B$ must be in normal form; otherwise, $uniform(B)$ will not be in normal form. Additionally, $B$ must not be equivalent to $true$, as $uniform(B)$ could be reduced to $true$ by the third rule.


%     % \begin{itemize}
%     %     \item Simplifies the module used for comparing natural numbers.
%     %     \item Fewer rules are required, thanks to the Innermost strategy and the three mentioned rules.
%     %     \item Fewer dependency pairs are generated due to the reduced number of rules.
%     % \end{itemize}
% \end{remark}
This section designs a hierarchical term rewriting system (HTRS) to simulate graph relabeling systems. It introduces two sets of modules: one specific to each graph relabeling system and the other common to all. These modules manage tasks such as generating and comparing graphs, checking for forbidden contexts, identifying applicable rewriting rules, and applying rewriting steps.

\begin{example}
    \label{example_fcpgrs}
    We present our translation using the graph relabeling system with priorities and forbidden contexts, $\mathcal{R} = \{r_1,\ r_2,\ r_3\}$
  
    \begin{tikzpicture}
      \draw (-1,0) node[left] {$r_1$:};
      \draw (1,0) node[above] {0};
      
      \node[draw,circle] (x_1) at (0,0) {$1$};
      \draw (0,0.3) node[above]{A};
      \node[draw,circle, minimum size = 1pt] (x_2) at (2,0) {$2$};
      \draw (2,0.3) node[above]{N};
      \draw[-] (x_1) -- (x_2);
      
      \draw (3,0) node[right] {$\longrightarrow$};
      \draw (6,0) node[above] {1};
      
      \node[draw,circle] (x3) at (5,0) {$1$};
      \draw (5,0.3) node[above]{A};
      \node[draw,circle, minimum size = 1pt] (x4) at (7,0) {$2$};
      \draw (7,0.3) node[above]{A'};
      \draw (7.5,-0.45) node[above]{$,$};
      \draw (8,-0.45) node[above]{$\Big\{$};
  
      % Context 2
      \node[draw,circle](c2_1) at (11,0.75) {$1$};
      \draw (11,1) node[above] {N};
      \draw (10.25,-1) node[below] {A};
      \draw (11.75,-1) node[below] {A'};
      \node[draw,circle](c2_2) at (10.25,-0.75) {$2$};
      \node[draw,circle](c2_3) at (11.75,-0.75) {$3$};
      \draw[-] (c2_1) -- (c2_2);
      \draw (10.6,0) node[left] {0};
      \draw (11.35,0) node[right] {0};
      \draw[-] (c2_1) -- (c2_3);
  
      \draw (12.5,-0.45) node[above]{$\Big\}$};
      \draw[-] (x3) -- (x4);
    \end{tikzpicture}
    
    \begin{tikzpicture}
      \draw (-1,0) node[left] {$r_2$:};
      \draw (1,0) node[above] {0};
      
      \node[draw,circle] (x_1) at (0,0) {$1$};
      \draw (0,0.3) node[above]{A'};
      \node[draw,circle, minimum size = 1pt] (x_2) at (2,0) {$2$};
      \draw (2,0.3) node[above]{N};
      \draw[-] (x_1) -- (x_2);
      
      \draw (3,0) node[right] {$\longrightarrow$};
      \draw (6,0) node[above] {1};
      
      \node[draw,circle] (x3) at (5,0) {$1$};
      \draw (5,0.3) node[above]{A'};
      \node[draw,circle, minimum size = 1pt] (x4) at (7,0) {$2$};
      \draw (7,0.3) node[above]{A'};
      \draw (7.5,-0.45) node[above]{$,$};
      \draw (8,-0.45) node[above]{$\Big\{$};
      \draw (8.5,-0.45) node[above]{$\Big\}$};
      \draw[-] (x3) -- (x4);
    \end{tikzpicture}
    
    \begin{tikzpicture}
      \draw (-1,0) node[left] {$r_3$:};
      
      \node[draw,circle] (x_1) at (0,0) {$1$};
      \draw (0,0.3) node[above]{A'};
      
      \draw (1,0) node[right] {$\longrightarrow$};
      
      \node[draw,circle] (x3) at (3,0) {$1$};
      \draw (3,0.3) node[above]{F};
      
      % Context 1
      \node[draw,circle](c1_1) at (5,0.75) {$1$};
      \draw (5,1) node[above] {A'};
      \node[draw,circle](c1_2) at (5,-0.75) {$2$};
      \draw (5,-1) node[below] {N};
      \draw (5,0) node[left] {0};
      \draw[-] (c1_1) -- (c1_2);
      \draw (5.5,-0.45) node[above]{$,$};
      
      % Context 2
      \node[draw,circle](c2_1) at (7,0.75) {$1$};
      \draw (7,1) node[above] {A'};
      \draw (6.25,-1) node[below] {A'};
      \draw (7.75,-1) node[below] {A'};
      \node[draw,circle](c2_2) at (6.25,-0.75) {$2$};
      \node[draw,circle](c2_3) at (7.75,-0.75) {$3$};
      \draw[-] (c2_1) -- (c2_2);
      \draw (6.6,0) node[left] {1};
      \draw (7.35,0) node[right] {1};
      \draw[-] (c2_1) -- (c2_3);
      \draw (8,-0.45) node[above]{$,$};
      
      % Context 3
      \node[draw,circle](c3_1) at (9.25, 0.75) {$1$};
      \node[draw,circle](c3_2) at (8.75, -0.75) {$2$};
      \node[draw,circle](c3_3) at (9.75, -0.75) {$3$};
      \draw (9.25,1) node[above] {A'};
      \draw (8.75,-1) node[below] {A'};
      \draw (9.75,-1) node[below] {A};
      \draw (9,0) node[left] {1};
      \draw (9.5,0) node[right] {1};
      \draw[-] (c3_1) -- (c3_2);
      \draw[-] (c3_1) -- (c3_3);
      
      \draw (3.5,-0.45) node[above]{$,$};
      \draw (4,-0.45) node[above]{$\Big\{$};
      \draw (10.5,-0.45) node[above]{$\Big\}$};
    \end{tikzpicture}
  \end{example} 
  




  \subsection{System-Specific Modules}


%   \subsection*{get\_context\_rule\_graphs}
%     \begin{flalign*}
%         \hspace{1cm}
%         1. \\
%         &get\_context\_rule\_graphs(n) &\\
%         &\longrightarrow &\\
%         &\overline{rule\_and\_context\_graphs(n)} &
%     \end{flalign*}

%     Given a rule number $n$, $get\_context\_rule\_graphs(n)$ returns the concatenantion of terms $graph(r,c,k)$ for all $c \in  where $r$ is the rule number, $c$ is the context number (0 for the rule graph), and $K$ is the size of the graph. Subterms are ordered according to the topological order of the priority.

\subsection*{get\_context\_rule\_graphs}
\begin{flalign*}
    \hspace{1cm}
    &get\_context\_rule\_graphs(R) &\\
    &\longrightarrow &\\
    &\circledast \{graph(R, c, k) \mid 0 \leq c \leq \operatorname{length}(f(r_R))
    %  \land (c = 0 \implies k = |lhs(r_N)|) \land 
    % (c \geq 1 \implies k = |\operatorname{codom}(f(r_N))|)
    \} &
\end{flalign*}

where $N$ is the rule number. The right-hand side (rhs) is the concatenation of terms in $\{ graph(R, c, k) \mid 0 \leq c \leq \operatorname{length}(f(r_N))\}$. If $c = 0$, then $k$ is the size of the left-hand side (lhs) graph of the $N$-th rule. Otherwise, $c \geq 1$ and $k$ is the size of the domain of the $c$-th forbidden context of the $N$-th rule.
    

    \subsection*{generate\_graph}
  
\begin{flalign*}
    \hspace{1cm}
    &generate\_graph(R, C, @ * X_n * X_{n-1} * \ldots * X_2 * X_1) & \\
    &\longrightarrow & \\
    & \lambda(node(x_1), l_{node(x_1)}) * \ldots * \lambda(node(x_n), l_{node(x_n)}) * & \\
    & \lambda(edge(x_1,x_2), l_{edge(x_1,x_2)}) * \lambda(edge(x_1,x_3), l_{edge(x_1,x_3)}) * \ldots * \lambda(edge(x_1,x_n), l_{edge(x_1,x_n)}) * & \\
    & \lambda(edge(x_2,x_1), l_{edge(x_2,x_1)}) * \lambda(edge(x_2,x_3), l_{edge(x_2,x_3)}) * \ldots * \lambda(edge(x_2,x_n), l_{edge(x_2,x_n)}) * & \\
    & \ldots & \\
    & \lambda(edge(x_n,x_1), l_{edge(x_n,x_1)}) * \lambda(edge(x_n,x_2), l_{edge(x_n,x_2)}) * \ldots * \lambda(edge(x_n,x_{n-1}), l_{edge(x_n,x_{n-1})}) &
\end{flalign*}
where $R$ is the rule number,  $C$ the number of graph in the rule $r_R$, $n$ the size of the lhs graph, if $C=0$, of the rule $r_R$ or the size of the codomain of the $C$-th of the rule. The context or lhs rule graph; subterms are ordered. The rhs of this rule is the representation of the corresponding labeled graph.

\subsection*{occurrence\_to\_rewrite\_step}

\noindent For each rule $R$ in the FC-P-GRS, there is a corresponding specific rule:
\begin{flalign*}
    \hspace{1cm}
    &occurrence\_to\_rewrite\_step(N, X) & \\
    &\longrightarrow & \\
    &\Gamma(\tau(rhs(r_N)), X)
\end{flalign*}
where $N$ is the rule numnber, $X$ is a term encoding all nodes participating to the rewriting step, and $\Gamma(\tau(rhs(r_N)), X)$ is the representation of the right-hand side of rule $r_N$ but every number $k$ is replaced by $X_k$.

    
    
    \subsection*{main}
    Combining with the innermost strategy, this module allows us to simply our translation by focusing only in the case where $s$ can be reduced to true for every $uniform(s)$.  
    \begin{flalign*}
        & main(
        \\ & \hspace{1cm} LabellingFunction, 
        \\ & \hspace{1cm} @
        \\ & \hspace{1cm} \vdots
        \\ & \hspace{1cm} @
        \\ & )
        \\
        &\longrightarrow
        \\
        & main(
            \\ & \hspace{1cm} LabellingFunction, 
            \\ & \hspace{1cm} identify\_applicable\_occs\_of\_rule(
                \\ & \hspace{2cm} LabellingFunction,
                \\ & \hspace{2cm} get\_context\_rule\_graphs(1), 
                \\ & \hspace{2cm} @, 
                \\ & \hspace{2cm} @,
            \\ & \hspace{1cm}),
            \\ & \hspace{1cm} identify\_applicable\_occs\_of\_rule(
                \\ & \hspace{2cm} LabellingFunction,
                \\ & \hspace{2cm} get\_context\_rule\_graphs(2), 
                \\ & \hspace{2cm} @, 
                \\ & \hspace{2cm} @,
            \\ & \hspace{1cm}),
            \\ & \hspace{1cm} \vdots
            \\ & \hspace{1cm} identify\_applicable\_occs\_of\_rule(
                \\ & \hspace{2cm} LabellingFunction,
                \\ & \hspace{2cm} get\_context\_rule\_graphs(n), 
                \\ & \hspace{2cm} @, 
                \\ & \hspace{2cm} @,
            \\ & \hspace{1cm}),
            \\ & )
    \end{flalign*} \\

    \begin{flalign*}
        & main(
            \\ & \hspace{1cm} LabellingFunction, 
            \\ & \hspace{1cm} identify\_applicable\_occs\_of\_rule(
                \\ & \hspace{2cm} X_1
                \\ & \hspace{2cm} X_2        
                \\ & \hspace{2cm} X_3
                \\ & \hspace{2cm} X_4
            \\ & \hspace{1cm}),
            \\ & \hspace{1cm} identify\_applicable\_occs\_of\_rule(
                \\ & \hspace{2cm} X_5
                \\ & \hspace{2cm} X_6
                \\ & \hspace{2cm} X_7 
                \\ & \hspace{2cm} X_8
            \\ & \hspace{1cm}),
            \\ & \hspace{1cm} \vdots
            \\ & \hspace{1cm} identify\_applicable\_occs\_of\_rule(
                \\ & \hspace{2cm} X_{4n-3}
                \\ & \hspace{2cm} X_{4n-2}
                \\ & \hspace{2cm} X_{4n-1}
                \\ & \hspace{2cm} X_{4n}
            \\ & \hspace{1cm}),
            \\ & )
        \\
        &\longrightarrow
        \\
        & LabellingFunction
    \end{flalign*} \\

    \begin{flalign*}
        1 \le k \le n
        \\ 
        & main(
        \\ & \hspace{1cm} LabellingFunction, 
        \\ & \hspace{1cm} X_1
        \\ & \hspace{1cm} X_2
        \\ & \hspace{1cm} \vdots
        \\ & \hspace{1cm} X_{k-1}
        \\ & \hspace{1cm} some(X_k)
        \\ & \hspace{1cm} X_{k-1}
        \\ & \hspace{1cm} \vdots
        \\ & \hspace{1cm} X_n
        \\ & )
        \\
        &\longrightarrow
        \\
        & main(
        \\ & \hspace{1cm} rewrite(
            \\ & \hspace{2cm} LabellingFunction,
            \\ & \hspace{2cm} generate\_occurrence\_to\_rewrite\_step(k, X_k)
        \\ & \hspace{1cm}) 
        \\ & \hspace{1cm} @
        \\ & \hspace{1cm} \vdots
        \\ & \hspace{1cm} @
        \\ & )
    \end{flalign*} \\


    \subsection*{Wrap}
    With this module and the innermost strategy, we can distinguish boolean expressions that can be evaluated to "true" and others.    

    \begin{flalign*}
        \\ & uniform(true) \longrightarrow true
    \end{flalign*}

    
    \subsection{General Modules}

    The following terminating modules are common for every FCGLS translation.

    \subsection*{if\_then\_else}

    \begin{flalign*}
        \hspace{1cm}
        &if\_then\_else(true, X, Y)&
        \\
        &\longrightarrow
        \\
        &X\\
        \hspace{1cm}
        &if\_then\_else(uniform(B), X, Y)&
        \\
        &\longrightarrow
         \\
         &Y
    \end{flalign*}


    \subsection*{eq\_list\_nat} 
    This module tests if two arguments encode the same list of natural numbers.

    \begin{flalign*}
        &eq\_list\_nat(X * Z, Y * Z)
        \\
        &\longrightarrow
        \\
        &eq\_list\_nat(X, Y)
        \\
        \hspace{1cm}
        &eq\_list\_nat(@, @)
        \\
        &\longrightarrow
        \\
        &true
    \end{flalign*}

    \subsection*{eq\_noe }
    This module compares two terms encoding nodes (or edges) and returns true if they encode the same node (or edge).

    \begin{flalign*}
        \hspace{1cm}
        & eq \_noe(node(X), node(X)) &
        \\
        & \longrightarrow
        \\
        & true
    \end{flalign*}
    
    % \begin{flalign*}
    %     4\\
    %     \hspace{1cm}
    %     & eq \_noe(edge( X, Y), edge( Z, W)) &\\
    %     & \longrightarrow\\
    %     & et(
    %        \\ & \hspace{1cm}eq\_nat(X, Z),
    %        \\ & \hspace{1cm}eq\_nat(Y,W)
    %     \\ & )
    % \end{flalign*}
    \begin{flalign*}
        & eq \_noe(edge( X, Y), edge( X, Y)) &\\
        & \longrightarrow\\
        & true
    \end{flalign*}

    \subsection*{mem} 
    This module tests if the first argument is an element of the second argument.

    \begin{flalign*}
        \hspace{1cm}
        &mem(X, Y * Z) & 
        \\
        &\longrightarrow
        \\
        &if\_then\_else(
            \\ & \hspace{1cm} uniform(eq\_list\_nat(@ * X, @ * Z)),
            \\ & \hspace{1cm} true,
            \\ & \hspace{1cm} mem(X, Y)
        \\ & )
    \end{flalign*}

    \subsection*{union}
    This module returns the concatenation of two lists.

    \begin{flalign*}
        \hspace{1cm}
        &union(L_1 * X, L_2)&
        \\
        &\longrightarrow
        \\
        &union(L_1, L_2 * X)&
        \\
        \\
        &union(@, L)&
        \\
        &\longrightarrow
        \\
        &L&
    \end{flalign*}

    
    % \subsection*{cons\_n}
    % cons\_n(x,l) locates the n-th list of $l_1$ and then adds x in the list :
    % \begin{flalign*}
    %     1\\
    %     \hspace{1cm}
    %     &cons \_n(
    %         \\ & \hspace{1cm} s(s(N)),
    %         \\ & \hspace{1cm} X,
    %         \\ & \hspace{1cm} L * Y
    %         \\ &)&
    %     \\
    %     &\longrightarrow
    %     \\
    %     &cons \_n(
    %         \\ & \hspace{1cm} s(N),
    %         \\ & \hspace{1cm} X,
    %         \\ & \hspace{1cm} L
    %         \\ &) * Y 
    % \end{flalign*}
    % Insert :   
    % \begin{flalign*}
    %     2\\
    %     \hspace{1cm}
    %     &cons \_n(s(o), X, L * Y)&
    %     \\
    %     &\longrightarrow
    %     \\
    %     &L * (Y * X)
    % \end{flalign*} 

    
    
    \subsection*{prefix}
    \noindent \texttt{prefix(X, Y) = true} if and only if the list \texttt{Y} is a prefix of the list \texttt{X}.
    
    
    \begin{flalign*}
        \hspace{1cm}
        &prefix(X, @) &
        \\
        &\longrightarrow 
        \\
        &true&
    \end{flalign*}
    % \begin{flalign*}
    %     \hspace{1cm}
    %     2\\
    %     &prefix(X * Y, W * Z)&
    %     \\
    %     &\longrightarrow 
    %     \\
    %     &and(
    %         \\ & \hspace{1cm} eq\_nat(Y,Z),
    %         \\ & \hspace{1cm} prefix(X, W),
    %         \\ &)
    % \end{flalign*}
    \begin{flalign*}
        \hspace{1cm}
        &prefix(X * Z, Y * Z)&
        \\
        &\longrightarrow 
        \\ &prefix(X, Y),
    \end{flalign*}



    \subsection*{graph\_to\_list\_nat}
    Given a graph labeling function encoded by $@ * \lambda(x_1,l_1) * \ldots * \lambda(x_n,l_n)$, this module returns the list $@ * l_1 * \ldots * l_n$.

    \begin{flalign*}
        \hspace{1cm}
        & graph\_to\_list\_nat(X * \lambda(Y,Z))&
        \\
        &\longrightarrow
        \\
        & graph\_to\_list\_nat(X) * Z
        \\
        \\
        & graph\_to\_list\_nat(@)&
        \\
        &\longrightarrow
        \\
        & @
    \end{flalign*}
    
    \subsection*{find\_label}
    This constructor's first argument is a node or an edge of a graph, and the second argument encodes the labeling function of the graph. It returns the label of the first argument.

        \begin{flalign*}
        \hspace{1cm}
        &find\_label(
            \\ & \hspace{1cm} X,
            \\ & \hspace{1cm} F\ *\ \lambda(Y, L)
            \\ &)
        \\
        &\longrightarrow
        \\
        &if\_then\_else(
            \\ & \hspace{1cm} uniform(eq\_noe(X,Y)),
            \\ & \hspace{1cm} L,   
            \\ & \hspace{1cm} find\_label(X, F)
            \\ &)
    \end{flalign*}
    % \noindent According to the result of key-comparison, either we return the label found or continue our recherche:
    % \begin{flalign*}
    %     2\\ 
    %     \hspace{1cm}
    %     &find\_label(X, F\ *\ \lambda(Y, L),true)&
    %     \\
    %     &\longrightarrow
    %     \\
    %     &L &
    % \end{flalign*}

    
    \subsection*{rewrite\_noe}
    Given a labeling function of a graph as the first argument and a node or an edge $K_2$ with label $L_2$ of the graph as the third argument, this module relabels the node or the edge.
    \begin{flalign*}
        \hspace{1cm}
        &rewrite\_noe(
            \\ & \hspace{1cm} F_1 * \lambda(K_1, L_1), 
            \\ & \hspace{1cm} F_2,
            \\ & \hspace{1cm} \lambda(K_2, L_2)
            \\ &)
        \\
        &\longrightarrow
        \\
        &if\_then\_else(
            \\ & \hspace{1cm} uniform(eq\_noe(K_1,K_2)),
            \\ & \hspace{1cm} union(F_1 * \lambda(K_1, L_2),F_2 ),
            \\ & \hspace{1cm} rewrite\_noe(
                \\ & \hspace{2cm} F_1,
                \\ & \hspace{2cm} F_2 * \lambda(K_1, L_1),
                \\ & \hspace{2cm} \lambda(K_2, L_2)
            \\ & \hspace{1cm})
            \\ & )
    \end{flalign*}
    
    % \begin{flalign*}
    %     2\\
    %     \hspace{1cm}
    %     &rewrite\_noe(
    %         \\ & \hspace{1cm} F_1 * \lambda(K_1, L_1), 
    %         \\ & \hspace{1cm} F_2,
    %         \\ & \hspace{1cm} \lambda(K_2, L_2)
    %         \\ &)
    %     \\
    %     &\longrightarrow
    %     \\
    %     &if\_then\_else(
    %         \\ & \hspace{1cm} eq\_noe(K_1,K_2),
    %         \\ & \hspace{1cm}rewrite\_noe(
    %             \\ & \hspace{2cm} F_1 * \lambda(K_1, L_2), 
    %             \\ & \hspace{2cm} F_2,
    %             \\ & \hspace{2cm} @
    %             \\ & \hspace{1cm}),
    %         \\ & \hspace{1cm} rewrite\_noe(
    %             \\ & \hspace{2cm} F_1,
    %             \\ & \hspace{2cm} F_2 * \lambda(K_1, L_1),
    %             \\ & \hspace{2cm} \lambda(K_2, L_2)
    %         \\ & \hspace{1cm})
    %         \\ & )
    % \end{flalign*}

    % \begin{flalign*}
    %     3\\
    %     \hspace{1cm}
    %     &rewrite\_noe(
    %         \\ & \hspace{1cm} F_1, 
    %         \\ & \hspace{1cm} F_2 * X,
    %         \\ & \hspace{1cm} @
    %         \\ &)
    %     \\
    %     &\longrightarrow
    %     \\
    %     &rewrite\_noe(
    %         \\ & \hspace{1cm} F_1 * X, 
    %         \\ & \hspace{1cm} F_2,
    %         \\ & \hspace{1cm} @
    %         \\ &)
    % \end{flalign*}

    % \begin{flalign*}
    %     4\\
    %     \hspace{1cm}
    %     &rewrite\_noe(
    %         \\ & \hspace{1cm} F_1, 
    %         \\ & \hspace{1cm} @,
    %         \\ & \hspace{1cm} @
    %         \\ &)
    %     \\
    %     &\longrightarrow
    %     \\
    %     &F_1
    % \end{flalign*}


    \subsection*{rewrite}
    Given a labeling function of a graph and a labeling function of a subgraph of the graph, this module relabels nodes and edges in the graph according to the subgraph labeling function.

    \begin{flalign*} 
        \hspace{1cm}
        \\ &rewrite(
            \\ & \hspace{1cm} F_1,
            \\ & \hspace{1cm} F_2 * X
            \\ &)&
        \\
        &\longrightarrow
        \\
        &rewrite(
            \\ & \hspace{1cm} rewrite\_noe(
                \\ & \hspace{2cm} F_1,
                \\ & \hspace{2cm} @,
                \\ & \hspace{2cm} X
                \\ & \hspace{1cm}),
            \\ & \hspace{1cm} F_2
        \\ &)
    \end{flalign*}
    \begin{flalign*}
        \hspace{1cm}
        \\ &rewrite(
            \\ & \hspace{1cm} F,
            \\ & \hspace{1cm} @
        \\ &)
        \\
        &\longrightarrow
        \\
        &F
    \end{flalign*}
    
    
    % \subsection*{apply\_rewrite\_steps}

    % \begin{flalign*}
    %     \hspace{1cm}
    %     1\\
    %     &apply\_rewrite\_steps(
    %         \\ & \hspace{1cm} X * Y,
    %         \\ & \hspace{1cm} F
    %     \\ & )
    %     \\
    %     &\longrightarrow
    %     \\
    %     &apply\_rewrite\_steps(
    %         \\ & \hspace{1cm} X,
    %         \\ & \hspace{1cm} rewrite(F, Y)
    %     \\ & )	
    % \end{flalign*}

    % \begin{flalign*}
    %     \hspace{1cm}
    %     2\\
    %     &apply\_rewrite\_steps(
    %         \\ & \hspace{1cm} @,
    %         \\ & \hspace{1cm} F
    %     \\ & )
    %     \\
    %     &\longrightarrow
    %     \\
    %     &F&
    % \end{flalign*}
    
    
    \subsection*{context\_inclusion\_check}
    This module is designed to make the translation easier to understand.
    \begin{flalign*}
        \hspace{1cm}
        1\\
        & context\_inclusion\_check( R, C)&
        \\
        &\longrightarrow
        \\
        & prefix(R, C)
    \end{flalign*}
    
    
    \subsection*{generate\_complete\_graph}
    Let $x_1,\ldots,x_n$ be natural numbers, $generate\_complete\_graph(@ * x_n * \ldots * x_1,\ @ * x_n * \ldots * x_1,\ @ * x_n * \ldots * x_1,\ @ * x_n * \ldots * x_1,\ @)$
       returns the following term encoding the unlabeled complete graph below:
        \begin{align*}
            & node(x_1) * \ldots * node(x_n) * \\
            & edge(x_1,x_2) * edge(x_1,x_3) * \ldots * edge(x_1,x_n) * \\
            & edge(x_2,x_1) * edge(x_2,x_3) * \ldots * edge(x_2,x_n) * \\
            & \ldots \\
            & edge(x_n,x_1) * edge(x_n,x_2) * \ldots * edge(x_n,x_{n-1})
        \end{align*}
    
    \begin{flalign*}
        1\\
        \hspace{1cm}
        &generate\_complete\_graph(
            \\ & \hspace{1cm} L_1,
            \\ & \hspace{1cm} L_2,
            \\ & \hspace{1cm} L_3, 
            \\ & \hspace{1cm} L_4 * X,
            \\ & \hspace{1cm} L_5
        \\ &)
        \\
        &\longrightarrow
        \\
        &generate\_complete\_graph(
            \\ & \hspace{1cm} L_1,
            \\ & \hspace{1cm} L_2,
            \\ & \hspace{1cm} L_3, 
            \\ & \hspace{1cm} L_4,
            \\ & \hspace{1cm} L_5 * node(X)
        \\ & )
    \end{flalign*}
    \begin{flalign*}
        2\\
        &generate\_complete\_graph(
            \\ & \hspace{1cm} L_1,
            \\ & \hspace{1cm} L_2 * X,
            \\ & \hspace{1cm} L_3 * Y,
            \\ & \hspace{1cm} @,
            \\ & \hspace{1cm} L_4
        \\ &)
        \\
        &\longrightarrow
        \\
        &if\_then\_else(
            \\ & \hspace{1cm} uniform(eq\_list\_nat(@ * X, @ * Y)), 
            \\ & \hspace{1cm} generate\_complete\_graph(
                \\ & \hspace{2cm} L_1,
                \\ & \hspace{2cm} L_2 * X, 
                \\ & \hspace{2cm} L_3,
                \\ & \hspace{2cm} @,
                \\ & \hspace{2cm} L_4
            \\ & \hspace{1cm}),
            \\ & \hspace{1cm} generate\_complete\_graph(
                \\ & \hspace{2cm} L_1,
                \\ & \hspace{2cm} L_2 * X, 
                \\ & \hspace{2cm} L_3,
                \\ & \hspace{2cm} @,
                \\ & \hspace{2cm} L_4 * edge(X,Y)
            \\ & \hspace{1cm})
        \\ &)
    \end{flalign*}
    \begin{flalign*}
        3\\
        &generate\_complete\_graph(
            \\ & \hspace{1cm} L_1, 
            \\ & \hspace{1cm} L_2 * X, 
            \\ & \hspace{1cm} @, 
            \\ & \hspace{1cm} @,
            \\ & \hspace{1cm} L3
        \\ &)
        \\
        &\longrightarrow 
        \\
        &generate\_complete\_graph(
            \\ & \hspace{1cm} L_1,   
            \\ & \hspace{1cm} L_2,
            \\ & \hspace{1cm} L_1,
            \\ & \hspace{1cm} @,
            \\ & \hspace{1cm} L3
        \\ & )
    \end{flalign*}
    \begin{flalign*}
        4\\
        &generate\_complete\_graph(
            \\ & \hspace{1cm} L_1, 
            \\ & \hspace{1cm} @, 
            \\ & \hspace{1cm} L_2, 
            \\ & \hspace{1cm} @, 
            \\ & \hspace{1cm} L_3
        \\ &)&
        \\
        &\longrightarrow
        \\
        &L_3&
    \end{flalign*}
    
    
    \subsection*{construct\_induced\_graph}
    Given a labeling function of a graph $F$ and a list $X$ of nodes in the graph, $generate\_complete\_graph(X,@,F)$ return the subgraph induced by nodes in $X$.  
    % \begin{description}
    %     \item[input] $@ * x_n * x_{n-1} * \ldots * x_1$
    %     \item[output]
    %     \begin{align*}
    %         & \lambda(node(x_1), l_{node(x_1)}) * \ldots * \lambda(node(x_n), l_{node(x_n)}) * \\
    %         & \lambda(edge(x_1,x_2), l_{edge(x_1,x_2)}) * \lambda(edge(x_1,x_3), l_{edge(x_1,x_3)}) * \ldots * \lambda(edge(x_1,x_n), l_{edge(x_1,x_n)}) * \\
    %         & \lambda(edge(x_2,x_1), l_{edge(x_2,x_1)}) * \lambda(edge(x_2,x_3), l_{edge(x_2,x_3)}) * \ldots * \lambda(edge(x_2,x_n), l_{edge(x_2,x_n)}) * \\
    %         & \ldots \\
    %         & \lambda(edge(x_n,x_1), l_{edge(x_n,x_1)}) * \lambda(edge(x_n,x_2), l_{edge(x_n,x_2)}) * \ldots * \lambda(edge(x_n,x_{n-1}), l_{edge(x_n,x_{n-1})})
    %     \end{align*}
    % \end{description}
    

    \begin{flalign*}
        \hspace{1cm}
        & construct\_induced\_graph(
            \\ & \hspace{1cm} Nodes * X,
            \\ & \hspace{1cm} @,
            \\ & \hspace{1cm} LabellingFunction
        \\ & )  &
        \\
        &\longrightarrow
        \\
        & construct\_induced\_graph(
            \\ & \hspace{1cm} @,
            \\ & \hspace{1cm} generate\_complete\_graph(
                \\ & \hspace{2cm} Nodes * X, 
                \\ & \hspace{2cm} Nodes * X,
                \\ & \hspace{2cm} Nodes * X,
                \\ & \hspace{2cm} Nodes * X , 
                \\ & \hspace{2cm} @
                \\ & \hspace{1cm}),
            \\ & \hspace{1cm} LabellingFunction
        \\ & )
    \end{flalign*}

    \begin{flalign*}
        \hspace{1cm}
        & construct\_induced\_graph(
            \\ & \hspace{1cm}@, 
            \\ & \hspace{1cm}Y * Z,
            \\ & \hspace{1cm}F
        \\ & ) 
        \\ &\longrightarrow
        \\ & construct\_induced\_graph(
            \\ & \hspace{1cm}@,
            \\ & \hspace{1cm}Y, 
            \\ & \hspace{1cm} F
        \\ &) * \lambda(Z, find\_label(Z, F))
        \\
        & construct\_induced\_graph(
            \\ & \hspace{1cm} @,
            \\ & \hspace{1cm} @,
            \\ & \hspace{1cm} F
        \\ &) &
        \\
        &\longrightarrow
        \\
        & @
    \end{flalign*}
    
\subsection*{get\_all\_nodes(f)}
Given a labeling function $F$ of a graph, this module reduces the term $get\_all\_nodes(F)$ to a term encoding the list of nodes in the graph.

    \begin{flalign*}
        \hspace{1cm}
        &get\_all\_nodes(
            \\ & \hspace{1cm}L * \lambda(node(X),Y)
        \\ &)
        \\ &\longrightarrow
        \\
        &get\_all\_nodes(
            \\ & \hspace{1cm} L
        \\ &) * X
        \\
        &get\_all\_nodes(
            \\ & \hspace{1cm} L * \lambda(edge(X,Y),W)
        \\ & )
        \\
        &\longrightarrow
        \\
        &get\_all\_nodes(
            \\ & \hspace{1cm} L
        \\ & )
        \\
        &get\_all\_nodes(
            \\ & \hspace{1cm} @
            \\ & )
        \\
        &\longrightarrow
        \\
        &@&
    \end{flalign*}
    
    \subsection*{isomorphism\_verification}
    Given two labeling functions $X$ and $Y$ of graphs sharing the same unlabeld graph, this module rewrites the term $isomorphism\_verification(X,Y)$ to true if the graphs are isomorphic.
    \begin{flalign*}
        \hspace{1cm}
        & isomorphism\_verification(X, Y)&
        \\
        &\longrightarrow
        \\
        % & if\_then\_else(
        %     \\ & \hspace{1cm} uniform( 
                \\ & \hspace{2cm} eq\_list\_nat(
                    \\ &\hspace{3cm}graph\_to\_list\_nat(X),
                    \\ &\hspace{3cm}graph\_to\_list\_nat(Y)
                 \\ & \hspace{2cm})
        %    \\ & \hspace{1cm})
        % \\ &)
    \end{flalign*}
    
    \subsection*{is\_included\_in\_fc}
    Let $X$ be a term encoding a occurrence of a rule lhs graph in a host graph. 
    Let $Y$ be a term encoding all occrrences of contexts in the host graph.
    This module rewrites $is\_included\_in\_fc(X,Y)$ to term $true$ if $X$ is included in a forbidden context in $Y$, and $false$ otherwise.

    \begin{flalign*}
        \hspace{1cm}
        & is\_included\_in\_fc(
            \\ & \hspace{1cm} Nodes,
            \\ & \hspace{1cm} Contexts\_occs * X,
        \\ & )
        \\
        &\longrightarrow
        \\
        &if\_then\_else(
            \\ & \hspace{1cm} uniform(context\_inclusion\_check(X, Nodes)),
            \\ & \hspace{1cm} true,
            \\ & \hspace{1cm} is\_included\_in\_fc(
                \\ & \hspace{2cm} Nodes,
                \\ & \hspace{2cm} Contexts\_occs
            \\ & \hspace{1cm})
        \\ & )
        \\
        & is\_included\_in\_fc(
            \\ & \hspace{1cm} Nodes,
            \\ & \hspace{1cm} @
        \\ & )
        \\
        &\longrightarrow
        \\
        & false&
    \end{flalign*}
    
    
\subsection*{choice}
 Arbitrarily select a rewrite step.

    \begin{flalign*}
        \\&choice(
            \\ & \hspace{1cm} L * X_1 * X_2
        \\ &)
        \\
        &\longrightarrow
        \\
        &choice(
            \\ & \hspace{1cm} L* X_1,
        \\ &)
    \end{flalign*}

    % If there is only one element, select that element:

    \begin{flalign*}
        \\&choice(
            \\ & \hspace{1cm} @ * X
        \\ &)
        \\
        &\longrightarrow
        \\
        & some(X)
    \end{flalign*}

    % An element can be selected at any moment:

    \begin{flalign*}
        \\&choice(
            \\ & \hspace{1cm} L * X
        \\ &)
        \\
        &\longrightarrow
        \\
        & some(X)
    \end{flalign*}
    % If the list is empty:

    \begin{flalign*}
        \\&choice(
            \\ & \hspace{1cm} @
        \\ &)
        \\
        &\longrightarrow
        \\
        & none
    \end{flalign*}


     
    
     

\subsection*{construct\_occs\_dfs}
Let $R$ be a term encoding a natural number between $1$ and $|\mathcal{R}|$, $C$ a natural number between $0$ and $|f(r_R)|$ the number of forbidden context of the rule $r_R$. Let $s(K)$ be the number of nodes in the host graph, $Y$ term encoding list of nodes in the host graph, $F$ a term encoding the labeling function of the host graph. This module rewrites $construct\_occs\_dfs(R,C,s(K),@, Y,F,@)$ the a term encoding the list of occurrences of the graph indicated by $(R,C)$ in the host graph.

    \begin{flalign*}
        \hspace{1cm}
        &construct\_occ\_dfs(
        \\ & \hspace{1cm} R, 
        \\ & \hspace{1cm} C,
        \\ & \hspace{1cm} nb\_nodes\_to\_choose(s(K)),
        \\ & \hspace{1cm} nodes\_chosen(X), 
        \\ & \hspace{1cm} nodes\_to\_consider(Y*V),
        \\ & \hspace{1cm} F,
        \\ & \hspace{1cm} Occs
        \\ & )
        \\
        &\longrightarrow
        \\ & union(
            \\ & \hspace{1cm} construct\_occ\_dfs(
                \\ & \hspace{2cm} R,
                \\ & \hspace{2cm} C,
                \\ & \hspace{2cm} nb\_nodes\_to\_choose( s(K)),
                \\ & \hspace{2cm} nodes\_chosen(X), 
                \\ & \hspace{2cm} nodes\_to\_consider(Y), 
                \\ & \hspace{2cm} F,
                \\ & \hspace{2cm} Occs
            \\ & \hspace{1cm}),
            \\ & \hspace{1cm} if\_then\_else(
                \\ & \hspace{2cm} uniform(mem(V, X)),
                \\ & \hspace{2cm} @, 
                \\ & \hspace{2cm} construct\_occ\_dfs(
                    \\ & \hspace{3cm} R,
                    \\ & \hspace{3cm} C,
                    \\ & \hspace{3cm} nb\_nodes\_to\_choose(K),
                    \\ & \hspace{3cm} nodes\_chosen(X * V), 
                    \\ & \hspace{3cm} nodes\_to\_consider(
                        \\ & \hspace{4cm} get\_all\_nodes(F)
                    \\ & \hspace{3cm})
                    \\ & \hspace{3cm}F,
                    \\ & \hspace{3cm}@
                \\ & \hspace{2cm})
        \\ & \hspace{1cm})
        \\ & )
    \end{flalign*}
    
    \begin{flalign*}
        \hspace{1cm}
        &construct\_occ\_dfs(
            \\ & \hspace{1cm} R,
            \\ & \hspace{1cm} C,
            \\ & \hspace{1cm} nb\_nodes\_to\_choose(s(K)),
            \\ & \hspace{1cm} X,
            \\ & \hspace{1cm} nodes\_to\_consider(@),
            \\ & \hspace{1cm} F,
            \\ & \hspace{1cm} Occs
            \\ & )
        \\ &\longrightarrow
        \\ & Occs
    \end{flalign*}


    \begin{flalign*}
        \hspace{1cm}
        &construct\_occ\_dfs(
            \\ & \hspace{1cm} R,
            \\ & \hspace{1cm} C,
            \\ & \hspace{1cm} nb\_nodes\_to\_choose(o),
            \\ & \hspace{1cm} nodes\_chosen(X),
            \\ & \hspace{1cm} Y,
            \\ & \hspace{1cm} F,
            \\ & \hspace{1cm} Occs,
            \\ & )
        \\ &\longrightarrow
        \\ &
        % if\_then\_else(
        %     \\ & \hspace{1cm}and(
        %         \\ & \hspace{2cm}not(has\_repetition(X)),
        %         \\ & \hspace{2cm} isomorphism\_verification(
        %             \\ & \hspace{3cm} construct\_induced\_graph(
        %                 \\ & \hspace{4cm} X,
        %                 \\ & \hspace{4cm} @,
        %                 \\ & \hspace{4cm} F
        %             \\ & \hspace{3cm}),
        %             \\ & \hspace{3cm} generate\_graph(
        %                 \\ & \hspace{4cm} R,
        %                 \\ & \hspace{4cm} C,
        %                 \\ & \hspace{4cm} X
        %             \\ & \hspace{3cm})
        %         \\ & \hspace{2cm})
        %     \\ & \hspace{1cm}),
        %     \\ & \hspace{1cm} Occs * X,
        % \\ & \hspace{1cm} Occs
        % \\ & )
        if\_then\_else(
                \\ & \hspace{1cm} uniform(isomorphism\_verification(
                    \\ & \hspace{2cm} construct\_induced\_graph(
                        \\ & \hspace{3cm} X,
                        \\ & \hspace{3cm} @,
                        \\ & \hspace{3cm} F
                    \\ & \hspace{2cm}),
                    \\ & \hspace{2cm} generate\_graph(
                        \\ & \hspace{3cm} R,
                        \\ & \hspace{3cm} C,
                        \\ & \hspace{3cm} X
                    \\ & \hspace{2cm})
                \\ & \hspace{1cm})),
            \\ & \hspace{1cm} Occs * X,
        \\ & \hspace{1cm} Occs
        \\ & )
    \end{flalign*}

    % \subsection*{construct\_rule\_occurrence\_dfs}		
    % When we have enough nodes for a rule or a context, the first node is the right-most node: $V_n * V_{n-1} * ... * V_1$. This is important when we proceed to compare *context-occurrence* and *rule-occurrence*.
    % \\
    % \\
    % \noindent While there are no enough nodes and there is a node that can be n-th node, we explore two possibilities: either we choose it as n-th node or not. (Termination: soit $nb\_node\_to\_chose$ diminue, soit $nodes\_to\_consider$ diminue en longueur)
    % \begin{flalign*}
    %     \hspace{1cm}
    %     1: \\
    %     &construct\_rule\_occurrence\_dfs(
    %     \\ & \hspace{1cm} R, 
    %     \\ & \hspace{1cm} nb\_nodes\_to\_choose(s(K)),
    %     \\ & \hspace{1cm} nodes\_chosen(X), 
    %     \\ & \hspace{1cm} nodes\_to\_consider(Y*V),
    %     \\ & \hspace{1cm} F,
    %     \\ & \hspace{1cm} Fcs,
    %     \\ & \hspace{1cm} Poccs,
    %     \\ & \hspace{1cm} Occs
    %     \\ & )
    %     \\
    %     &\longrightarrow
    %     \\
    %     &construct\_rule\_occurrence\_dfs(
    %         \\ & \hspace{1cm} R,
    %         \\ & \hspace{1cm} nb\_nodes\_to\_choose( s(K)),
    %         \\ & \hspace{1cm} nodes\_chosen(X), 
    %         \\ & \hspace{1cm} nodes\_to\_consider(Y), 
    %         \\ & \hspace{1cm} F,
    %         \\ & \hspace{1cm} Fcs,
    %         \\ & \hspace{1cm} Poccs,
    %         \\ & \hspace{1cm} construct\_rule\_occurrence\_dfs(
    %             \\ & \hspace{2cm} R,
    %             \\ & \hspace{2cm} nb\_nodes\_to\_choose(K),
    %             \\ & \hspace{2cm} nodes\_chosen(X * V), 
    %             \\ & \hspace{2cm} nodes\_to\_consider(
    %                 \\ & \hspace{3cm} get\_all\_nodes(F)
    %             \\ & \hspace{2cm})
    %             \\ & \hspace{2cm} F,
    %             \\ & \hspace{2cm} Fcs,
    %             \\ & \hspace{2cm} Poccs,
    %             \\ & \hspace{2cm} Occs
    %         \\ & \hspace{1cm}),
    %     \\ &)
    % \end{flalign*}
    
    % \begin{flalign*}
    %     2: \hspace{1cm}
    %     &construct\_rule\_occurrence\_dfs(
    %         \\ & \hspace{1cm} R,
    %         \\ & \hspace{1cm} nb\_nodes\_to\_choose(s(K)),
    %         \\ & \hspace{1cm} X,
    %         \\ & \hspace{1cm} nodes\_to\_consider(@),
    %         \\ & \hspace{1cm} F,
    %         \\ & \hspace{1cm} Fcs,
    %         \\ & \hspace{1cm} Poccs,
    %         \\ & \hspace{1cm} Occs
    %         \\ & )
    %     \\ &\longrightarrow
    %     \\ & Occs
    % \end{flalign*}

    % \begin{flalign*}
    %     3: \hspace{1cm}
    %     &construct\_rule\_occurrence\_dfs(
    %         \\ & \hspace{1cm} R,
    %         \\ & \hspace{1cm} nb\_nodes\_to\_choose(o),
    %         \\ & \hspace{1cm} nodes\_chosen(X),
    %         \\ & \hspace{1cm} Y,
    %         \\ & \hspace{1cm} F,
    %         \\ & \hspace{1cm} Fcs,
    %         \\ & \hspace{1cm} Poccs,
    %         \\ & \hspace{1cm} Occs
    %     \\ &)
    %     \\  &\longrightarrow
    %     \\ &
    %     if\_then\_else(
    %         \\ & \hspace{1cm}and (
    %             \\ & \hspace{2cm}and(
    %                 \\ & \hspace{3cm} not(has\_repetition(X)),
    %                 \\ & \hspace{3cm} isomorphism\_verification(
    %                     \\ & \hspace{4cm} construct\_induced\_graph(X,@, F),
    %                     \\ & \hspace{4cm} generate\_graph(R,o,X)
    %                 \\ & \hspace{3cm})
    %                 \\ & \hspace{2cm}),
    %             \\ & \hspace{2cm} and(
    %                 \\ & \hspace{3cm} not(
    %                     \\ & \hspace{4cm} is\_included\_in\_fc(
    %                         \\ & \hspace{5cm} R,
    %                         \\ & \hspace{5cm} X,
    %                         \\ & \hspace{5cm} Fcs,
    %                         \\ & \hspace{5cm} @
    %                     \\ & \hspace{4cm})
    %                 \\ & \hspace{3cm}),
    %                 \\ & \hspace{3cm} not(
    %                     \\ & \hspace{4cm} is\_dominated(
    %                         \\ & \hspace{5cm} R,
    %                         \\ & \hspace{5cm} X,
    %                         \\ & \hspace{5cm} @,
    %                         \\ & \hspace{5cm} @, 
    %                         \\ & \hspace{5cm} Poccs
    %                     \\ & \hspace{4cm})
    %                 \\ & \hspace{3cm})
    %             \\ & \hspace{2cm})
    %     \\ & \hspace{1cm}),
    %         \\ & \hspace{1cm} Occs * X,
    %     \\ & \hspace{1cm}  Occs
    %     \\ & )
    % \end{flalign*}

    
    % \subsection*{context\_rule\_enumeration}

    % If there is a context or a rule, we proceed to find its occurrences: todotodo continue {bar @ * bar @}
    % \begin{flalign*}
    %     \hspace{1cm}
    %     1\\
    %     & context\_rule\_enumeration( 
    %         \\ & \hspace{1cm} L * graph(R, C, K),
    %         \\ & \hspace{1cm} F,
    %         \\ & \hspace{1cm} Occs
    %     \\ & )
    %     \\
    %     &\longrightarrow
    %     \\
    %     & context\_rule\_enumeration( 
    %         \\ & \hspace{1cm} L,
    %         \\ & \hspace{1cm} F,
    %         \\ & \hspace{1cm} construct\_occurrence\_dfs(
    %             \\ & \hspace{2cm} R, 
    %             \\ & \hspace{2cm} C,
    %             \\ & \hspace{2cm} nb\_nodes\_to\_choose(K), 
    %             \\ & \hspace{2cm} nodes\_chosen(@),
    %             \\ & \hspace{2cm} nodes\_to\_consider(
    %                 \\ & \hspace{3cm} get\_all\_nodes(F)
    %             \\ & \hspace{2cm}),
    %             \\ & \hspace{2cm} F,
    %             \\ & \hspace{2cm} Occs
    %         \\ & \hspace{1cm})
    %     \\ & )
    % \end{flalign*}

    % \noindent
    % \begin{flalign*}
    %     \hspace{1cm}
    %     2\\
    %     & context\_rule\_enumeration( 
    %         \\ & \hspace{1cm} @,
    %         \\ & \hspace{1cm} F,
    %         \\ & \hspace{1cm} Cos * Ros
    %     \\ & )
    %     \\
    %     &\longrightarrow
    %     \\
    %     & Ros
    % \end{flalign*}

    % \subsection*{get\_poccs}
    % \begin{flalign*}
    %     \hspace{1cm}
    %     1\\
    %     & get\_poccs( 
    %         \\ & \hspace{1cm} R,
    %         \\ & \hspace{1cm} Occs,
    %         \\ & \hspace{1cm} none, 
    %         \\ & \hspace{1cm} none 
    %     \\ & )
    %     \\
    %     &\longrightarrow
    %     \\
    %     & get\_poccs( 
    %         \\ & \hspace{1cm} R,
    %         \\ & \hspace{1cm} Occs,
    %         \\ & \hspace{1cm} get\_rules\_with\_priority(R), 
    %         \\ & \hspace{1cm} @ 
    %     \\ & )
    % \end{flalign*}

    % \begin{flalign*}
    %     \hspace{1cm}
    %     2\\
    %     & get\_poccs( 
    %         \\ & \hspace{1cm} R,
    %         \\ & \hspace{1cm} Occs,
    %         \\ & \hspace{1cm} Rules * Rule, 
    %         \\ & \hspace{1cm} Poccs
    %     \\ & )
    %     \\
    %     &\longrightarrow
    %     \\
    %     & get\_poccs( 
    %         \\ & \hspace{1cm} R,
    %         \\ & \hspace{1cm} Occs,
    %         \\ & \hspace{1cm} Rules, 
    %         \\ & \hspace{1cm} union(Poccs, nth(Rule, Occs)) 
    %     \\ & )
    % \end{flalign*}

    % \begin{flalign*}
    %     \hspace{1cm}
    %     2\\
    %     & get\_poccs( 
    %         \\ & \hspace{1cm} R,
    %         \\ & \hspace{1cm} Occs,
    %         \\ & \hspace{1cm} @, 
    %         \\ & \hspace{1cm} Poccs
    %     \\ & )
    %     \\
    %     &\longrightarrow
    %     \\
    %     & Poccs
    % \end{flalign*}


\subsection*{cal\_occs\_not\_included\_in\_fc}
Let $X$ be a term encoding a list of occurrences of the left-hand-side graph of a rule in a host graph. Let $Y$ be a term encoding the list of occurrences of the forbidden contexts in the host graph. This module rewrites the term $cal\_occs\_not\_included\_in\_fc(X,Y)$ to a term encoding the list of occurrences in $Y$ which is allowed.
\begin{flalign*}
& cal\_occs\_not\_included\_in\_fc(
        \\ & \hspace{1cm} Occs * X,
        \\ & \hspace{1cm} Contexts\_occs
    \\ &)
    \\ & \longrightarrow \\
    \\ & if\_then\_else(
        \\ & \hspace{1cm} uniform(is\_included\_in\_fc(X, Contexts\_occs)),
        \\ & \hspace{1cm} cal\_occs\_not\_included\_in\_fc(
            \\ & \hspace{2cm} Occs,
            \\ & \hspace{2cm} Contexts\_occs
        \\ & \hspace{1cm}),
        \\ & \hspace{1cm} cal\_occs\_not\_included\_in\_fc(
            \\ & \hspace{2cm} Occs,
            \\ & \hspace{2cm} Contexts\_occs
            \\ & \hspace{1cm}) * X
    \\ & )
\end{flalign*}
\begin{flalign*}
    & cal\_occs\_not\_included\_in\_fc(
        \\ & \hspace{1cm} @,
        \\ & \hspace{1cm} Contexts\_occs
    \\ &)
    \\ & \longrightarrow \\
    \\ & @
\end{flalign*}


\subsection*{identify\_applicable\_occs\_of\_rule}
    Let $F$ be a term encoding the relabeling function of a graph.
    Let $X$ be a term encoding the list of informations of all left-hand-side graphs and codomains of all forbidden contexts.
    This module rewrites $identify\_applicable\_occs\_of\_rule(F,X,@,@)$ to 
    $identify\_applicable\_occs\_of\_rule(F,@,Y,Z)$ where $Y$ is a term encoding a list of lists of occurrences of left-hand-side graphs in the host graph, and $Z$ a term encoding all occurrences of forbidden contexts of all rules in the host graph.
    \begin{flalign*}
        \hspace{1cm}
        & identify\_applicable\_occs\_of\_rule( 
            \\ & \hspace{1cm} F,
            \\ & \hspace{1cm} Graphs * graph(R,s(C),K),
            \\ & \hspace{1cm} Context\_occs,
            \\ & \hspace{1cm} Rule\_occs,
        \\ & )
        \\
        &\longrightarrow
        \\
        & identify\_applicable\_occs\_of\_rule( 
            \\ & \hspace{1cm} F,
            \\ & \hspace{1cm} Graphs,
            \\ & \hspace{1cm} union(
                \\ & \hspace{2cm} Context\_occs, 
                \\ & \hspace{2cm} construct\_occs\_dfs(
                    \\ & \hspace{3cm} R,
                    \\ & \hspace{3cm} s(C),
                    \\ & \hspace{3cm} nb\_nodes\_to\_choose(K),
                    \\ & \hspace{3cm} nodes\_chosen(@),
                    \\ & \hspace{3cm} nodes\_to\_consider(
                        \\ & \hspace{4cm} get\_all\_nodes(F)
                    \\ & \hspace{3cm} ),
                    \\ & \hspace{3cm} F,
                    \\ & \hspace{3cm} @
                \\ & \hspace{2cm} )
            \\ & \hspace{1cm} ),
            \\ & \hspace{1cm} Rule\_occs
        \\ & )
    \end{flalign*}
    \begin{flalign*}
        \hspace{1cm}
        & identify\_applicable\_occs\_of\_rule( 
            \\ & \hspace{1cm} F
            \\ & \hspace{1cm} @ * graph(R,o,K),
            \\ & \hspace{1cm} Context\_occs,
            \\ & \hspace{1cm} Rule\_occs
        \\ & )
        \\
        &\longrightarrow
        \\ & identify\_applicable\_occs\_of\_rule(
            \\ & \hspace{1cm} F
            \\ & \hspace{1cm} @,
            \\ & \hspace{1cm} Context\_occs,
            \\ & \hspace{1cm} construct\_occs\_dfs(
                \\ & \hspace{2cm} R,
                \\ & \hspace{2cm} o,
                \\ & \hspace{2cm} nb\_nodes\_to\_choose(K),
                \\ & \hspace{2cm} nodes\_chosen(@),
                \\ & \hspace{2cm} nodes\_to\_consider(
                    \\ & \hspace{3cm} get\_all\_nodes(F)
                \\ & \hspace{2cm} ),
                \\ & \hspace{2cm} F,
                \\ & \hspace{2cm} @
            \\ & \hspace{1cm})
        \\ &)
    \end{flalign*}

    \begin{flalign*}
        & identify\_applicable\_occs\_of\_rule(
            \\ & \hspace{1cm} F,
            \\ & \hspace{1cm} @, 
            \\ & \hspace{1cm} Context\_occs,
            \\ & \hspace{1cm} Rule\_occs
        \\ & )
        \\ & \longrightarrow \\
        \\ & choice(\\ & \hspace{1cm}cal\_occs\_not\_included\_in\_fc(
                \\ & \hspace{2cm} Rule\_occs,
                \\ & \hspace{2cm} Context\_occs
            \\ & \hspace{1cm})
        \\ & )
    \end{flalign*}
    
    The translation into an innermost HTRS without an AC symbol is interesting because it allows us to apply termination methods designed for TRS or Innermost TRS. Additionally, we can apply termination methods designed for hierachical TRS. 

\begin{theorem}
    \label{thm:fcgls_preservation_termination}
    If $(G,\Lambda) \to_\mathcal{R} (G,\Lambda')$ then $\tau_{fcgls}(G,\Lambda) \to_{\mathcal{R}'} \tau_{fcgls}(G,\Lambda')$.
\end{theorem}
\begin{proof}
    to do
\end{proof}

\begin{corollary}
    If  $\mathcal{R}'$ terminates then so does $\mathcal{R}$.
\end{corollary}
\begin{proof}
    by~\autoref{thm:fcgls_preservation_termination}.
\end{proof}

\part{Termination of graph relabeling systems via morphisms to graph rewriting relations}

\chapter{Introduction}


\chapter{Preliminaries} 
    \label{sec:grs}

    \section{Hypergraph}
        \label{sec:hypergraph}
        We brively recall the definitions of hypergraphs from~\cite{plump2018modular}.
Intuitively, a hypergraph is a mathematical structure consisting of a set of nodes and a set of hyperedges, where nodes and hyperedges are labelled. Each hyperedge can connect any number of nodes which form a finite sequences of nodes. The label of a hyperedge decides the number of nodes it connects and the labels of the nodes.
\begin{definition}[Signature of a hypergraph~\cite{plump2018modular}]
    A signature $\Sigma = \langle \Sigma_V, \Sigma_E, \text{Type} \rangle$ consists of 
    \begin{itemize}
        \item a set $\Sigma_V$ of node labels,
        \item a set $\Sigma_E$ of hyperedge labels, and
        \item a function $\text{Type}$ assigning to each $l \in \Sigma_E$ a set of strings $\text{Type}(l) \subseteq \Sigma_V^*$.
    \end{itemize}
\end{definition}
Unless stated otherwise, we denote by $\Sigma$ an arbitrary but fixed signature over which all hypergraphs are labelled.
\begin{remark}
    The extension $f^* : A^* \to B^*$ of a function $f : A \to B$ maps the empty string to itself and $a_1 \ldots a_n$ to $f(a_1) \ldots f(a_n)$.
\end{remark}
\begin{definition}[Hypergraph~\cite{plump2018modular}]
    A hypergraph over $\Sigma$ is a system $G = \langle V_G, E_G, \text{mark}_G, \text{lab}_G, \text{att}_G \rangle$ consisting of 
    \begin{itemize}
        \item a finite set $V_G$ of nodes,
        \item a finite set $E_G$ of hyperedges,
        \item two labelling functions $\text{mark}_G : V_G \to \Sigma_V$ and $\text{lab}_G : E_G \to \Sigma_E$, and
        \item an attachment function $\text{att}_G : E_G \to V_G^*$ such that $\text{mark}_G^*(\text{att}_G(e)) \in \text{Type}(\text{lab}_G(e))$ for each hyperedge $e$.
    \end{itemize} 
\end{definition} 

\begin{example}
    We adopt the notation from~\cite[\textsection 2.1]{plump2018modular} to visualize hypergraphs. \enquote{nodes and hyperedges are drawn as circles and boxes, respectively, with labels inside. Lines represent the attachment of hyperedges to nodes, where numbers specify the left-to-right order in the attachment string.} The following shows a hypergraph with three nodes all labeled with $\epsilon$ and three hyperedges (labeled with $s$ and $0$).

    \begin{tikzpicture}[scale=0.75]
        \node[draw,circle] (n1) at (-3,3) {1};
        \node[draw,circle] (n2) at (3,3) {2};
        \node[draw,circle] (n3) at (0,0) {3};
        \node[draw,rectangle] (e1) at (-1.5,1.5) {$s$};
        \node[draw,rectangle] (e2) at (1.5,1.5) {$s$};
        \node[draw,rectangle] (e3) at (0, -1.5) {0};
        \draw (e1) -- (n1) node[midway,left] {1};
        \draw (e1) -- (n3) node[midway,right] {2};
        \draw (e2) -- (n2) node[midway,left] {1};
        \draw (e2) -- (n3) node[midway,right] {2};
        \draw (e3) -- (n3) node[midway,left] {1};
    \end{tikzpicture}
\end{example}

\begin{definition}[Hypergraph homomorphism~\cite{plump2018modular}]
    Give hypergraphs $G$ and $H$, a hypergraph homomorphism $h : G \to H$ consists of functions $h_V : V_G \to V_H$ and $h_E : E_G \to E_H$ that preserve labels and attachment to nodes, i.e., $\opn{mark}_H \circ h_V = \opn{mark}_G$ and $\opn{lab}_H \circ h_E = \opn{lab}_G$ and $\opn{att}_H \circ h_E = g_V \circ \opn{att}_G$.
\end{definition}
    
    \section{Pushout and pullback}
        \label{sec:category_theory}
        \begin{definition}[Category~\cite{pierce1991basic,barr1990category}]
    \label{def:cat}
    A \textbf{category} is an unlabeled graph \( C \) together with a total function \( u : C_0 \to C_1 \) and a partial function \( \star: C_1 \times C_1 \to C_1 \) such that 
        (i) for all edges \( f:X \to Y \) and \( g:Y \to Z \), the edge \( f \star g :X \to Z \) is defined; 
        (ii) for every node \( X \), \( u(X) \) is an edge from \( X \) to \( X \);
        (iii) for every \( f:X \to Y \), we have \(u(X) \star f = f = f \star u(Y)\);
        (iv) for all edges \( f \), \( g \) and \(h\), we have \( (f \star g) \star h = f \star (g \star h) \) whenever either side is defined.
    Edges are called \textbf{morphisms}. The function $\star$ is called \textbf{composition}. For all \( X \in C_0 \), the edge \( u(X) \) is denoted \( \operatorname{id}_X \) and is called the \textbf{identity} of the object \( X \).
    % \( C \) is called the \textbf{underlying graph} of the category \( \mathcal{C} \).
\end{definition} 

% cat notation * 
\begin{notation}
    The composition of morphisms \( f : X \to Y \) and \( g : Y \to Z \) is written in diagrammatic order as \( f \star g \), rather than in functional order \( g \circ f \). 
    % The advantage is that, when reading from left to right, the morphisms appear in the same order as in the corresponding diagram, making the notation more intuitive for visual reasoning.
\end{notation}  

\begin{definition}[Monomorphism~\cite{pierce1991basic,barr1990category}]
    \label{def:cat:homo}
    A morphism \( f : X \to Y \) is \textbf{monic} (or a \textbf{monomorphism}) if for all morphisms \( g \) and \( h \), if \( g \star f = h \star f \), then \( g = h \). A monomorphism is denoted by \( f : X \rightarrowtail Y \). The set of all monomorphisms from graph \( X \) to graph \( Y \) is denoted by \( \operatorname{Mono}(X, Y) \).
\end{definition} 
\begin{example}
     Finite edge-labeled directed multigraphs and their homomorphisms form a category, hereafter denoted \textbf{Graph}. Its objects are labeled graphs, its morphisms are graph homomorphisms, and the monomorphisms are injective homomorphisms.
\end{example}
\begin{definition}[Category of Rewriting Systems]
    Rewriting systems and homomorphisms between them form a category, denoted $\mathcal{RS}$.   
  \end{definition}
In category theory, a ordered pair \((\alpha : A \to B,\, \beta : A \to C)\) of morphisms with a common domain is called a \textbf{span} \cite{lowe2010graph}, denoted by
\(
B \overset{\alpha}{\leftarrow} A \overset{\beta}{\rightarrow} C.
\)
Likewise, an ordered pair \((\beta' : B \to D,\, \alpha' : C \to D)\) of morphisms with a common codomain is called a \textbf{cospan}, denoted by
\(
B \overset{\beta'}{\rightarrow} D \overset{\alpha'}{\leftarrow} C.
\)

\begin{definition}[Diagram \cite{barr1990category}]
    \label{def:cat:diagram}
    Let \( G \) be an unlabeled graph. A \textbf{diagram} (in \( \mathcal{C} \) of shape \( G \)) is a homomorphism of unlabeled graphs $ h : G \to C $ where \( C \) is the underlying unlabeled graph of the category \( \mathcal{C} \). A diagram is \textbf{commutative} if, for all nodes \( u \), \( v \), and any two paths from \( u \) to \( v \) in the unlabeled graph \( G \):

    \begin{center}
    \resizebox{12cm}{!}{
        \begin{tikzpicture}
        \node (u) at (0,0) {\( u \)};
        \node (k1) at (2,0.5) {\( k_1 \)};
        \node (k2) at (4,0.5) {\( k_2 \)};
        \node (ketc) at (6,0.5) {\( \dots \)};
        \node (knm2) at (8,0.5) {\( k_{n-2} \)};
        \node (knm1) at (10,0.5) {\( k_{n-1} \)};
        \node (v) at (12,0) {\( v \)};
        \node (l1) at (2,-0.5) {\( l_1 \)};
        \node (l2) at (4,-0.5) {\( l_2 \)};
        \node (letc) at (6,-0.5) {\( \dots \)};
        \node (lnm2) at (8,-0.5) {\( l_{m-2} \)};
        \node (lnm1) at (10,-0.5) {\( l_{m-1} \)};
        \draw[->] (u) -- (k1) node [midway,above] {\( s_1 \)};
        \draw[->] (k1) -- (k2) node [midway,above] {\( s_2 \)};
        \draw[->] (k2) -- (ketc);
        \draw[->] (ketc) -- (knm2); 
        \draw[->] (knm2) -- (knm1) node[midway,above] {\( s_{n-2} \)}; 
        \draw[->] (knm1) -- (v) node[midway,above] {\( s_{n-1} \)}; 
        \draw[->] (u) -- (l1) node[midway,below] {\( t_1 \)};
        \draw[->] (l1) -- (l2) node[midway,below] {\( t_2 \)};
        \draw[->] (l2) -- (letc);
        \draw[->] (letc) -- (lnm2); 
        \draw[->] (lnm2) -- (lnm1) node[midway,below] {\( t_{m-2} \)}; 
        \draw[->] (lnm1) -- (v) node[midway,below] {\( t_{m-1} \)}; 
        \end{tikzpicture}
    }
    \end{center}
    \noindent
    the equality \( h(s_1) \star h(s_2) \star \dots  \star h(s_{n-1}) = h(t_1) \star h(t_2) \star \dots  \star h(t_{m-1}) \) holds.
\end{definition}

\begin{definition}[Pushout \cite{barr1990category}]
    \label{def:cat:po}
    \ \newline
\noindent
\begin{minipage}{0.7\textwidth}  
    A \textbf{pushout} of a span \( B \overset{\alpha}{\leftarrow} A \overset{\beta}{\rightarrow} C \), as shown on the right, is defined as a cospan \( B \overset{\beta'}{\rightarrow} D \overset{\alpha'}{\leftarrow} C \) such that \( \alpha \star \beta' = \beta \star \alpha' \), and for every cospan \( B \overset{\gamma'}{\rightarrow} E \overset{\gamma}{\leftarrow} C \), if \( \alpha \star \gamma' = \beta \star \gamma \) holds, then there exists a unique morphism \(\delta : D \to E\) such that \( \gamma' = \beta' \star \delta \) and \( \gamma = \alpha' \star \delta \).
\end{minipage}
\hfill
\begin{minipage}{0.299\textwidth}
    \hfill
\resizebox{0.9\textwidth}{!}{
           \begin{tikzpicture}
                \node (i) at (0,0) {A};
                \node (r) at (1,1) {B};
                \node (c) at (1,-1) {C};
                \node (h) at (2,0) {D};
                % \node () at (1,-1) {\( \Delta \)};
                \draw[->]  (i) -- (r) node [midway,left] {$ \alpha $};
                \draw[->] (c) -- (h) node [midway,left] {$ \alpha' $};
                \draw[->] (r) -- (h) node[midway, left] {$ \beta' $};
                \draw[->] (i) -- (c) node[midway, left] {$ \beta $};
                \node (d') at (4,0) {E};
                \draw[->] (c) -- (d') node [midway,below]{$ \gamma $};
                \draw[->] (r) -- (d') node [midway,above]{$ \gamma' $};
                \draw[->,dashed] (h) -- (d') node [midway]{$ \delta $};
            \end{tikzpicture}
}
\end{minipage}
The diagram involving \( (\alpha, \beta, \alpha', \beta') \) is called the \textbf{pushout square}, with \(D\) as the \textbf{pushout object}. The existence of a unique morphism is known as the universal mapping property of the pushout.
\end{definition} 
\begin{example}
    \label{ex:cat:po}
    Consider the square in \textbf{Graph} illustrated below.
    A pushout of the span \( B \overset{\alpha}{\leftarrowtail} A \overset{\beta}{\rightarrowtail} C \) is the cospan \( C \overset{\alpha'}{ 
    \rightarrowtail} D \overset{\beta'}{\leftarrowtail} B \), where pushout object \( D \) is the graph union of $B$ and $C$ and the morphisms \( \alpha' \) and \( \beta' \) are inclusion functions. 

    \begin{center} 
        \resizebox{0.45\textwidth}{!}{
        \begin{tikzpicture} 
            \graphbox{\( B \)}{40mm}{20mm}{34mm}{12mm}{2mm}{2mm}{
                \coordinate (o) at (0mm,-8mm); 
                \node[draw,circle] (l1) at ($(o)+(-10mm,0mm)$) {1};
                \node[draw,circle] (l2) at ($(l1)+(2,0)$) {2};
                \node[draw,circle,red] (l3) at ($(l1) + (1,0)$) {3};
                \draw[->,red] (l1) -- (l3) node[midway,above] {$a$};
                \draw[->,red] (l3) -- (l2) node[midway,above] {$a$};
            } 
    
            \graphbox{\( A \)}{0mm}{0mm}{34mm}{12mm}{2mm}{2mm}{
                \coordinate (o) at (0mm,-8mm); 
                \node[draw,circle] (l1) at ($(o)+(-10mm,0mm)$) {1};
                \node[draw,circle] (l2) at ($(l1)+(2,0)$) {2};
            }  
            \graphbox{\( D \)}{90mm}{5mm}{34mm}{22mm}{2mm}{-3mm}{
                \coordinate (o) at (0mm,-3mm); 
                \node[draw,circle] (l1) at ($(o)+(-10mm,0mm)$) {1};
                \node[draw,circle] (l2) at ($(l1)+(2,0)$) {2};
                \node[draw,circle,red] (l3) at ($(l1) + (1,0)$) {3};
                \node[draw,circle,blue] (l4) at ($(l2) + (0,-1)$) {6};
                \draw[->,red] (l1) -- (l3) node[midway,above] {$a$};
                \draw[->,red] (l3) -- (l2) node[midway,above] {$a$};
                \draw[->,blue] (l2) -- (l4) node[midway,right] {$a$};
                \node[draw,circle,blue] (l6) at ($(l1) + (0,-1)$) {7};
                \draw[<-,blue] (l1) -- (l6) node[midway,left] {$a$};
                \draw[->,blue] (l2) edge[out=-135,in=-45]node[midway,below] {$a$} (l1) ;
            }    
     
            \graphbox{\( C  \)}{40mm}{-20mm}{34mm}{22mm}{2mm}{-3mm}{
                \coordinate (o) at (0mm,-3mm); 
                \node[draw,circle] (l1) at ($(o)+(-10mm,0mm)$) {1};
                \node[draw,circle] (l2) at ($(l1)+(2,0)$) {2};
                \node[draw,circle,blue] (l4) at ($(l2) + (0,-1)$) {6};
                \draw[->,blue] (l2) -- (l4) node[midway,right] {$a$};
                \draw[->,blue] (l2) edge[out=-135,in=-45]node[midway,below] {$a$} (l1) ;
                \node[ draw,circle,blue] (l6) at ($(l1) + (0,-1)$) {7};
                \draw[<-,blue] (l1) -- (l6) node[midway,left] {$a$};
            }      
            % K to L
            \draw[>->] (17mm,5mm) -- node[above] {$\alpha$} (37mm,15mm);
            % C to G
            \draw[>->] (76mm,-28mm)-- node[below] {$\alpha'$} (104mm,-20mm) ;
            % K to C
            \draw[>->] (17mm,-17mm) -- node[below] {$\beta$} (37mm,-28mm);
            % L to G
            \draw[>->] (76mm,16mm) -- node[above] {$\beta'$} (104mm,7mm);
            \node () at (57mm,-6mm) {$PO$};
        \end{tikzpicture}
        }
    \end{center}
\end{example}
\begin{definition}[Pullback \cite{pierce1991basic}]
    \label{def:cat:pb}
    \ \newline
\noindent
\begin{minipage}{0.7\textwidth}  
   A \textbf{pullback} of a cospan \(B \overset{\beta'}{\rightarrow} D \overset{\alpha'}{\leftarrow} C \), as shown on the right, is defined as a span \( B \overset{\alpha}{\leftarrow} A \overset{\beta}{\rightarrow} C \) such that \( \alpha \star \beta' = \beta \star \alpha' \), and for every span \( B \overset{\gamma'}{\leftarrow} E \overset{\gamma}{\rightarrow} C \) if \(\gamma' \star \beta' = \gamma \star \alpha'\) holds, then there exists a unique morphism \(\delta: E \to A\) such that $\gamma' = \delta \star \alpha$ and $\gamma = \delta \star \beta$. 
\end{minipage}
\hfill
\begin{minipage}{0.299\textwidth}
    \hfill
\resizebox{0.9\textwidth}{!}{
            \begin{tikzpicture}
                \node (i) at (0,0) {A};
                \node (r) at (1,1) {B};
                \node (c) at (1,-1) {C};
                \node (h) at (2,0) {D};
                % \node () at (1,-1) {\( \Delta \)};
                \draw[->]  (i) -- (r) node [midway,left] {$\alpha$};
                \draw[->] (c) -- (h) node [midway,left] {$\alpha'$};
                \draw[->] (r) -- (h) node[midway, left] {$\beta'$};
                \draw[->] (i) -- (c) node[midway, left] {$\beta$};
                \node (d') at (-2,0) {E};
                \draw[<-] (c) -- (d') node [midway,below]{$\gamma$};
                \draw[<-] (r) -- (d') node [midway,above]{$\gamma'$};
                \draw[->, dashed] (d') -- (i) node [midway]{$\delta$};
            \end{tikzpicture}
}
\end{minipage}
The diagram involving \( (\alpha, \beta, \alpha', \beta') \) is called the \textbf{pullback square}, with \(A\) as the \textbf{pullback object}. The existence of a unique morphism is known as the universal mapping property of the pullback.
\end{definition} 
 
\begin{definition}
    \label{def:cat:pb}
    Consider the square in \autoref{ex:cat:po}. A pullback of the cospan \( C \overset{\alpha'}{\rightarrowtail} D \overset{\beta'}{\leftarrowtail} B \) is the span \( B \overset{\alpha}{\leftarrowtail} A \overset{\beta}{\rightarrowtail} C \), where pullback object \( A \) is the graph intersection of $B$ and $C$ and the morphisms \( \alpha \) and \( \beta \) are inclusion functions.
\end{definition}
\begin{notation}
    When the context makes it clear, a morphism \( h : A \to B \) will be denoted by \( h_{AB} \), and diagrams will be referred to by their nodes, as is standard in geometry. For example, the diagram involving morphisms \( ( \alpha, \beta, \alpha', \beta' ) \) in \autoref{def:cat:po} will be denoted by \( ACDB \) or \( ABDC \).
\end{notation}   

    \section{Algebraic Rewriting Systems} 
    \subsection{DPO Rewriting Systems}

    \begin{itemize}
      \item the most dominant graph rewriting method, called the Double Pushout (DPO) approach.
      \item It can  specify the specification of rewrite steps with deletion, identification and addition features, for matches m that satisfy the gluing condition.
      \item Alternatives to the DPO approach avoid the construction of pushout complements. For instance, the Single Pushout (SPO) approach [17] relies on a single pushout construction, but uses partial graph homomorphisms instead of total morphisms, in order to specify deletion. In this approach, the gluing condition no longer needs to be checked either: all edges incident to a removed vertex are simply deleted.
      \item As another example, the Sesqui-Pushout (corradini2006sesqui) approach [8] replaces the first PO square of DPO by what is called a final pullback complement square. This square allows duplication with deterministic behavior, and like SPO, deletes any edges that would be left dangling.
      \item !!! implicite deletion of nodes or edges is a problem when we want to prove the termination 
      \item In DPO, we are only able to delete nodes if they do not leave any edges “dangling” (called the gluing condition). This can be considered a pleasant safety feature, but also a limitation. However, since GLS never delete or duplicate nodes, DPO GRS is a well studied approach with advanced termination method.
    \end{itemize}

    \begin{definition}[Rewriting rule, Match~\cite{ehrig1997algebraic}]
        \label{def:grs:dpo_rule}
      A \textbf{DPO rewriting rule} $\rho$ is a span \( L \overset{l}{\leftarrow} K \overset{r}{\rightarrow} R \), where       
      the object \( K \) is the \textbf{interface}, 
      the morphism \( l \) is the \textbf{left-hand-side}, denoted \( \operatorname{lhs}(\rho) \),
      the morphism \( r \) is the \textbf{right-hand-side}, denoted \( \operatorname{rhs}(\rho) \),
      the object \( L \) is the \textbf{left-hand-side object} and the object \( R \) is the \textbf{right-hand-side object}. The rule is \textbf{left-monic} if \( l \) is monic, \textbf{right-monic} if \( r \) is monic, \textbf{monic} if it is left- and right-monic.
      A \textbf{match} of the rule in a graph \( G \) is a morphism \( m: L \mathop{\rightarrow} G \).   
      \end{definition}
    
      \begin{example}
        \label{ex:grsaa}
        The injective DPO rule from \cite[Example 6]{bruggink2014termination} will be used as a running example throughout this paper to illustrate the concepts disccussed.
        \begin{center} 
            \resizebox{0.7\textwidth}{!}{
            \begin{tikzpicture}
                \graphbox{$L$}{0mm}{0mm}{34mm}{15mm}{2mm}{-5mm}{
                    \coordinate (o) at (0mm,-3mm); 
                    \node[draw,circle] (l1) at ($(o)+(-10mm,0mm)$) {1};
                    \node[draw,circle] (l2) at ($(l1)+(2,0)$) {2};
                    \node[draw,circle] (l3) at ($(l1)+(1,0)$) {3};
                    \draw[->] (l1) -- (l3) node[midway,above] {$a$};
                    \draw[->] (l3) -- (l2) node[midway,above] {$a$};
                }     
                \graphbox{$K$}{40mm}{0mm}{24mm}{15mm}{2mm}{-5mm}{
                    \coordinate (o) at (5mm,-3mm); 
                    \node[draw,circle] (l1) at ($(o)+(-10mm,0mm)$) {1};
                    \node[draw,circle] (l2) at ($(l1)+(1,0)$) {2};
                    % \node[draw,circle] (l3) at ($(l1)+(1,0)$) {$\ $};
                    % \draw[->] (l1) -- (l3) node[midway,above] {$a$};
                    % \draw[->] (l3) -- (l2) node[midway,above] {$a$};
                }    
                \graphbox{$R$}{70mm}{0mm}{45mm}{15mm}{2mm}{-5mm}{
                    \coordinate (o) at (-5mm,-3mm); 
                    \node[draw,circle] (l1) at ($(o)+(-10mm,0mm)$) {1};
                    \node[draw,circle] (l2) at ($(l1)+(3,0)$) {2};
                    \node[draw,circle] (l3) at ($(l1)+(1,0)$) {4};
                    \node[draw,circle] (l4) at ($(l1)+(2,0)$) {5};
                    \draw[->] (l1) -- (l3) node[midway,above] {$a$};
                    \draw[->] (l3) -- (l4) node[midway,above] {$b$};
                    \draw[->] (l4) -- (l2) node[midway,above] {$a$};
                }    
                \node () at (37mm,-8mm) {$\leftarrowtail$};
                \node () at (67mm,-8mm) {$\rightarrowtail$};
                % \draw[>->] (51mm,2mm) -- (52mm,3mm);
            \end{tikzpicture}
            }
        \end{center}
      \end{example}
      
      \begin{definition}[Witness, Rewriting step \cite{endrullis2024generalized}]
        \label{def:rewriting_step}
          \ \newline
          \noindent
          \begin{minipage}{0.72\textwidth}
            A DPO diagram $\delta$ is a diagram as shown on the right.
            This diagram $\delta$ is a \textbf{witness} for the \textbf{rewriting step} from \( G \) to \( H \) using the rule \( \rho \) and match \( m \), denoted \( G \mathop{\Rightarrow}_\rho^m H \) or \( G \mathop{\Rightarrow}_\rho^\delta H \). We denote $\operatorname{left}(\delta)$ and $\operatorname{right}(\delta)$ the pushout squares $KLGC$ and $KRHC$, respectively.
          \end{minipage}
          \hfill
          \begin{minipage}{0.28\textwidth}
                % \begin{center}
                \hfill
                \resizebox{\textwidth}{!}{
                \begin{tikzpicture}
                  % [node distance=11mm]
                  \node (I) {$K$};
                  \node (L) [left of=I] {$L$};
                  \node (R) [right of=I] {$R$};
                  \node (G) [below of=L] {$G$};
                  \node (C) [below of=I] {$C$};
                  \node (H) [below of=R] {$H$};
                  \draw [->] (I) to  node [midway,above] {$l$} (L);
                  \draw [->] (I) to  node [midway,above] {$r$} (R);
                  \draw [->] (L) to node [midway,left] {$m$} (G);
                  \draw [->] (I) to (C);
                  \draw [->] (R) to node [midway,right] {$m'$} (H);
                  \draw [->] (C) to node [midway,below] {$l'$} (G);
                  \draw [->] (C) to node [midway,below] {$r'$} (H);
                  \node [at=($(I)!.5!(G)$)] {\normalfont PO};
                  \node [at=($(I)!.5!(H)$)] {\normalfont PO};
                \end{tikzpicture}
              % \end{center}
              }
              \end{minipage}
        \end{definition}
    

        \begin{example}
            \label{ex:rewriting_step_grs_aa}
            The DPO diagram below defines a rewriting step using the rule from Example~\ref{ex:grsaa}.
            \begin{center} 
                \resizebox{0.7\textwidth}{!}{
                \begin{tikzpicture}
                    \graphbox{\( L \)}{0mm}{-3mm}{34mm}{12mm}{2mm}{2mm}{
                        \coordinate (o) at (0mm,-8mm); 
                        \node[draw,circle] (l1) at ($(o)+(-10mm,0mm)$) {1};
                        \node[draw,circle] (l2) at ($(l1)+(2,0)$) {2};
                        \node[draw,circle] (l3) at ($(l1)+(1,0)$) {3};
                        \draw[->] (l1) -- (l3) node[midway,above] {$a$};
                        \draw[->] (l3) -- (l2) node[midway,above] {$a$};
                    } 
            
                    \graphbox{\( K \)}{40mm}{-3mm}{34mm}{12mm}{2mm}{2mm}{
                        \coordinate (o) at (0mm,-8mm); 
                        \node[draw,circle] (l1) at ($(o)+(-10mm,0mm)$) {1};
                        \node[draw,circle] (l2) at ($(l1)+(2,0)$) {2};
                    }  
            
                    \graphbox{\( R \)}{80mm}{-3mm}{45mm}{12mm}{2mm}{2mm}{
                        \coordinate (o) at (-5mm,-8mm); 
                        \node[draw,circle] (l1) at ($(o)+(-10mm,0mm)$) {1};
                        \node[draw,circle] (l2) at ($(l1)+(3,0)$) {2};
                        \node[draw,circle] (l3) at ($(l1)+(1,0)$) {4};
                        \node[draw,circle] (l4) at ($(l1)+(2,0)$) {5};
                        \draw[->] (l1) -- (l3) node[midway,above] {$a$};
                        \draw[->] (l3) -- (l4) node[midway,above] {$b$};
                        \draw[->] (l4) -- (l2) node[midway,above] {$a$};
                    }    
            
                    \graphbox{\( G \)}{0mm}{-22mm}{34mm}{22mm}{2mm}{-3mm}{
                        \coordinate (o) at (0mm,-3mm); 
                        \node[draw,circle] (l1) at ($(o)+(-10mm,0mm)$) {1};
                        \node[draw,circle] (l2) at ($(l1)+(2,0)$) {2};
                        \node[draw,circle] (l3) at ($(l1)+(1,0)$) {3};
                        \node[draw,circle] (l4) at ($(l2)+(0,-1)$) {6};
                        \draw[->] (l1) -- (l3) node[midway,above] {$a$};
                        \draw[->] (l3) -- (l2) node[midway,above] {$a$};
                        \draw[->] (l2) -- (l4) node[midway,right] {$a$};
                        \node[draw,circle] (l6) at ($(l1)+(0,-1)$) {7};
                        \draw[<-] (l1) -- (l6) node[midway,left] {$a$};
                        \draw[->] (l2) edge[out=-135,in=-45]node[midway,below] {$a$} (l1) ;
                    }    
            
                    \graphbox{\( C  \)}{40mm}{-22mm}{34mm}{22mm}{2mm}{-3mm}{
                        \coordinate (o) at (0mm,-3mm); 
                        \node[draw,circle] (l1) at ($(o)+(-10mm,0mm)$) {1};
                        \node[draw,circle] (l2) at ($(l1)+(2,0)$) {2};
                        \node[draw,circle] (l4) at ($(l2)+(0,-1)$) {6};
                        \draw[->] (l2) -- (l4) node[midway,right] {$a$};
                        \draw[->] (l2) edge[out=-135,in=-45]node[midway,below] {$a$} (l1) ;
                        \node[ draw,circle] (l6) at ($(l1)+(0,-1)$) {7};
                        \draw[<-] (l1) -- (l6) node[midway,left] {$a$};
                    }    
            
                    \graphbox{\( H \)}{80mm}{-22mm}{45mm}{22mm}{2mm}{-3mm}{
                        \coordinate (o) at (-5mm,-3mm); 
                        \node[draw,circle] (l1) at ($(o)+(-10mm,0mm)$) {1};
                        \node[draw,circle] (l2) at ($(l1)+(3,0)$) {2};
                        \node[draw,circle] (l3) at ($(l1)+(1,0)$) {4};
                        \node[draw,circle] (l4) at ($(l1)+(2,0)$) {5};
                        \node[ draw,circle] (l5) at ($(l2)+(0,-1)$) {6};
                        \node[ draw,circle] (l6) at ($(l1)+(0,-1)$) {7};
                        \draw[<-] (l1) -- (l6) node[midway,left] {$a$};
                        \draw[->] (l1) -- (l3) node[midway,above] {$a$};
                        \draw[->] (l3) -- (l4) node[midway,above] {$b$};
                        \draw[->] (l4) -- (l2) node[midway,above] {$a$};
                        \draw[->] (l2) -- (l5) node[midway,right] {$a$};
                        \draw[->] (l2) edge[out=-135,in=-45]node[midway,below] {$a$} (l1) ;
                    }    
            
                    \node () at (37mm,-8mm) {\( \leftarrowtail \)}; % K -> L
                    \node () at (77mm,-8mm) {\( \rightarrowtail \)}; % K -> R
                    \node () at (15mm,-18mm) {\( m\ \downarrowtail \)};
                    \node () at (37mm,-33mm) {\( \leftarrowtail \)};
                    \node () at (58mm,-18mm) {\( \downarrowtail \)};
                    \node () at (102mm,-18mm) {\( \downarrowtail \)};
                    \node () at (77mm,-33mm) {\( \rightarrowtail \)}; % C -> H
                \end{tikzpicture}
                }
            \end{center}
          \end{example}
    
          When a direct graph transformation with a production p and a match m is performed, all the vertices and edges which are matched by $L \mathop{\setminus} K$ are removed from G. The removed part is not a graph, in general, but the remaining structure $D := (G \mathop{\setminus} m(L)) \mathop{\cup}  m(K)$ still has to be a graph. This means that the match $m$ has to satisfy a suitable gluing condition, which makes sure that the gluing of $L\mathop{\setminus} K$ and $D$ is equal to G. In the second step of a graph rewriting step, the graph $D$ is glued together with $R \mathop{\setminus} K$ to obtain the derived graph $H$. For gluing newly created vertices and edges into D, the graph $m(K)$ and $r(K)$ are used: nodes and edges with the same pre-image will be unified.

    \begin{definition}[Rewriting framework \cite{endrullis2024generalized}]
        A \textbf{DPO rewriting framework} $\mathfrak{F}$ is a mapping of DPO rewriting rules to classes of DPO diagrams such that, for every rule $\rho$, $\mathfrak{F}(\rho)$ is a class of DPO diagrams with top-span $\rho$.
      \end{definition}

      
      \begin{definition}[Rewriting relation]
        The \textbf{rewriting relation $\mathop{\Rightarrow}_{\mathfrak{dpo},\mathfrak{F},\rho}$ induced by a rule $\rho$ in $\mathfrak{F}$} is defined as follows: $G \mathop{\Rightarrow}_{\mathfrak{dpo},\mathfrak{F},\rho} H$ iff $G \mathop{\Rightarrow}_\rho^\delta H$ for some $\delta \mathop{\in} \mathfrak{F}(\rho)$. 
          % for some $\delta \mathop{\in} \mathfrak{F}(\rho)$
          The \textbf{DPO rewriting relation $\mathop{\Rightarrow}_{\mathcal{R},\mathfrak{F}}$ induced by a set $\mathcal{R}$ of DPO rewriting rules in $\mathfrak{F}$} is given by: $G \mathop{\Rightarrow}_{\mathfrak{dpo}, \mathfrak{F},\mathcal{R}} H$ iff $G \mathop{\Rightarrow}_{\mathfrak{dpo},\mathfrak{F}, \rho} H$ for some $\rho \mathop{\in} \mathcal{R}$. When $\mathfrak{F}$ is clear from the context, we 
          suppress $\mathfrak{F}$ and 
          write $\mathop{\Rightarrow}_{\mathfrak{dpo},\rho}$ and $\mathop{\Rightarrow}_{\mathfrak{dpo},\mathcal{R}}$.
      \end{definition} 

    % \begin{definition}[DPO Graph Rewriting System]
    %   Let $M, \mathfrak{M}, W, \mathfrak{W}, \mathfrak{I}$ be the class of DPO matches, the match mechanism, the class of DPO witnesses, the witness mechanism and the interpretation function, respectively, defined in Definition~\ref{def:grs:dpo_rule}.
    %   Let $\mathfrak{F}$ be a DPO rewriting framework as defined in Definition~\ref{def:rewriting_framework}.
    %   The rewriting system $(\mathbf{Graph}, \mathcal{R}, M, \mathfrak{M}, W, \mathfrak{W}, \mathfrak{I}, \mathfrak{F})$ is called a \textbf{DPO Graph Rewriting System in $\mathfrak{F}$}.
    % \end{definition}

    
    \subsection{DPO GRS with Negative Application Condition}
    \begin{definition}[Rewriting rule, Match~\cite{bottoni2010atermination}]
        \label{def:grs:dpo_nac_rule}
        A \textbf{DPO rewriting rule with a negative application condition} $\varphi$ is an ordered pair \( \left( n , \rho \right) \) where $n$ is a monomorphism and $\rho$ is an monic DPO rewriting rule
         such that $\operatorname{codom}(\operatorname{lhs}(\rho)) \mathop{=} \operatorname{dom}(n)$.
         
         A rule is denoted \( N \overset{n}{\leftarrowtail} L \overset{l}{\leftarrowtail} K \overset{r}{\rightarrowtail} R \) with $\rho \mathop{=} (L \overset{l}{\leftarrowtail} K \overset{r}{\rightarrowtail} R) $.
    
        A match of the rule in an object \( G \) is a monomorphism \( m: L \mathop{\rightarrow} G \) such that there is no monomorphism \( q: N \mathop{\rightarrow} G \) such that \( n \mathop{\star} q \mathop{=} m \), as illustrated below.
    
        A DPO diagram as shown below is a \textbf{witness} for a \textbf{rewriting step} from \( G \) to \( H \) using the rule $\varphi$ and match \( m \), denoted \( G \mathop{\Rightarrow}_\rho^m H \) or \( G \mathop{\Rightarrow}_\rho^\delta H \). 
        % We denote $\operatorname{left}(\delta)$ and $\operatorname{right}(\delta)$ the pushout squares $KLGC$ and $KRHC$, respectively.
    
    \begin{figure}[H]
        \begin{tikzpicture} 
            % Define nodes
            \node (K)  {K};
            \node (L) [left=of K] {L};
            \node (R) [right=of K] {R};
            \node (N) [left=of L] {N};
            \node (G) [below=of L] {G};
            \node (D) [below=of K] {D};
            \node (H) [below=of R] {H};
    
            % Draw arrows
            % Top row
            \draw[>->] (L) -- node[above] {n} (N);
            \draw[>->] (K) -- node[above] {$l$} (L);
            \draw[>->] (K) -- node[above] {$r$} (R);
          
            % Vertical arrows
            \draw[>->] (L) -- node[left] {$m$} (G);
            \draw[>->] (K) --  (D);
            \draw[>->] (R) --  (H);
          
            % Bottom row
            \draw[>->] (G) --  (D);
            \draw[>->] (D) --   (H);
          
            % Bent arrow from G to N
            \draw[>->, bend right,red,dashed] (N) to node[below] {q} (G);
    
            \node at ($(N)!0.5!(G)$) {\textcolor{red}{$\mathop{\neq}$}};
            \node at ($(K)!0.5!(G)$) {$\mathrm{PO}$};
            \node at ($(K)!0.5!(H)$) {$\mathrm{PO}$};
          \end{tikzpicture}
          \label{fig:grs:dpo_nac_rule}
    \end{figure}
    \end{definition}
    \begin{example}
      The following diagram illustrates a DPO rewriting rule with a negative application condition:
      \begin{center} 
        \resizebox{0.7\textwidth}{!}{
        \begin{tikzpicture}
          \graphbox{$N$}{-40mm}{0mm}{34mm}{20mm}{2mm}{-5mm}{
            \coordinate (o) at (0mm,-3mm); 
            \node[draw,circle] (l1) at ($(o)+(-10mm,0mm)$) {1};
            \node[draw,circle] (l2) at ($(l1)+(2,0)$) {2};
            \node[draw,circle] (l3) at ($(l1)+(1,0)$) {3};
            \draw[->] (l1) -- (l3) node[midway,above] {$a$};
            \draw[->] (l3) -- (l2) node[midway,above] {$a$};
            \draw[->] (l2) edge[out=-135,in=-45]node[midway,below] {$a$} (l1) ;
        }  
            \graphbox{$L$}{0mm}{0mm}{34mm}{15mm}{2mm}{-5mm}{
                \coordinate (o) at (0mm,-3mm); 
                \node[draw,circle] (l1) at ($(o)+(-10mm,0mm)$) {1};
                \node[draw,circle] (l2) at ($(l1)+(2,0)$) {2};
                \node[draw,circle] (l3) at ($(l1)+(1,0)$) {3};
                \draw[->] (l1) -- (l3) node[midway,above] {$a$};
                \draw[->] (l3) -- (l2) node[midway,above] {$a$};
            }     
            \graphbox{$K$}{40mm}{0mm}{24mm}{15mm}{2mm}{-5mm}{
                \coordinate (o) at (5mm,-3mm); 
                \node[draw,circle] (l1) at ($(o)+(-10mm,0mm)$) {1};
                \node[draw,circle] (l2) at ($(l1)+(1,0)$) {2};
                % \node[draw,circle] (l3) at ($(l1)+(1,0)$) {$\ $};
                % \draw[->] (l1) -- (l3) node[midway,above] {$a$};
                % \draw[->] (l3) -- (l2) node[midway,above] {$a$};
            }    
            \graphbox{$R$}{70mm}{0mm}{45mm}{15mm}{2mm}{-5mm}{
                \coordinate (o) at (-5mm,-3mm); 
                \node[draw,circle] (l1) at ($(o)+(-10mm,0mm)$) {1};
                \node[draw,circle] (l2) at ($(l1)+(3,0)$) {2};
                \node[draw,circle] (l3) at ($(l1)+(1,0)$) {4};
                \node[draw,circle] (l4) at ($(l1)+(2,0)$) {5};
                \draw[->] (l1) -- (l3) node[midway,above] {$a$};
                \draw[->] (l3) -- (l4) node[midway,above] {$b$};
                \draw[->] (l4) -- (l2) node[midway,above] {$a$};
            }    
            \node () at (-3mm,-8mm) {$\leftarrowtail$};
            \node () at (37mm,-8mm) {$\leftarrowtail$};
            \node () at (67mm,-8mm) {$\rightarrowtail$};
            % \draw[>->] (51mm,2mm) -- (52mm,3mm);
        \end{tikzpicture}
        }
    \end{center}
      The rewriting step in Example~\ref{ex:rewriting_step_grs_aa} is not possible with this rule in DPO rewriting with negative application conditions, because there exists a monomorphism \( q: N \mathop{\rightarrow} G \) such that \( n \mathop{\star} q \mathop{=} m \) as shown in the diagram below.
      \begin{center} 
        \resizebox{0.7\textwidth}{!}{
        \begin{tikzpicture}
          \graphbox{$N$}{-40mm}{0mm}{34mm}{20mm}{2mm}{-5mm}{
            \coordinate (o) at (0mm,-3mm); 
            \node[draw,circle] (l1) at ($(o)+(-10mm,0mm)$) {1};
            \node[draw,circle] (l2) at ($(l1)+(2,0)$) {2};
            \node[draw,circle] (l3) at ($(l1)+(1,0)$) {3};
            \draw[->] (l1) -- (l3) node[midway,above] {$a$};
            \draw[->] (l3) -- (l2) node[midway,above] {$a$};
            \draw[->] (l2) edge[out=-135,in=-45]node[midway,below] {$a$} (l1) ;
        }  
            \graphbox{\( L \)}{0mm}{-3mm}{34mm}{12mm}{2mm}{2mm}{
                \coordinate (o) at (0mm,-8mm); 
                \node[draw,circle] (l1) at ($(o)+(-10mm,0mm)$) {1};
                \node[draw,circle] (l2) at ($(l1)+(2,0)$) {2};
                \node[draw,circle] (l3) at ($(l1)+(1,0)$) {3};
                \draw[->] (l1) -- (l3) node[midway,above] {$a$};
                \draw[->] (l3) -- (l2) node[midway,above] {$a$};
            } 
    
            \graphbox{\( K \)}{40mm}{-3mm}{34mm}{12mm}{2mm}{2mm}{
                \coordinate (o) at (0mm,-8mm); 
                \node[draw,circle] (l1) at ($(o)+(-10mm,0mm)$) {1};
                \node[draw,circle] (l2) at ($(l1)+(2,0)$) {2};
            }  
    
            \graphbox{\( R \)}{80mm}{-3mm}{45mm}{12mm}{2mm}{2mm}{
                \coordinate (o) at (-5mm,-8mm); 
                \node[draw,circle] (l1) at ($(o)+(-10mm,0mm)$) {1};
                \node[draw,circle] (l2) at ($(l1)+(3,0)$) {2};
                \node[draw,circle] (l3) at ($(l1)+(1,0)$) {4};
                \node[draw,circle] (l4) at ($(l1)+(2,0)$) {5};
                \draw[->] (l1) -- (l3) node[midway,above] {$a$};
                \draw[->] (l3) -- (l4) node[midway,above] {$b$};
                \draw[->] (l4) -- (l2) node[midway,above] {$a$};
            }    
    
            \graphbox{\( G \)}{0mm}{-22mm}{34mm}{22mm}{2mm}{-3mm}{
                \coordinate (o) at (0mm,-3mm); 
                \node[draw,circle] (l1) at ($(o)+(-10mm,0mm)$) {1};
                \node[draw,circle] (l2) at ($(l1)+(2,0)$) {2};
                \node[draw,circle] (l3) at ($(l1)+(1,0)$) {3};
                \node[draw,circle] (l4) at ($(l2)+(0,-1)$) {6};
                \draw[->] (l1) -- (l3) node[midway,above] {$a$};
                \draw[->] (l3) -- (l2) node[midway,above] {$a$};
                \draw[->] (l2) -- (l4) node[midway,right] {$a$};
                \node[draw,circle] (l6) at ($(l1)+(0,-1)$) {7};
                \draw[<-] (l1) -- (l6) node[midway,left] {$a$};
                \draw[->] (l2) edge[out=-135,in=-45]node[midway,below] {$a$} (l1) ;
            }    
    
            \graphbox{\( C  \)}{40mm}{-22mm}{34mm}{22mm}{2mm}{-3mm}{
                \coordinate (o) at (0mm,-3mm); 
                \node[draw,circle] (l1) at ($(o)+(-10mm,0mm)$) {1};
                \node[draw,circle] (l2) at ($(l1)+(2,0)$) {2};
                \node[draw,circle] (l4) at ($(l2)+(0,-1)$) {6};
                \draw[->] (l2) -- (l4) node[midway,right] {$a$};
                \draw[->] (l2) edge[out=-135,in=-45]node[midway,below] {$a$} (l1) ;
                \node[ draw,circle] (l6) at ($(l1)+(0,-1)$) {7};
                \draw[<-] (l1) -- (l6) node[midway,left] {$a$};
            }    
    
            \graphbox{\( H \)}{80mm}{-22mm}{45mm}{22mm}{2mm}{-3mm}{
                \coordinate (o) at (-5mm,-3mm); 
                \node[draw,circle] (l1) at ($(o)+(-10mm,0mm)$) {1};
                \node[draw,circle] (l2) at ($(l1)+(3,0)$) {2};
                \node[draw,circle] (l3) at ($(l1)+(1,0)$) {4};
                \node[draw,circle] (l4) at ($(l1)+(2,0)$) {5};
                \node[ draw,circle] (l5) at ($(l2)+(0,-1)$) {6};
                \node[ draw,circle] (l6) at ($(l1)+(0,-1)$) {7};
                \draw[<-] (l1) -- (l6) node[midway,left] {$a$};
                \draw[->] (l1) -- (l3) node[midway,above] {$a$};
                \draw[->] (l3) -- (l4) node[midway,above] {$b$};
                \draw[->] (l4) -- (l2) node[midway,above] {$a$};
                \draw[->] (l2) -- (l5) node[midway,right] {$a$};
                \draw[->] (l2) edge[out=-135,in=-45]node[midway,below] {$a$} (l1) ;
            }    
            \draw[>->, bend right] (-21mm,-22mm) to node[below] {q} (-4mm,-33mm);
            \node () at (-3mm,-8mm) {$\leftarrowtail$};
            \node () at (37mm,-8mm) {\( \leftarrowtail \)}; % K -> L
            \node () at (77mm,-8mm) {\( \rightarrowtail \)}; % K -> R
            \node () at (15mm,-18mm) {\( m\ \downarrowtail \)};
            \node () at (37mm,-33mm) {\( \leftarrowtail \)};
            \node () at (58mm,-18mm) {\( \downarrowtail \)};
            \node () at (102mm,-18mm) {\( \downarrowtail \)};
            \node () at (77mm,-33mm) {\( \rightarrowtail \)}; % C -> H
        \end{tikzpicture}
        }
    \end{center}
    \end{example}
    
    \begin{definition}[Rewriting framework]
      \label{def:dponac:rewriting_framework}
      A framework of rewriting with a negative application condition $\mathfrak{F}$ is a mapping of DPO rewriting rules with a negative condition to classes of diagrams as illustrated in Definition~\ref{def:grs:dpo_nac_rule} such that, for every rule $\varphi =(n, \rho)$, every diagram in $\mathfrak{F}(\varphi)$ has $\rho$ as the top-span of the DPO subdiagram.
    \end{definition}

    \trackedtext{
      \begin{definition}[Rewriting relation]
        The \textbf{rewriting relation $\mathop{\Rightarrow}_{\mathfrak{dponac},\mathfrak{F},\rho}$ induced by a rule $\rho$ in $\mathfrak{F}$} is defined as follows: $G \mathop{\Rightarrow}_{\mathfrak{dponac},\mathfrak{F},\rho} H$ iff $G \mathop{\Rightarrow}_\rho^\delta H$ for some $\delta \mathop{\in} \mathfrak{F}(\rho)$. 
          % for some $\delta \mathop{\in} \mathfrak{F}(\rho)$
          The \textbf{DPO rewriting relation $\mathop{\Rightarrow}_{\mathfrak{dponac},\mathfrak{F},\mathcal{R}}$ induced by a set $\mathcal{R}$ of DPO rewriting rules in $\mathfrak{F}$} is given by: $G \mathop{\Rightarrow}_{\mathfrak{dponac}, \mathfrak{F},\mathcal{R}} H$ iff $G \mathop{\Rightarrow}_{\mathfrak{dponac},\mathfrak{F}, \rho} H$ for some $\rho \mathop{\in} \mathcal{R}$. When $\mathfrak{F}$ is clear from the context, we 
          suppress $\mathfrak{F}$ and 
          write $\mathop{\Rightarrow}_{\mathfrak{dponac},\rho}$ and $\mathop{\Rightarrow}_{\mathfrak{dponac},\mathcal{R}}$.
      \end{definition}
    }

    \begin{definition}
      Let $M, \mathfrak{M}, W, \mathfrak{W}, \mathfrak{I}$ be the class of DPO matches, the match mechanism, the class of DPO witnesses, the witness mechanism and the interpretation function, respectively, defined in Definition~\ref{def:grs:dpo_nac_rule}.

      Let $\mathfrak{F}$ be a DPO rewriting framework with a negative application condition as defined in Definition~\ref{def:dponac:rewriting_framework}.

      The rewriting system $(\mathbf{Graph}, \mathcal{R}, M, \mathfrak{M}, W, \mathfrak{W}, \mathfrak{I}, \mathfrak{F})$ is called a \textbf{DPO Graph Rewriting System with a negative application condition in $\mathfrak{F}$}.
    \end{definition}
    
    \subsection{\pbpop~Rewriting Systems}
    \begin{definition}[PBPO+ Rewriting Rule~\cite{overbeek2023pbpo_JLAMP}]
      \label{def:pbpop:rule}
      A \pbpop~\textbf{rewriting rule} \(\rho\) is a diagram of the form
    \begin{center}
    
    \begin{tikzpicture}[node distance=12mm, auto]
    
    \node (L) {$L$};
    \node (K) [right=of L] {$K$};
    \node (R) [right=of K] {$R$};
    
    \node (L') [below=of L] {$L'$};
    \node (K') [below=of K] {$K'$};
    
    % Arrows for the upper row
    \draw[<-] (L) to node [above] {$l$} (K);
    \draw[->] (K) to node [above] {$r$} (R);
    
    % Vertical arrows 
    \draw[>->] (L) to node [left] {$t_L$} (L');
    \draw[>->] (K) to node [right] {$t_K$} (K');
    
    % Arrows for the lower row
    \draw[<-] (L') to node [below] {$l'$} (K');
    
    % Pullback label
    \node at ($(L)!0.5!(K')$) {$PB$};
    \end{tikzpicture}
    \end{center}
    
    \noindent 
    where \(L\), \(K\) and \(R\) are the input, interface and output patterns; \(L'\) is the context type of $L$, and \(t_L\) is the context typing of $L$; \(K'\) is the context type of $K$, and \(t_K\) is the context typing of $K$. 
    \end{definition} 
    \textcolor{red}{intuition:
    L is the pattern; L' restricts the shape of the context; pullback: duplication and filtering(extraction); pushout:merging and extension; strong matching: leftmost pullback square ensures that the context in $G_L$ is not mapped onto the pattern L in L'
    //
    The morphism \(l' : L' \mathop{\leftarrow} K'\) is used to define deletion and duplication operations on the type graph. The morphism \(r : K \mathop{\to} R\) subsequently specifies the identification and addition operations, particularly focused on \(l'^{-1}(t_L)\).
    }
    % For PBPO\(^*\) matches, it is required that \(\alpha^{-1}(\operatorname{Im} t_L) \overset{iso}{=} L\). Such a match is called a \textit{strong match}, which distinguishes $PBPO+$ from $PBPO$, and is uniquely determined by the means of an adherence morphism.
    \begin{example}
        \url{https://www.youtube.com/watch?v=_gwz64o1eBQ&t=1352s} 30:20 an example of a rule which remove a single loop from a graph.
    \end{example}

    \begin{proposition}[\text{\cite[Proposition 12]{overbeek2023pbpo_JLAMP}}]
      Consider the following pullback square:

      \begin{tikzpicture}
          \node (L1) at (0,0) {$L$}; 
          \node (L2) at (0,-2) {$L$};
          \node (G) at (2,0) {$G$};
          \node (L') at (2,-2) {$L'$};
          \draw[->] (L1) to node [above] {$m$} (G);
          \draw[->] (L2) to node [above] {$t_L$} (L'); 
          \draw[->] (G) to node [right] {$\alpha$} (L');
          \draw[->] (L1) to node [left] {$\opn{id}_L$} (L2);
          \node () at (1,-1) {$\mathrm{PB}$};
      \end{tikzpicture}

      The morphism \(m\) is monic if and only if $t_L$ is monic.
    \end{proposition}

    \begin{definition}[Match]
      \label{def:pbpop:match}
      A \textbf{match} of a \pbpop~rewriting rule \(\rho\) as illustrated in Definition~\ref{def:pbpop:rule} in an object \(G\) is an ordered pair of morphisms \((m:L \rightarrowtail G, \alpha: G \mathop{\rightarrow} L')\) such that the following is a pullback square:
    
      \begin{tikzpicture}
        \node (L1) at (0,0) {$L$}; 
        \node (L2) at (0,-2) {$L$};
        \node (G) at (2,0) {$G$};
        \node (L') at (2,-2) {$L'$};
        \draw[>->] (L1) to node [above] {$m$} (G);
        \draw[>->] (L2) to node [above] {$t_L$} (L'); 
        \draw[->] (G) to node [right] {$\alpha$} (L');
        \draw[->] (L1) to node [left] {$\opn{id}_L$} (L2);
        \node () at (1,-1) {$\mathrm{PB}$};
    \end{tikzpicture}

      \end{definition}
    
\begin{remark}[\text{\cite[Remark 11]{overbeek2023pbpo_JLAMP}}]
In both \textbf{Set} and \textbf{Graph} the match diagram of Definition~\ref{def:pbpop:match} states that the preimage of $t_L(L)$ under $\alpha$ is $L$ itself. 
\end{remark}
Intuitively, to rewrite a graph \(G\) using a \pbpop~rewriting rule \(\rho\), we need to find an occurrence $m(L)$ of \(L\) in \(G\) such that $m(L)$ is mapped to the occurrence $t_L(L)$ of $L$ in \(L'\), and the other part of \(G\) is mapped to the other part of \(L'\) by \(\alpha\).

\begin{example}
  todo
\end{example}

    \begin{proposition}
      In any category, if diagrams
    \[
    \begin{tikzcd}
    L \arrow[r, "m'"] \arrow[d, equals] & G \arrow[d, "m"] \\
    L \arrow[r, "t_L"] & L'
    \arrow[phantom, from=1-1, to=2-2, "\mathrm{PB}" description]
    \end{tikzcd}
    \quad \text{and} \quad
    \begin{tikzcd}
    L \arrow[r, "m''"] \arrow[d, equals] & G \arrow[d, "m"] \\
    L \arrow[r, "t_L"] & L'
    \arrow[phantom, from=1-1, to=2-2, "\mathrm{=}" description]
    \end{tikzcd}
    \]
    hold then $m' \mathop{=} m''$.
    \end{proposition}
    
    \begin{definition}[\pbpop~rewrite step~\cite{overbeek2023pbpo_JLAMP}]
      \label{def:pbpop:rewriting_step}
      Let $\rho$ be a \pbpop~rewriting rule as illustrated in Definition~\ref{def:pbpop:rule}. 
    
      Let $G_L$ be an object and \(m, \alpha\) a match. 
      
     A diagram as illustrated below, where \(u : K \mathop{\rightarrow} G_K\) is the unique morphism satisfying \(t_K \mathop{=} u' \circ u\), is a \textbf{witness} of a \pbpop~\textbf{rewrite step} from $G_L$ to $G_R$ using rule $\rho$ and match $m$, denoted \(G_L \mathop{\Rightarrow}^m_{\text{\pbpop},\rho} G_R\).
    
      \begin{tikzpicture} 
        % Nodes of the diagram
          \node (L) {$L$};
          \node (GL) [right=of L] {$G_L$};
          \node (GK) [right=of GL] {$G_K$};
          \node (K) [above=of GK] {$K$};
          \node (R) [right=of K] {$R$};
          \node (GR) [below=of R] {$G_R$};
    
          \node (L2) [below=of L] {$L$};
          \node (L') [right=of L2] {$L'$};
          \node (K') [right=of L'] {$K'$};
    
          % Arrows on the top row
          \draw[>->] (L) to node [above] {$m$} (GL);
          \draw[->] (K) to node [above] {$r$} (R);
          \draw[<-] (GL) to node [above] {$g_L$} (GK);
          \draw[->] (GK) to node [left] {$u'$} (K');
          \draw[<-,dashed] (GK) to node [left] {$!u$} (K);
          \draw[->] (R) to node [right] {$w$} (GR);
          \draw[->] (GK) to node [above] {$g_R$} (GR);
    
          % Arrows on the bottom row
          \draw[>->] (L2) to node [below] {$t_L$} (L');
          \draw[<-] (L') to node [below] {$l'$} (K');
    
          % Vertical arrows
          \draw[->] (GL) to node [right] {$\alpha$} (L');
          \draw[=] (L) to node [left] {} (L2);
    
          % Curved arrow from K to GK
          \draw[>->] (K) to [bend left] node [pos=0.75,right] {$t_K$} (K');
    
          % PB and PO labels
          \node at ($(L)!0.5!(L')$) {$PB$};
          \node at ($(GK)!0.5!(L')$) {$PB$};
          \node at ($(K)!0.5!(GR)$) {$PO$}; 
      \end{tikzpicture}
    \end{definition}
    
    \textcolor{red}{
    \begin{definition}[Rewriting framework]
      A \pbpop~\textbf{rewriting framework} $\mathfrak{F}$ is a mapping associating to each \pbpop~rule $\rho$ as depicted in Definition~\ref{def:pbpop:rule} 
       the class of \pbpop~witnesses of form as illustrated in Definition~\ref{def:pbpop:rule}.
    \end{definition}
    }

    \trackedtext{
      \begin{definition}[Rewriting relation]
        The \textbf{rewriting relation $\mathop{\Rightarrow}_{\text{\pbpop},\mathfrak{F},\rho} $ induced by a rule $\rho$ in $\mathfrak{F}$} is defined as follows: $G_L \mathop{\Rightarrow}_{\text{\pbpop},\mathfrak{F},\rho} G_R$ iff $G_L \mathop{\Rightarrow}^m_{\text{\pbpop},\rho} G_R$ for some match \(m\) and some witness \(w\) in \(\mathfrak{F}(\rho)\).
        \\
          The \textbf{rewriting relation $\mathop{\Rightarrow}_{\text{\pbpop},\mathfrak{F},\mathcal{R}}$ induced by a set $\mathcal{R}$ of \pbpop~rewriting rules in $\mathfrak{F}$} is given by: $G_L \mathop{\Rightarrow}_{\text{\pbpop},\mathfrak{F},\mathcal{R}} G_R$ iff $G_L \mathop{\Rightarrow}_{\text{\pbpop},\mathfrak{F},\rho} G_R$ for some rule $\rho \mathop{\in} \mathcal{R}$. When $\mathfrak{F}$ is clear from the context, we suppress $\mathfrak{F}$ and write $\mathop{\Rightarrow}_{\text{\pbpop},\rho}$ and $\mathop{\Rightarrow}_{\text{\pbpop},\mathcal{R}}$.
      \end{definition}
    }
    
    \begin{definition}
      Let $\mathcal{R}$ be a set of \pbpop~rewriting rules.
      Let $M, \mathfrak{M}, W, \mathfrak{W}, \mathfrak{I}$ be the class of \pbpop~matches, the match mechanism, the class of \pbpop~witnesses, the witness mechanism and the interpretation function, respectively, defined in Definition~\ref{def:pbpop:rule} and Definition~\ref{def:pbpop:rewriting_step}.
      \\
      The rewriting system $(\mathbf{Graph}, \mathcal{R}, M, \mathfrak{M}, W, \mathfrak{W}, \mathfrak{I})$ is called a \textbf{\pbpop~graph rewriting system}.
    \end{definition}

    % \section{Termination of Graph Rewriting Systems}
    %     \label{sec:grs_termination}
    %     todo
    %     \subsection{Modular Termination of DPO GRS}
    %     \label{sec:modular_termination}
    %     Plump \cite{plump2018modular} later proposed a modular critical pair-based strategy for left-injective DPO hypergraph rewriting with monic matches. 
Our method complements this: while modularity reduces global complexity, each subsystem requires individual termination proofs. 

We brively recall the basic definitions of hypergraphs and show that a graph is a hypergraph.

 For a more comprehensive introduction to DPO hypergraph rewriting, we refer the reader to \cite{plump2018modular}. 

\begin{definition}
    A signature $\Sigma = \langle \Sigma_V, \Sigma_E, \text{Type} \rangle$ consists of 
    \begin{itemize}
        \item a set $\Sigma_V$ of node labels,
        \item a set $\Sigma_E$ of hyperedge labels, and
        \item a function $\text{Type}$ assigning to each $l \in \Sigma_E$ a set of strings $\text{Type}(l) \subseteq \Sigma_V^*$.
    \end{itemize}
\end{definition}
Unless stated otherwise, we denote by $\Sigma$ an arbitrary but fixed signature over which all hypergraphs are labelled.

\begin{remark}
    The extension $f^* : A^* \to B^*$ of a function $f : A \to B$ maps the empty string to itself and $a_1 \ldots a_n$ to $f(a_1) \ldots f(a_n)$.
\end{remark}
\begin{definition}
    A hypergraph over $\Sigma$ is a system $G = \langle V_G, E_G, \text{mark}_G, \text{lab}_G, \text{att}_G \rangle$ consisting of 
    \begin{itemize}
        \item a finite set $V_G$ of nodes,
        \item a finite set $E_G$ of hyperedges,
        \item two labelling functions $\text{mark}_G : V_G \to \Sigma_V$ and $\text{lab}_G : E_G \to \Sigma_E$, and
        \item an attachment function $\text{att}_G : E_G \to V_G^*$ such that $\text{mark}_G^*(\text{att}_G(e)) \in \text{Type}(\text{lab}_G(e))$ for each hyperedge $e$.
    \end{itemize} 
\end{definition} 

\begin{definition}[Sequential independence]
    \todo{rewrite}
   Direct derivations $G \Rightarrow_{r_1} H \Rightarrow_{r_2} M$ as in Figure 4 are sequentially independent if there are morphisms $R_1 \to D_2$ and $L_2 \to D_1$ such that the following holds.
    \begin{itemize}
        \item Commutativity: $R_1 \to D_2 \to H = R_1 \to H$ and $L_2 \to D_1 \to H = L_2 \to H$.
        \item Injectivity: $R_1 \to D_2 \to M$ is injective.
    \end{itemize}
    
    \begin{tikzpicture} 
        % Define nodes
        \node (L1) at (0,2) {$L_1$};
        \node (K1) at (2,2) {$K_1$};
        \node (R1) at (4,2) {$R_1$};
        \node (L2) at (6,2) {$L_2$};
        \node (K2) at (8,2) {$K_2$};
        \node (R2) at (10,2) {$R_2$};
        \node (G) at (0,0) {$G$};
        \node (T) at (5,0) {$H$};
        \node (M) at (10,0) {$M$};
        \node (D1) at (2,0) {$D_1$};
        \node (D2) at (8,0) {$D_2$};
      
        % Draw arrows
        % Top row
        \draw[->] (L1) -- node[above] {} (K1);
        \draw[->] (K1) -- node[above] {} (R1);
        \draw[->] (L2) -- node[above] {} (K2);
        \draw[->] (K2) -- node[above] {} (R2);
      
        % Vertical arrows
        \draw[->] (L1) -- node[left] {} (G);
        \draw[->] (R1) -- node[right] {} (T);
        \draw[->] (R2) -- node[right] {} (M);
        \draw[->] (K1) -- node[right] {} (D1);
        \draw[->] (K2) -- node[right] {} (D2);
      
        % Bottom row
        \draw[->] (G) -- node[below] {} (D1);
        \draw[->] (D1) -- node[below] {} (T);
        \draw[->] (T) -- node[below] {} (D2);
        \draw[->] (D2) -- node[below] {} (M);

        \draw[->,dashed] (R1) -- node[below] {} (D2);
        \draw[->,dashed] (L2) -- node[below] {} (D1);
      
        % Diagonal arrows
        \draw[->] (R1) -- node[below left] {$h_1$} (T);
        \draw[->] (L2) -- node[below right] {$g_2$} (T);
      \end{tikzpicture}
\end{definition}

\begin{definition}[Sequential critical pair]
    Let $r_i \colon \langle L_i \leftarrow K_i \rightarrow R_i \rangle$ be rules, for $i = 1, 2$. A pair of direct derivations $G \Rightarrow_{r_1} H \Rightarrow_{r_2} M$ as depicted below is a sequential critical pair if the following holds.
    \begin{itemize}
        \item Minimality:$H = h_1(R_1) \cup g_2(L_2)$.
        \item Conflict: The steps are not sequentially independent.
    \end{itemize}
    \begin{tikzpicture} 
        % Define nodes
        \node (L1) at (0,2) {$L_1$};
        \node (K1) at (2,2) {$K_1$};
        \node (R1) at (4,2) {$R_1$};
        \node (L2) at (6,2) {$L_2$};
        \node (K2) at (8,2) {$K_2$};
        \node (R2) at (10,2) {$R_2$};
        \node (G) at (0,0) {$G$};
        \node (T) at (5,0) {$H$};
        \node (M) at (10,0) {$M$};
        \node (D1) at (2,0) {$D_1$};
        \node (D2) at (8,0) {$D_2$};
      
        % Draw arrows
        % Top row
        \draw[->] (L1) -- node[above] {} (K1);
        \draw[->] (K1) -- node[above] {} (R1);
        \draw[->] (L2) -- node[above] {} (K2);
        \draw[->] (K2) -- node[above] {} (R2);
      
        % Vertical arrows
        \draw[->] (L1) -- node[left] {} (G);
        \draw[->] (R1) -- node[right] {} (T);
        \draw[->] (R2) -- node[right] {} (M);
        \draw[->] (K1) -- node[right] {} (D1);
        \draw[->] (K2) -- node[right] {} (D2);
      
        % Bottom row
        \draw[->] (G) -- node[below] {} (D1);
        \draw[->] (D1) -- node[below] {} (T);
        \draw[->] (T) -- node[below] {} (D2);
        \draw[->] (D2) -- node[below] {} (M);

        % \draw[->,dashed] (R1) -- node[below] {} (D2);
        % \draw[->,dashed] (L2) -- node[below] {} (D1);
      
        % Diagonal arrows
        \draw[->] (R1) -- node[below left] {$h_1$} (T);
        \draw[->] (L2) -- node[below right] {$g_2$} (T);
      \end{tikzpicture}
\end{definition}

\begin{theorem}
    \todo{rewrite}
    Let $\langle \Sigma, \mathcal{R} \rangle$ and $\langle \Sigma, \mathcal{S} \rangle$ be terminating hypergraph transformation systems. If there are no sequential critical pairs of shape $S \Rightarrow_{\mathcal{R}} T \Rightarrow_{\mathcal{S}} U$, then the combined system $\langle \Sigma, \mathcal{R} \cup \mathcal{S} \rangle$ is terminating.
\end{theorem}

    %     \subsection{Forward Closure Method of DPO GRS}  
    %     \label{sec:forward_closure}
    %     Plump \cite{Plump1995} introduced a necessary and sufficient termination condition for left-injective DPO hypergraph rewriting via forward closure, though verifying this condition is undecidable.


\begin{definition}[Instance]
    \todo{rewrite}
    Consider a finite or infinite derivation
    

    \begin{tikzpicture}
      % Nodes
      \node (pi1) at (0,0) {$\pi_1$};
      \node (pip1) at (2,0) {$\pi'_1$}; 
      \node (pi2) at (4,0) {$\pi_2$};
      \node (pip2) at (6,0) {$\pi'_2$};
      \node (pi3) at (8,0) {$\pi_3$};
      \node (pip3) at (10,0) {$\pi'_3$};
      \node (dots) at (12,0) {$\dots$};
    
      % Horizontal arrows in the top row
      \draw[->] (pi1) -- ++(-2,0);
      \draw[->] (pip1) -- ++(2,0);
      \draw[->] (pi2) -- ++(-2,0);
      \draw[->] (pip2) -- ++(2,0);
      \draw[->] (pi3) -- ++(-2,0);
      \draw[->] (pip3) -- ++(2,0);
    
      % Vertical arrows
      \draw[->] (pi1) -- ++(0,-1);
      \draw[->] (pip1) -- ++(0,-1);
      \draw[->] (pi2) -- ++(0,-1);
      \draw[->] (pip2) -- ++(0,-1);
      \draw[->] (pi3) -- ++(0,-1);
      \draw[->] (pip3) -- ++(0,-1);
    
      % Diagonal arrows
      \draw[->] (pip1) -- ++(1,-1);
      \draw[->] (pip2) -- ++(1,-1);
      \draw[->] (pip3) -- ++(1,-1);
    
      % Horizontal arrows in the bottom row
      \draw[->] (pi1) -- ++(2,0);
      \draw[->] (pip1) -- ++(-2,0);
      \draw[->] (pi2) -- ++(2,0);
      \draw[->] (pip2) -- ++(-2,0);
      \draw[->] (pi3) -- ++(2,0);
      \draw[->] (pip3) -- ++(-2,0);
    \end{tikzpicture}
    
    A derivation is an \textit{instance} (or \textit{embedding}) of this derivation if it can be written
    

    
    where for each $i \geq 1$, $\phi_i$ and $\phi'_i$ are pushouts with injective vertical morphisms. (So the pushouts constituting the derivation are composed of the pushouts $\pi_i$ and $\phi_i$, respectively $\pi'_i$ and $\phi'_i$.)
\end{definition}

\begin{definition}[forward closure]
    Forward closures are inductively defined as follows:
    \begin{enumerate}
        \item Every direct derivation $L \Rightarrow_{r,g} R$ with surjective $g$ is a forward closure.
        \item A derivation $G \Rightarrow^+ H \Rightarrow M$ is a forward closure if $G \Rightarrow^+ H$ is an instance of a forward closure, and $H \Rightarrow M$ is a direct derivation such that
        \begin{itemize}
            \item $G \Rightarrow^+ H$ and $H \Rightarrow M$ are dependent,
            \item $\text{New}(H) \subseteq \text{Redex}(H \Rightarrow M)$.
        \end{itemize}
    \end{enumerate}
\end{definition}

\begin{definition}[infinite forward closure]
    \todo{rewrite}
    An \textit{infinite forward closure} is an infinite derivation $G_0 \Rightarrow G_1 \Rightarrow G_2 \Rightarrow \dots$ that contains a forward closure as a prefix. That is, there is some $n \geq 1$ such that $G_0 \Rightarrow^+ G_n$ is a forward closure.
\end{definition}

\begin{definition}[Consumptive System]
    \todo{rewrite}
    A graph rewriting system is \textit{consumptive} if each rule $(L \leftarrow K \rightarrow R)$ has a non-surjective morphism $K \to L$.
\end{definition}
    
\begin{theorem}[Main Theorem]
    \todo{rewrite}
    A graph rewriting system is terminating if, and only if, it is consumptive and does not admit an infinite forward closure.
\end{theorem}
    %     \subsection{Termination Criterion of injective DPO GRS based on E-dependency relation \textcolor{red}{todo}} 
    % \begin{definition}[E-graph]
\end{definition}

\begin{definition}[size of a graph]
\end{definition}

\begin{definition}[production]
\end{definition}

\begin{definition}[E-concurrent production~\text{\cite[Def~3.1]{LEVENDOVSZKY200787}}]
    An E-concurrent production $p^*$ is an E-based composition if there is at least one input graph $G_0$ with an E-related transformation $G_0 \xrightarrow{p^*} H$.
\end{definition}

\begin{definition}[E-dependency relation~\text{\cite[Def~3.2]{LEVENDOVSZKY200787}}]
    Consider a possibly infinite sequence of graph productions $p_i$, ($i \mathop{=} 1, 2, \ldots$) and a sequence of E-dependency relations $((E_i, e_i^*, e_{i+1}))$ leading to a sequence of their E-based compositions $(p_i^* \mathop{=} (L_i^* \mathop{\leftarrow} K_i^* \mathop{\rightarrow} R_i^*))$ with $p_1^* \mathop{=} p_1$ and $p_n^* \mathop{=} (p_1 *_{E_1} p_2) *_{E_2} \ldots *_{E_n} p_n$.
    
    A cumulative LHS series of this sequence is the graph series $L_n^*$ consisting of the left-hand side graphs of $p_n^*$. Moreover, a cumulative size series of a production sequence is the nonnegative integer series $|L_n^*|$.
    \end{definition}

\begin{theorem}[Termination~\cite{LEVENDOVSZKY200787}]
    A \text{GTS} $= (P)$ (Def.??) terminates if for all infinite cumulative LHS sequences $(L_i^*)$ of the graph productions created from the members of $P$, it holds that
    \[
    \lim_{i \mathop{\to} \infty} |L_i^*| \mathop{=} \infty.
    \]
    Note that we assume finite input graphs and injective matches.
\end{theorem} 
    %     \subsection{Type Graph Method for DPO GRS}
    %     \label{sec:type_graph_method}
    %         \subsubsection{Well-founded semirings}
    %         \begin{definition}[Monoid]
    \label{def:monoid}
    \todo{to check and add reference}
    A \textbf{monoid} is a triple \(\langle S, \cdot, e \rangle\) where \(S\) is a set, \(\cdot : S \times S \to S\) is a binary operation and \(e \in S\) is constant
    satisfying for all \(a, b, c \in S\):
    \begin{itemize}
        \item \textbf{Associativity}: \((a \cdot b) \cdot c = a \cdot (b \cdot c)\),
        \item \textbf{Identity}: \(a \cdot e = a = e \cdot a\).
    \end{itemize}
    The monoid is \textbf{commutative} if \(a \cdot b = b \cdot a\) for all \(a, b \in S\).
\end{definition} 

\begin{definition}[Semiring]
    \label{def:semiring}
    \todo{to check and to add some references}
    A \textbf{semiring} is a tuple \(\langle S, \oplus, \odot, 0, 1 \rangle\) where:
    \begin{itemize}
        \item  \(\langle S, \oplus, 0 \rangle\) forms a commutative monoid,
        \item  \(\langle S, \odot, 1 \rangle\) forms a monoid,
        \item  $0$ is an annihilator for $\odot$ : For all \(a \in S\),
              \(
                  0 \odot a = 0 = a \odot 0
              \),
        \item $\odot$ distributes over $\oplus$: For all \(a, b, x \in S\),
              \begin{itemize}
                \item $(a \oplus b) \odot x = (a \odot x) \oplus (b \odot x)$
                \item $x \odot (a \oplus b) = (x \odot a) \oplus (x \odot b)$
              \end{itemize}                 
    \end{itemize}
    The semiring is \textbf{commutative} if its multiplicative monoid \(\langle S, \odot, 1 \rangle\) is commutative.
\end{definition}


\begin{definition}[Well-founded semiring]
    \label{def:well_founded_semiring}
    A \textbf{well-founded semiring} is a tuple \(\langle S, \oplus, \odot, 0, 1, \prec, \leq \rangle\) where:
    \begin{itemize}
        \item  \(\langle S, \oplus, \odot, 0, 1 \rangle\) is a semiring,
        \item \textbf{Order Relations}: 
            \begin{itemize}
                \item \(\prec, \leq \subseteq S \times S\) are non-empty orders,
                \item \(\prec \subseteq \leq\) (i.e., \(x \prec y \implies x \leq y\)),
                \item \(\leq\) is reflexive.
            \end{itemize}
    \end{itemize}
    The structure satisfies:
    \begin{itemize}
        \item \(\text{SN}(\succ / \geq)\) (strong normalization modulo \(\geq\)),
        \item \(0 \neq 1\),
        \item For all \(x, x', y, y' \in S\):
            \begin{align*}
                x \leq x' \land y \leq y' &\implies x \oplus y \leq x' \oplus y' \tag{S1} \label{ax:S1} \\
                x \prec x' \land y \prec y' &\implies x \oplus y \prec x' \oplus y' \tag{S2} \label{ax:S2} \\
                x \leq x' \land 1 \leq y &\implies x \odot y \leq x' \odot y \tag{S3} \label{ax:S3} \\
                x \prec x' \land 1 \leq y \neq 0 &\implies x \odot y \prec x' \odot y \tag{S4} \label{ax:S4}
            \end{align*}
    \end{itemize}
    The semiring is \textbf{strictly monotonic} if it additionally satisfies:
    \[
        x \prec x' \land y \leq y' \implies x \oplus y \prec x' \oplus y' \tag{S5} \label{ax:S5}
    \]
\end{definition}

\begin{example}[Semiring examples~\text{\cite[Ex. 2.7]{endrullis2024generalized}}]
    \label{ex:semiring_examples}
    \textbf{Arithmetic semiring}: \(\langle \mathbb{N}, +, \cdot, 0, 1, <, \leq \rangle\) is strictly monotonic and well-founded.

    \textbf{Tropical semiring}: \(\langle \mathbb{N} \cup \{\infty\}, \min, +, \infty, 0, <, \leq \rangle\) is well-founded but not strictly monotone.

    \textbf{Arctic semiring}: \(\langle \mathbb{N} \cup \{-\infty\}, \max, +, -\infty, 0, <, \leq \rangle\) is well-founded but not strictly monotone.
\end{example}
    %         \subsubsection{Weighted Type Graphs}
    %         The following definition of a weighted type graph is obtained from~\cite[\textdef~3.1]{endrullis2024generalized}.
\begin{definition}[Weighted Type Graph]
    \label{def:weighted_type_graph}
    A \textbf{weighted type graph} \(\mathcal{T} \mathop{=} (T, \mathbb{E}, \mathcal{S}, w)\) consists of:
    \begin{itemize} 
        \item An object \(T \mathop{\in} \mathcal{C}_0\), called the \textbf{type graph},
        \item A set \(\mathbb{E}\) of arrows \(e \mathop{\in} \mathcal{C}_1\) with \(\operatorname{codom}(e) \mathop{=} T\), called the \textbf{\(T\)-valued elements},
        \item A well-founded, commutative semiring \(\mathcal{S}=\langle S, \mathop{\oplus}, \mathop{\odot}, 0, 1, \prec, \leq \rangle \),
        \item A weight function \(w : \mathbb{E} \mathop{\to} S \mathop{\setminus} \{0_S\}\).
    \end{itemize}
    \(\mathcal{T}\) is \textbf{finitary} if for every \((e:X \mathop{\to} T) \mathop{\in} \mathbb{E}\) and every \(G \mathop{\in} \mathcal{C}_0\), the sets \(\operatorname{Hom}(X, G)\) and \(\operatorname{Hom}(G, T)\) are finite.
\end{definition}
\todo{include : $weighted_type_graph_remark_geq1$}
\todo{include : $weighted_type_graph_remark_neq0$}
    %         \subsubsection{Weighing Morphisms and Objects}
    %         % \begin{notation} For morphisms \( \alpha : A \to B \), \( \beta : B \to C \), \( \gamma : A \to C \) and $E \subseteq \operatorname{Hom}(A,B)$, define:
%     \begin{flalign*}
%                \set{ \alpha \star - = \gamma } &\overset{\operatorname{def}}{=} \{ \beta \in \operatorname{Hom}(B, C) \mid \alpha \star \beta = \gamma \}.
%    \\
%                \set{ - \star \beta = \gamma }  &\overset{\operatorname{def}}{=} \{ \alpha \in \operatorname{Hom}(A, B) \mid \alpha \star \beta = \gamma \}.
%    \\
%                E \star \beta                   &\overset{\operatorname{def}}{=} \set{ \alpha \star \beta \mid \alpha \in E }.
%     \end{flalign*}
%    \end{notation} 
   
   \begin{definition} 
       \label{def:bigodot}
   Let $(S, \oplus, \odot, 0_s, 1_s)$ be a semiring. We define 
    \begin{flalign*}
       \bigodot \emptyset &\overset{\operatorname{def}}{=} 1_s
   \\
       \bigodot \left( E \cup \{x\} \right) &\overset{\operatorname{def}}{=} \left( \bigodot E \right) \odot x
       \\
       \bigoplus \emptyset &\overset{\operatorname{def}}{=} 0_s
       \\
           \bigoplus \left( E \cup \{x\} \right) &\overset{\operatorname{def}}{=} \left( \bigoplus E \right) \oplus x
   \end{flalign*}
   \end{definition}
    
   
   \begin{definition}[\text{\cite[\textdef~3.4]{endrullis2024generalized}}] 
       \label{def:weight}
       Let $\mathcal{T} = (T,\mathbb{E},S, w)$ be a finitary type
       \newline
       \noindent
       \begin{minipage}{0.6\textwidth}graph.~The \textbf{weight of a morphism $h: G \rightarrow T$ relative to $(e:X \to T) \in \mathbb{E}$} is defined as the weight $w(e)$ raised to the power of the number of $\iota$ making the diagram (shown on the right) commutative.
       \end{minipage}
       \begin{minipage}{0.29\textwidth}
           \begin{center}
                   \begin{tikzpicture}
                       \node (a) at (0,0) {X};
                       \node (c) at (4,0) {G};
                       \node (d) at (2,-1) {T};
                       \draw[->] (a) -- (d) node [midway,below] {e};
                       \draw[->] (a) -- (c) node [midway,above] {$\iota$};
                       \draw[->] (c) -- (d) node[midway, below] {$h$};
                       % \node (d) at (2,-0.5) {=};
                   \end{tikzpicture}
           \end{center} 
       \end{minipage}
                   \[
                   w_e(h) 
                       \overset{\operatorname{def}}{=}
                   \underset{\alpha \in \{- \star h = e\}}{\bigodot}w(e) 
                   \]
           % which can be visualized as $w(e)$ raised to the power of the number of possible morphisms $g$ which make the following commutative diagram hold
   
           \noindent
           The \textbf{weight of a morphism $h: G \rightarrow T$ relative to \(\mathcal{T}\)} is defined as the semiring product of $w(e)$ for $e \in \mathbb{E}$:
           \[  w_\mathcal{T}(h) \overset{\operatorname{def}}{=} \underset{e \in \mathbb{E}}{\bigodot} 
                   w_e(h) \]
   
           \noindent
          The \textbf{weight of an object \( G \in \mathcal{C}_0 \) relative to \( \mathcal{T}\)} is defined as the semiring sum of $w_\mathcal{T}(h)$ for $h \in \operatorname{Hom}(G,T)$:
           \[w_\mathcal{T}(G) \overset{\operatorname{def}}{=} \underset{h \in \operatorname{Hom}(G,T)}{\bigoplus}  w_\mathcal{T}(h) \]
   \end{definition}
   
   
   % In certain scenarios, we want to exclude specific morphisms when calculating morphism weight relative to a give T-valued element $e:X \to T$. For this reason, for all set \( \Gamma \subseteq \operatorname{Hom}(A, G) \), we define 
   
   
   
   % The weight of the morphism \( \phi : G \to T\) excluding morphisms in \( \Gamma' \) is defined as the number of morphisms from \( X \) to \( G \) that do not factor through any \( \alpha \in \Gamma \).
   
   \begin{definition}[\text{\cite[\textdef~3.4]{endrullis2024generalized}}]
       \label{def:weight_excluding}
       Let \( A \in \mathcal{C}_0 \) and $\Gamma \subseteq \operatorname{Hom}(A,G)$. We de-
       \newline
       \noindent
       \begin{minipage}{0.6\textwidth}
         fine $\Gamma'$ as the set consisting of all morphisms \( \iota : X \to G \) admitting morphisms \( \zeta \colon X \to A \) and \( \alpha \in \Gamma \) such that the diagram illustrated on the right is commutative. Formally, 
       \end{minipage}
       \begin{minipage}{0.4\textwidth}
           \hfill 
           % \begin{center}
               \begin{tikzpicture}[scale=0.8]
                   \node (X) at (0,0) {\(X \)};
                   \node (A) at (2,1) {\( A \)};
                   \node (G) at (4,0) {\( G \)}; 
                   \draw[->] (X) -- (A) node[midway, above] {\( \zeta \)};
                   \draw[ ->] (X) -- (G) node[midway, below] {\( \iota \)};
                   \draw[->] (A) -- (G) node[midway, above] {\( \alpha \)};
                   % \node at (2,0.5) {\( = \)};
                   % \node (T) at (2,-1) {\( T \)};
                   % \draw[ ->] (X) -- (T) node[midway, below] {e};
                   % \draw[<-] (T) -- (G) node[midway, above] {};
               \end{tikzpicture}
           % \end{center} 
       \end{minipage}
   
       \[
       \Gamma' \overset{\operatorname{def}}{=} \left\{ \iota \in \operatorname{Hom}(X, G)~\middle|~\exists \alpha \in \Gamma,~\exists \zeta:X \to A,~\zeta \star \alpha = \iota \right\}.
       \]
   
       \noindent
       The \textbf{weight of a morphism \(h : G \to T\) excluding morphisms in \( \Gamma' \) relative to $(e:X \to T) \in \mathbb{E}$} is defined as $w(e)$ raised to the power of the number of morphisms \( (\iota : X \to G) \notin \Gamma' \).
           \[
           w_e(h - \Gamma) \overset{\operatorname{def}}{=} \underset{
               \substack{\alpha \in \set{- \star h = e} \\
                           \alpha \notin \Gamma'}}{\bigodot} w(e)\] 
           The \textbf{weight of a morphism $h: G \to T$ excluding morphisms in \( \Gamma' \) relative to \(\mathcal{T}\)} is defined as the semiring product of $w_e(h-\Gamma)$ for $e \in \mathbb{E}$:
           \[ 
               w_\mathcal{T}(h-\Gamma) \overset{\operatorname{def}}{=} \underset{e \in \mathbb{E}}{\bigodot} 
           w_e(h-\Gamma)
                   \]
   \end{definition} 
   % \begin{definition}[Weight excluding morphisms \cite{endrullis2024generalized}]
   %     \label{def:weight_excluding}
   %     Let \(\Gamma \subseteq \operatorname{Hom}(A, G)\). Define:
   %     \[
   %     \Gamma' \overset{\text{def}}{=} \left\{ \iota \in \operatorname{Hom}(X, G) \,\middle|\, \exists \alpha \in \Gamma, \exists \zeta : X \to A, \, \zeta \star \alpha = \iota \right\}.
   %     \]
   %     The \textbf{weight of \(h : G \to T\) excluding \(\Gamma'\)} is:
   %     \[
   %     w_e(h - \Gamma) \overset{\text{def}}{=} \bigodot_{\substack{\iota \in \{- \star h = e\} \\ \iota \notin \Gamma'}} w(e), \quad
   %     w_\mathcal{T}(h - \Gamma) \overset{\text{def}}{=} \bigodot_{e \in \mathbb{E}} w_e(h - \Gamma).
   %     \]
   % \end{definition}
       If \( \Gamma \) is a singleton \( \{ \alpha \} \), we denote \( w_{\mathcal{T}_\Sigma^X}(h - \alpha) \) instead of \( w_{\mathcal{T}_\Sigma^X}(h - \{ \alpha \}) \).
    %         \subsubsection{Estimating Weight Changes}
    %         In this section, we recall the concepts of traceability, relative monicity and weighable pushout square. These concepts are used to constrain \(T\)-valued elements and the DPO rewriting framework.  
For further details, we refer the reader to \cite[Section 4.1 and 4.2]{endrullis2024generalized}.
\begin{definition}[Traceability~\text{\cite[\textdef~4.1]{endrullis2024generalized}}]
    \label{def:traceability}
    \ \newline
\noindent 
\begin{minipage}{0.5\textwidth}
An object $X \mathop{\in} \mathcal{C}$ is called \emph{traceable} along a class of pushout squares $\Delta$ if for every diagram $\delta \mathop{\in} \Delta$ (shown on the right), the following hold
\end{minipage}
\begin{minipage}{0.5\textwidth}
\begin{center}
    \begin{tikzpicture}[node distance=12mm]
      \node (A) {$A$};
      \node [right of=A] (B) {$B$}; 
      \node [below of=A] (C) {$C$}; 
      \node [below of=B] (D) {$D$}; 
      \node [right of=D] (X) {$X$};
      \begin{scope}[nodes=rectangle]          
      \draw [->] (A) to node [above,label,pos=0.5] {$\alpha$} (B);
      \draw [->] (A) to node [left,label,pos=0.5] {$\beta$} (C);
      \draw [->] (B) to node [right,label,pos=0.45] {$\beta'$} (D); 
      \draw [->] (C) to node [below,label,pos=0.45] {$\alpha'$} (D);
      \draw [->] (X) to node [below,label,pos=0.4] {$f$} (D);
      \end{scope}
      \node at ($(A)!.5!(D)$) {$\delta$};
    \end{tikzpicture}
\end{center}
\end{minipage}
    \begin{enumerate}[label=(\alph*)]
        \item\label{traceable:a} there exists a morphism $h : X \mathop{\to} B$ such that $f \mathop{=} h \mathop{\star} \beta'$, or
        \item\label{traceable:b} there exists an $h : X \mathop{\to} C$ such that $f \mathop{=} h \mathop{\star} \alpha'$.
    \end{enumerate}
    If additionally, 
    \textcolor{red}{whenever \ref{traceable:a} and \ref {traceable:b} hold,} \todo{This starts to be quite hard to follow without intuition or example}
    \begin{enumerate}[label=(\alph*),resume]
        \item 
        there exists a morphism $h : X \mathop{\to} A$ such that $f \mathop{=} h \mathop{\star} \alpha \mathop{\star} \beta' $,
    \end{enumerate}
    then we say that $X$ is \emph{strongly traceable} along $\Delta$.
\end{definition}
\begin{definition}[Relative Monicity~\text{\cite[\textdef~4.6]{endrullis2024generalized}}]
    Let $A,B,C$ and $X$ be objects, $\Gamma$ a set of objects, $S$ a set of morphisms to $A$ and $u:C\to A$ a morphism. A morphism $f : A \to B$ is called 
    \begin{enumerate}[label=(\roman*)] 
        \item 
            \emph{monic for $S$} 
            if $g \star f = h \star f$ implies $g = h$ for all $g, h \in S$;
        \item 
            \emph{$X$-monic} if $f$ is monic for $\operatorname{Hom}(X, A)$.
        \item \emph{$X$-monic outside of $u$}, if $f$ is monic for \( \operatorname{Hom}(X,A) \setminus \left ( \operatorname{Hom}(X,C) \star u \right ) \).
        \item  \emph{$\Gamma$-monic} if $f$ is $X$-monic for every $X \in \Gamma$.
    \end{enumerate}
\end{definition} 
\begin{definition}[Weighable Pushout Square~\text{\cite[\textdef~4.9]{endrullis2024generalized}}]
    \label{def:weighable}
    \ \newline
    \noindent
    \begin{minipage}{0.7\textwidth}
        Let  $\mathcal{T} = (T,\mathbb{E}, S, w)$ be a finitary weighted type graph.
        Consider the pushout square $\delta$ shown on the right. We say that $\delta$ is
    \end{minipage}
    \begin{minipage}{0.3\textwidth}
        \begin{center}
            \begin{tikzpicture}[node distance=12mm]
                \node (A) {$A$};
                \node (B) [right of=A] {$B$};
                \node (C) [below of=A] {$C$};
                \node (D) [right of=C] {$D$};
                \draw [->] (A) to node [above, label] {$\alpha$} (B);
                \draw [->] (A) to node [left, label] {$\beta$} (C);
                \draw [->] (B) to node [right, label] {$\beta'$} (D);
                \draw [->] (C) to node [below, label] {$\alpha'$} (D);
                \node [at=($(A)!.5!(D)$)] {$\delta$};
            \end{tikzpicture}
        \end{center}
    \end{minipage}
    \begin{enumerate}[label=(\alph*)]
        \item \emph{weighable} with $\mathcal{T}$ if the following hold:
            \begin{enumerate}[label=(\roman*)]
                \item $dom(\mathbb{E})$ is strongly traceable along $\delta$,
                \item $\beta'$ is $dom(\mathbb{E})$-monic,
                \item $\alpha'$ is $dom(\mathbb{E})$-monic outside of $\beta$.
            \end{enumerate}
        \item \emph{bounded-above} by $\mathcal{T}$ if $dom(\mathbb{E})$ is traceable along $\delta$.
    \end{enumerate}
\end{definition}
The existence of a context closure guarantees the existence of morphisms from any object subject to rewriting to the type graph—a critical requirement as emphasized in \hyperref[remark:semiring_0_unpredictable]{\textsection~\ref{remark:semiring_0_unpredictable}}. Further details can be found in~\cite[\textsection 5.3]{endrullis2024generalized}.
\begin{definition}[Context Closure~\text{\cite[\textdef~5.3]{endrullis2024generalized}}]
    \label{def:context_closure} 
    \ \newline 
\begin{minipage}{0.65\textwidth}
    Let $\mathcal{T} = (T,\mathbb{E}, S, w)$ be a type graph, \(\rho = (L \overset{l}{\leftarrow} K \overset{r}{\rightarrow} R ) \) a DPO rewriting rule and $\mathfrak{F}$ a rewriting framework. 
    A \textbf{context closure} for $\rho$ and $\mathcal{T}$ in $\mathfrak{F}$ is a morphism $c:L \rightarrow T$ such that for every DPO diagram in $\mathfrak{F}(\rho)$ (shown on the right) 
    there exists $\alpha : G \rightarrow T$ such that $m \star \alpha = c$.
\end{minipage}
\begin{minipage}{0.35\textwidth}
    \begin{center}
        \begin{tikzpicture}[rotate=90]
          \node (I) {$K$}; 
          \node (L) [left of=I] {$L$};
          \node (R) [right of=I] {$R$};
          \node (G) [below of=L] {$G$};
          \node (C) [below of=I] {$C$};
          \node (H) [below of=R] {$H$};
          \node (T) [left=of $(L)!0.5!(G)$] {$T$};
          \draw [->] (L) to  node [label, above] {$c$}  (T);
          \draw [->] (G) to  node [label, below] {$\alpha$} (T);
          \draw [->] (I) to node [label, above] {$l$} (L);
          \draw [->] (I) to node [label,above] {$r$} (R);
          \draw [->] (L) to node [label, right] {$m$} (G);
          \draw [->] (I) to (C);
          \draw [->] (R) to (H);
          \draw [->] (C) to (G);
          \draw [->] (C) to (H);
        \end{tikzpicture}
      \end{center}
\end{minipage}
\end{definition}
    %         % \subsection{Decreasing Rule}
    %         \todo{to be weakened to the original version by endrullis}
\todo{maybe add the core lemma for estimating pushout object weight, but adapt to estimating DPO objects weight difference}
\begin{definition}[Decreasing Rules~\text{\cite[\textdef~5.9]{endrullis2024generalized}}]
    \label{def:decreasing_rule}
    Let $\mathcal{T} = (T,\mathbb{E},\mathcal{S}, w)$ be a finitary weighted type graph where $\mathcal{S}=\langle S, \oplus, \odot, 0, 1, \prec, \leq \rangle$ is a well-founded, commutative, \(\mathfrak{F}\) a DPO rewriting framework, $\rho = (L \overset{l}{\leftarrow} K \overset{r}{\rightarrow} R)$ a DPO rewriting rule. 

    \noindent
    The rule $\rho$ is \textbf{weakly decreasing} w.r.t. $\mathcal{T}$ in $\mathfrak{F}$ if 
            for every $t_K : K \to T$,
                $$ 
                  w_\mathcal{T}(\{l \star - = t_K\}) \succeq w_\mathcal{T}(\{r\star - = t_K\})$$
           
    \noindent
    The rule $\rho$ is \textbf{uniformly decreasing} w.r.t. $\mathcal{T}$ in $\mathfrak{F}$ if
        \begin{itemize}
            \item[]- there exists a context closure $c_\rho$ for $\rho$ and $\mathcal{T}$ in $\mathfrak{F}$, and
            \item[]- for every $t_K : K \to T$,
            \begin{itemize}
                \item[] $\bullet$ $\{l \star - = t_K\} = \emptyset = \{r \star - = t_K\}$, or
                \item[] $\bullet$ $w_\mathcal{T}(\{l \star - = t_K\}) \succ    w_\mathcal{T}(\{r \star - = t_K\}) + \delta$.
            \end{itemize}
        \end{itemize}  
         
    \noindent
    The rule $\rho$ is
            \textbf{$\delta$-closure decreasing} w.r.t. $\mathcal{T}$ in $\mathfrak{F}$ if the following hold:
            \begin{itemize}
                \item[]- $S$ is strictly monotonic semiring,
                \item[]- $\rho$ is weakly decreasing,
                \item[]- there exists a context closure $c_\rho$ for $\rho$ and $\mathcal{T}$ in $\mathfrak{F}$,
                \item[]- $w_\mathcal{T}(\{l \star - = t_K\}) \succ  w_\mathcal{T}(\{r \star - = t_K\})$ for $t_K = l \star c_\rho$.
            \end{itemize}
\end{definition}

The following lemma states that decreasing rules reduce the weights of host graphs, provided specific constraints are satisfied.
\begin{lemma}[Decreasing steps~\text{\cite[Theorem C.3]{endrullis2024generalized}}]
    \label{lem:decreasing_step}
    \ \newline 
\begin{minipage}{0.7\textwidth}
    Let $\mathcal{T} = (T,\mathbb{E}, \mathcal{S}, w)$ be a finitary weighted type graph where $\mathcal{S} = \langle S, \oplus, \odot, 0, 1, \prec, \leq \rangle$ is a strong monotonic, commutative semiring, $\rho$ a rewriting rule and $\Delta \in \mathfrak{F}(\rho)$ a DPO diagram
    (shown on the right)   such that the following conditions hold:
\end{minipage}  
\begin{minipage}{0.3\textwidth}
    \begin{center}
        \begin{tikzpicture}[node distance=11mm]
          \node (I) {$K$};
          \node (L) [left of= I] {$L$};
          \node (R) [right of=I] {$R$}; 
          \node (G) [below of=L] {$G$};
          \node (C) [below of=I] {$C$};
          \node (H) [below of=R] {$H$};
        %   \node (T) [left=of $(L)!0.5!(G)$] {$T$};
        %   \draw [->] (L) to  node [label, above] {$c$}  (T);
        %   \draw [->] (G) to  node [label, below] {$\alpha$} (T);
          \draw [->] (I) to node [label, below] {$l$} (L);
          \draw [->] (I) to node [label, below] {$r$} (R);
          \draw [->] (L) to  (G);
          \draw [->] (I) to (C);
          \draw [->] (R) to (H);
          \draw [->] (C) to (G);
          \draw [->] (C) to (H);
        \end{tikzpicture}
      \end{center}
\end{minipage}
   \begin{itemize}
       \item $\operatorname{left}(\Delta)$ is weighable with \(\mathcal{T}\), and
       \item $\operatorname{right}(\Delta)$ is bounded above by \(\mathcal{T}\), and
       \item $w(e) \succeq 1_S$ for all $e \in \mathbb{E}$.
   \end{itemize}

   \noindent
  We have:
   \begin{itemize}
       \item $w_\mathcal{T}(G) \succeq w_\mathcal{T}(H)$ if $\rho$ is weakly decreasing,
       \item $w_\mathcal{T}(G) \succ w_\mathcal{T}(H)$ if $w(e) \succeq 1_S$ for all $e \in \mathbb{E}$, and $\rho$ is uniformly or closure decreasing.
   \end{itemize}
\end{lemma} 
\begin{proof}
   See the \hyperref[proof:decreasing_step]{Appendix}.
\end{proof} 
    %         \subsubsection{Termination Criterion} 
    %         \todo{to be weakened to the original version by endrullis}
Finally, a termination criterion can be established.
\begin{theorem}[Termination~\text{\cite[Theorem 5.11]{endrullis2024generalized}}] 
    \label{thm:termination_grs}
    Let $\mathcal{A}$ and $\mathcal{B}$ be sets of DPO rewriting rules, $\mathcal{T} \mathop{=} (T,\mathbb{E}, \mathcal{S}, w)$ be a finitary weighted type graph where $\mathcal{S} \mathop{=} \langle S, \mathop{\oplus}, \mathop{\odot}, 0, 1, \prec, \leq \rangle$ is a strong monotonic, commutative semiring and $\mathfrak{F}$ a DPO rewriting framework such that

     \begin{enumerate}
        \item\label{thm1:hyp3} $w(e) \mathop{\succeq} 1_S$ for all $e \mathop{\in} \mathbb{E}$,
        % \item\label{thm1:hyp4} $\{s \mathop{\in} S\mid 1_S \leq s \mathop{\neq} 0_S\} \mathop{\subseteq} \mathbb{R}_{>0}$ 
        % \item\label{thm1:hyp4} for all $x \mathop{\in} S$, if $ 1_S \mathop{\preceq} x \mathop{\neq} 0_S$ then $\mu(x) \mathop{\geq} \mu(1_S)$ and $\mu(x) \mathop{\in} \mathbb{R}$,
        % \item\label{thm1:hyp4} for all $x \mathop{\in} S$, if $ 1_S \mathop{\preceq} x \mathop{\neq} 0_S$ then $\mu(x) \mathop{\geq} \mu(1_S)$ and $\mu(x) \mathop{\in} \mathbb{R}$,
        \item for every rule $\rho \mathop{\in} (\mathcal{A }\mathop{\cup} \mathcal{B })$ and every double pushout diagram  
        $\Delta \mathop{\in} \mathfrak{F}(\rho)$ 
        \begin{itemize}
            \item \(\operatorname{left}(\Delta)\) is weighable with \(\mathcal{T}\), and
            \item \(\operatorname{right}(\Delta)\) is bounded above by \(\mathcal{T}\). 
        \end{itemize}
    \end{enumerate}       

    \noindent If the following conditions hold:
    \begin{enumerate}
        \item either every $\rho \mathop{\in} \mathcal{A}$ is uniformly decreasing, or every $\rho \mathop{\in} \mathcal{A}$ is closure decreasing,
        \item every rule $\rho \mathop{\in} \mathcal{B}$ is weakly decreasing,
    \end{enumerate}
    then $\mathop{\Rightarrow}_{\mathcal{A},\mathfrak{F}}$ is \textbf{terminating} relative to $\mathop{\Rightarrow}_{\mathcal{B},\mathfrak{F}}$.
\end{theorem} 
\begin{proof}
    See the 
    % \hyperref[proof_termination_grs]{Appendix}.
    ref
\end{proof}
    %     \subsection{Subgraph Counting}  
    %     The subgraph-counting method in 2024 lmcs by Overbeek and Endrullis is designed for the PBPO+ 2023 pbpo+ ref —a rewriting formalism capable of simulating left-injective DPO rewriting.
    %     \subsection{Termination Criterion for DPO GRS with NAC}
    %     \begin{definition}
    Let $V_L$ be the set of nodes in $L$, and $M_L \mathop{=} \{m_1, \dots, m_r\}$ the set of matches $m_i : V_L \mathop{\to} V_L$ of $V_L$ into itself, including the identity $id_{V_L} \mathop{=} m_1$. We construct a Labelled Transition System $\mathcal{L}^p \mathop{=} (S, \Lambda, \longrightarrow)$ as follows:

    \begin{enumerate}
        \item $S$ contains a state $s_i$ for each graph 
        $H_i^p \mathop{\in} \mathfrak{H}_p^p$. 
        % $H_i^p \mathop{\in} \mathscr{H}$.
         Each $s_i$ induces a classification function $c_i$ for matches of $p$ such that $c_i(m_l) \mathop{=} \text{true}$ iff $m_l$ can be extended to a match on $H_i^p$, but not to a match on any other 
         $H_j^p \mathop{\in} \mathfrak{H}_p^p \mathop{\setminus} (\{H_1^p, H_i^p\} \mathop{\cup} \{H_t^p \mid \exists h_t^i : H_t^p \mathop{\to} H_i^p\})$,
          i.e., $H_i^p$ is the biggest graph to which $m_l$ can be extended.
        
        \item $\Lambda$ contains a label $p^i$, $i \mathop{=} 1, \dots, r$ for each morphism in $M_L$.
        
        \item The transition relation $\longrightarrow \mathop{\subseteq} S \mathop{\times} \Lambda \mathop{\times} S$ is such that $(s_i, p^l, s_j) \mathop{\in} \longrightarrow$ (denoted by $s_i \xrightarrow{p^l} s_j$) if applying $p$ with match $m_l$ on graph $H_i^p$ produces a graph for which $c_j(m_l) \mathop{=} \text{true}$.
    \end{enumerate}
\end{definition} 

\begin{definition}
    In order to consider the variation on the number of possible matches induced by the application of $p$ on the minimal context for a given state, let $@|\cdot|: S \mathop{\to} \mathbb{N}$ be a function which associates with each state $s_i$ the number of matches for $p$ on the graph $H_i^p$ prior to application of $p$, and with $|\cdot|: S \mathop{\to} \mathbb{N}$ the function defining the number of matches on $H_i^p$ after the application of $p$.
\end{definition}

\begin{definition}
    Let $\rho$ be a graph rewriting rule. We say that $\rho$
    \begin{itemize}
        \item \emph{must terminate} if it should terminate simply and for all transitions $s_i \xrightarrow{p_l} s_j$ and all states $s_h \mathop{\neq} s_i, s_j, s_k$, we have $\left| s_i \right| < @ | s_i |$, $\left| s_j \right| \leq @ | s_j |\mathop{+}1$ and $\left| s_h \right| \leq @ | s_h |$.
        \item \emph{may terminates} if it may terminate simply and for all states on the path the same condition on matches as above applies.
        \item \emph{does not terminate} if (it does not terminate simply AND for all states from which $s_k$ is reachable, there is a state for which the number of matches increases for some transition leading to $s_k$ increases) OR (it should or may terminate simply, but there is at least one state $s_i$ on a path from $s_1$ to $s_k$ for which $\left| s_i \right| \mathop{\geq} @ | s_i |$, for a transition reachable from state $s_i$).
    \end{itemize}
\end{definition}

\begin{theorem}[Termination~\text{\cite[Theorem~2]{bottoni2010atermination}}]
    \todo{rewrite}
    Let $@|\cdot|: S \mathop{\to} \mathbb{N}$ and $|\cdot|: S \mathop{\to} \mathbb{N}$ be counting functions as defined above, let $p$ be a rule, and let $G$ be a finite graph. Then the following holds:
    \begin{enumerate}
        \item If rule $p$ is of type \textbf{must terminate}, then the application of \texttt{asLongAsPossible $p$ end} on the starting graph $G$ terminates after a finite number of steps.
        
        \item If $p$ is of type \textbf{does not terminate}, then the application of \texttt{asLongAsPossible $p$ end} on the starting graph $G$ does not terminate.
    \end{enumerate}
\end{theorem}   
    \section{Morphisms of Graph Rewriting Systems}
        \label{sec:morphisms_from_dpo_to_pbpop}
        \subsection{Morphism from left-injective DPO GRS with injective match to PBPO+}
        \label{sec:morphism_from_dpo_grs_to_pbpop}
        \begin{example}[ \cite{overbeek2023pbpo_JLAMP}]
            In the category of unlabeled directed graphs, there exists a mono-partial morphism classifier $(T,\eta)$ where 
            \begin{itemize}
                \item for all unlabeled graph $G=(V,A,s,t)$, we have $T(G) \mathop{=} (V_*,E_*,s_*,t_*)$ where $V_* \mathop{=} V \uplus \{*\}, E_* \mathop{=} E \uplus (V_* \mathop{\times} V_x), s_*(e) \mathop{=} s(e)$ if $e \mathop{\in} E$ and $\pi_1(e)$ otherwise, and $t_*(e) \mathop{=} t(e)$ if $e \mathop{\in} E$ and $\pi_2(e)$ otherwise.
            \end{itemize}
            An example is given by 
            \begin{center}
                $G \; \mathop{=} $
                \begin{tikzcd}
                v \arrow[loop, distance=2em, in=215, out=145] \arrow[rr] &  & w
                \end{tikzcd}
                \quad and \quad 
                $T(G) \; \mathop{=} \hspace{-3mm} $ % https://tikzcd.yichuanshen.de/#N4Igdg9gJgpgziAXAbVABwnAlgFyxMJZABgBpiBdUkANwEMAbAVxiVpAF9T1Nd9CUAJnJVajFmwDunbiAzY8BIgEZSy0fWatEIADq6ARhBydRMKAHN4RUADMAThAC2SMiBwRX1BhAhoiggDsZLaMcDCiDHQGMAwACryKAiD2WBYAFiZcdo4uiG4eSMrZIA7OXu6eiMIgDFhg2iBQxjjmINQxYFBIALQAzMQlZXk1hYiqtfWNzTit3UO5SKNVNT5+RCFhEd7RsQkK-GypGSbeU2wzczI55ePUYxNr-igAnJsM4ZG78YmHOseZdqTBoXFptBa3ApVNxPALBUihD7bWrffZ8JT-NKAs4gnSXcGyYYVB4dGBdJADHHTMHzQmLO6VCqdbqISnA6mzAk3Eb3aFU0Gc2ncoq8pak8msmHnPE00wcIA
                \begin{tikzcd}
                v \arrow[loop, distance=2em, in=215, out=145] \arrow[rr] \arrow[rd, dotted, bend right] \arrow[dotted, loop, distance=4em, in=240, out=140, looseness=3] \arrow[rr, dotted, bend left] &                                                                                                & w \arrow[dotted, loop, distance=2em, in=35, out=325] \arrow[ll, dotted, bend left] \arrow[ld, dotted, bend left] \\
                                                                                                                                                                  & \mathop{\star} \arrow[ru, dotted] \arrow[dotted, loop, distance=2em, in=305, out=235] \arrow[lu, dotted] &                                                                                                       
                \end{tikzcd},
            \end{center}
        
            where the dotted edges represent the edges $e \mathop{\in} V_\star \mathop{\times} V_\star$.
            It can be seen that for any partial homomorphism $\psi: H \mathop{\to} G$ defined on subgraph $H' \mathop{\subseteq} H$, there exists exactly one homomorphism $\psi_\star : H \mathop{\to} T(G)$ such that $\psi_\star(x) \mathop{=} \psi(x)$ for $x \mathop{\in} V_{H'} \mathop{\cup} E_{H'}$ and $\psi_\star(x) \notin V_{G} \mathop{\cup} E_{G}$ for $x \notin V_{H'} \mathop{\cup} E_{H'}$. Equivalently, $\psi_\star$ is the unique morphism $\varphi_{m,f} : H \mathop{\to} T(G)$ making
                \begin{center}
                    % https://tikzcd.yichuanshen.de/#N4Igdg9gJgpgziAXAbVABwnAlgFyxMJZABgBpiBdUkANwEMAbAVxiRAA0QBfU9TXfIRQBGclVqMWbAELdeIDNjwEiZYePrNWiEAEE5fJYKKj11TVJ0AVABTSAlN3EwoAc3hFQAMwBOEALZIoiA4EEgAzNQ4dFgMbAAWEBAA1iDmktogADpZMNEA+rLUDHQARjAMAAr8ykIgWGDYsAYgvgFIZCFhiABMUTFxOokpaRJabIHFZRXVRio6DU2sPN5+gYidoUHp4zpeoyXlVTXGC41YzSuta0h9XRE7ltlZaHQ+AMYlcHDA-lzAXi4B2mxzmdUWF2WFC4QA
                    \begin{tikzcd}[column sep=15mm]
                    H' \arrow[d, "m" description, hook] \arrow[r, "f" description] & G \arrow[d, "\eta_G" description, hook] \\
                    H \arrow[r, "\varphi_{m,f}" description]                    & T(G)                                   
                    \end{tikzcd}
                \end{center}
            a pullback square, where $H \stackrel{m}{\hookleftarrow} H' \stackrel{f}{\to} G$ a partial map span representation of $\psi$, and $\eta_G$ and $m$ are inclusions.
            
            The generalization to labeled graphs is straightforward: between any two nodes $u,v \mathop{\in} V_\star$ and for every label $l$, there is one $l$-labeled edge representing an undefined $l$-edge between $u$ and $v$.
        \end{example}
        
        The following theorem is an instance of \cite[Definition 71]{overbeek2023pbpo_JLAMP}.
        \begin{definition}
            Let $L \overset{l}{\leftarrowtail} K \overset{r}{\rightarrowtail} R$ be a DPO graph rewriting rule. The diagram depicted below where the left square a pushout (which is also a pullback) is a $\mathbf{PBPO}^+$ rule, denoted $|\rho|$.
        \[ 
        \begin{tikzcd}
        L \arrow[r, "l"] & K \arrow[r, "r"] & R \\
        L' \arrow[u, "\eta_L", hook] \arrow[r, "l'"'] & T(K) \arrow[u, "\eta_K"', hook]
        \end{tikzcd}
        \]
        
        
        \end{definition} 
        
        The following theorem is an instance of \cite[Theorem 72]{overbeek2023pbpo_JLAMP}.
        \begin{theorem}
            In $\mathbf{Graph}$.
            
            let $\mathrm{DPO}$ be the DPO graph rewriting system with injective matches.
        
            For any left-injective DPO rule $\rho$, $|\rho|$ is a well-defined $\mathrm{PBPO}^+$ rule. 
            
            For any $G \mathop{\to} ^{\rho}_{\mathrm{DPO}} H$, we have $G \mathop{\to} ^{|\rho|}_{\mathrm{PBPO}^+} H$.  
        \end{theorem}
        \begin{proof}
            \todo{todo}
        \end{proof}
        
\chapter{Homomorphisms}
to do
    \section{Homomorphism from GLS rewriting relations to DPO graph rewriting relations}
    \label{sec:morphism_from_gls_to_dpo_grs}
    % \begin{theorem}
%     Let $\mathcal{T}_{gls \mathop{\to} dpo}$ be the functor from the category of Graph relabeling systems to the category of DPO Graph rewriting systems which assocates $\set{\rho' \mathop{\mid} \rho \mathop{\in} \mathcal{R}}$.
% \end{theorem}

\begin{definition}
    \label{def:gls_to_dpo}
    Let $\rho \mathop{=} ((R,\lambda) \mathop{\rightarrow} (R,\lambda'))$ be a graph relabeling rule.
    We define $|\rho|$ as the following injective DPO rewriting rule 
    $$  ((R,\lambda) \overset{l}{\leftarrowtail} (K, \emptyset) \overset{r}{\rightarrowtail} (R,\lambda'))$$ 
    where $K$ is the discrete graph obtained by removing all edges from $R$, and $l : (K,\emptyset) \rightarrowtail (R,\lambda)$ and $r : (K, \emptyset) \rightarrowtail (R,\lambda')$ are the inclusion functions.
\end{definition}
For any graph relabeling rule $\rho$, $|\rho|$ is a well-defined DPO rewriting rule.

\begin{theorem}
    Let $\rho$ be a graph relabeling rule. 
    If there is a graph relabeling step from a graph $G$ to a graph $H$ using the rule $\rho$, then there is a DPO rewriting step from $G$ to $H$ using the rule $|\rho|$.
\end{theorem}
\begin{proof}
    Suppose that there is a graph relabeling step from a graph $G$ to a graph $H$ using the rule $\rho$.
    By definition, there is an occurrence of $(R,\lambda)$ in $G$. Therefore, there is a match $m$ of DPO rule $|\rho|$ in $G$ such that $\opn{Im}(m)$ is the occurrence. 
    There is a DPO rewriting step from $G$ using rule $|\rho|$ and match $m$ to a graph $H'$ where $H'$ is $G$ with the occurrence of $(R,\lambda)$ replaced by an occurrence of $(R,\lambda')$. 
    Since $H$ is obtained from $G$ by replacing the labels of the occurrence of $(R,\lambda)$ by the labels of $(R,\lambda')$, we have $H' \mathop{=} H$.
    We conclude that there is a DPO rewriting step from $G$ to $H$ using the rule $|\rho|$.
\end{proof}

% \begin{proposition}[Termination Preservation]
%     Let $\mathcal{R}$ be a graph relabeling system, let $\mathcal{R}' \isdef \set{\rho' \mathop{\mid} \rho \mathop{\in} \mathcal{R}}$ the resulting DPO graph rewriting system. If $(G,\lambda) \mathop{\to} _\mathcal{R} (G,\lambda')$ then $(G,\lambda) \mathop{\to} _\mathcal{R} (G,\lambda')$.
% \end{proposition} 

% \begin{theorem}
%     Let $\mathcal{R}$ be a graph relabeling system. If the resulting DPO graph rewriting system terminates then  $\mathcal{R}$ terminates.
% \end{theorem}

% A graph relabelling system, as described in \cite{Litovsky1999gls}, is a rule-based graph transformation system of \textcolor{red}{single graphs}. A graph relabelling rule consists of a couple of simple graphs $(L,R)$ and an isomorphism between their underlying unlabelled graphs.
  
  
  
% Given a graph,  the graph transformation consists in finding an \todo{occurrence} of $L$ and then replacing the labels on the occurrence by corresponding labels on $R$ according to the isomorphism. 
% Graph relabeling systems are very intuitive and many distributed algorithms can be modelled in this way.
% A graph relabelling system is a graph rewriting system which repalces arrows with arrows of the same extremities.
% \begin{definition}{\cite{Litovsky1999gls}}
%    A graph rewriting rule $(l: K \rightarrowtail L,r : K \rightarrowtail R)$ is said to be a \textbf{graph relabelling rule} if
%    \begin{itemize}
%        \item $K$ is discrete graph,
%        \item $L$ and $R$ are simple graphs,
%        \item there is an isomorphism $\varphi :L' \mathop{\to} R'$ such that $l \mathop{\star} \varphi \mathop{=} r$ and $r \mathop{\star} \varphi^{-1} \mathop{=} l$ where $L'$ (reps. $R'$) is the underlying unlabelled graph of $L$ (reps. $R$).
%    \end{itemize}  
% \end{definition}

    \section{Homomorphism from FCGLS rewriting relations to rewriting relations of DPO graph rewriting systems with Negative Application Conditions}
    \label{sec:morphism_from_pgls_to_nac_dpo_grs}
    When a graph relabeling system consists of one rule and the rule has a unique forbidden context, we can transform the graph relabeling system into a DPO graph rewriting system with one negative application condition. For the latter, the termination technique that Bottoni et al.~\cite{bottoni2010atermination} proposed applies. 
\begin{definition}
    Let $\rho$ be a graph relabeling rule with a forbidden context $n:\opn{lhs}(\rho) \rightarrowtail N$.
      We define $|(n,\rho)|$ as the DPO rewriting rule with a negative application condition $(n, |\rho|)$
    where $|\rho|$ is the DPO rewriting rule defined in \autoref{def:gls_to_dpo}.
\end{definition}  

\begin{theorem}
    \label{prop:fcgls_to_dpogrs_with_one_nac}
   Let $\rho$ be a graph relabeling rule with a forbidden context $n:\opn{lhs}(\rho) \rightarrowtail N$.
    For any graph relabeling step from $G$ to $H$ using rule $\rho$ with forbidden context $n$, there is a DPO rewriting step with a negative application condition from $G$ to $H$ using the rule $(n, |\rho|)$.
\end{theorem}
The proof is straightforward by definition.
    \section{Homomorphism from DPO graph rewriting relations to DPO hypergraph rewriting relations}
    \label{sec:morphism_from_dpo_grs_to_hypergraph}
    \begin{definition}
    Let $(G,\lambda)$ be a graph over edge label set $\Sigma_E$.
    We define the signature $\Sigma$ of the hypergraph as
    $(\{\epsilon\}, \Sigma_E, \text{Type})$ where $\text{Type}(l) \mathop{=} \{\epsilon\epsilon\}$ for all $l \mathop{\in} \Sigma_E$. 
    Furthermore, we define $|(G,\lambda)|$ as the hypergraph $(V_G, E_G,\opn{mark}, \opn{lab},\text{att})$ over the signature $\Sigma$ where
    \begin{itemize}
        \item $\opn{mark}(v) \mathop{=} \epsilon$ for all $v \mathop{\in} V_G$,
        \item $\opn{lab}(e) \mathop{=} \lambda(e)$ for all $e \mathop{\in} E_G$,
        \item $\text{att}(e) \mathop{=} st \mathop{\in} V_G^*$ for all $e \mathop{\in} E_G$ with source $s$ and target $t$,
    \end{itemize}
\end{definition}

\begin{example}
    We show a graph $G$, on the left,and its hypergraph $|G|$,on the right.

    \begin{tikzpicture}[scale=0.75]
        \node[draw,circle] (l1) at (-3,3) {1};
        \node[draw,circle] (l2) at (3 ,3) {2};
        \node[draw,circle] (l3) at (0,0) {3};
        \draw[->] (l1) -- (l3) node[midway,left] {$s$};
        \draw[->] (l2) -- (l3) node[midway,right] {$s$};
        \draw[->] (l3) edge [loop below] node {0} (l3);
    \end{tikzpicture}
    \hspace{1cm}
    \begin{tikzpicture}[scale=0.75]
        \node[draw,circle] (n1) at (-3,3) {1};
        \node[draw,circle] (n2) at (3,3) {2};
        \node[draw,circle] (n3) at (0,0) {3};
        \node[draw,rectangle] (e1) at (-1.5,1.5) {$s$};
        \node[draw,rectangle] (e2) at (1.5,1.5) {$s$};
        \node[draw,rectangle] (e3) at (0, -1.5) {0};
        \draw (e1) -- (n1) node[midway,left] {1};
        \draw (e1) -- (n3) node[midway,right] {2};
        \draw (e2) -- (n2) node[midway,left] {1};
        \draw (e2) -- (n3) node[midway,right] {2};
        \draw (e3) -- (n3) node[midway,left] {1};
    \end{tikzpicture}
\end{example}

Since a graph homomorphism is a hypergraph homomorphism, we transform graph homomorphisms into hypergraph homomorphisms in a trivial way.
\begin{definition}
    Let $h: A \rightarrowtail B$ be a graph homomorphism. We define the hypergraph homomorphism $|h| : |A| \rightarrowtail |B|$ such that for all $v \mathop{\in} V_A$, $|h|(v) \mathop{=} h(v)$, for all $e \mathop{\in} E_A$, $|h|(e) \mathop{=} h(e)$.
\end{definition}

\begin{definition}
    Let $\rho \mathop{=} (L \overset{l}{\leftarrowtail} K \overset{r}{\rightarrowtail} R)$ be a DPO graph relabeling rule. We define the hypergraph rewriting rule $|\rho|$ as $|L| \overset{|l|}{\leftarrowtail} |K| \overset{|r|}{\rightarrowtail} |R|$. 
\end{definition}


\begin{theorem}
    Let $\rho$ be a DPO graph rewriting rule.
    For any DPO rewriting step from a graph $G$ to a graph $H$ using the rule $\rho$, there is a DPO hypergraph rewriting step from $|G|$ to $|H|$ using the rule $|\rho|$.
\end{theorem}
The proof is straightforward by definition.
    % \chapter{Morphism from FCGLS to PBPO+ GRS: todo maybe } 
    % 


\begin{corollary}
    Let $\mathcal{R}$ be a graph relabeling system(GLS). There is a PBPO+ graph rewriting system $\mathcal{R}'$ and a homomorphism $h: \mathcal{R}\to \mathcal{R}'$.
\end{corollary}
\begin{proof}
    by~\autoref{thm:dpo_instanceOf_pbpo}
\end{proof}

Let $\mathcal{R}$ be a FCGLS. Let $\rho = ((G,\lambda) \to (G,\lambda'), \set{C_1,...,C_n}) \in \mathcal{R}$. We can translate this graph relabeling rule with forbidden contexts to a PBPO+ graph rewriting:



\part{Termination techniques for DPO graph rewriting systems}
\chapter{Termination of Graph Rewriting using Weighted Type Graphs over Non-well-founded Semirings}
\label{part:morphisms_from_gls_to_grs}

  
\chapter{Termination of Injective DPO Graph Rewriting Systems using Subgraph Counting}
 \label{part:terminating_injective_dpo_grs}
 \begin{itemize}
    \item This part of the thesis is dedicated to the study of terminating injective DPO graph rewriting systems.
    \item Since~\autoref{part:morphisms_from_gls_to_grs} has introduced morphisms from graph relabeling systems to injective DPO graph rewriting systems on graphs, the terminating methods developed in this part can be applied to graph relabeling systems.
    \item This part is organized as follows: todo todo \todo{todo}
\end{itemize}
 
\section{Introduction}
\label{sec:intro}

In this chap, we present an new automated method for relative termination~\cite{geser1990relative} of double-pushout (DPO) graph rewriting systems with injective rules~\cite{corradini1997algebraic,habel2001double,konig2018atutorial}. The method, based on the interpretation method framework~\cite{nipkow1998term}, works as follows. Let \( \mathbb{X} \) denote a fixed set of graphs. We assign weights (natural numbers) to monomorphisms from graphs in \( \mathbb{X} \), and define the weight of a graph $G$ as the sum of the weights of all monomorphisms from graphs in \( \mathbb{X} \) to $G$. 

Let $\mathcal{A}$ and $\mathcal{B}$ be sets of DPO rewriting rules. Let $X \in \mathbb{X}$ and $G \Rightarrow H$ be a rewriting step defined by the following double-pushouts diagram.
\begin{center}
 \resizebox{0.4\textwidth}{!}{
     \begin{tikzpicture}
            % [node distance=15mm]
            \node (I) at (0,0) {$K$};
            \node (L)  at (-2,0) {$L$};
            \node (R)  at (2,0) {$R$};
            \node (G)  at (-2,-2) {$G$};
            \node (C)  at (0,-2) {$C$};
            \node (H)  at (2,-2) {$H$};
            \draw [>->] (I) to  node [midway,below] {$l$} (L);
            \draw [>->] (I) to  node [midway,below] {$r$} (R);
            \draw [>->] (L) to node [midway,right] {$m$} (G);
            \draw [>->] (I) to  node [midway,right] 
            % {$u$}
            {} (C);
            \draw [>->] (R) to  node [midway,right] 
            {$m'$}
            (H);
            \draw [>->] (C) to node [midway,above] {$l'$} (G);
            \draw [>->] (C) to node [midway,above] 
            {$r'$} 
            (H);
            % \node [at=($(I)!.5!(G)$)] {\normalfont PO};
            % \node [at=($(I)!.5!(H)$)] {\normalfont PO};
        \end{tikzpicture}
    }
\end{center}

For a monomorphism from $X$ to $G$, its image falls into exactly one of the following three mutually exclusive cases. First, the image may be fully included in $m(L)$. Second, the image may be fully included in $l'(C)$; Third, the image may be neither fully included in $m(L)$ nor fully included in $l'(C)$. Similary, for a monomorphism from $X$ to $H$, its image falls into exactly one of the following three mutually exclusive cases. First, the image may be fully included in $m'(R)$. Second, the image may be fully included in $r'(C)$; Third, the image may be neither fully included in $m'(R)$ nor fully included in $r'(C)$. The number of monomorphisms from $X$ to $G$ whose image is fully included in $m(L)$ can be precisely calculated, as can the number of monomorphisms from $X$ to $H$ whose image is fully included in $m'(R)$. Moreover, the number of monomorphisms from $X$ to $G$ equals to the that of monomorphisms to $H$. 

To estimate the change in weight when applying a rewriting rule, we focus on 
the monomorphisms in the third case. Specifically, our technique requires that, for every graph $X$ in $\mathbb{X}$ and every rewriting step $G \Rightarrow H$ using a rule from $\mathcal{A}$ or $\mathcal{B}$, there exists an injective mapping from the set of monomorphisms $x : X \to H$ (where $\opn{Im}(x)$ is neither fully included in the right-hand side graph nor fully included in the context) to the set of monomorphisms from $y : X \to G$ (where $\opn{Im}(y)$ is neither fully included in the left-hand side graph nor fully included in the context). 

Suppose, additionally, that the following conditions hold: (i) For every rule in \( \mathcal{A} \), the left-hand side graph's weight is strictly greater than the right-hand side graph's weight; (ii) For every rule in \( \mathcal{B} \), the left-hand side graph's weight is greater than or equal to the right-hand side graph's weight. 

Under these conditions, rewriting steps using rules in \( \mathcal{A} \) strictly decrease the weights of the host graphs, while rewriting steps using rules in \( \mathcal{B} \) do not increase them.
Consequently, the rewriting rules in \( \mathcal{A} \) can only be applied finitely many times, since the weight of a finite graph is a natural number and strictly decreases with each rewriting step using a rule in \( \mathcal{A} \), the process must terminate after finitely many iterations.  
 
Our method proves termination for systems such as the one in~\autoref{subgraph_counting:ex_contrib_variant}, which prior interpretation-based approaches~\cite{zantema2014termination,bruggink2014termination,bruggink2015proving,
endrullis2024generalized_arxiv_v2,
overbeek2024termination_lmcs} cannot handle. 
We also provide an implementation of our technique.  
   
The remainder of this chapter is organized as follows.
We first define the weight of a graph in~\textsection~\ref{subgraph_counting:sec:interpretation}. 
Using this definition, we outline the general approach and identify the key challenge in~\textsection~\ref{subgraph_counting:sec:general_idea} and introduce the concept of non-increasing rules in \textsection~\ref{subgraph_counting:sec:non-increasing}. 
A solution to the challenge is then proposed in~\textsection~\ref{subgraph_counting:sec:solution_to_the_key_challenge}, culminating in our termination criterion in~\textsection~\ref{subgraph_counting:sec:termination}.
\textsection~\ref{subgraph_counting:sec:examples} illustrates the method with some examples.
\textsection~\ref{subgraph_counting:sec:related_work} compares our approach with some existing methods.
Finally,~\textsection~\ref{subgraph_counting:sec:conclusion} concludes with remarks and future research directions. The missing proofs are available in 
\iflongversion
the appendix.
\else
the preprint version~\cite{qiu2025termination}.
\fi

\section{Preliminaries}
\label{sec:pre} 
While Definition~\ref{def:rewriting_inclusion} of injective rewriting rules and rewriting steps may appear more restrictive, they are equivalent to those introduced in \textsection~\ref{sec:dpo}. This equivalence arises because algebraic graph rewriting via double pushouts is defined up to isomorphism. Our preference for these definitions stems from their practical advantage: they provide a clearer view of how the number of monomorphisms from specific graphs changes after a rewriting step. 
 
\begin{definition}
    \label{def:rewriting_inclusion}
    An \textbf{inclusion} is a function $f$ such that $\opn{dom}(f) \subseteq \opn{codom}(f)$ and $f(x) = x$ for every $x \in \opn{dom}(f)$.
    A rewriting rule $\varphi = (L \overset{l}{\leftarrowtail} K \overset{r}{\rightarrowtail} R)$ consists of two inclusions $l$ and $r$.  
    Two rules $\varphi' = (L' \overset{l'}{\leftarrowtail} K' \overset{r'}{\rightarrowtail} R')$ and $\varphi$ are \textbf{equivalent} if there exist isomorphisms $a,b,c$ such that 
    $l' \star a = b \star l$ and $r' \star c = b \star r$, i.e. the following diagram is commutative:    
    % the commutative diagram illustrated in Figure~\ref{fig:subgraph_counting:general_rewriting_step_sdfsl_sdfjlsdkjfl} can be constructed:
            \begin{center}
                    \resizebox{0.40\textwidth}{!}{
                        \begin{tikzpicture}
                                \node (I) at (0,0) {$K'$};
                                \node (L) at (-2,0) {$L'$};
                                \node (R) at (2,0) {$R'$};
                                \node (G) at (-2,-2) {$L$};
                                \node (C) at (0,-2) {$K$};
                                \node (H) at (2,-2) {$R$};
                                \draw [>->] (I) to  node [midway,below] {$l'$} (L);
                                \draw [>->] (I) to  node [midway,below] {$r'$} (R);
                                \draw [>->] (L) to node [midway,right] {\textbf{a}} (G);
                                \draw [>->] (I) to node [midway,right] {\textbf{b}} (C);
                                \draw [>->] (R) to node [midway,left] {\textbf{c}} (H);
                                \draw [>->] (C) to node [midway,above] {$l$} (G);
                                \draw [>->] (C) to node [midway,above] {$r$} (H);
                                \node [at=($(I)!.5!(G)$)] {$\mathop{=}$};
                                \node [at=($(I)!.5!(H)$)] {$\mathop{=}$};
                            \end{tikzpicture}
                    }
                    % \caption{}
                    % \label{fig:subgraph_counting:general_rewriting_step_sdfsl_sdfjlsdkjfl}
                \end{center}
    A rewriting step $G \Rightarrow_\varphi H$ is a double-pushout:
    %  in Figure~\ref{fig:subgraph_counting:gedfadsdsfdsfds}
%   with injective match $m :L' \rightarrowtail G$:
         \begin{center}
          \resizebox{0.40\textwidth}{!}{
              \begin{tikzpicture}
                    \node (I) at (0,0) {$K'$};
                    \node (L) at (-2,0) {$L'$};
                    \node (R) at (2,0) {$R'$};
                    \node (G) at (-2,-2) {$G$};
                    \node (C) at (0,-2) {$C$};
                    \node (H) at (2,-2) {$H$};
                    \draw [>->] (I) to  node [midway,below] {$l'$} (L);
                    \draw [>->] (I) to  node [midway,below] {$r'$} (R);
                    \draw [>->] (L) to node [midway,left] {} (G);
                    \draw [>->] (I) to node [midway,right] { } (C);
                    \draw [>->] (R) to node [midway,left] { } (H);
                    \draw [>->] (C) to node [midway,above] { } (G);
                    \draw [>->] (C) to node [midway,above] { } (H);
                    \node [at=($(I)!.5!(G)$)] {\normalfont PO};
                    \node [at=($(I)!.5!(H)$)] {\normalfont PO};
                  \end{tikzpicture}          }
      \end{center}
      where all arrows are inclusions and $L' \overset{l'}{\leftarrowtail} K' \overset{r'}{\rightarrowtail} R'$ is equivalent to $\varphi$.
\end{definition}

\subsection{pregraphs}
Graphs in double-pushouts that define rewriting steps in \textbf{Graph} can be decomposed into unions of pregraphs. pregraphs model finite directed edge-labeled multigraphs that may contain dangling edges.
\begin{definition}
    \label{def:pregraph}
    A pregraph $G$ consists of
    \begin{itemize}
        \item a finite set of nodes, denoted by $V(G)$,
        \item a subset $V_{\nexists}(G) \subseteq V(G)$ of missing nodes
        \item a finite set of edges, denoted by $E(G)$,
        \item a source function, denoted by \( \textbf{src}_G \), (resp. a target function, denoted by \( \textbf{dst}_G \)) that assigns to each edge an existing or missing node,
        \item an edge-labeling function, denoted by $l_G$, that assigns labels to edges.
    \end{itemize}
\end{definition} 
The set of existing nodes, $V(G) \setminus V_{\nexists}(G)$, is denoted by $V_\exists(G)$.

 \begin{example}
    \label{def:pregraph}
     A pregraph with node-set $\{1,2,3\}$, a missing node $2$, and 3 dangling edges is illustrated below:
     \begin{center}
        \resizebox{0.5\textwidth}{!}{
        \begin{tikzpicture}
            \graphbox{\(\)}{00mm}{-20mm}{45mm}{20mm}{2mm}{-5mm}{ 
                \coordinate (o) at (-5mm,-8mm); 
                \node[draw,circle] (l1) at ($(o)+(-10mm,0mm)$) {1};
                % \node[draw,circle] (l2) at ($(l1)+(3,0)$) {};
                \node[draw,circle,dashed,red] (l3) at ($(l1)+(1,0)$) {2};
                \node[draw,circle] (l4) at ($(l1)+(2,0)$) {3};
                \draw[->] (l1) edge[bend right]  node[midway,below] {$a$} (l3);
                \draw[->] (l1) edge[bend left] node[midway,above] {$a$}  (l3);
                \draw[->] (l3) -- (l4) node[midway,above] {$b$};
                % \draw[->] (l4) -- (l2) node[midway,above] {$a$};
            }  
        \end{tikzpicture} 
    }
    \end{center}
 \end{example} 
 
 A homomorphism of pregraphs is a mapping between two pregraphs that preserves structure and respects missing and existing nodes.
 \begin{definition} 
    \label{def:pregraph:homomorphism}
    Let \( G \) and \( H \) be pregraphs. A \textbf{homomorphism of pregraphs} $h: G \to H$ consists a pair of functions $h_V: V(G) \to V(H) $ and $h_E: E(G) \to E(H)$ such that
    \begin{itemize}
        \item for every edge \( a \colon s \to t \) in \( G \), the edge \( h_E(a) \) in $H$ has source \( h_V(s) \) and target \( h_V(t) \),
        \item for all $s \in V_{\nexists}(G)$, we have $h_V(s) \in V_{\nexists}(H)$,
        \item for all $t \in V_\exists(G)$, we have $h_V(t) \in V_\exists(H)$.
    \end{itemize} 
\end{definition}

\begin{definition} 
    \label{def:pregraph:composition}
    Let $f: (G,\lambda) \to (G',\lambda')$ and $g: (G',\lambda') \to (G'',\lambda'')$ be homomorphisms of pregraphs. Their composition, denoted by $g \circ f$, is the homomorphism defined by the component-wise composition of functions $g_V \circ f_V$ and $g_E \circ f_E$.
\end{definition}
\begin{definition} 
    \label{def:subpregraph} Let $G$ be a pregraph. A pregraph $H$ is a \textbf{subpregraph} of $G$, written $H \rightarrowtail G$, if all of the following conditions hold:
     \begin{itemize} 
            \item $V(H) \subseteq V(G)$, \item $V_{\nexists}(H) = V(H) \cap V_{\nexists}(G)$, 
            \item $E(H) \subseteq E(G)$, 
            \item the source and target functions of $H$ are the restrictions of those of $G$ to $E(H)$,
            \item the edge-labeling function of $H$ is the restriction $l_H = {l_G}_{|_{E(H)}}$. 
        \end{itemize} 
\end{definition}
We define two operations on pregraphs that will be used in the following discussion.
The \textbf{union} $C \cup R$ of two subpregraphs $C \subseteq G$ and $R \subseteq G$ is a pregraph with 
\begin{itemize}
    \item Node set: $ V(C) \cup V(R) $,
    \item Missing nodes:
        $(
            V_{\nexists}(C) \cup V_{\nexists}(R)) \setminus 
            (   V_\exists(C) \cup 
             V_\exists(R)  
                   )$
    \item Edge set: $ E(C) \cup E(R) $, 
    \item Source function: $ \opn{src}_C \cup \opn{src}_R $,  
    \item Target function: $ \opn{dst}_C \cup \opn{dst}_R $,  
    \item Edge-labeling function: $ l_C \cup l_R $.   
\end{itemize} 
The definition is illustrated below:
% in Figure~\ref{fig:subgraph_counting:union_of_pregraphs}.
% \begin{figure}[H]
%     \centering
\begin{center}
    \resizebox{\textwidth}{!}{
        \begin{tikzpicture}  
            \graphbox{\( R \)}{91mm}{-22mm}{45mm}{25mm}{2mm}{2mm}{
                \coordinate (o) at (-5mm,-8mm); 
                \node[draw,circle] (l1) at ($(o)+(-10mm,0mm)$) {1};
                \node[draw,circle] (l2) at ($(l1)+(3,0)$) {2};
                \node[draw,circle] (l3) at ($(l1) + (1,0)$) {4};
                \node[draw,circle] (l4) at ($(l1) + (2,0)$) {5};
                \draw[->] (l1) -- (l3) node[midway,above] {$a$};
                \draw[->] (l3) -- (l4) node[midway,above] {$b$};
                \draw[->] (l4) -- (l2) node[midway,above] {$a$};
            }    
        
            \graphbox{\( C  \)}{0mm}{-22mm}{34mm}{25mm}{2mm}{-3mm}{
                \coordinate (o) at (0mm,-6mm); 
                \node[draw,circle,dashed] (l1) at ($(o)+(-10mm,0mm)$) {1};
                \node[draw,circle,dashed] (l2) at ($(l1)+(2,0)$) {2};
                \node[draw,circle] (l4) at ($(l2) + (0,-1)$) {6};
                \draw[->] (l2) -- (l4) node[midway,right] {$a$};
                \draw[->] (l2) edge[out=-135,in=-45]node[midway,below] {$a$} (l1) ;
                % \node[ draw,circle] (l6) at ($(l1) + (0,-1)$) {7};
                % \draw[<-] (l1) -- (l6) node[midway,left] {$a$};
            }    

            \graphbox{\(C \cup R \)}{40mm}{-22mm}{45mm}{25mm}{2mm}{-3mm}{
                \coordinate (o) at (-5mm,-6mm); 
                \node[draw,circle] (l1) at ($(o)+(-10mm,0mm)$) {1};
                \node[draw,circle] (l2) at ($(l1)+(3,0)$) {2};
                \node[draw,circle] (l3) at ($(l1) + (1,0)$) {4};
                \node[draw,circle] (l4) at ($(l1) + (2,0)$) {5};
                \node[ draw,circle] (l5) at ($(l2) + (0,-1)$) {6};
                % \node[ draw,circle] (l6) at ($(l1) + (0,-1)$) {7};
                % \draw[<-] (l1) -- (l6) node[midway,left] {$a$};
                \draw[->] (l1) -- (l3) node[midway,above] {$a$};
                \draw[->] (l3) -- (l4) node[midway,above] {$b$};
                \draw[->] (l4) -- (l2) node[midway,above] {$a$};
                \draw[->] (l2) -- (l5) node[midway,right] {$a$};
                \draw[->] (l2) edge[out=-135,in=-45]node[midway,below] {$a$} (l1) ;
            }    
             
        \end{tikzpicture}
    }
    % \caption{Union of pregraphs}
    % \label{fig:subgraph_counting:union_of_pregraphs}
    % \end{figure}
    \end{center}

    % \textcolor{red}{Intersection} $G \cap H$: 
    % \vspace{1mm}

    % % \begin{center}
    %     \resizebox{0.95\textwidth}{!}{
    %     \begin{tikzpicture}
    %         \graphbox{\( G \)}{0mm}{-22mm}{34mm}{26mm}{2mm}{-3mm}{
    %             \coordinate (o) at (0mm,-6mm); 
    %             \node[draw,circle] (l1) at ($(o)+(-10mm,0mm)$) {1};
    %             \node[draw,circle] (l2) at ($(l1)+(2,0)$) {2};
    %             \node[draw,circle] (l3) at ($(l1) + (1,0)$) {3};
    %             \node[draw,circle] (l4) at ($(l2) + (0,-1)$) {6};
    %             \draw[->] (l1) -- (l3) node[midway,above] {$a$};
    %             \draw[->] (l3) -- (l2) node[midway,above] {$a$};
    %             \draw[->] (l2) -- (l4) node[midway,right] {$a$};
    %             \node[draw,circle] (l6) at ($(l1) + (0,-1)$) {7};
    %             \draw[<-] (l1) -- (l6) node[midway,left] {$a$};
    %             \draw[->] (l2) edge[out=-135,in=-45]node[midway,below] {$a$} (l1) ;
    %         }    
    
    %         \graphbox{ \textcolor{red}{\(G \cap H\)}  }{40mm}{-22mm}{34mm}{25mm}{2mm}{-3mm}{
    %             \coordinate (o) at (0mm,-6mm); 
    %             \node[draw,circle] (l1) at ($(o)+(-10mm,0mm)$) {1};
    %             \node[draw,circle] (l2) at ($(l1)+(2,0)$) {2};
    %             \node[draw,circle] (l4) at ($(l2) + (0,-1)$) {6};
    %             \draw[->] (l2) -- (l4) node[midway,right] {$a$};
    %             \draw[->] (l2) edge[out=-135,in=-45]node[midway,below] {$a$} (l1) ;
    %             \node[ draw,circle] (l6) at ($(l1) + (0,-1)$) {7};
    %             \draw[<-] (l1) -- (l6) node[midway,left] {$a$};
    %         }    
    
    %         \graphbox{\( H \)}{80mm}{-22mm}{45mm}{25mm}{2mm}{-6mm}{
    %             \coordinate (o) at (-5mm,-3mm); 
    %             \node[draw,circle] (l1) at ($(o)+(-10mm,0mm)$) {1};
    %             \node[draw,circle] (l2) at ($(l1)+(3,0)$) {2};
    %             \node[draw,circle] (l3) at ($(l1) + (1,0)$) {4};
    %             \node[draw,circle] (l4) at ($(l1) + (2,0)$) {5};
    %             \node[ draw,circle] (l5) at ($(l2) + (0,-1)$) {6};
    %             \node[ draw,circle] (l6) at ($(l1) + (0,-1)$) {7};
    %             \draw[<-] (l1) -- (l6) node[midway,left] {$a$};
    %             \draw[->] (l1) -- (l3) node[midway,above] {$a$};
    %             \draw[->] (l3) -- (l4) node[midway,above] {$b$};
    %             \draw[->] (l4) -- (l2) node[midway,above] {$a$};
    %             \draw[->] (l2) -- (l5) node[midway,right] {$a$};
    %             \draw[->] (l2) edge[out=-135,in=-45]node[midway,below] {$a$} (l1) ;
    %         }    
           
    %         % \node () at (37mm,-33mm) {\( \leftarrowtail \)};
    %         % \node () at (77mm,-33mm) {\( \rightarrowtail \)}; % C -> H
    %     \end{tikzpicture}
    %     }         
    % \end{center}
The \textbf{relative complement} of $R$ in $H$ where $R \subseteq H$, denoted by $H \setminus R$, is the pregraph with
\begin{itemize}
    \item Edge set: $X \isdef E(H) \setminus E(R)$,
    \item Node set: $Y \isdef \bigcup_{e \in X} \{\opn{src}(e), \opn{dst}(e)\}$,
    \item Set of missing nodes: $
   ( V_{\nexists}(H)  \cup
        (
                  V_\exists(H) 
                  \cap 
                  V_\exists(R)
                  )
   )
        \cap  
        Y
                  $,
    \item Source function: ${\opn{src}_H}$ restricted to $X$,
    \item Target function: ${\opn{dst}_H}$ restricted to $X$,
    \item Edge-labeling function : $l_H$ restricted to $X$.
\end{itemize}
An example is shown below where the graph $H$, $R$ and the relative complement $H \setminus R$ are illustrated:
% in Figure~\ref{fig:subgraph_counting:relative_complement_of_pregraphs}.
\begin{center}  
% \begin{figure}[H]
    %     \centering
  \resizebox{\textwidth}{!}{
        \begin{tikzpicture}  
            \phantom{
                \graphbox{\( G \)}{0mm}{-22mm}{34mm}{22mm}{2mm}{-3mm}{
                }  
            }
            \graphbox{\( R \)}{80mm}{-22mm}{45mm}{25mm}{2mm}{2mm}{
                \coordinate (o) at (-5mm,-8mm); 
                \node[draw,circle] (l1) at ($(o)+(-10mm,0mm)$) {1};
                \node[draw,circle] (l2) at ($(l1)+(3,0)$) {2};
                \node[draw,circle] (l3) at ($(l1) + (1,0)$) {4};
                \node[draw,circle] (l4) at ($(l1) + (2,0)$) {5};
                \draw[->] (l1) -- (l3) node[midway,above] {$a$};
                \draw[->] (l3) -- (l4) node[midway,above] {$b$};
                \draw[->] (l4) -- (l2) node[midway,above] {$a$};
            }    
        
            \graphbox{\( H \setminus R  \)}{40mm}{-22mm}{34mm}{25mm}{2mm}{-3mm}{
                \coordinate (o) at (0mm,-6mm); 
                \node[draw,dashed,circle] (l1) at ($(o)+(-10mm,0mm)$) {1};
                \node[draw,dashed,circle] (l2) at ($(l1)+(2,0)$) {2};
                \node[draw,circle] (l4) at ($(l2) + (0,-1)$) {6};
                \draw[->] (l2) -- (l4) node[midway,right] {$a$};
                \draw[->] (l2) edge[out=-135,in=-45]node[midway,below] {$a$} (l1) ;
                % \node[ draw,circle] (l6) at ($(l1) + (0,-1)$) {7};
                % \draw[<-] (l1) -- (l6) node[midway,left] {$a$};
            }    

            \graphbox{\( H \)}{-11mm}{-22mm}{45mm}{25mm}{2mm}{-3mm}{
                \coordinate (o) at (-5mm,-6mm); 
                \node[draw,circle] (l1) at ($(o)+(-10mm,0mm)$) {1};
                \node[draw,circle] (l2) at ($(l1)+(3,0)$) {2};
                \node[draw,circle] (l3) at ($(l1) + (1,0)$) {4};
                \node[draw,circle] (l4) at ($(l1) + (2,0)$) {5};
                \node[ draw,circle] (l5) at ($(l2) + (0,-1)$) {6};
                % \node[ draw,circle] (l6) at ($(l1) + (0,-1)$) {7};
                % \draw[<-] (l1) -- (l6) node[midway,left] {$a$};
                \draw[->] (l1) -- (l3) node[midway,above] {$a$};
                \draw[->] (l3) -- (l4) node[midway,above] {$b$};
                \draw[->] (l4) -- (l2) node[midway,above] {$a$};
                \draw[->] (l2) -- (l5) node[midway,right] {$a$};
                \draw[->] (l2) edge[out=-135,in=-45]node[midway,below] {$a$} (l1) ;
            }    
            % \node () at (102mm,-18mm) {\( \downarrowtail \)};
            % \node () at (77mm,-33mm) {\( \rightarrowtail \)}; % C -> H
        \end{tikzpicture}
    }
    % \caption{Relative complement of $R$ in $H$.}
    % \label{fig:subgraph_counting:relative_complement_of_pregraphs}
    % \end{figure}
\end{center}

 
\subsection{$X$-occurrences}

    Consider the following double pushout diagram in the category \textbf{Graph}, where all arrows are inclusions:
    % illustrated in Figure~\ref{fig:subgraph_counting:pushoutsquare} can be decomposed into the diagram illustrated in Figure~\ref{fig:subgraph_counting:decomposition_of_rewriting_step}.
    % \begin{figure}[H]
    %     \centering
    \begin{center}
       \resizebox{0.49\textwidth}{!}{
        \begin{tikzpicture}
            % [node distance=15mm]
            \node (I) at (0,0) {$K$};
            \node (L)  at (-2,0) {$L$};
            \node (R)  at (2,0) {$R$};
            \node (G)  at (-2,-2) {$G$};
            \node (C)  at (0,-2) {$C$};
            \node (H)  at (2,-2) {$H$};
            \draw [>->] (I) to  node [midway,below] { } (L);
            \draw [>->] (I) to  node [midway,below] { } (R);
            \draw [>->] (L) to node [midway,right] { } (G);
            \draw [>->] (I) to  node [midway,right] 
            % {$u$}
            {} (C);
            \draw [>->] (R) to  node [midway,right] 
            {}
            (H);
            \draw [>->] (C) to node [midway,above] {} (G);
            \draw [>->] (C) to node [midway,above] 
            {}
            (H);
            % \node [at=($(I)!.5!(G)$)] {\normalfont PO};
            % \node [at=($(I)!.5!(H)$)] {\normalfont PO};
        \end{tikzpicture}
    }
    % \caption{}
    % \label{fig:subgraph_counting:pushoutsquare}
    % \end{figure}
    \end{center}
    Since all arrows are inclusions, by Definition~\ref{def:rewriting_inclusion}, graphs can be decomposed as unions of pregraphs as illustrated below, 
     where $K, L', R',C'$ are pregraphs and $L = L' \cup K$, $R = R' \cup K$, $C = C' \cup K$, $G = L' \cup K \cup C'$, $H = R' \cup K \cup C'$.
    % \begin{figure}[H]
    %     \centering
    \begin{center}
    \resizebox{0.49\textwidth}{!}{
            \begin{tikzpicture} 
                \coordinate (k) at (0, 0);
                \draw[fill=white] ($(k)+(0,0)$) rectangle ($(k)+(0.5,0.5)$);
                \node () at ($(k)+(0.25,0.25)$) {\( \mathrm{K} \)};
            
                \coordinate (c) at (0, -2.2);
                \draw[fill=blue!20]
                ($(c)+(0,-0.5)$)
                -- ($(c)+(0,0.5)$) 
                -- ($(c)+(1,0.5)$) 
                arc[start angle=0, end angle=-90, radius=1]
                -- cycle;
                \node () at ($(c)+(0.75,0.25)$) {\( \mathrm{C'} \)};
                \draw[fill=white] ($(c)+(0,0)$) rectangle ($(c)+(0.5,0.5)$);
                \node () at ($(c)+(0.25,0.25)$) {\( \mathrm{K} \)};
            
                \coordinate (l) at (-3, 0);
                \draw[fill=orange!20] ($(l)+(-0.5,0)$) rectangle ($(l)+(0.5,1)$);
                \node () at ($(l)+(-0.23,0.25)$) {\( \mathrm{L'} \)};
                \draw[fill=white] ($(l)+(0,0)$) rectangle ($(l)+(0.5,0.5)$);
                \node () at ($(l)+(0.25,0.25)$) {\( \mathrm{K} \)};
            
                \coordinate (g) at (-3, -2.2);
                \draw[fill=blue!20]
                ($(g)+(0,-0.5)$)
                -- ($(g)+(0,0.5)$)
                -- ($(g)+(1,0.5)$) 
                arc[start angle=0, end angle=-90, radius=1]
                -- cycle;
                \draw[fill=orange!20] ($(g)+(-0.5,0)$) rectangle ($(g)+(0.5,1)$);
                \node () at ($(g)+(0.75,0.25)$) {\( \mathrm{C'} \)};
                \node () at ($(g)+(-0.23,0.25)$) {\( \mathrm{L'} \)};
                \draw[fill=white] ($(g)+(0,0)$) rectangle ($(g)+(0.5,0.5)$);
                \node () at ($(g)+(0.25,0.25)$) {\( \mathrm{K} \)};
            
                \coordinate (r) at (3,0);
                \draw[fill=red!20] ($(r)+(-0.5,0)$)
                -- ($(r)+(-0.5,0.5)$)
                -- ($(r)+(0,1)$)
                --  ($(r)+(0.5,1)$)
                -- ($(r)+(0.5,0)$)
                -- cycle;
                \node () at ($(r)+(-0.23,0.25)$) {\( \mathrm{R'} \)};
                \draw[fill=white] ($(r)+(0,0)$) rectangle ($(r)+(0.5,0.5)$);
                \node () at ($(r)+(0.25,0.25)$) {\( \mathrm{K} \)};
            
                \coordinate (h) at (3, -2.2);
                \draw[fill=blue!20]
                ($(h)+(0,-0.5)$)
                -- ($(h)+(0,0.5)$)
                -- ($(h)+(1,0.5)$) 
                arc[start angle=0, end angle=-90, radius=1]
                -- cycle;
                \draw[fill=red!20] ($(h)+(-0.5,0)$)
                -- ($(h)+(-0.5,0.5)$)
                -- ($(h)+(0,1)$)
                --  ($(h)+(0.5,1)$)
                -- ($(h)+(0.5,0)$)
                -- cycle;
            \node () at ($(h)+(0.75,0.25)$) {\( \mathrm{C'} \)};
            \draw[fill=white] ($(h)+(0,0)$) rectangle ($(h)+(0.5,0.5)$);
            \node () at ($(h)+(0.25,0.25)$) {\( \mathrm{K} \)};
            \node () at ($(h)+(-0.23,0.25)$) {\( \mathrm{R'} \)};
            
                \node[ font=\huge] (kl) at ($(k)!0.5!(l)+(0.25,0.25)$)
            %    {\( \overset{l}{\leftarrowtail} \)}
                {\( \leftarrowtail \)}
                ; 
                \node[ font=\huge] (kr) at ($(k)!0.5!(r)+(0.25,0.25)$)
                {\( \rightarrowtail \)}
                ;  
                \node[ font=\huge] (cg) at ($(c)!0.5!(g)+(0.25,0.25)$) 
                {\( \leftarrowtail \)}
            ;  
                \node[ font=\huge] (ch) at ($(c)!0.5!(h)+(0.25,0.25)$)
                {\( \rightarrowtail \)}
            ; 
                \node[ font=\huge] (kc) at ($(k)!0.5!(c)+(0.2,0.4)$) {\( \downarrowtail \)}; 
            %   \node[ font=\LARGE] () at ($(l)!0.5!(g)+(0.5,0.4)$) {$m$}; 
                \node[ font=\huge] (lg) at ($(l)!0.5!(g)+(0.1,0.4)$) {\( \downarrowtail \)}; 
                \node[ font=\huge] (rh) at ($(r)!0.5!(h)+(0.1,0.4)$) {\( \downarrowtail \)}; 
            %   \node[ font=\LARGE] () at ($(r)!0.5!(h)+(0.55,0.4)$) {$m'$}; 
            %   \node[ font=\LARGE] () at ($(k)!0.5!(c)+(0.5,0.4)$) {$u$}; 
            \end{tikzpicture}
          }
%         \caption{}
%         \label{fig:subgraph_counting:decomposition_of_rewriting_step}
%    \end{figure}
    \end{center}
\begin{definition}
    \label{def:x_occurrence}
     Let $X$ be a subgraph of $L$.  An \textbf{$X$-occurrence} in a graph $G$ is a \textbf{monomorphism} from $X$ to $G$. An $X$-occurrence in $G$ is 
    \begin{itemize}
        \item \textbf{shared} if its image is included in $C$,
        \item \textbf{explicit} if its image is  included in $L$,
        \item \textbf{implicit} if its image has elements in both $C'$ and $L'$.
    \end{itemize}  
    An $X$-occurrence in $H$ is 
    \begin{itemize}
        \item \textbf{shared} if its image is included in $C$,
        \item \textbf{explicit} if its image is included in $R$,
        \item \textbf{implicit} if its image has elements in both $C'$ and $R'$.
    \end{itemize} 

    Intuitively, the rewriting step implicitly creates the implicit $X$-occurrences in $H$ and implicitly destroys implicit $X$-occurrences in $G$. 
    
    Given an implicit $X$-occurrence $x$ in $H$, its image can be decomposed as the union $R_x \cup C_x$, where $R_x = \opn{Im}(x) \cap R$ (the part of the image overlapping the right-hand side R), and $C_x = \opn{Im}(x) \cap C$ (the part overlapping the context C).
    
    For every graph monomorphism $h : R_x \rightarrowtail L$ that preserves interface elements (i.e., \( h(k) = k \) for all \( k \in K \)), we define the \textbf{corresponding $X$-occurrence of $x$ relative to $h$} as the $X$-occurrence $y$ in $G$ given by:
        $$
        y(z) = 
        \begin{cases}
        z & \text{if } z \in x^{-1}(C_x), \\
        h(x(z)) & \text{if } z \in x^{-1}(R_x \setminus C_x).
        \end{cases}
        $$
   \end{definition} 



% \begin{definition}
%     \label{def:corresponding_occurrence}
%      An \textbf{$X$-occurrence} in $G$ is a subgraph of $G$ isomorphic to $X$. An \textbf{$X$-occurrence $x$ implicitly created by the rewriting step} is an \( X \)-occurrence in $H$ which is included neither in \( R \) nor in $C$. 
%       The occurrence $x$ can be decomposed as the union of graphs $x = R_x \cup C_x$ where $R_x = x \cap R$ and $C_x = x \cap C$. For every graph monomorphism $h : R_x \rightarrowtail L$ that preserves interface elements (i.e., \( h(k) = k \) for all \( k \in K \)), we define the \textbf{corresponding $X$-occurrence of $x$ relative to $h$} as the $X$-occurrence in $G$ given by $h(R_x) \cup C_x$.
%     An \textbf{$X$-occurrence implicitly destroyed by the rewriting step} is an \( X \)-occurrence in $G$ which is included neither in \( L \) nor in $C$.
% \end{definition}


\begin{example}
    Let $X$ be the graph
\tikz[baseline=-0.5ex]{ 
        \node[draw,circle] (x) at (0,0) {};  
        \node[draw,circle] (y) at (1,0) {};
        \node[draw,circle] (z) at (2,0) {};
        \draw[->] (x) -- node[midway,above] {$a$} (y) ;
        \draw[->] (y) -- node[midway,above] {$a$} (z) ;
}.
Consider the $X$-occurrences in the decomposition of a rewriting step using the rule in Example~\ref{ex:grsaa} and recalled below:
% in Figure~\ref{fig:subgraph_counting:occurrdlfdjslfkjsa}.
    % \begin{figure}[H]
    %   \centering
    \begin{center}
      \resizebox{0.7\textwidth}{!}{
      \begin{tikzpicture}
          \graphbox{$L$}{0mm}{0mm}{34mm}{15mm}{2mm}{-5mm}{
              \coordinate (o) at (0mm,-3mm); 
              \node[draw,circle] (l1) at ($(o)+(-10mm,0mm)$) {1};
              \node[draw,circle] (l2) at ($(l1)+(2,0)$) {2};
              \node[draw,circle] (l3) at ($(l1) + (1,0)$) {3};
              \draw[->] (l1) -- (l3) node[midway,above] {$a$};
              \draw[->] (l3) -- (l2) node[midway,above] {$a$};
          }     
          \graphbox{$K$}{40mm}{0mm}{24mm}{15mm}{2mm}{-5mm}{
              \coordinate (o) at (5mm,-3mm); 
              \node[draw,circle] (l1) at ($(o)+(-10mm,0mm)$) {1};
              \node[draw,circle] (l2) at ($(l1)+(1,0)$) {2};
              % \node[draw,circle] (l3) at ($(l1) + (1,0)$) {$\ $};
              % \draw[->] (l1) -- (l3) node[midway,above] {$a$};
              % \draw[->] (l3) -- (l2) node[midway,above] {$a$};
          }    
          \graphbox{$R$}{70mm}{0mm}{45mm}{15mm}{2mm}{-5mm}{
              \coordinate (o) at (-5mm,-3mm); 
              \node[draw,circle] (l1) at ($(o)+(-10mm,0mm)$) {1};
              \node[draw,circle] (l2) at ($(l1)+(3,0)$) {2};
              \node[draw,circle] (l3) at ($(l1) + (1,0)$) {4};
              \node[draw,circle] (l4) at ($(l1) + (2,0)$) {5};
              \draw[->] (l1) -- (l3) node[midway,above] {$a$};
              \draw[->] (l3) -- (l4) node[midway,above] {$b$};
              \draw[->] (l4) -- (l2) node[midway,above] {$a$};
          }    
          \node () at (37mm,-8mm) {$\overset{l}{\leftarrowtail}$};
          \node () at (67mm,-8mm) {$\overset{r}{\rightarrowtail}$};
          % \draw[>->] (51mm,2mm) -- (52mm,3mm);
      \end{tikzpicture}
      }
%       \caption{}
%       \label{fig:subgraph_counting:occurrdlfdjslfkjsa}
%   \end{figure}
    \end{center}
   We illustrate the aforementioned decomposition
    % in Figure~\ref{fig:subgraph_counting:occurrendfsdadsgces_of_x} 
    using the visual notation for graph homomorphisms introduced in Notation~\ref{notation:graph_homomorphism}, with distinct pregraphs highlighted in different colors.
% \begin{figure}[H]
%     \centering
    \begin{center}
    \resizebox{0.8\textwidth}{!}{
      \begin{tikzpicture}
              \graphbox{\( L = L' \cup K \)}{0mm}{5mm}{34mm}{20mm}{2mm}{-5mm}{
                  \coordinate (o) at (0mm,-8mm); 
                  \node[draw,circle] (l1) at ($(o)+(-10mm,0mm)$) {1};
                  \node[draw,circle] (l2) at ($(l1)+(2,0)$) {2};
                  \node[orange,draw,circle] (l3) at ($(l1) + (1,0)$) {3};
                  \draw[orange,->] (l1) -- (l3) node[midway,above] {$a$};
                  \draw[orange,->] (l3) -- (l2) node[midway,above] {$a$};
              } 

              \graphbox{\( K \)}{40mm}{5mm}{34mm}{20mm}{2mm}{-5mm}{
                  \coordinate (o) at (0mm,-8mm); 
                  \node[draw,circle] (l1) at ($(o)+(-10mm,0mm)$) {1};
                  \node[draw,circle] (l2) at ($(l1)+(2,0)$) {2};
              }  

              \graphbox{\( R = R' \cup K \)}{80mm}{5mm}{45mm}{20mm}{2mm}{-5mm}{
                  \coordinate (o) at (-5mm,-8mm); 
                  \node[draw,circle] (l1) at ($(o)+(-10mm,0mm)$) {1};
                  \node[draw,circle] (l2) at ($(l1)+(3,0)$) {2};
                  \node[red,draw,circle] (l3) at ($(l1) + (1,0)$) {4};
                  \node[red,draw,circle] (l4) at ($(l1) + (2,0)$) {5};
                  \draw[red,->] (l1) -- (l3) node[midway,above] {$a$};
                  \draw[red,->] (l3) -- (l4) node[midway,above] {$b$};
                  \draw[red,->] (l4) -- (l2) node[midway,above] {$a$};
              }    

              \graphbox{\( G = L' \cup K \cup C' \)}{0mm}{-22mm}{34mm}{30mm}{2mm}{-10mm}{
                  \coordinate (o) at (0mm,-3mm); 
                  \node[draw,circle] (l1) at ($(o)+(-10mm,0mm)$) {1};
                  \node[draw,circle] (l2) at ($(l1)+(2,0)$) {2};
                  \node[draw,circle,orange] (l3) at ($(l1) + (1,0)$) {3};
                  \node[blue, draw,circle] (l4) at ($(l2) + (0,-1)$) {6};
                  \draw[orange,->] (l1) -- (l3) node[midway,above] {$a$};
                  \draw[orange,->] (l3) -- (l2) node[midway,above] {$a$};
                  \draw[blue,->] (l2) -- (l4) node[midway,right] {$a$};
                %   \node[blue,draw,circle] (l6) at ($(l1) + (0,-1)$) {7};
                %   \draw[blue,<-] (l1) -- (l6) node[midway,left] {$a$};
                \node[blue,draw,circle] (l7) at ($(l4) + (-1,0)$) {7};
                  \draw[blue,->] (l4) -- (l7) node[midway,below] {$a$};
              }    

              \graphbox{\( C = K \cup C' \)}{40mm}{-22mm}{34mm}{30mm}{2mm}{-10mm}{
                  \coordinate (o) at (0mm,-3mm); 
                  \node[draw,circle] (l1) at ($(o)+(-10mm,0mm)$) {1};
                  \node[draw,circle] (l2) at ($(l1)+(2,0)$) {2};
                  \node[blue,draw,circle] (l4) at ($(l2) + (0,-1)$) {6};
                  \draw[blue,->] (l2) -- (l4) node[midway,right] {$a$};
                %   \node[blue,draw,circle] (l6) at ($(l1) + (0,-1)$) {7};
                %   \draw[blue,<-] (l1) -- (l6) node[midway,left] {$a$};
                    \node[blue,draw,circle] (l7) at ($(l4) + (-1,0)$) {7};
                  \draw[blue,->] (l4) -- (l7) node[midway,below] {$a$};
              }    

              \graphbox{\( H = R' \cup K \cup C' \)}{80mm}{-22mm}{45mm}{30mm}{2mm}{-10mm}{
                  \coordinate (o) at (-5mm,-3mm); 
                  \node[draw,circle] (l1) at ($(o)+(-10mm,0mm)$) {1};
                  \node[draw,circle] (l2) at ($(l1)+(3,0)$) {2};
                  \node[draw,circle,red] (l3) at ($(l1) + (1,0)$) {4};
                  \node[draw,circle,red] (l4) at ($(l1) + (2,0)$) {5};
                  \node[blue,draw,circle] (l5) at ($(l2) + (0,-1)$) {6};
                %   \node[blue,draw,circle] (l6) at ($(l1) + (0,-1)$) {7};
                %   \draw[blue,<-] (l1) -- (l6) node[midway,left] {$a$};
                  \draw[red,->] (l1) -- (l3) node[midway,above] {$a$};
                  \draw[red,->] (l3) -- (l4) node[midway,above] {$b$};
                  \draw[red,->] (l4) -- (l2) node[midway,above] {$a$};
                  \draw[blue,->] (l2) -- (l5) node[midway,right] {$a$};
                        \node[blue,draw,circle] (l7) at ($(l5) + (-1,0)$) {7};
                  \draw[blue,->] (l5) -- (l7) node[midway,below] {$a$};
              }    

              \node () at (37mm,-8mm) {\( \leftarrowtail \)}; % K -> L
              \node () at (77mm,-8mm) {\( \rightarrowtail \)}; % K -> R
              \node () at (17mm,-18mm) {\( m\ \downarrowtail \)};
              \node () at (37mm,-33mm) {\( \leftarrowtail \)};
              \node () at (52mm,-18mm) {\( \downarrowtail \)};
              \node () at (92mm,-18mm) {\( \downarrowtail \)};
              \node () at (77mm,-33mm) {\( \rightarrowtail \)}; % C -> H
      \end{tikzpicture}
      }
%       \caption{}
%       \label{fig:subgraph_counting:occurrendfsdadsgces_of_x}
%    \end{figure}
    \end{center}
    In this example, the graph $X$ has 4 distinct $X$-occurrences in $G$ and $H$.
   There are 
\begin{itemize}
    \item 1 occurrence shared by $G$ and $H$: 
     \tikz[baseline=-0.5ex]{ 
        \node[draw,circle] (x) at (0,0) {2};  
        \node[blue,draw,circle] (y) at (1,0) {6};
        \node[blue,draw,circle] (z) at (2,0) {7};
        \draw[blue,->] (x) -- node[midway,above] {$a$} (y) ;
        \draw[blue,->] (y) -- node[midway,above] {$a$} (z) ;
     }
    \item 1 explicit occurrence in $G$:
     \tikz[baseline=-0.5ex]{ 
        \node[draw,circle] (x) at (0,0) {1};  
        \node[orange,draw,circle] (y) at (1,0) {3};
        \node[draw,circle] (z) at (2,0) {2};
        \draw[orange,->] (x) -- node[midway,above] {$a$} (y) ;
        \draw[orange,->] (y) -- node[midway,above] {$a$} (z) ;
    }
 
    \item 1 implicit occurrences in $G$:
    \tikz[baseline=-0.5ex]{ 
        \node[orange,draw,circle] (x) at (0,0) {3};  
        \node[draw,circle] (y) at (1,0) {2};
        \node[blue,draw,circle] (z) at (2,0) {6};
        \draw[orange,->] (x) -- node[midway,above] {$a$} (y) ;
        \draw[blue,->] (y) -- node[midway,above] {$a$} (z) ;
    } 
 
    \item 0 explicit occurrence in $H$,
    \item 1 implicit occurrence in $H$: 
    \tikz[baseline=-0.5ex]{ 
        \node[red,draw,circle] (x) at (0,0) {5};  
        \node[draw,circle] (y) at (1,0) {2};
        \node[blue,draw,circle] (z) at (2,0) {6};
        \draw[red,->] (x) -- node[midway,above] {$a$} (y) ;
        \draw[blue,->] (y) -- node[midway,above] {$a$} (z) ;
}
 
\end{itemize}
    The corresponding X-occurrence of the implicit occurrences \tikz[baseline=-0.5ex]{ 
        \node[red,draw,circle] (x) at (0,0) {5};  
        \node[draw,circle] (y) at (1,0) {2};
        \node[blue,draw,circle] (z) at (2,0) {6};
        \draw[red,->] (x) -- node[midway,above] {$a$} (y) ;
        \draw[blue,->] (y) -- node[midway,above] {$a$} (z) ;
} in $H$ relative to the unique monomorphism from 
    \tikz[baseline=-0.5ex]{
      \node[red,draw,circle] (x) at (0,0) {5};
      \node[draw,circle] (y) at (1,0) {2};
      \draw[red,->] (x) -- (y) node[midway, above] {$a$};
  }
     to 
    \tikz[baseline=-0.5ex]{
      \node[orange,draw,circle] (x) at (0,0) {3};
      \node[draw,circle] (y) at (1,0) {2};
      \draw[orange,->] (x) -- (y) node[midway, above] {$a$};
  }
    is the implicit occurrences 
       \tikz[baseline=-0.5ex]{ 
        \node[orange,draw,circle] (x) at (0,0) {3};  
        \node[draw,circle] (y) at (1,0) {2};
        \node[blue,draw,circle] (z) at (2,0) {6};
        \draw[orange,->] (x) -- node[midway,above] {$a$} (y) ;
        \draw[blue,->] (y) -- node[midway,above] {$a$} (z) ;
    } 
    in $G$.
\end{example}
\subsection{Graphs}
 \begin{definition}[\cite{barr1990category}]
    \label{def:graph:unlabeled}
    An \textbf{unlabeled graph} \( G \) consists of a collection of \textbf{nodes} (also called \textbf{objects}) and a collection of \textbf{edges}, each equipped with a \textbf{source} (or \textbf{domain}) node and a \textbf{target} (or \textbf{codomain}) node. 
    For an unlabeled graph \( G \), we denote by \( G_0 \) its collection of nodes, \( G_1 \) its collection of edges, \( \operatorname{dom}:G_1{\to}G_0 \) the domain function, and \( \operatorname{cod}:G_1{\to}G_0 \) the codomain function. An unlabeled graph is \textbf{finite} if \( G_0 \) and \( G_1 \) are both finite sets.
    We write \( a: s \to t \) to indicate that \( a \) is a directed edge from \( s \) to \( t \).
\end{definition}   

\begin{example}
    In an unlabeled graph, edges are directed and will be visualized as arrows; between any two nodes, there can be multiple edges, and loops are allowed. The following is an example of an unlabeled graph:
     \begin{center}
        \begin{tikzpicture}
            \graphbox{}{0mm}{0mm}{32mm}{28mm}{-10mm}{-14mm}{
                \node[draw,circle] (1) at (0,0) {1};
                \node[draw,circle] (2) at (2,0) {2};
                \draw[->] (1) edge[loop above] (1) ;
                \draw[->] (1) edge[loop below] (1) ;
                \draw[->] (1) edge[bend left] (2)  ;
                \draw[->] (2) edge[bend left] (1)   ;
            }
        \end{tikzpicture}
    \end{center}
\end{example}
% A homomorphism of unlabeled graphs is a mapping between the nodes and edges of two graphs that preserves the graph structure.
\begin{definition}
    \label{def:unlabeled_graph:homomorphism}
    Let \( G \) and \( H \) be unlabeled graphs. A \textbf{homomorphism of unlabeled graphs} $h: G \to H$ is a pair of functions $h_V: V(G) \to V(H) $ and $h_E: E(G) \to E(H)$ such that for every edge \( a: s \to t \) in \( G \), the edge \( h_E(a) \) is from \( h_V(s) \) to \( h_V(t) \).
\end{definition}


\begin{definition}
    \label{def:graph:composition}
    Let $f: (G,\lambda) \to (G',\lambda')$ and $g: (G',\lambda') \to (G'',\lambda'')$ be two homomorphisms of unlabeled graphs. Their composition, denoted $g \circ f$, is defined as the pair of functions $g_V \circ f_V$ and $g_E \circ f_E$.
\end{definition}
\begin{definition} 
    \label{def:graph}
     Let \(\Sigma\) be a finite set of labels. A \textbf{directed edge-labeled multigraph} is a structure \((V, E, \opn{dom}, \opn{cod},\lambda)\) where
    \begin{itemize}
        \item $V$ is a collection of distinct elements called \textbf{nodes} (or \textbf{objects}),
        \item $E$ is a collection of distinct elements, disjoint from $V$, called \textbf{edges} (or \textbf{arrows}),
        \item $\opn{dom}: E \to V$ is the \textbf{domain} (or \textbf{source}) function assigning to each edge its source node,
        \item $\opn{cod}: E \to V$ is the \textbf{codomain} (or \textbf{target}) function assigning to each edge its target node,
        \item $\lambda: E \to \Sigma$ is the \textbf{labeling} function assigning to each edge a label from $\Sigma$.
    \end{itemize}
    A directed edge-labeled multigraph is said to be \textbf{finite} if $V$ and $E$ are finite, and \textbf{discrete} if its edge set \(E\) is empty.
    For a directed edge-labeled multigraph \( G \), we write \( V(G) \) for its set of nodes, \( E(G) \) for its set of edges, \( \opn{dom}_G \) for its domain function, \( \opn{cod}_G \) for its codomain function, and \( \lambda_G \) for its labeling function. $a : s\overset{l}{\rightarrow} t$ denotes an arrow $a$ labeled by $l$ from $s$ to $t$.
\end{definition}
A directed edge-labeled multigraph may have multiple edges from one node to another, as well as loops (edges from a node to itself), because different edges $u,v$ can have the same source and target nodes, i.e. $\opn{dom}(u) = \opn{dom}(v)$ and $\opn{cod}(u) = \opn{cod}(v)$.

    Throughout, a directed edge-labeled multigraph will be simply referred to as a \enquote{labeled graph} or \enquote{graph} when the context makes it clear, and we fix a set $\Sigma$ of labels.

% \begin{definition} 
%     \label{def:graph}
%     Let \(\Sigma\) be a finite set of labels. A \textbf{directed edge-labeled multigraph} is an ordered pair \((G,\lambda)\) where \( G \) is an unlabeled graph and \( \lambda : E(G) \rightarrow \Sigma\) is an edge-labeling function. 
%     It is called \textbf{finite} if its underlying unlabeled graph is finite.
%     Throughout, a directed edge-labeled multigraph will be simply referred to as a \enquote{labeled graph} when the context makes it clear.
%     An arrow $a$ labeled by $l$ from $s$ to $t$ is denoted by $a : s\overset{l}{\rightarrow} t$.
% \end{definition}
%  The definition of labeled graphs extends the definition of unlabeled graphs, as unlabeled graphs can be seen as labeled graphs with a single label for all edges, i.e., the label function is a constant function.


\begin{example}
    Consider the graph shown in Figure~\ref{fig:preliminaries:labeled_graph}.
    It has two nodes which are marked with numbers to facilitate discussion, but in general, nodes are not labeled.
    There are five edges: two loops on node $1$ labeled \(a\) and \(b\), two edges from node $1$ to node $2$ both labeled \(a\), and one edge from node $2$ to node $1$ labeled \(a\). Note that the two edges from node $1$ to node $2$ are distinct edges, even though they share the same source, target, and label.

    \begin{figure}[H]
       \centering
       \resizebox{0.4\textwidth}{!}{
        \begin{tikzpicture}
            \graphbox{}{0mm}{0mm}{32mm}{28mm}{-10mm}{-14mm}{
                \node[draw,circle] (1) at (0,0) {1};
                \node[draw,circle] (2) at (2,0) {2};
                \draw[->] (1) edge[loop above] node[midway, above] {$a$} (1) ;
                \draw[->] (1) edge[loop below] node[midway, below] {$b$} (1) ;
                \draw[->] (1) edge[bend left] node[midway, above] {$a$}  (2)  ;
                \draw[->] (1) edge node[midway, above] {$a$}  (2)  ;
                \draw[->] (2) edge[bend left] node[midway, below] {$a$} (1)   ;
            }
        \end{tikzpicture}
       }
        \caption{Labeled graph.}
        \label{fig:preliminaries:labeled_graph}
    \end{figure}

\end{example}
  A homomorphism of labeled graphs is a homomorphism of unlabeled graphs that preserves the labels assigned to the edges. 
% \begin{definition}
%     \label{def:graph:homomorphism}
%     Let \( (G,\lambda) \) and \( (H,\lambda') \) be labeled graphs. A \textbf{homomorphism of labeled graphs} $h:(G,\lambda) \mathop{\rightarrow} (H,\lambda')$ is a homomorphism of unlabeled graphs such that for each edge \( a \) in \( G \), we have \( \lambda (a) \mathop{=} \lambda' (h_E (a)) \).
% \end{definition}

% \begin{definition}
%     \label{def:unlabeled_graph:homomorphism}
%     Let \( G \) and \( H \) be unlabeled graphs. A \textbf{homomorphism of unlabeled graphs} $h: G \mathop{\to} H$ 
% \end{definition}
\begin{definition}
    \label{def:graph:homomorphism}
    Let \( G \) and \( H \) be labeled graphs. A \textbf{homomorphism of labeled graphs}\index{Homomorphism of labeled graphs} $h:G \mathop{\rightarrow} H$ is a pair of functions $(h_V: V(G) \mathop{\to} V(H), h_E: E(G) \mathop{\to} E(H))$ such that for every edge \( a: s \mathop{\to} t \) in \( G \), the edge \( h_E(a) \) in $H$ is from node \( h_V(s) \) to node \( h_V(t) \), and \( \lambda_H(h_E(a)) \mathop{=} \lambda_G(a) \).
\end{definition}

\begin{notation}
    \label{notation:graph_homomorphism}
    We use the notation from Overbeek and Endrullis~\cite[Notation 1]{overbeek2023apbpotutorial} to visualize edge-labeled graph homomorphisms. Labeled graphs are enclosed in boxes with their names displayed in the top-left corner. Nodes and edges are assigned subsets of \(\mathbb{N}\) as identifiers, and these identifiers are chosen such that: (i) Each node \( y \) (resp. edge \( y \)) in the codomain graph is assigned the union of the identifiers of all nodes (resp. edges) in the domain graph that are mapped to \( y \); (ii) The graph homomorphism is uniquely determined by this assignment. To improve readability, we represent sets by listing their elements. Additionally, we omit identifiers when doing so does not cause confusion. An example of a homomorphism of labeled graphs is shown below.
    % in Figure~\ref{fig:preliminaries:graph_homomorphism}.
    In this example, the sets \(\{1\}\), \(\{2\}\), \(\{3\}\), \(\{4\}\), and \(\{3,4\}\) are represented as \(1\), \(2\), \(3\), \(4\), and \(3\ 4\), respectively. Edge identifiers are omitted.
    % \begin{figure}[H]
    %     \centering
    \begin{center}
        \resizebox{0.8\textwidth}{!}{
        \begin{tikzpicture}
            \graphbox{\( G \)}{00mm}{-20mm}{45mm}{24mm}{2mm}{-8mm}{
                \coordinate (o) at (-5mm,-8mm); 
                \node[draw,circle] (l1) at ($(o)+(-10mm,0mm)$) {1};
                \node[draw,circle] (l2) at ($(l1)+(3,0)$) {2};
                \node[draw,circle] (l3) at ($(l1)+(1,0)$) {3};
                \node[draw,circle] (l4) at ($(l1)+(2,0)$) {4};
                \draw[->] (l1) -- (l3) node[midway,above] {$a$};
                \draw[->] (l3) -- (l4) node[midway,above] {$b$};
                \draw[->] (l4) -- (l2) node[midway,above] {$a$};
            }  
            \graphbox{\( H \)}{50mm}{-20mm}{34mm}{24mm}{2mm}{-8mm}{
                \coordinate (o) at (0mm,-8mm); 
                \node[draw,circle] (l1) at ($(o)+(-10mm,0mm)$) {1};
                \node[draw,circle] (l2) at ($(l1)+(2,0)$) {2};
                \node[draw,circle] (l3) at ($(l1)+(1,0)$) {3\ 4};
                \draw[->] (l1) -- (l3) node[midway,above] {$a$};
                \draw[->] (l3) edge[loop above] node[midway,above] {$b$} (l3) ;
                \draw[->] (l3) -- (l2) node[midway,above] {$a$};
            }      
            \node () at (48mm,-29mm) {$\rightarrow$};
        \end{tikzpicture}
    }
    \end{center}
    % \caption{}
    % \label{fig:preliminaries:graph_homomorphism}
    % \end{figure}
\end{notation} 


\begin{definition}
    \label{def:graph:composition}
    Let $G, G'$ and $G''$ be labeled graphs, and $f: G \mathop{\to} G'$ and $g: G' \mathop{\to} G''$ homomorphisms of labeled graphs. Their composition, denoted by $g \circ f$, is defined as the component-wise composition of their node and edge mapping functions $(g_V \circ f_V, g_E \circ f_E)$.
\end{definition}

\begin{definition}
    \label{def:graph:unlabeled}
    An \textbf{unlabeled multigraph}\index{Unlabeled multigraph} \( G \) is a labeled graph such that the labeling function is constant, i.e., there is a label \( l \mathop{\in} \Sigma \) such that \( \lambda_G(e) \mathop{=} l \) for all edges \( e \) in $G$.
\end{definition}
    Hereafter, we omit the labeling function when referring to an unlabeled graph, and \( a: s \mathop{\to} t \) denotes an arrow \( a \) from \( s \) to \( t \). 
\subsection{Pushout and pullback}  
 The concepts and notation in this section follow the treatments in~\cite{pierce1991basic,barr1990category}.

\begin{definition}
    \label{def:cat}
    A \textbf{category} is an unlabeled graph \( C \) together with a total function \( u : V(C) \to E(C) \) and a partial function \( \star: E(C) \times E(C) \to E(C) \) such that 
        \begin{itemize}
            \item for all edges \( f:X \to Y \) and \( g:Y \to Z \), the edge \( f \star g :X \to Z \) is defined; 
            \item  for every node \( X \), \( u(X) \) is an edge from \( X \) to \( X \);
            \item for every \( f:X \to Y \), we have \(u(X) \star f = f = f \star u(Y)\);
            \item for all edges \( f \), \( g \) and \(h\), we have \( (f \star g) \star h = f \star (g \star h) \) whenever either side is defined.
        \end{itemize}
    Edges are called \textbf{morphisms}. The function $\star$ is called \textbf{composition}. For all \( X \in V(C) \), the edge \( u(X) \) is denoted by \( \operatorname{id}_X \) and is called the \textbf{identity} of the object \( X \).
    % \( C \) is called the \textbf{underlying graph} of the category \( \mathcal{C} \).
\end{definition} 
\begin{definition}
    A category \(\mathcal{C}\) is said to be \textbf{locally small} if for all objects \(X,Y\) in \(\mathcal{C}\), the collection $\opn{Hom}(X,Y)$ of morphisms from \(X\) to \(Y\) is a set (called a \textbf{hom-set}). For a locally small category, $\opn{Mono}(X,Y)$ denotes the set of all monomorphisms from $X$ to $Y$.
\end{definition}
 Throughout this section, we fix a locally small category \( \mathcal{C} \).
\begin{example}
    Consider the unlabeled graph shown in~\autoref{fig:preliminaries:category}. 
        \begin{figure}[H]
        \centering
        \begin{tikzpicture}
            \graphbox{}{0mm}{0mm}{32mm}{16mm}{-10mm}{-8mm}{
                \node[draw,circle] (1) at (0,0) {};
                \node[draw,circle] (2) at (2,0) {};
                \draw[->] (1) edge[loop above] node[midway, above] { } (1) ;
                \draw[->] (2) edge[loop above] node[midway, above] { } (2) ;
                \draw[->] (1) edge[bend left] node[midway, above] {}  (2)  ;
                \draw[->] (2) edge[bend left] node[midway, below] {} (1)   ;
            }
        \end{tikzpicture}
        \caption{}
        \label{fig:preliminaries:category}
    \end{figure}
    It can be considered as a category where the objects are the nodes and the morphisms are the paths between nodes; composition is path concatenation. The identity of a node is the self-loop of the node. There are at least three morphisms from the left node to itself: the identity morphism (the self-loop), the path that traverses the self-loop twice, and the path that goes to the right node and back.
\end{example}

% cat notation * 
\begin{notation}
    The composition of morphisms \( f : X \to Y \) and \( g : Y \to Z \) is written in diagrammatic order as \( f \star g \), rather than in functional order \( g \circ f \) as is common in litterature. The advantage is that, when reading from left to right, the morphisms appear in the same order as in the corresponding diagram, making the reasoning accompanying diagrams more intuitive. 
\end{notation}  

\begin{definition} 
    \label{def:cat:homo}
    A morphism \( f : X \to Y \) is said to be a \textbf{monomorphism} (is \textbf{monic}) if given any morphisms \( g,h: Z \to X  \), \( g \star f = h \star f \) implies \( g = h \). 
    In this case, we write $f : X \rightarrowtail Y$ to indicate that $f$ is a monomorphism.
\end{definition} 

\begin{example}
The category \(\mathbf{Set}\) has sets as objects and total functions between them as morphisms. For \(f\colon A\to B\) and \(g\colon B\to C\), composition is given by \(g\circ f\), and the identity morphism on a set \(A\) is the identity function \(\mathrm{id}_A\).
\end{example}

\begin{example} 
    Finite labeled graphs and their homomorphisms form a category, hereafter denoted by \textbf{Graph}. Its objects are labeled graphs, its morphisms are graph homomorphisms, and the monomorphisms are homomorphisms. 
    $\textbf{Graph}$ is locally small. 
\end{example}

% \begin{definition}[Span \cite{lowe2010graph}]
%     A pair \( (\alpha : A \to B,~\beta : A \to C) \) of morphisms with a common domain is called a \textbf{span}, denoted by \( B \overset{\alpha}{\leftarrow} A \overset{\beta}{\rightarrow} C \).
% \end{definition}
 
% \begin{definition}[Cospan]
%     A pair \( (\beta' : B \to D,~\alpha' : C \to D) \) of morphisms with a common codomain is called a \textbf{cospan}, denoted by \( B \overset{\beta'}{\rightarrow} D \overset{\alpha'}{\leftarrow} C \). 
% \end{definition} 
A span (resp. cospan) is a couple of morphisms with a common domain (resp. codomain).
\begin{definition}
An ordered pair \((\alpha : A \to B,\, \beta : A \to C)\) of morphisms with a common domain is called a \textbf{span} \cite{lowe2010graph}, denoted by
\(
B \overset{\alpha}{\leftarrow} A \overset{\beta}{\rightarrow} C
\). 
% An example of a span $(\alpha, \beta)$ is shown below.
Likewise, an ordered pair \((\beta' : B \to D,\, \alpha' : C \to D)\) of morphisms with a common codomain is called a \textbf{cospan}, denoted by
\(
B \overset{\beta'}{\rightarrow} D \overset{\alpha'}{\leftarrow} C
\). 
\end{definition}
\begin{example}
Consider the diagram in the category \textbf{Graph} shown in~\autoref{fig:preliminaries:span_cospan}. $(\alpha, \beta)$ is a span, and $(\beta', \alpha')$ is a cospan.
\end{example}


\begin{figure}[H]
   \centering
        \resizebox{0.7\textwidth}{!}{
        \begin{tikzpicture} 
            \graphbox{\( L \)}{40mm}{20mm}{34mm}{12mm}{2mm}{2mm}{
                \coordinate (o) at (0mm,-8mm); 
                \node[draw,circle] (l1) at ($(o)+(-10mm,0mm)$) {1};
                \node[draw,circle] (l2) at ($(l1)+(2,0)$) {2};
                \node[draw,circle,red] (l3) at ($(l1) + (1,0)$) {3};
                \draw[->,red] (l1) -- (l3) node[midway,above] {$a$};
                \draw[->,red] (l3) -- (l2) node[midway,above] {$a$};
            } 
    
            \graphbox{\( K \)}{0mm}{0mm}{34mm}{12mm}{2mm}{2mm}{
                \coordinate (o) at (0mm,-8mm); 
                \node[draw,circle] (l1) at ($(o)+(-10mm,0mm)$) {1};
                \node[draw,circle] (l2) at ($(l1)+(2,0)$) {2};
            }  
            \graphbox{\(G = L \cup C \)}{90mm}{5mm}{34mm}{27mm}{2mm}{-3mm}{
                \coordinate (o) at (0mm,-8mm); 
                \node[draw,circle] (l1) at ($(o)+(-10mm,0mm)$) {1};
                \node[draw,circle] (l2) at ($(l1)+(2,0)$) {2};
                \node[draw,circle,red] (l3) at ($(l1) + (1,0)$) {3};
                \node[draw,circle,blue] (l4) at ($(l2) + (0,-1)$) {6};
                \draw[->,red] (l1) -- (l3) node[midway,above] {$a$};
                \draw[->,red] (l3) -- (l2) node[midway,above] {$a$};
                \draw[->,blue] (l2) -- (l4) node[midway,right] {$a$};
                \node[draw,circle,blue] (l6) at ($(l1) + (0,-1)$) {7};
                \draw[<-,blue] (l1) -- (l6) node[midway,left] {$a$};
                \draw[->,blue] (l2) edge[out=-135,in=-45]node[midway,below] {$a$} (l1) ;
            }   
     
            \graphbox{\( C \)}{40mm}{-20mm}{34mm}{27mm}{2mm}{-3mm}{
                \coordinate (o) at (0mm,-8mm); 
                \node[draw,circle] (l1) at ($(o)+(-10mm,0mm)$) {1};
                \node[draw,circle] (l2) at ($(l1)+(2,0)$) {2};
                \node[draw,circle,blue] (l4) at ($(l2) + (0,-1)$) {6};
                \draw[->,blue] (l2) -- (l4) node[midway,right] {$a$};
                \draw[->,blue] (l2) edge[out=-135,in=-45]node[midway,below] {$a$} (l1) ;
                \node[ draw,circle,blue] (l6) at ($(l1) + (0,-1)$) {7};
                \draw[<-,blue] (l1) -- (l6) node[midway,left] {$a$};
            }      
            % K to L
            \draw[>->] (17mm,5mm) -- node[above] {$\alpha$} (37mm,15mm);
            % C to G
            \draw[>->] (76mm,-28mm)-- node[below] {$\alpha'$} (104mm,-24mm) ;
            % K to C
            \draw[>->] (17mm,-17mm) -- node[below] {$\beta$} (37mm,-28mm);
            % L to G
            \draw[>->] (76mm,16mm) -- node[above] {$\beta'$} (104mm,7mm);
            % \node () at (57mm,-6mm) {$PO$};
        \end{tikzpicture}
        }
        \caption{}
        \label{fig:preliminaries:span_cospan}
\end{figure}

\begin{definition}[\cite{barr1990category}]
    \label{def:cat:diagram}
    Let $\mathcal{C}$ be a category, and \( G \) an unlabeled graph. A \textbf{diagram} (of shape \( G \)) is a homomorphism of unlabeled graphs \( h : G \to \mathcal{C} \) where \( \mathcal{C} \) is considered as an unlabeled graph. A diagram is said to be \textbf{commutative} if, for all nodes \( u \), \( v \), and any two paths from \( u \) to \( v \) in the unlabeled graph \( G \):

    \begin{center}
    \resizebox{12cm}{!}{
        \begin{tikzpicture}
        \node (u) at (0,0) {\( u \)};
        \node (k1) at (2,0.5) {\( k_1 \)};
        \node (k2) at (4,0.5) {\( k_2 \)};
        \node (ketc) at (6,0.5) {\( \dots \)};
        \node (knm2) at (8,0.5) {\( k_{n-2} \)};
        \node (knm1) at (10,0.5) {\( k_{n-1} \)};
        \node (v) at (12,0) {\( v \)};
        \node (l1) at (2,-0.5) {\( l_1 \)};
        \node (l2) at (4,-0.5) {\( l_2 \)};
        \node (letc) at (6,-0.5) {\( \dots \)};
        \node (lnm2) at (8,-0.5) {\( l_{m-2} \)};
        \node (lnm1) at (10,-0.5) {\( l_{m-1} \)};
        \draw[->] (u) -- (k1) node [midway,above] {\( s_1 \)};
        \draw[->] (k1) -- (k2) node [midway,above] {\( s_2 \)};
        \draw[->] (k2) -- (ketc);
        \draw[->] (ketc) -- (knm2); 
        \draw[->] (knm2) -- (knm1) node[midway,above] {\( s_{n-2} \)}; 
        \draw[->] (knm1) -- (v) node[midway,above] {\( s_{n-1} \)}; 
        \draw[->] (u) -- (l1) node[midway,below] {\( t_1 \)};
        \draw[->] (l1) -- (l2) node[midway,below] {\( t_2 \)};
        \draw[->] (l2) -- (letc);
        \draw[->] (letc) -- (lnm2); 
        \draw[->] (lnm2) -- (lnm1) node[midway,below] {\( t_{m-2} \)}; 
        \draw[->] (lnm1) -- (v) node[midway,below] {\( t_{m-1} \)}; 
        \end{tikzpicture}
    }
    \end{center}
    \noindent
    the equality \( h(s_1) \star h(s_2) \star \dots  \star h(s_{n-1}) = h(t_1) \star h(t_2) \star \dots  \star h(t_{m-1}) \) holds.
\end{definition}

\begin{example}
    A commutative diagram in the category \textbf{Graph} of finite directed edge-labeled multigraphs is illustrated in~\autoref{fig:preliminaries:commutative_diagram}. The symbol $\mathop{=}$ in the center of the diagram
    is used to indicate that the two paths from node 1 to node 2 are equal.
 
    \begin{figure}[H]
        \centering
        \resizebox{0.7\textwidth}{!}{
        \begin{tikzpicture} 
            \graphbox{\( L \)}{40mm}{20mm}{34mm}{12mm}{2mm}{2mm}{
                \coordinate (o) at (0mm,-8mm); 
                \node[draw,circle] (l1) at ($(o)+(-10mm,0mm)$) {1};
                \node[draw,circle] (l2) at ($(l1)+(2,0)$) {2};
                \node[draw,circle,red] (l3) at ($(l1) + (1,0)$) {3};
                \draw[->,red] (l1) -- (l3) node[midway,above] {$a$};
                \draw[->,red] (l3) -- (l2) node[midway,above] {$a$};
            } 
    
            \graphbox{\( K \)}{0mm}{0mm}{34mm}{12mm}{2mm}{2mm}{
                \coordinate (o) at (0mm,-8mm); 
                \node[draw,circle] (l1) at ($(o)+(-10mm,0mm)$) {1};
                \node[draw,circle] (l2) at ($(l1)+(2,0)$) {2};
            }  
            \graphbox{\(G = L \cup C \)}{90mm}{5mm}{34mm}{27mm}{2mm}{-3mm}{
                \coordinate (o) at (0mm,-8mm); 
                \node[draw,circle] (l1) at ($(o)+(-10mm,0mm)$) {1};
                \node[draw,circle] (l2) at ($(l1)+(2,0)$) {2};
                \node[draw,circle,red] (l3) at ($(l1) + (1,0)$) {3};
                \node[draw,circle,blue] (l4) at ($(l2) + (0,-1)$) {6};
                \draw[->,red] (l1) -- (l3) node[midway,above] {$a$};
                \draw[->,red] (l3) -- (l2) node[midway,above] {$a$};
                \draw[->,blue] (l2) -- (l4) node[midway,right] {$a$};
                \node[draw,circle,blue] (l6) at ($(l1) + (0,-1)$) {7};
                \draw[<-,blue] (l1) -- (l6) node[midway,left] {$a$};
                \draw[->,blue] (l2) edge[out=-135,in=-45]node[midway,below] {$a$} (l1) ;
            }   
     
            \graphbox{\( C \)}{40mm}{-20mm}{34mm}{27mm}{2mm}{-3mm}{
                \coordinate (o) at (0mm,-8mm); 
                \node[draw,circle] (l1) at ($(o)+(-10mm,0mm)$) {1};
                \node[draw,circle] (l2) at ($(l1)+(2,0)$) {2};
                \node[draw,circle,blue] (l4) at ($(l2) + (0,-1)$) {6};
                \draw[->,blue] (l2) -- (l4) node[midway,right] {$a$};
                \draw[->,blue] (l2) edge[out=-135,in=-45]node[midway,below] {$a$} (l1) ;
                \node[ draw,circle,blue] (l6) at ($(l1) + (0,-1)$) {7};
                \draw[<-,blue] (l1) -- (l6) node[midway,left] {$a$};
            }      
            % K to L
            \draw[>->] (17mm,5mm) -- node[above] {$\alpha$} (37mm,15mm);
            % C to G
            \draw[>->] (76mm,-28mm)-- node[below] {$\alpha'$} (104mm,-24mm) ;
            % K to C
            \draw[>->] (17mm,-17mm) -- node[below] {$\beta$} (37mm,-28mm);
            % L to G
            \draw[>->] (76mm,16mm) -- node[above] {$\beta'$} (104mm,7mm);
            \node () at (57mm,-6mm) {$\mathop{=}$};
        \end{tikzpicture}
        }
        \caption{}
        \label{fig:preliminaries:commutative_diagram}
    \end{figure}
\end{example}

\begin{notation}
    When the context makes it clear, \( h_{AB} \) denotes a morphism \( h : A \to B \), and we refer to diagrams by listing their nodes, as is standard in geometry. 
    For example, we denote the diagram shown in \autoref{def:cat:po} by \( ACDB \) or \( ABDC \).
\end{notation}   

The pushout is a construction in category theory that can often be thought of as the construction of a new structure from two given structures by gluing them along a common interface structure.
\begin{definition}
    \label{def:cat:po} 
    A \textbf{pushout} of a span \( B \overset{\alpha}{\leftarrow} A \overset{\beta}{\rightarrow} C \), as shown in~\autoref{fig:preliminaries:pushout_sdfkjasdlgjfl}
    is defined as a cospan \( B \overset{\beta'}{\rightarrow} D \overset{\alpha'}{\leftarrow} C \) such that \( \alpha \star \beta' = \beta \star \alpha' \), and for every cospan \( B \overset{\gamma'}{\rightarrow} E \overset{\gamma}{\leftarrow} C \), if \( \alpha \star \gamma' = \beta \star \gamma \) holds, then there exists a unique morphism \(\delta : D \to E\) such that \( \gamma' = \beta' \star \delta \) and \( \gamma = \alpha' \star \delta \).
    \begin{figure}[H]
        \centering
        \resizebox{0.4\textwidth}{!}{
            \begin{tikzpicture}
                    \node (i) at (0,0) {A};
                    \node (r) at (1,1) {B};
                    \node (c) at (1,-1) {C};
                    \node (h) at (2,0) {D};
                    % \node () at (1,-1) {\( \Delta \)};
                    \draw[->]  (i) -- (r) node [midway,left] {$ \alpha $};
                    \draw[->] (c) -- (h) node [midway,left] {$ \alpha' $};
                    \draw[->] (r) -- (h) node[midway, left] {$ \beta' $};
                    \draw[->] (i) -- (c) node[midway, left] {$ \beta $};
                    \node (d') at (4,0) {E};
                    \draw[->] (c) -- (d') node [midway,below]{$ \gamma $};
                    \draw[->] (r) -- (d') node [midway,above]{$ \gamma' $};
                    \draw[->,dashed] (h) -- (d') node [midway]{$ \delta $};
                \end{tikzpicture}
        }
        \caption{}
        \label{fig:preliminaries:pushout_sdfkjasdlgjfl}
    \end{figure}
The diagram involving \( (\alpha, \beta, \alpha', \beta') \) is called the \textbf{pushout square}, with \(D\) as the \textbf{pushout object}. The existence of a unique morphism is known as the \textbf{universal mapping property of the pushout}.
\end{definition} 

\begin{example}
    \label{ex:cat:posfjsdlkgja}
     Pushouts of two morphisms always exist in \(\mathbf{Set}\), and (up to isomorphism) can be described as follows. Let
    \( B \overset{\alpha}{\leftarrow} A \overset{\beta}{\rightarrow} C \) be a span. Its pushout is the cospan \( B \overset{\beta'}{\rightarrow} D \overset{\alpha'}{\leftarrow} C \) where the pushout object of \((\alpha,\beta)\) is the set
    \[
    D \;=\; (B + C)/{\sim}
    \]
    where \(B + C\) denotes the disjoint union of $B$ and $C$ and \(\sim\) is the smallest equivalence relation that includes \(\set{(\alpha(a),\beta(b))\mid a \in A }\). The maps
    \(\beta' \colon B\to D\) and \(\alpha' \colon C\to D\) send each element to its equivalence class.

    %  Consider the functions \(b\) and \(c\) shown in~\autoref{fig:preliminaries:a_rewriting_step_dfjalsdkjflg} in the category \(\mathbf{Set}\). Each set is drawn as a box; elements of a set are represented by circles. The numbers inside circles indicate how the functions map those elements. 
    %  The element labeled by $n$ in set $X$ is denoted by $n_X$ to avoid ambiguity.
    %  The equivalence relation \(\sim\) is the set \(\set{(1_B,1_C),(2_B,2_C)}\). The disjoint union of \(B\) and \(C\) is $D = \set{1_B,2_B,4_B,1_C,2_C,3_C}$. The quotient set $D/{\sim}$ is $\set{[1_B],[2_B],[3_C],[4_B]}$, where $[x]$ denotes the equivalence class of $x$. Note that $[1_C] = [1_B]$ and $[2_C] = [2_B]$, and $\set{[1_B],[2_B],[3_C],[4_B]}$ is isomorphic to $\set{1_D,2_D,3_D,4_D}$.
    For example, consider the functions \(\alpha\) and \(\beta\) shown in Figure~\ref{fig:preliminaries:a_rewriting_step_dfjalsdkjflg} in the category \(\mathbf{Set}\) where each set is drawn as a box and its elements are represented by circles. The numbers inside circles indicate how the functions map those elements.
    \begin{figure}[H]
      \centering 
      \resizebox{0.7\textwidth}{!}{
      \begin{tikzpicture}
          \graphbox{\( A\)}{40mm}{-3mm}{34mm}{12mm}{2mm}{2mm}{
              \coordinate (o) at (0mm,-8mm); 
              \node[draw,circle] (l1) at ($(o)+(-10mm,0mm)$) {1};
              \node[draw,circle] (l2) at ($(l1)+(2,0)$) {2};
          }  
          \graphbox{\( B \)}{80mm}{-3mm}{45mm}{12mm}{2mm}{2mm}{
              \coordinate (o) at (-5mm,-8mm); 
              \node[draw,circle] (l1) at ($(o)+(-10mm,0mm)$) {1};
              \node[draw,circle] (l2) at ($(l1)+(3,0)$) {2};
              \node[draw,circle] (l3) at ($(l1) + (1,0)$) {4};
            %   \node[draw,circle] (l4) at ($(l1) + (2,0)$) {5};
            %   \draw[ ] (l1) -- (l3) node[midway,above] {$a$};
            %   \draw[ ] (l3) -- (l4) node[midway,above] {$b$};
            %   \draw[ ] (l4) -- (l2) node[midway,above] {$a$};
          }     
          \graphbox{\( C  \)}{40mm}{-22mm}{34mm}{22mm}{2mm}{-3mm}{
              \coordinate (o) at (0mm,-3mm); 
              \node[draw,circle] (l1) at ($(o)+(-10mm,0mm)$) {1};
              \node[draw,circle] (l2) at ($(l1)+(2,0)$) {2};
            %   \node[draw,circle] (l4) at ($(l2) + (0,-1)$) {6};
              \node[ draw,circle] (l6) at ($(l1) + (0,-1)$) {3};
            %   \draw[ ] (l1) -- (l6) node[midway,left] {$a$};
            %   \draw[ ] (l2) -- (l4) node[midway,right] {$a$};
          }    
          \graphbox{\( D \)}{80mm}{-22mm}{45mm}{22mm}{2mm}{-3mm}{
              \coordinate (o) at (-5mm,-3mm); 
              \node[draw,circle] (l1) at ($(o)+(-10mm,0mm)$) {1};
              \node[draw,circle] (l2) at ($(l1)+(3,0)$) {2};
              \node[draw,circle] (l3) at ($(l1) + (1,0)$) {4};
            %   \node[draw,circle] (l4) at ($(l1) + (2,0)$) {5};
            %   \node[ draw,circle] (l5) at ($(l2) + (0,-1)$) {6};
              \node[ draw,circle] (l6) at ($(l1) + (0,-1)$) {3};
            %   \draw[ ] (l1) -- (l6) node[midway,left] {$a$};
            %   \draw[] (l1) -- (l3) node[midway,above] {$a$};
            %   \draw[] (l3) -- (l4) node[midway,above] {$b$};
            %   \draw[ ] (l4) -- (l2) node[midway,above] {$a$};
            %   \draw[ ] (l2) -- (l5) node[midway,right] {$a$};
          }    
          \node () at (77mm,-8mm) {\( \overset{\alpha}{\rightarrowtail} \)}; % K -> R
          \node () at (58mm,-18mm) {\( \beta\downarrowtail \)};
          \node () at (102mm,-18mm) {\( \beta'\downarrowtail \)};
          \node () at (77mm,-33mm) {\( \overset{\alpha'}{\rightarrowtail} \)}; % C -> H
      \end{tikzpicture}
      }
      \caption{}
      \label{fig:preliminaries:a_rewriting_step_dfjalsdkjflg}
  \end{figure}
    Throughout this example, an element labeled \(n\) in a set \(X\) is denoted by \(n_X\) to avoid ambiguity.        
    The binary relation \(\sim\) is the reflexive, symmetric and transitive closure of the binary relation $\{(1_B,1_C),(2_B,2_C)\}$.
 
    The disjoint union of \(B\) and \(C\) is
    \[ 
    D' = \{1_B,2_B,4_B,1_C,2_C,3_C\},
    \]
    and the quotient \(D'/\sim\) is
    \[
    D'/\sim = \{[1_B],[2_B],[3_C],[4_B]\},
    \]
    We have \([1_C]=[1_B]\) and \([2_C]=[2_B]\), and the maps
    \(\beta'' \colon B\to D'\) and \(\alpha'' \colon C\to D'\) send each element to its equivalence class. Note that \(\{[1_B],[2_B],[3_C],[4_B]\}\) is isomorphic to \(D\), shown in \autoref{fig:preliminaries:a_rewriting_step_dfjalsdkjflg}, which is expected because the pushout of a span is unique up to isomorphism.
\end{example}

\begin{proposition}{\cite{corradini1997algebraic}[p.188]}
    \label{prop:pushout_graph_always_exists}
    In category \textbf{Graph} the pushout of two arrows always exists: It can be computed componentwise (as a pushout in \textbf{Set}) for the nodes and for the edges, and the source, target, and labeling mappings are uniquely determined.
\end{proposition}

\begin{example}
    The diagram in~\autoref{fig:preliminaries:pushout_injective} is a pushout square in the category \textbf{Graph}. In this example, both $\alpha$ and $\beta$ are injective morphisms. Therefore, the pushout object $G$ can be constructed easily by taking the interface graph $K$ and adding elements from $L$ and $C$ which are not present in $K$.
    \begin{figure}[H]
        \centering
        \resizebox{0.7\textwidth}{!}{
        \begin{tikzpicture} 
            \graphbox{\( L \)}{40mm}{20mm}{34mm}{12mm}{2mm}{2mm}{
                \coordinate (o) at (0mm,-8mm); 
                \node[draw,circle] (l1) at ($(o)+(-10mm,0mm)$) {1};
                \node[draw,circle] (l2) at ($(l1)+(2,0)$) {2};
                \node[draw,circle,red] (l3) at ($(l1) + (1,0)$) {3};
                \draw[->,red] (l1) -- (l3) node[midway,above] {$a$};
                \draw[->,red] (l3) -- (l2) node[midway,above] {$a$};
            } 
    
            \graphbox{\( K \)}{0mm}{0mm}{34mm}{12mm}{2mm}{2mm}{
                \coordinate (o) at (0mm,-8mm); 
                \node[draw,circle] (l1) at ($(o)+(-10mm,0mm)$) {1};
                \node[draw,circle] (l2) at ($(l1)+(2,0)$) {2};
            }  
            \graphbox{\(G  \)}{90mm}{5mm}{34mm}{27mm}{2mm}{-3mm}{
                \coordinate (o) at (0mm,-8mm); 
                \node[draw,circle] (l1) at ($(o)+(-10mm,0mm)$) {1};
                \node[draw,circle] (l2) at ($(l1)+(2,0)$) {2};
                \node[draw,circle,red] (l3) at ($(l1) + (1,0)$) {3};
                \node[draw,circle,blue] (l4) at ($(l2) + (0,-1)$) {6};
                \draw[->,red] (l1) -- (l3) node[midway,above] {$a$};
                \draw[->,red] (l3) -- (l2) node[midway,above] {$a$};
                \draw[->,blue] (l2) -- (l4) node[midway,right] {$a$};
                \node[draw,circle,blue] (l6) at ($(l1) + (0,-1)$) {7};
                \draw[<-,blue] (l1) -- (l6) node[midway,left] {$a$};
                \draw[->,blue] (l2) edge[out=-135,in=-45]node[midway,below] {$a$} (l1) ;
            }   
     
            \graphbox{\( C \)}{40mm}{-20mm}{34mm}{27mm}{2mm}{-3mm}{
                \coordinate (o) at (0mm,-8mm); 
                \node[draw,circle] (l1) at ($(o)+(-10mm,0mm)$) {1};
                \node[draw,circle] (l2) at ($(l1)+(2,0)$) {2};
                \node[draw,circle,blue] (l4) at ($(l2) + (0,-1)$) {6};
                \draw[->,blue] (l2) -- (l4) node[midway,right] {$a$};
                \draw[->,blue] (l2) edge[out=-135,in=-45]node[midway,below] {$a$} (l1) ;
                \node[ draw,circle,blue] (l6) at ($(l1) + (0,-1)$) {7};
                \draw[<-,blue] (l1) -- (l6) node[midway,left] {$a$};
            }      
            % K to L
            \draw[>->] (17mm,5mm) -- node[above] {$\alpha$} (37mm,15mm);
            % C to G
            \draw[>->] (76mm,-28mm)-- node[below] {$\alpha'$} (104mm,-24mm) ;
            % K to C
            \draw[>->] (17mm,-17mm) -- node[below] {$\beta$} (37mm,-28mm);
            % L to G
            \draw[>->] (76mm,16mm) -- node[above] {$\beta'$} (104mm,7mm);
            \node () at (57mm,-6mm) {$PO$};
        \end{tikzpicture}
        }
        \caption{Pushout square with injective morphisms.}
        \label{fig:preliminaries:pushout_injective}
    \end{figure}
\end{example}

\begin{example}
    The diagram in~\autoref{fig:preliminaries:pushout_non_injective} is a pushout square in the category \textbf{Graph}. In this example, $\beta$ is not injective. Therefore, some elements are merged in the pushout object $G$.

    \begin{figure}[H]
        \centering 
        \resizebox{0.7\textwidth}{!}{
        \begin{tikzpicture} 
            \graphbox{\( L \)}{40mm}{20mm}{34mm}{20mm}{2mm}{2mm}{
                \coordinate (o) at (0mm,-8mm); 
                \node[draw,circle] (l1) at ($(o)+(-10mm,0mm)$) {1};
                \node[draw,circle] (l2) at ($(l1)+(2,0)$) {2};
                \draw[->,red] (l2) edge[out=-135,in=-45]node[midway,below] {$a$} (l1) ;
                \node[draw,circle,red] (l3) at ($(l1) + (1,0)$) {3};
                \draw[->,red] (l1) -- (l3) node[midway,above] {$a$};
                \draw[->,red] (l3) -- (l2) node[midway,above] {$a$};
            } 
    
            \graphbox{\( K \)}{0mm}{0mm}{34mm}{12mm}{2mm}{2mm}{
                \coordinate (o) at (0mm,-8mm); 
                \node[draw,circle] (l1) at ($(o)+(-10mm,0mm)$) {1};
                \node[draw,circle] (l2) at ($(l1)+(2,0)$) {2};
            }  
            \graphbox{\(G  \)}{110mm}{5mm}{34mm}{40mm}{2mm}{-3mm}{
                \coordinate (o) at (0mm,-20mm); 
                 \node[draw,circle] (l1) at ($(o)+(0,0)$) {1\ 2};
                % \node[draw,circle] (l1) at ($(o)+(-10mm,0mm)$) {1};
                % \node[draw,circle] (l2) at ($(l1)+(2,0)$) {2};
                \draw[->,red] (l1) edge[loop below] node[midway, below] {$a$} (l1) ;
                \node[draw,circle,red] (l3) at ($(l1) + (0,1.4)$) {3};
                \node[draw,circle,blue] (l4) at ($(l1) + (1,-1)$) {6};
                \draw[->,red] (l1) edge[bend left] node[midway,left] {$a$} (l3);
                \draw[->,red] (l3) edge[bend left] node[midway,right] {$a$} (l1);
                \draw[->,blue] (l1) edge  node[midway,right] {$a$} (l4);
                \node[draw,circle,blue] (l6) at ($(l1) + (-1,-1)$) {7};
                \draw[<-,blue] (l1) edge node[midway,left] {$a$} (l6) ;
                % \draw[->,blue] (l1) edge[out=-135,in=-45]node[midway,below] {$a$} (l1) ;
            }   
     
            \graphbox{\( C \)}{40mm}{-20mm}{34mm}{35mm}{2mm}{-3mm}{
                \coordinate (o) at (0mm,-14mm); 
                \node[draw,circle] (l1) at ($(o)+(0,0)$) {1\ 2};
                % \node[draw,circle] (l2) at ($(l1)+(2,0)$) {2};
                \node[draw,circle,blue] (l4) at ($(l1) + (1,-1)$) {6};
                \draw[->,blue] (l1) -- (l4) node[midway,right] {$a$};
                % \draw[->,blue] (l1) edge[loop above] node[midway, above] {$a$} (l1) ;
                \node[ draw,circle,blue] (l6) at ($(l1) + (-1,-1)$) {7};
                \draw[<-,blue] (l1) -- (l6) node[midway,left] {$a$};
            }
            % K to L  
            \draw[>->] (17mm,5mm) -- node[above] {$\alpha$} (37mm,15mm);
            % C to G
            \draw[>->] (76mm,-28mm)-- node[below] {$\alpha'$} (104mm,-24mm) ;
            % K to C
            \draw[->] (17mm,-17mm) -- node[below] {$\beta$} (37mm,-28mm);
            % L to G
            \draw[->] (76mm,16mm) -- node[above] {$\beta'$} (104mm,7mm);
            \node () at (57mm,-6mm) {$PO$};
        \end{tikzpicture}
        }
        \caption{Pushout square with a non-injective morphism.}
        \label{fig:preliminaries:pushout_non_injective}
    \end{figure} 
\end{example}
The \emph{pullback}, defined below, is the dual construction of the pushout in category theory, and it can be thought of as construction of the interface structure along which two structures are glued together. 
\begin{definition} 
    \label{def:cat:pb}
   A \textbf{pullback} of a cospan \(B \overset{\beta'}{\rightarrow} D \overset{\alpha'}{\leftarrow} C \), as shown in~\autoref{fig:preliminaries:pullback_ssdsfd},
is a span \( B \overset{\alpha}{\leftarrow} A \overset{\beta}{\rightarrow} C \) such that \( \alpha \star \beta' = \beta \star \alpha' \), and for every span \( B \overset{\gamma'}{\leftarrow} E \overset{\gamma}{\rightarrow} C \) if \(\gamma' \star \beta' = \gamma \star \alpha'\) holds, then there exists a unique morphism \(\delta: E \to A\) such that $\gamma' = \delta \star \alpha$ and $\gamma = \delta \star \beta$. 
    \begin{figure}[H]
        \centering  
        \resizebox{0.4\textwidth}{!}{
                    \begin{tikzpicture}
                        \node (i) at (0,0) {A};
                        \node (r) at (1,1) {B};
                        \node (c) at (1,-1) {C};
                        \node (h) at (2,0) {D}; 
                        % \node () at (1,-1) {\( \Delta \)};
                        \draw[->]  (i) -- (r) node [midway,left] {$\alpha$};
                        \draw[->] (c) -- (h) node [midway,left] {$\alpha'$};
                        \draw[->] (r) -- (h) node[midway, left] {$\beta'$};
                        \draw[->] (i) -- (c) node[midway, left] {$\beta$};
                        \node (d') at (-2,0) {E};
                        \draw[<-] (c) -- (d') node [midway,below]{$\gamma$};
                        \draw[<-] (r) -- (d') node [midway,above]{$\gamma'$};
                        \draw[->, dashed] (d') -- (i) node [midway]{$\delta$};
                    \end{tikzpicture}
        }
        \caption{}
        \label{fig:preliminaries:pullback_ssdsfd}
    \end{figure}
The diagram involving \( (\alpha, \beta, \alpha', \beta') \) is called the \textbf{pullback square}, with \(A\) as the \textbf{pullback object}. The existence of a unique morphism is known as the \textbf{universal mapping property of the pullback}.
\end{definition} 
In the category \textbf{Graph} pullbacks always exist and are computed pointwise: take the pullback in \textbf{Set} of the node sets and of the edge sets, and equip the resulting graph with source, target and labeling maps induced componentwise; the projection graph homomorphisms are the corresponding componentwise projections and satisfy the universal property.
\begin{example} 
    In the category \textbf{Set},
    the pullback 
    of a cospan \(B \overset{\beta'}{\rightarrow} D \overset{\alpha'}{\leftarrow} C \) is the span \( B \overset{\alpha}{\leftarrow} A \overset{\beta}{\rightarrow} C \) where
    the pullback object $A$ is the set $A = \set{(b,c) \mid \beta (c) = \alpha (b)} \subseteq B \times C$; $\alpha$ and $\beta$ are defined as the corresponding projections, e.g., $\alpha((b, c)) = b$ and $\beta((b, c)) = c$.

    Consider the cospan \(B \overset{\beta'}{\rightarrow} D \overset{\alpha'}{\leftarrow} C \) shown in~\autoref{fig:preliminaries:a_rewriting_step_dfjalsdkdfsdfjflg}. 
    The span \( B \overset{\alpha}{\leftarrow} A \overset{\beta}{\rightarrow} C \) shown in~\autoref{fig:preliminaries:a_rewriting_step_dfjalsdkdfsdfjflg} is the pullback of this cospan.
    Indeed, the pullback object can be taken as the set $A' = \set{(1_B,1_C),(2_B,2_C)} \subseteq B \times C$, and the morphisms $\alpha'' \colon A' \to B$ and $\beta'' \colon A' \to C$ are defined as the corresponding projections, e.g., $\alpha''((x,y)) = x$ and $\beta''((x,y)) = y$. Note that $A$ is isomorphic to $A'$, which is expected because the pullback of a cospan is unique up to isomorphism.
      \begin{figure}[H]
      \centering 
      \resizebox{0.7\textwidth}{!}{
      \begin{tikzpicture}
          \graphbox{\( A\)}{40mm}{-3mm}{34mm}{12mm}{2mm}{2mm}{
              \coordinate (o) at (0mm,-8mm); 
              \node[draw,circle] (l1) at ($(o)+(-10mm,0mm)$) {1};
              \node[draw,circle] (l2) at ($(l1)+(2,0)$) {2};
          }  
          \graphbox{\( B \)}{80mm}{-3mm}{45mm}{12mm}{2mm}{2mm}{
              \coordinate (o) at (-5mm,-8mm); 
              \node[draw,circle] (l1) at ($(o)+(-10mm,0mm)$) {1};
              \node[draw,circle] (l2) at ($(l1)+(3,0)$) {2};
              \node[draw,circle] (l3) at ($(l1) + (1,0)$) {4};
            %   \node[draw,circle] (l4) at ($(l1) + (2,0)$) {5};
            %   \draw[ ] (l1) -- (l3) node[midway,above] {$a$};
            %   \draw[ ] (l3) -- (l4) node[midway,above] {$b$};
            %   \draw[ ] (l4) -- (l2) node[midway,above] {$a$};
          }     
          \graphbox{\( C  \)}{40mm}{-22mm}{34mm}{22mm}{2mm}{-3mm}{
              \coordinate (o) at (0mm,-3mm); 
              \node[draw,circle] (l1) at ($(o)+(-10mm,0mm)$) {1};
              \node[draw,circle] (l2) at ($(l1)+(2,0)$) {2};
            %   \node[draw,circle] (l4) at ($(l2) + (0,-1)$) {6};
              \node[ draw,circle] (l6) at ($(l1) + (0,-1)$) {3};
            %   \draw[ ] (l1) -- (l6) node[midway,left] {$a$};
            %   \draw[ ] (l2) -- (l4) node[midway,right] {$a$};
          }    
          \graphbox{\( D \)}{80mm}{-22mm}{45mm}{22mm}{2mm}{-3mm}{
              \coordinate (o) at (-5mm,-3mm); 
              \node[draw,circle] (l1) at ($(o)+(-10mm,0mm)$) {1};
              \node[draw,circle] (l2) at ($(l1)+(3,0)$) {2};
              \node[draw,circle] (l3) at ($(l1) + (1,0)$) {4};
            %   \node[draw,circle] (l4) at ($(l1) + (2,0)$) {5};
            %   \node[ draw,circle] (l5) at ($(l2) + (0,-1)$) {6};
              \node[ draw,circle] (l6) at ($(l1) + (0,-1)$) {3};
            %   \draw[ ] (l1) -- (l6) node[midway,left] {$a$};
            %   \draw[] (l1) -- (l3) node[midway,above] {$a$};
            %   \draw[] (l3) -- (l4) node[midway,above] {$b$};
            %   \draw[ ] (l4) -- (l2) node[midway,above] {$a$};
            %   \draw[ ] (l2) -- (l5) node[midway,right] {$a$};
          }    
          \node () at (77mm,-8mm) {\( \overset{\alpha}{\rightarrowtail} \)};
          \node () at (58mm,-18mm) {\( \beta\downarrowtail \)};
          \node () at (102mm,-18mm) {\( \beta'\downarrowtail \)};
          \node () at (77mm,-33mm) {\( \overset{\alpha'}{\rightarrowtail} \)};
      \end{tikzpicture}
      }
      \caption{}
      \label{fig:preliminaries:a_rewriting_step_dfjalsdkdfsdfjflg}
  \end{figure}
\end{example}

\begin{example}
    Consider the diagram shown in~\autoref{fig:preliminaries:pullback} in the category \textbf{Graph}.
    The span $L \overset{\alpha}{\leftarrow} K \overset{\beta}{\rightarrow} C$ is a pullback of the cospan $L \overset{\beta'}{\rightarrow} G \overset{\alpha'}{\leftarrow} C$.
    \begin{figure}[H]
        \centering
        \resizebox{0.9\textwidth}{!}{
        \begin{tikzpicture} 
            \graphbox{\( L \)}{40mm}{20mm}{34mm}{20mm}{2mm}{2mm}{
                \coordinate (o) at (0mm,-8mm); 
                \node[draw,circle] (l1) at ($(o)+(-10mm,0mm)$) {1};
                \node[draw,circle] (l2) at ($(l1)+(2,0)$) {2};
                \draw[->,red] (l2) edge[out=-135,in=-45]node[midway,below] {$a$} (l1) ;
                \node[draw,circle,red] (l3) at ($(l1) + (1,0)$) {3};
                \draw[->,red] (l1) -- (l3) node[midway,above] {$a$};
                \draw[->,red] (l3) -- (l2) node[midway,above] {$a$};
            } 
    
            \graphbox{\( K \)}{0mm}{0mm}{34mm}{12mm}{2mm}{2mm}{
                \coordinate (o) at (0mm,-8mm); 
                \node[draw,circle] (l1) at ($(o)+(-10mm,0mm)$) {1};
                \node[draw,circle] (l2) at ($(l1)+(2,0)$) {2};
            }  
            \graphbox{\(G  \)}{110mm}{5mm}{34mm}{40mm}{2mm}{-3mm}{
                \coordinate (o) at (0mm,-20mm); 
                 \node[draw,circle] (l1) at ($(o)+(0,0)$) {1\ 2};
                % \node[draw,circle] (l1) at ($(o)+(-10mm,0mm)$) {1};
                % \node[draw,circle] (l2) at ($(l1)+(2,0)$) {2};
                \draw[->,red] (l1) edge[loop below] node[midway, below] {$a$} (l1) ;
                \node[draw,circle,red] (l3) at ($(l1) + (0,1.4)$) {3};
                \node[draw,circle,blue] (l4) at ($(l1) + (1,-1)$) {6};
                \draw[->,red] (l1) edge[bend left] node[midway,left] {$a$} (l3);
                \draw[->,red] (l3) edge[bend left] node[midway,right] {$a$} (l1);
                \draw[->,blue] (l1) edge  node[midway,right] {$a$} (l4);
                \node[draw,circle,blue] (l6) at ($(l1) + (-1,-1)$) {7};
                \draw[<-,blue] (l1) edge node[midway,left] {$a$} (l6) ;
                % \draw[->,blue] (l1) edge[out=-135,in=-45]node[midway,below] {$a$} (l1) ;
            }   
     
            \graphbox{\( C \)}{40mm}{-20mm}{34mm}{35mm}{2mm}{-3mm}{
                \coordinate (o) at (0mm,-14mm); 
                \node[draw,circle] (l1) at ($(o)+(0,0)$) {1\ 2};
                % \node[draw,circle] (l2) at ($(l1)+(2,0)$) {2};
                \node[draw,circle,blue] (l4) at ($(l1) + (1,-1)$) {6};
                \draw[->,blue] (l1) -- (l4) node[midway,right] {$a$};
                % \draw[->,blue] (l1) edge[loop above] node[midway, above] {$a$} (l1) ;
                \node[ draw,circle,blue] (l6) at ($(l1) + (-1,-1)$) {7};
                \draw[<-,blue] (l1) -- (l6) node[midway,left] {$a$};
            }
            % K to L  
            \draw[>->] (17mm,5mm) -- node[above] {$\alpha$} (37mm,15mm);
            % C to G
            \draw[>->] (76mm,-28mm)-- node[below] {$\alpha'$} (104mm,-24mm) ;
            % K to C
            \draw[->] (17mm,-17mm) -- node[below] {$\beta$} (37mm,-28mm);
            % L to G 
            \draw[->] (76mm,16mm) -- node[above] {$\beta'$} (104mm,7mm);
            \node () at (57mm,-6mm) {$PB$};
        \end{tikzpicture}
        }
        \caption{Pullback.}
        \label{fig:preliminaries:pullback}
    \end{figure} 

 
    % Specifically, for a possible span \( B \overset{\gamma'}{\leftarrow} E \overset{\gamma}{\rightarrow} C \) to satisfy \(\gamma' \star \beta' = \gamma \star \alpha'\), graph $E$ must not contain any edges as $\alpha'$ and $\beta'$ do not map any edges from $C$ and $L$ to a common edge in $G$.
    %  For all graph $E$ without edges and all morphisms $\gamma': E \rightarrow L $ and $\gamma : E \rightarrow C$, we have \(\gamma' \star \beta' = \gamma \star \alpha'\). Three examples of such spans are illustrated in~\autoref{fig:preliminaries:three_squares}.
    % \begin{figure}[H]
    %     \centering
    %         \resizebox{0.7\textwidth}{!}{
    %         \begin{tikzpicture} 
    %             \graphbox{\( E_1 \)}{-50mm}{0mm}{34mm}{12mm}{2mm}{2mm}{
    %                 \coordinate (o) at (0mm,-8mm); 
    %                 \node[draw,circle] (l1) at ($(o)+(-10mm,0mm)$) {1};
    %                 \node[draw,circle] (l2) at ($(l1)+(2,0)$) {2};
    %             }  
    %             \graphbox{\( L \)}{40mm}{20mm}{34mm}{20mm}{2mm}{2mm}{
    %                 \coordinate (o) at (0mm,-8mm); 
    %                 \node[draw,circle] (l1) at ($(o)+(-10mm,0mm)$) {1};
    %                 \node[draw,circle] (l2) at ($(l1)+(2,0)$) {2};
    %                 \draw[->,red] (l2) edge[out=-135,in=-45]node[midway,below] {$a$} (l1) ;
    %                 \node[draw,circle,red] (l3) at ($(l1) + (1,0)$) {3};
    %                 \draw[->,red] (l1) -- (l3) node[midway,above] {$a$};
    %                 \draw[->,red] (l3) -- (l2) node[midway,above] {$a$};
    %             } 
        
    %             \graphbox{\( K \)}{0mm}{0mm}{34mm}{12mm}{2mm}{2mm}{
    %                 \coordinate (o) at (0mm,-8mm); 
    %                 \node[draw,circle] (l1) at ($(o)+(-10mm,0mm)$) {1};
    %                 \node[draw,circle] (l2) at ($(l1)+(2,0)$) {2};
    %             }  

    %             \graphbox{\( C \)}{40mm}{-20mm}{34mm}{35mm}{2mm}{-3mm}{
    %                 \coordinate (o) at (0mm,-14mm); 
    %                 \node[draw,circle] (l1) at ($(o)+(0,0)$) {1\ 2};
    %                 % \node[draw,circle] (l2) at ($(l1)+(2,0)$) {2};
    %                 \node[draw,circle,blue] (l4) at ($(l1) + (1,-1)$) {6};
    %                 \draw[->,blue] (l1) -- (l4) node[midway,right] {$a$};
    %                 % \draw[->,blue] (l1) edge[loop above] node[midway, above] {$a$} (l1) ;
    %                 \node[ draw,circle,blue] (l6) at ($(l1) + (-1,-1)$) {7};
    %                 \draw[<-,blue] (l1) -- (l6) node[midway,left] {$a$};
    %             }

    %             \graphbox{\(G  \)}{110mm}{5mm}{34mm}{40mm}{2mm}{-3mm}{
    %                 \coordinate (o) at (0mm,-20mm); 
    %                 \node[draw,circle] (l1) at ($(o)+(0,0)$) {1\ 2};
    %                 % \node[draw,circle] (l1) at ($(o)+(-10mm,0mm)$) {1};
    %                 % \node[draw,circle] (l2) at ($(l1)+(2,0)$) {2};
    %                 \draw[->,red] (l1) edge[loop below] node[midway, below] {$a$} (l1) ;
    %                 \node[draw,circle,red] (l3) at ($(l1) + (0,1.4)$) {3};
    %                 \node[draw,circle,blue] (l4) at ($(l1) + (1,-1)$) {6};
    %                 \draw[->,red] (l1) edge[bend left] node[midway,left] {$a$} (l3);
    %                 \draw[->,red] (l3) edge[bend left] node[midway,right] {$a$} (l1);
    %                 \draw[->,blue] (l1) edge  node[midway,right] {$a$} (l4);
    %                 \node[draw,circle,blue] (l6) at ($(l1) + (-1,-1)$) {7};
    %                 \draw[<-,blue] (l1) edge node[midway,left] {$a$} (l6) ;
    %                 % \draw[->,blue] (l1) edge[out=-135,in=-45]node[midway,below] {$a$} (l1) ;
    %             }   
    %             % E to K
    %             \draw[>->] (-15mm,-6mm) -- node[above] {$\delta_1$} (-1mm,-6mm);
    %             % E to L
    %             \draw[>->] (-33mm,5mm) -- node[above] {$\gamma'_1$} (33mm,15mm);
    %             % K to L  
    %             \draw[>->] (17mm,5mm) -- node[above] {$\alpha$} (37mm,15mm);
    %             % C to G
    %             \draw[>->] (76mm,-28mm)-- node[below] {$\alpha'$} (104mm,-24mm) ;
    %             % K to C
    %             \draw[->] (17mm,-17mm) -- node[below] {$\beta$} (37mm,-28mm);
    %             % E to C
    %             \draw[->] (-33mm,-17mm) -- node[below] {$\gamma_1$} (33mm,-28mm);
    %             % L to G
    %             \draw[->] (76mm,16mm) -- node[above] {$\beta'$} (104mm,7mm);
    %             % \node () at (57mm,-6mm) {$PB$};
    %         \end{tikzpicture}
    %         }

    %         \vspace{5mm}
    %         \resizebox{0.7\textwidth}{!}{
    %         \begin{tikzpicture} 
    %             \graphbox{\( E_2 \)}{-50mm}{0mm}{34mm}{12mm}{2mm}{2mm}{
    %                 \coordinate (o) at (0mm,-8mm); 
    %                 \node[draw,circle] (l1) at ($(o)+(-10mm,0mm)$) {1};
    %                 % \node[draw,circle] (l2) at ($(l1)+(2,0)$) {2};
    %             }  
    %             \graphbox{\( L \)}{40mm}{20mm}{34mm}{20mm}{2mm}{2mm}{
    %                 \coordinate (o) at (0mm,-8mm); 
    %                 \node[draw,circle] (l1) at ($(o)+(-10mm,0mm)$) {1};
    %                 \node[draw,circle] (l2) at ($(l1)+(2,0)$) {2};
    %                 \draw[->,red] (l2) edge[out=-135,in=-45]node[midway,below] {$a$} (l1) ;
    %                 \node[draw,circle,red] (l3) at ($(l1) + (1,0)$) {3};
    %                 \draw[->,red] (l1) -- (l3) node[midway,above] {$a$};
    %                 \draw[->,red] (l3) -- (l2) node[midway,above] {$a$};
    %             } 
        
    %             \graphbox{\( K \)}{0mm}{0mm}{34mm}{12mm}{2mm}{2mm}{
    %                 \coordinate (o) at (0mm,-8mm); 
    %                 \node[draw,circle] (l1) at ($(o)+(-10mm,0mm)$) {1};
    %                 \node[draw,circle] (l2) at ($(l1)+(2,0)$) {2};
    %             }  

    %             \graphbox{\( C \)}{40mm}{-20mm}{34mm}{35mm}{2mm}{-3mm}{
    %                 \coordinate (o) at (0mm,-14mm); 
    %                 \node[draw,circle] (l1) at ($(o)+(0,0)$) {1\ 2};
    %                 % \node[draw,circle] (l2) at ($(l1)+(2,0)$) {2};
    %                 \node[draw,circle,blue] (l4) at ($(l1) + (1,-1)$) {6};
    %                 \draw[->,blue] (l1) -- (l4) node[midway,right] {$a$};
    %                 % \draw[->,blue] (l1) edge[loop above] node[midway, above] {$a$} (l1) ;
    %                 \node[ draw,circle,blue] (l6) at ($(l1) + (-1,-1)$) {7};
    %                 \draw[<-,blue] (l1) -- (l6) node[midway,left] {$a$};
    %             }

    %             \graphbox{\(G  \)}{110mm}{5mm}{34mm}{40mm}{2mm}{-3mm}{
    %                 \coordinate (o) at (0mm,-20mm); 
    %                 \node[draw,circle] (l1) at ($(o)+(0,0)$) {1\ 2};
    %                 % \node[draw,circle] (l1) at ($(o)+(-10mm,0mm)$) {1};
    %                 % \node[draw,circle] (l2) at ($(l1)+(2,0)$) {2};
    %                 \draw[->,red] (l1) edge[loop below] node[midway, below] {$a$} (l1) ;
    %                 \node[draw,circle,red] (l3) at ($(l1) + (0,1.4)$) {3};
    %                 \node[draw,circle,blue] (l4) at ($(l1) + (1,-1)$) {6};
    %                 \draw[->,red] (l1) edge[bend left] node[midway,left] {$a$} (l3);
    %                 \draw[->,red] (l3) edge[bend left] node[midway,right] {$a$} (l1);
    %                 \draw[->,blue] (l1) edge  node[midway,right] {$a$} (l4);
    %                 \node[draw,circle,blue] (l6) at ($(l1) + (-1,-1)$) {7};
    %                 \draw[<-,blue] (l1) edge node[midway,left] {$a$} (l6) ;
    %                 % \draw[->,blue] (l1) edge[out=-135,in=-45]node[midway,below] {$a$} (l1) ;
    %             }   
    %             % E to K
    %             \draw[>->] (-15mm,-6mm) -- node[above] {$\delta_2$} (-1mm,-6mm);
    %             % E to L
    %             \draw[>->] (-33mm,5mm) -- node[above] {$\gamma'_2$} (33mm,15mm);
    %             % K to L  
    %             \draw[>->] (17mm,5mm) -- node[above] {$\alpha$} (37mm,15mm);
    %             % C to G
    %             \draw[>->] (76mm,-28mm)-- node[below] {$\alpha'$} (104mm,-24mm) ;
    %             % K to C
    %             \draw[->] (17mm,-17mm) -- node[below] {$\beta$} (37mm,-28mm);
    %             % E to C
    %             \draw[->] (-33mm,-17mm) -- node[below] {$\gamma_2$} (33mm,-28mm);
    %             % L to G
    %             \draw[->] (76mm,16mm) -- node[above] {$\beta'$} (104mm,7mm);
    %             % \node () at (57mm,-6mm) {$PB$};
    %         \end{tikzpicture}
    %         }

    %         \vspace{5mm}
    %         \resizebox{0.7\textwidth}{!}{
    %         \begin{tikzpicture} 
    %             \graphbox{\( E_3 \)}{-50mm}{0mm}{34mm}{12mm}{2mm}{2mm}{
    %                 \coordinate (o) at (0mm,-8mm); 
    %                 % \node[draw,circle] (l1) at ($(o)+(-10mm,0mm)$) {1};
    %                 \node[draw,circle] (l2) at ($(o)+(10mm,0mm)$) {2};
    %             }  
    %             \graphbox{\( L \)}{40mm}{20mm}{34mm}{20mm}{2mm}{2mm}{
    %                 \coordinate (o) at (0mm,-8mm); 
    %                 \node[draw,circle] (l1) at ($(o)+(-10mm,0mm)$) {1};
    %                 \node[draw,circle] (l2) at ($(l1)+(2,0)$) {2};
    %                 \draw[->,red] (l2) edge[out=-135,in=-45]node[midway,below] {$a$} (l1) ;
    %                 \node[draw,circle,red] (l3) at ($(l1) + (1,0)$) {3};
    %                 \draw[->,red] (l1) -- (l3) node[midway,above] {$a$};
    %                 \draw[->,red] (l3) -- (l2) node[midway,above] {$a$};
    %             } 
        
    %             \graphbox{\( K \)}{0mm}{0mm}{34mm}{12mm}{2mm}{2mm}{
    %                 \coordinate (o) at (0mm,-8mm); 
    %                 \node[draw,circle] (l1) at ($(o)+(-10mm,0mm)$) {1};
    %                 \node[draw,circle] (l2) at ($(l1)+(2,0)$) {2};
    %             }  

    %             \graphbox{\( C \)}{40mm}{-20mm}{34mm}{35mm}{2mm}{-3mm}{
    %                 \coordinate (o) at (0mm,-14mm); 
    %                 \node[draw,circle] (l1) at ($(o)+(0,0)$) {1\ 2};
    %                 % \node[draw,circle] (l2) at ($(l1)+(2,0)$) {2};
    %                 \node[draw,circle,blue] (l4) at ($(l1) + (1,-1)$) {6};
    %                 \draw[->,blue] (l1) -- (l4) node[midway,right] {$a$};
    %                 % \draw[->,blue] (l1) edge[loop above] node[midway, above] {$a$} (l1) ;
    %                 \node[ draw,circle,blue] (l6) at ($(l1) + (-1,-1)$) {7};
    %                 \draw[<-,blue] (l1) -- (l6) node[midway,left] {$a$};
    %             }

    %             \graphbox{\(G  \)}{110mm}{5mm}{34mm}{40mm}{2mm}{-3mm}{
    %                 \coordinate (o) at (0mm,-20mm); 
    %                 \node[draw,circle] (l1) at ($(o)+(0,0)$) {1\ 2};
    %                 % \node[draw,circle] (l1) at ($(o)+(-10mm,0mm)$) {1};
    %                 % \node[draw,circle] (l2) at ($(l1)+(2,0)$) {2};
    %                 \draw[->,red] (l1) edge[loop below] node[midway, below] {$a$} (l1) ;
    %                 \node[draw,circle,red] (l3) at ($(l1) + (0,1.4)$) {3};
    %                 \node[draw,circle,blue] (l4) at ($(l1) + (1,-1)$) {6};
    %                 \draw[->,red] (l1) edge[bend left] node[midway,left] {$a$} (l3);
    %                 \draw[->,red] (l3) edge[bend left] node[midway,right] {$a$} (l1);
    %                 \draw[->,blue] (l1) edge  node[midway,right] {$a$} (l4);
    %                 \node[draw,circle,blue] (l6) at ($(l1) + (-1,-1)$) {7};
    %                 \draw[<-,blue] (l1) edge node[midway,left] {$a$} (l6) ;
    %                 % \draw[->,blue] (l1) edge[out=-135,in=-45]node[midway,below] {$a$} (l1) ;
    %             }   
    %             % E to K
    %             \draw[>->] (-15mm,-6mm) -- node[above] {$\delta_3$} (-1mm,-6mm);
    %             % E to L
    %             \draw[>->] (-33mm,5mm) -- node[above] {$\gamma'_3$} (33mm,15mm);
    %             % K to L  
    %             \draw[>->] (17mm,5mm) -- node[above] {$\alpha$} (37mm,15mm);
    %             % C to G
    %             \draw[>->] (76mm,-28mm)-- node[below] {$\alpha'$} (104mm,-24mm) ;
    %             % K to C
    %             \draw[->] (17mm,-17mm) -- node[below] {$\beta$} (37mm,-28mm);
    %             % E to C
    %             \draw[->] (-33mm,-17mm) -- node[below] {$\gamma_3$} (33mm,-28mm);
    %             % L to G
    %             \draw[->] (76mm,16mm) -- node[above] {$\beta'$} (104mm,7mm);
    %             % \node () at (57mm,-6mm) {$=$};
    %         \end{tikzpicture}
    %         }
    %     \caption{}
    %     \label{fig:preliminaries:three_squares} 
    % \end{figure}
\end{example} 
An important phenomenon about pushouts and pullbacks in the category \textbf{Graph} is stated in the following proposition. For examples, consider the diagrams in~\autoref{fig:preliminaries:pushout_non_injective} and~\autoref{fig:preliminaries:a_rewriting_step_dfjalsdkdfsdfjflg}
\begin{proposition}[\text{\cite[Lemma 13]{lack2004adhesive}}]
    \label{prop:pb_eq_po}
    In the category \textbf{Graph}, pushouts along monomorphisms are also pullbacks. 
\end{proposition}

\subsection{DPO rewriting} 
\begin{definition}[Rewriting rule and match~\cite{corradini1997algebraic}]
  \label{def:grs:dpo_rule}
A \textbf{DPO rewriting rule} $\rho$ is a span \( L \overset{l}{\leftarrow} K \overset{r}{\rightarrow} R \), where \( K \) is the \textbf{interface}, \( L \) is the \textbf{left-hand-side graph}, denoted \( \operatorname{lhs}(\rho) \), and \( R \) is the \textbf{right-hand-side graph}, denoted \( \operatorname{rhs}(\rho) \). The rule is \textbf{monic} if $l$ and $r$ are both monic.
A match of the rule in an graph \( G \) is a morphism \( m: L \rightarrow G \).   
\end{definition}
   In this paper, we use examples from the category \textbf{Graph} of edge-labeled directed multigraphs (see~\cite{konig2018atutorial}) to illustrate the discussed concepts. 
\begin{example}
  \label{ex:grsaa}
  The injective DPO rule from \cite[Example 6]{bruggink2014termination} will be used to illustrate the concepts discussed throughout this paper.
  The rule can be visualized as follows:
  \begin{center} 
      \resizebox{0.7\textwidth}{!}{
      \begin{tikzpicture}
          \graphbox{$L$}{0mm}{0mm}{34mm}{15mm}{2mm}{-5mm}{
              \coordinate (o) at (0mm,-3mm); 
              \node[draw,circle] (l1) at ($(o)+(-10mm,0mm)$) {1};
              \node[draw,circle] (l2) at ($(l1)+(2,0)$) {2};
              \node[draw,circle] (l3) at ($(l1) + (1,0)$) {3};
              \draw[->] (l1) -- (l3) node[midway,above] {a};
              \draw[->] (l3) -- (l2) node[midway,above] {a};
          }     
          \graphbox{$K$}{40mm}{0mm}{24mm}{15mm}{2mm}{-5mm}{
              \coordinate (o) at (5mm,-3mm); 
              \node[draw,circle] (l1) at ($(o)+(-10mm,0mm)$) {1};
              \node[draw,circle] (l2) at ($(l1)+(1,0)$) {2};
              % \node[draw,circle] (l3) at ($(l1) + (1,0)$) {$\ $};
              % \draw[->] (l1) -- (l3) node[midway,above] {a};
              % \draw[->] (l3) -- (l2) node[midway,above] {a};
          }    
          \graphbox{$R$}{70mm}{0mm}{45mm}{15mm}{2mm}{-5mm}{
              \coordinate (o) at (-5mm,-3mm); 
              \node[draw,circle] (l1) at ($(o)+(-10mm,0mm)$) {1};
              \node[draw,circle] (l2) at ($(l1)+(3,0)$) {2};
              \node[draw,circle] (l3) at ($(l1) + (1,0)$) {4};
              \node[draw,circle] (l4) at ($(l1) + (2,0)$) {5};
              \draw[->] (l1) -- (l3) node[midway,above] {a};
              \draw[->] (l3) -- (l4) node[midway,above] {b};
              \draw[->] (l4) -- (l2) node[midway,above] {a};
          }    
          \node () at (37mm,-8mm) {$\overset{l}{\leftarrowtail}$};
          \node () at (67mm,-8mm) {$\overset{r}{\rightarrowtail}$};
          % \draw[>->] (51mm,2mm) -- (52mm,3mm);
      \end{tikzpicture}
      }
  \end{center}
\end{example}
\begin{definition}[DPO Rewriting step \cite{endrullis2024generalized_arxiv_v2}]
  \label{def:rewriting_step}
    \ \newline
    \noindent
    \begin{minipage}{0.72\textwidth}
      A DPO diagram $\delta$ is a diagram as shown on the right.
      This diagram $\delta$ is a witness for the \textbf{rewriting step} from \( G \) to \( H \) using the rule \( \rho \) and \textbf{match} \( m \), denoted \( G \Rightarrow_\rho^m H \) or \( G \Rightarrow_\rho^\delta H \). We denote $\operatorname{left}(\delta)$ and $\operatorname{right}(\delta)$ the pushout squares $KLGC$ and $KRHC$, respectively.
    \end{minipage}
    \hfill
    \begin{minipage}{0.28\textwidth}
          % \begin{center}
          \hfill
          \resizebox{0.85\textwidth}{!}{
          \begin{tikzpicture}
            % [node distance=11mm]
            \node (I) at (0,0) {$K$};
            \node (L) at (-2,0) {$L$};
            \node (R) at (2,0) {$R$};
            \node (G) at (-2,-2) {$G$};
            \node (C) at (0,-2) {$C$};
            \node (H) at (2,-2) {$H$};
            \draw [->] (I) to  node [midway,below] {$l$} (L);
            \draw [->] (I) to  node [midway,below] {$r$} (R);
            \draw [->] (L) to node [midway,right] {$m$} (G);
            \draw [->] (I) to node [midway,right] {$u$} (C);
            \draw [->] (R) to node [midway,left] {$m'$} (H);
            \draw [->] (C) to node [midway,above] {$l'$} (G);
            \draw [->] (C) to node [midway,above] {$r'$} (H);
            \node [at=($(I)!.5!(G)$)] {\normalfont PO};
            \node [at=($(I)!.5!(H)$)] {\normalfont PO};
          \end{tikzpicture}
        % \end{center}
        }
        \end{minipage}
  \end{definition}
\begin{example}
  \label{ex:rewriting_step_grs_aa}
  The DPO diagram below defines a rewriting step using the rule from Example~\ref{ex:grsaa}.
  \begin{center} 
      \resizebox{0.7\textwidth}{!}{
      \begin{tikzpicture}
          \graphbox{\( L \)}{0mm}{-3mm}{34mm}{12mm}{2mm}{2mm}{
              \coordinate (o) at (0mm,-8mm); 
              \node[draw,circle] (l1) at ($(o)+(-10mm,0mm)$) {1};
              \node[draw,circle] (l2) at ($(l1)+(2,0)$) {2};
              \node[draw,circle] (l3) at ($(l1) + (1,0)$) {3};
              \draw[] (l1) -- (l3) node[midway,above] {$a$};
              \draw[] (l3) -- (l2) node[midway,above] {$a$};
          } 
          \graphbox{\( K \)}{40mm}{-3mm}{34mm}{12mm}{2mm}{2mm}{
              \coordinate (o) at (0mm,-8mm); 
              \node[draw,circle] (l1) at ($(o)+(-10mm,0mm)$) {1};
              \node[draw,circle] (l2) at ($(l1)+(2,0)$) {2};
          }  
          \graphbox{\( R \)}{80mm}{-3mm}{45mm}{12mm}{2mm}{2mm}{
              \coordinate (o) at (-5mm,-8mm); 
              \node[draw,circle] (l1) at ($(o)+(-10mm,0mm)$) {1};
              \node[draw,circle] (l2) at ($(l1)+(3,0)$) {2};
              \node[draw,circle] (l3) at ($(l1) + (1,0)$) {4};
              \node[draw,circle] (l4) at ($(l1) + (2,0)$) {5};
              \draw[ ] (l1) -- (l3) node[midway,above] {$a$};
              \draw[ ] (l3) -- (l4) node[midway,above] {$b$};
              \draw[ ] (l4) -- (l2) node[midway,above] {$a$};
          }    
          \graphbox{\( G \)}{0mm}{-22mm}{34mm}{22mm}{2mm}{-3mm}{
              \coordinate (o) at (0mm,-3mm); 
              \node[draw,circle] (l1) at ($(o)+(-10mm,0mm)$) {1};
              \node[draw,circle] (l2) at ($(l1)+(2,0)$) {2};
              \node[draw,circle] (l3) at ($(l1) + (1,0)$) {3};
              \node[draw,circle] (l4) at ($(l2) + (0,-1)$) {6};
              \draw[] (l1) -- (l3) node[midway,above] {$a$};
              \draw[] (l3) -- (l2) node[midway,above] {$a$};
              \draw[ ] (l2) -- (l4) node[midway,right] {$a$};
              \node[draw,circle] (l6) at ($(l1) + (0,-1)$) {7};
              \draw[] (l1) -- (l6) node[midway,left] {$a$};
          }    
          \graphbox{\( C  \)}{40mm}{-22mm}{34mm}{22mm}{2mm}{-3mm}{
              \coordinate (o) at (0mm,-3mm); 
              \node[draw,circle] (l1) at ($(o)+(-10mm,0mm)$) {1};
              \node[draw,circle] (l2) at ($(l1)+(2,0)$) {2};
              \node[draw,circle] (l4) at ($(l2) + (0,-1)$) {6};
              \draw[ ] (l2) -- (l4) node[midway,right] {$a$};
              \node[ draw,circle] (l6) at ($(l1) + (0,-1)$) {7};
              \draw[ ] (l1) -- (l6) node[midway,left] {$a$};
          }    
          \graphbox{\( H \)}{80mm}{-22mm}{45mm}{22mm}{2mm}{-3mm}{
              \coordinate (o) at (-5mm,-3mm); 
              \node[draw,circle] (l1) at ($(o)+(-10mm,0mm)$) {1};
              \node[draw,circle] (l2) at ($(l1)+(3,0)$) {2};
              \node[draw,circle] (l3) at ($(l1) + (1,0)$) {4};
              \node[draw,circle] (l4) at ($(l1) + (2,0)$) {5};
              \node[ draw,circle] (l5) at ($(l2) + (0,-1)$) {6};
              \node[ draw,circle] (l6) at ($(l1) + (0,-1)$) {7};
              \draw[ ] (l1) -- (l6) node[midway,left] {$a$};
              \draw[] (l1) -- (l3) node[midway,above] {$a$};
              \draw[] (l3) -- (l4) node[midway,above] {$b$};
              \draw[ ] (l4) -- (l2) node[midway,above] {$a$};
              \draw[ ] (l2) -- (l5) node[midway,right] {$a$};
          }    
          \node () at (37mm,-8mm) {\( \leftarrowtail \)}; % K -> L
          \node () at (77mm,-8mm) {\( \rightarrowtail \)}; % K -> R
          \node () at (15mm,-18mm) {\( m\ \downarrowtail \)};
          \node () at (37mm,-33mm) {\( \leftarrowtail \)};
          \node () at (58mm,-18mm) {\( u\downarrowtail \)};
          \node () at (102mm,-18mm) {\( \downarrowtail \)};
          \node () at (77mm,-33mm) {\( \rightarrowtail \)}; % C -> H
      \end{tikzpicture}
      }
  \end{center}
\end{example} 
\subsection{Relative termination}
\begin{definition}[Rewriting sequence]
    Let \(\mathcal{R}\) be a set of DPO rewriting rules and $\mathfrak{F}$ a DPO rewriting framework in \textbf{Graph}.
    A \textbf{$(\mathcal{R},\mathfrak{F})$-rewriting sequence} is a finite sequence \(s_0,s_1,\hdots, s_m\) of objects such that for each \( 0 \leq n \leq m-1\), we have \(s_n \Rightarrow_{\mathcal{R},\mathfrak{F}} s_{n+1}\), or an infinite sequence \(s_0,s_1,\hdots\) of objects such that \(s_n \Rightarrow_{\mathcal{R},\mathfrak{F}} s_{n+1}\) for each \(n \in \mathbb{N}\).
\end{definition}
We adapt the concept of relative termination presented in \cite{klop1987term,geser1990relative}.
\begin{definition}[Relative termination]
    \label{termination:def:relative_termination}
     Let $\mathcal{A}$ and $\mathcal{B}$ be sets of rewriting rules, and $\mathfrak{F}$ a DPO rewriting framework in \textbf{Graph}. We say that $\Rightarrow_{\mathcal{A},\mathfrak{F}}$ is \textbf{terminating relative to} $\Rightarrow_{\mathcal{B}, \mathfrak{F}}$ if any infinite $(\mathcal{A} \cup \mathcal{B},\mathfrak{F})$-rewriting sequence has only a finite number of rewriting steps with rules in $\mathcal{A}$.
\end{definition}
   
\section{A Termination Criterion} 
\label{sec:termination_criterion}
% In this section, we present a method for termination analysis of DPO graph rewriting systems with injective rules and monic matches. 
% Since it is an interpretation method \cite{nipkow1998term}, we first define the weight of a graph in \autoref{def:graph_weight}. 
% With this definition, we explain the general idea of the method and identify the key challenge in \autoref{sec:general_idea}. 
% In \autoref{sec:non-increasing}, we introduce the concept of non-increasing rules. 
% Using this concept, a solution to this challenge is then proposed in \autoref{sec:solution_to_the_key_challenge}. 
% Finally, in \autoref{sec:termination}, we state our termination criterion.

\section{Measuring graphs by counting subgraphs}
\label{sec:interpretation}
Let $\mathcal{R} = \mathcal{A} \cup \mathcal{B}$ be a set of injective rules.
Define $\mathrm{Sub}(\mathrm{left}(\mathcal{R}))$ as the set of all subgraphs of left-hand sides graphs of rules in $\mathcal{R}$.
We call elements of $\mathrm{Sub}(\mathrm{left}(\mathcal{R}))$ \textbf{ruler-graphs}. Analogous to measuring physical objects with rulers, we measure graphs by counting occurrences of these ruler-graphs.

\begin{definition}
    % Let \( X \) be a ruler-graph and \( G \) a graph. The \textbf{measurement} of \( G \) with respect to \( X \), denoted \( m_X(G) \), is 
    For a ruler-graph \( X \) and graph \( G \), the \textbf{measurement} \( m_X(G) \) is:
    \(
        m_X(G) = \card{\operatorname{Mono}(X, G)}
    \).
\end{definition} 
The occurrence count of ruler-graphs before and after a rewriting step provides a termination criterion: 
    If every rewriting step using a rule in $\mathcal{A}$ strictly decreases the number of $X$'s occurrences, and every rewriting step using a rule in $\mathcal{B}$ never increases it, then the rewriting relation induced by $\mathcal{A}$ is terminating relative to the rewriting relation induced by $\mathcal{B}$.
 
For some ruler-graphs, occurrence changes during a rewriting step using an injective rule are precisely computable by simply counting their occurrences in the left- and right-hand side graphs. Examples include: the ruler-graph \tikz[baseline=-0.5ex]{
        \node (x) at (0,0) {$\bullet$};
    } consisting of a single node; the ruler-graph \tikz[baseline=-0.5ex]{
        \node (x) at (0,0) {$\bullet$};
        \draw[->] (x) edge[loop right] node[midway, right] {$x$} (x) ;
    } consisting of a single node and a single loop labeled by $x\in \Sigma$; the ruler-graph \tikz[baseline=-0.5ex]{
        \node (x) at (0,0) {$\bullet$};
        \node (y) at (1,0) {$\bullet$ };
        \draw[->] (x) -- (y) node[midway, above] {$x$};
    } consisting of two nodes and a single labeled edge between them.  
    
However, systems like \autoref{ex:grsaa_rx} require more complex ruler-graphs such as \tikz[baseline=-0.5ex]{
        \node (x) at (0,0) {$\bullet$};
        \node (y) at (1,0) {$\bullet$ };
        \node (z) at (2,0) { $\bullet$};
        \draw[->] (x) -- (y) node[midway, above] {$a$};
        \draw[->] (y) -- (z) node[midway, above] {$a$};
    } as shown in \autoref{ex:termination:grsaa}.
    For such ruler-graph, occurrence changes during rewriting is unpredictable in general, as their occurrence creation and destruction depend on the context graph\textemdash a challenge we address in later sections.
    
To combine measurements from multiple ruler-graphs (analogous to combining measurements of physical objects from multiple rulers), we introduce:
\begin{definition}[Weight function]
    \label{def:weight_function}
    A \textbf{weight function} for a set of ruler-graphs \( \mathbb{X} \) is a map \( s_{\mathbb{X}} \colon \mathbb{X} \to \mathbb{N} \) assigning a weight in $\mathbb{N}$ to each ruler-graph.
\end{definition}
\autoref{ex:overbeek_5d6} demonstrates the usefulness of a weight function that assigns different weight to measurements from different ruler-graphs.
\begin{definition}[Graph weight]
    \label{def:graph_weight}  
    Let \( \mathbb{X} \) be a set of ruler-graphs, 
    \( s_{\mathbb{X}} \) a weight function, and \( G \) a graph. The \textbf{weight of graph} \( G \) relative to \( s_{\mathbb{X}} \), denoted \( w_{s_{\mathbb{X}}}(G) \), is defined as: 
    \[
        w_{s_{\mathbb{X}}}(G) = \sum_{X \in \mathbb{X}} s_{\mathbb{X}}(X) \cdot m_X(G).  
    \]  
\end{definition}


% \section{Measuring graphs with ruler-Graphs}
% \label{sec:interpretation}
% 
%     Let $\mathcal{R}$ be a set of injective DPO graph rewriting rules and $\mathrm{Sub}(\mathrm{left}(\mathcal{R}))$ the set of subgraphs of the left-hand-side graphs of the rules in $\mathcal{R}$. 
%     Analogous to measuring vectors using a basis within a vector space, 
%     or measuring physical objects using rulers, 
%     we measure graphs using ruler-graphs\textemdash special graphs in $\mathrm{Sub}(\mathrm{left}(\mathcal{R}))$ that act as reference structures. 
%     A ruler-graph measures a graph $G$, and provides a measurement of $G$ relative to that ruler-graph.  
%     Similarly, analogous to aggregating the coordinates of a vector to compute a weight, or combining measurements of physical objects provided by multiple rulers to obtain a weight, we summarize the measurements across multiple ruler-graphs as the weight of a graph.
% }
% \todo{A: all of this is a little bit over-formalized}
% \begin{definition}[Ruler-graph]
%     A \textbf{ruler-graph} is a graph in $\mathrm{Sub}(\mathrm{left}(\mathcal{R}))$.
% \end{definition} 
% \begin{definition}[Measurement]
%     Let \( X \) be a ruler-graph and \( G \) a graph. The \textbf{measurement} of \( G \) with respect to \( X \), denoted \( m_X(G) \), is 
%     \(
%         m_X(G) = \big|\!\operatorname{Mono}(X, G)\!\big| 
%     \).
% \end{definition}
% \begin{definition}[Weight function]
%     \label{def:weight_function}
%     A \textbf{weight function} for a set of ruler-graphs \( \mathbb{X} \) is a map \( s_{\mathbb{X}} \colon \mathbb{X} \to \mathbb{N} \) assigning a weight in $\mathbb{N}$ to each ruler-graph.
% \end{definition}
% \begin{definition}[Graph weight]
%     \label{def:graph_weight}  
%     % Let \( \mathcal{R} \) be a set of graph rewriting rules.
%     % Let $\mathbb{X}$ be a set of ruler-graphs,  
%     % \( s_\mathbb{X} \colon \mathbb{X} \rightarrow \mathbb{N} \) a weight function, and \( G \) a graph. 
%     % The \textbf{weight \( w_{s_\mathbb{X}}(G) \)} of \( G \) is the weighted sum of the measurements of \( G \) provided by rulers in \( \mathbb{X} \):
%     % \[
%     %     w_{s_\mathbb{X}}(G) = \sum_{X \in \mathbb{X}} s_\mathbb{X}(X) \cdot \big|\!\operatorname{Mono}(X, G)\!\big|.
%     % \]
%     Let \( \mathbb{X} \) be a set of ruler-graphs, 
%     \( s_{\mathbb{X}} \) a weight function, and \( G \) a graph. The \textbf{weight of graph} \( G \) relative to \( s_{\mathbb{X}} \), denoted \( w_{s_{\mathbb{X}}}(G) \) (or simply $w(G)$ when $s_{\mathbb{X}}$ is clear from context), is defined as: 
%     \[
%         w_{s_{\mathbb{X}}}(G) = \sum_{X \in \mathbb{X}} s_{\mathbb{X}}(X) \cdot m_X(G).  
%     \]  
% \end{definition}
% % \begin{remark}
% %     This work focuses on deriving sufficient conditions for termination, deferring the construction of ruler-graphs to future work.
% %     %  however, considering a set of subgraphs of the left-hand-side graphs of the rewriting rules is already fruitful, as demonstrated in \autoref{sec:related_work}.
% % \end{remark}
% % \begin{remark}
%     Trivial ruler-graphs include graphs with a single node and no edges, graphs with a single node and a single labeled loop, and graphs with two nodes and a single labeled edge between them. 
%     For any such trivial ruler-graph, the change in the number of its occurrences in a graph rewriting step using a rule in $\mathcal{R}$ can be precisely computed by simply counting its occurrences in the left- and right-hand-side graphs. }\todo{reviewer 1: The choice of the set of ruler graphs is rather opaque. Intuitively, it should be related to the left-hand sides of the rules in $A$ and $B$. A discussion on this would be valuable before introducing ruler graphs.}




\section{General idea and key challenge}
\label{sec:general_idea}
For the remainder of this section, we assume: a set \( \mathcal{R} = \mathcal{A} \cup \mathcal{B} \) of injective DPO rewriting rules, a set \( \mathbb{X} \) of ruler-graphs, and a weight function \( s_{\mathbb{X}} \).  Let \( \rho = (L \overset{l}{\leftarrowtail} K \overset{r}{\rightarrowtail} R) \) denote a rule in \( \mathcal{R} \) and \( X \in \mathbb{X} \).

The rewriting relation \( \Rightarrow_{\mathcal{A},\mathfrak{M}} \) is terminating relative to $\Rightarrow_{\mathcal{B},\mathfrak{M}}$ if the following hold: (i) for all \(G \Rightarrow_{\mathcal{A},\mathfrak{M}} H\), \( w_{s_\mathbb{X}}(G) > w_{s_\mathbb{X}}(H) \), and (ii) for all \(G \Rightarrow_{\mathcal{B},\mathfrak{M}} H\), \( w_{s_\mathbb{X}}(G) \geq w_{s_\mathbb{X}}(H) \). However, directly verifying all rewriting steps is infeasible due to their potential infiniteness. 
Thus, we aim to provide a rule-based sufficient condition for termination.
% \begin{remark}
%     To simplify our presentation, we treat the following three notions interchangeably when the monomorphism \( h: A \to B \) is understood: the domain graph \( A \), the monomorphism \( h: A \to B \), the image \( h(A) \subseteq B \) (an occurrence of \( A \) in \( B \)).
%     Consequently, in a rewriting step defined by the diagram in \autoref{def:rewriting_step}, we streamline descriptions involving subgraph images. Rather than stating that there is an occurrence of a ruler-graph \( X \) formed by the image of subgraphs \( R' \subseteq R \) and \( C' \subseteq C \) around the image of \( K' \subseteq K \), we simply say \( X \) is formed by \( R' \) and \( C' \) around \( K' \). Rather than stating that $b \notin h(A)$, we simply say that $b \notin A$.
% \end{remark}
\begin{notation}
    % The disjoint union of two sets \( S \) and \( S' \) will be denoted by \( S \uplus S' \). Let \( X, A, B, G \) be graphs, and let \( \alpha \colon A \to G \) and \( \beta \colon B \to G \) be morphisms. We define the following sets of monomorphisms of $X$ in $G$: 
\ \newline
\noindent
\begin{minipage}{0.69\textwidth}
\vspace{2mm}
\noindent  The disjoint union of two sets \( S \) and \( S' \) will be denoted by \( S \uplus S' \). Let \( X, A, B, G \) be graphs, and let \( \alpha \colon A \to G \) and \( \beta \colon B \to G \) be morphisms as illustrated on the right. We define the following sets of monomorphisms from $X$ in $G$ based on their interaction with $\alpha$ and $\beta$:    
\end{minipage}%
\begin{minipage}{0.30\textwidth}
    \hfill
        \resizebox{0.9\textwidth}{!}{
        % \begin{tikzpicture}
        %     \node (B) at (2,0) {$A$}; 
        %     \node (C) at (0,2) {$B$}; 
        %     \node (D) at (0,0) {$G$}; 
        %     \draw [>->] (B) to node [below,label,pos=0.45] {$\alpha$} (D); 
        %     \draw [>->] (C) to node [left,label,pos=0.45] {$\beta$} (D);
        % \end{tikzpicture}
            \begin{tikzpicture}
            \node (X) at (-1.5,-1) {$X$};
            \node (B) at (2,0) {$A$}; 
            \node (C) at (0,2) {$B$}; 
            \node (D) at (0,0) {$G$}; 
            \draw [>->] (B) to node [above,label,pos=0.45] {$\alpha$} (D); 
            \draw [>->] (C) to node [right,label,pos=0.45] {$\beta$} (D);
            \draw [>->,dashed] (X) to node [above,label,pos=0.45] {$\iota$} (D);
            \draw [>->,dashed] (X) to node [below,label,pos=0.45] {$\zeta$} (B);
            \draw [>->,dashed] (X) to node [above,label,pos=0.45] {$\eta$} (C);
        \end{tikzpicture}
        }
\end{minipage} 
    \begin{align*}
        \operatorname{Mono}(X,G,\alpha) &= \left\{ \iota \colon X \rightarrowtail G \;\middle|\; \exists \zeta \colon X \rightarrowtail A.\, \iota = \zeta \star \alpha \right\}
        \\
        \operatorname{Mono}(X,G,\lnot \alpha) &= \left\{ \iota \colon X \rightarrowtail G \;\middle|\; \nexists \zeta \colon X \rightarrowtail A.\, \iota = \zeta \star \alpha \right\}
        \\
        \operatorname{Mono}(X,G,\lnot \alpha, \beta) &= \left\{ 
            \iota \colon X \rightarrowtail G \;\middle|\; 
                \begin{aligned}  
                    &(\nexists \zeta \colon X \rightarrowtail A.\, \iota = \zeta \star \alpha) \\ 
                    &\land (\exists \eta \colon X \rightarrowtail B.\, \iota = \eta \star \beta)
                \end{aligned}
        \right\}
        \\
        \operatorname{Mono}(X,G,\lnot \alpha, \lnot \beta) &= \left\{ 
            \iota \colon X \rightarrowtail G \;\middle|\; 
                \begin{aligned}
                    &(\nexists \zeta \colon X \rightarrowtail A.\, \iota = \zeta \star \alpha) \\
                    &\land (\nexists \eta \colon X \rightarrowtail B.\, \iota = \eta \star \beta)
                \end{aligned}
        \right\}
    \end{align*}
\end{notation}
\noindent
\begin{minipage}{0.7\textwidth}\setlength{\parindent}{1em}
\noindent Consider a rewriting step \( G \Rightarrow_{\rho,\mathfrak{M}} H \) defined by the DPO diagram illustrated on the right. The sets of monomorphisms \( \operatorname{Mono}(X,G) \) and \( \operatorname{Mono}(X,H) \) can be decomposed into the following disjoint unions:
\end{minipage}%
\begin{minipage}{0.29\textwidth}
    \hfill
    % \begin{tikzpicture}
    %     % [node distance=15mm]
    %     \node (I) {$K$};
    %     \node (L) [left of=I] {$L$};
    %     \node (R) [right of=I] {$R$};
    %     \node (G) [below of=L] {$G$};
    %     \node (C) [below of=I] {$C$};
    %     \node (H) [below of=R] {$H$};
    %     \draw [>->] (I) to  node [midway,below] {$l$} (L);
    %     \draw [>->] (I) to  node [midway,below] {$r$} (R);
    %     \draw [>->] (L) to node [midway,right] {$m$} (G);
    %     \draw [>->] (I) to  node [midway,right] {} (C);
    %     \draw [>->] (R) to  node [midway,right] {$m'$} (H);
    %     \draw [>->] (C) to node [midway,above] {$l'$} (G);
    %     \draw [>->] (C) to node [midway,above] {$r'$} (H);
    %     % \node [at=($(I)!.5!(G)$)] {\normalfont PO};
    %     % \node [at=($(I)!.5!(H)$)] {\normalfont PO};
    %   \end{tikzpicture}
    \resizebox{0.9\textwidth}{!}{
     \begin{tikzpicture}
            % [node distance=15mm]
            \node (I) at (0,0) {$K$};
            \node (L)  at (-2,0) {$L$};
            \node (R)  at (2,0) {$R$};
            \node (G)  at (-2,-2) {$G$};
            \node (C)  at (0,-2) {$C$};
            \node (H)  at (2,-2) {$H$};
            \draw [>->] (I) to  node [midway,below] {$l$} (L);
            \draw [>->] (I) to  node [midway,below] {$r$} (R);
            \draw [>->] (L) to node [midway,right] {$m$} (G);
            \draw [>->] (I) to  node [midway,right] 
            % {$u$}
            {} (C);
            \draw [>->] (R) to  node [midway,right] 
            {$m'$}
            (H);
            \draw [>->] (C) to node [midway,above] {$l'$} (G);
            \draw [>->] (C) to node [midway,above] 
            {$r'$} 
            (H);
            % \node [at=($(I)!.5!(G)$)] {\normalfont PO};
            % \node [at=($(I)!.5!(H)$)] {\normalfont PO};
        \end{tikzpicture}
    }
\end{minipage}


\begin{flalign*}
    \operatorname{Mono}(X,G) &= 
    \operatorname{Mono}(X,G,l')
    \uplus
    \operatorname{Mono}(X,G,\lnot l',m) 
    \uplus
    \operatorname{Mono}(X,G,\lnot l',\lnot m)
    \\
    \operatorname{Mono}(X,H) &= 
    \operatorname{Mono}(X,H,r')
    \uplus
    \operatorname{Mono}(X,H,\lnot r',m') 
    \uplus
    \operatorname{Mono}(X,H,\lnot r',\lnot m')
\end{flalign*}
Thus, we have:
\begin{flalign*}
    \card{\operatorname{Mono}(X,G)} ={} &
    \card{\operatorname{Mono}(X,G,l')}
    +
    \card{\operatorname{Mono}(X,G,\lnot l',m)} \\
    &+
    \card{\operatorname{Mono}(X,G,\lnot l',\lnot m)}
    \\
    \card{\operatorname{Mono}(X,H)} ={} &
    \card{\operatorname{Mono}(X,H,r')} 
    +
    \card{\operatorname{Mono}(X,H,\lnot r',m')} \\
    &+\card{\operatorname{Mono}(X,H,\lnot r',\lnot m')}
\end{flalign*}
\noindent Using the following \autoref{lem:decomp_w_u} 
(see \iflongversion
Appendix \autoref{proof:dcomp_w_u} 
\else
\cite[Lemma 27]{qiu2025termination}
\fi for a proof), we derive:
\begin{flalign*}
&\card{\operatorname{Mono}(X,G)} =
\card{\operatorname{Mono}(X,C)}
+
\card{\operatorname{Mono}(X,L, \lnot l)}
+
\card{\operatorname{Mono}(X,G, \lnot l',\lnot m)} 
 \\
 &\card{\operatorname{Mono}(X,H)} =
 \card{\operatorname{Mono}(X,C)}
 +
 \card{ \operatorname{Mono}(X,R, \lnot r) }
 +
 \card{ \operatorname{Mono}(X,H, \lnot r',\lnot m')}  
\end{flalign*}
\begin{lemma}
    \label{lem:decomp_w_u}
    \ \newline
    \noindent
    \begin{minipage}{0.69\textwidth}
        Let $X$ be a ruler-graph. For a pushout square as shown on the right, we have 
        % $\card{\operatorname{Mono}(X, B)} = \card{\operatorname{Mono}(X, D, \beta')}$ and $\card{\operatorname{Mono}(X, C, \lnot \beta)} = \card{\operatorname{Mono}(X, D, \lnot \beta', \alpha')}$.
        \begin{flalign*}
            \card{\operatorname{Mono}(X, B)} &= \card{\operatorname{Mono}(X, D, \beta')}
            \\
            \card{\operatorname{Mono}(X, C, \lnot \beta)} &= \card{\operatorname{Mono}(X, D, \lnot \beta', \alpha')}
        \end{flalign*}
    \end{minipage}
    \hfill
    \begin{minipage}{0.29\textwidth}
        \hfill
        \begin{tikzpicture}
            \node (A) {$A$};
            \node [below of=A] (B) {$B$}; 
            \node [left of=A] (C) {$C$}; 
            \node [left of=B] (D) {$D$}; 
            \begin{scope}[nodes=rectangle]          
            \draw [>->] (A) to node [right,label,pos=0.5] {$\alpha$} (B);
            \draw [>->] (A) to node [above,label,pos=0.5] {$\beta$} (C);
            \draw [>->] (B) to node [below,label,pos=0.45] {$\beta'$} (D); 
            \draw [>->] (C) to node [left,label,pos=0.45] {$\alpha'$} (D);
            \end{scope}
        \end{tikzpicture}
    \end{minipage} 
\end{lemma}
Therefore, we obtain:
\begin{flalign*}
      &\card{\operatorname{Mono}(X,G)} -  \card{\operatorname{Mono}(X,H)} \\
     =& \card{\operatorname{Mono}(X,L, \lnot l)}- \card{\operatorname{Mono}(X,R, \lnot r)}  + \\ 
     &\card{\operatorname{Mono}(X,G, \lnot l',\lnot m)} - 
     \card{\operatorname{Mono}(X,H, \lnot r',\lnot m') } 
\end{flalign*}
By Lemma~\ref{lem:xlnlmxrnr} (see \iflongversion
    Appendix \autoref{proof:lem:xlnlmxrnr}
    \else
    \cite[Lemma 28]{qiu2025termination}
    \fi for a proof), it can be simplified to:
\begin{flalign*}
     &\card{\operatorname{Mono}(X,G)} -  \card{\operatorname{Mono}(X,H)} \\
    = & \card{\operatorname{Mono}(X,L)}- \card{\operatorname{Mono}(X,R)}  + \\ 
      & \card{\operatorname{Mono}(X,G, \lnot l',\lnot m)} - 
    \card{\operatorname{Mono}(X,H, \lnot r',\lnot m') } 
\end{flalign*}
\begin{lemma}
    \label{lem:xlnlmxrnr}
    Let $X$ be a graph. Let $L \overset{l}{\leftarrowtail} K \overset{r}{\rightarrowtail} R$ be an injective DPO graph rewriting rule. We have 
    \[
       \card{\operatorname{Mono}(X, L, \lnot l)}  - \card{\operatorname{Mono}(X, R, \lnot r)} 
       = 
       \card{\operatorname{Mono}(X, L)}  - \card{\operatorname{Mono}(X, R)} 
        \]
\end{lemma}
Consequently, the weight difference satisfies:
\begin{flalign*}
     &w_{s_\mathbb{X}}(G) - w_{s_\mathbb{X}}(H)\\
    \overset{\operatorname{def}}{=}&\sum_{X \in \mathbb{X}}^{}s_\mathbb{X}(X) * m_X(G) - \sum_{X \in \mathbb{X}}^{}s_\mathbb{X}(X) * m_X(H)\\
    \overset{\operatorname{def}}{=}&\sum_{X \in \mathbb{X}}^{}s_\mathbb{X}(X) * |\operatorname{Mono}(X,G)| - \sum_{X \in \mathbb{X}}^{}s_\mathbb{X}(X) * |\operatorname{Mono}(X,H)|\\
    =&\sum_{X \in \mathbb{X}}^{}s_\mathbb{X}(X) * \left( \card{\operatorname{Mono}(X,G)} -  \card{\operatorname{Mono}(X,H)} \right)\\
    =&\sum_{X \in \mathbb{X}}^{}s_\mathbb{X}(X) * \left(  \card{\operatorname{Mono}(X,L)}- \card{\operatorname{Mono}(X,R)} \right)+\\
       & \sum_{X \in \mathbb{X}}^{}s_\mathbb{X}(X) * \left( 
     \card{\operatorname{Mono}(X,G, \lnot l',\lnot m)} - 
    \card{\operatorname{Mono}(X,H, \lnot r',\lnot m') } \right)
\end{flalign*}
Since $\sum_{X \in \mathbb{X}}^{}s(X) * \left(  \card{\operatorname{Mono}(X,L)}- \card{\operatorname{Mono}(X,R)} \right)$ can be precisely computed, the key challenge is to establish a lower bound for
\begin{flalign*}
    % \label{ineq:key_challenge}
    \sum_{X \in \mathbb{X}}^{} s_\mathbb{X}(X) * 
         \left ( \card{\operatorname{Mono}(X,G, \lnot l',\lnot m)} - 
     \card{\operatorname{Mono}(X,H, \lnot r',\lnot m') } \right ) 
    % \label{ineq:ge0_sss}
\end{flalign*} 
\section{Non-increasing rules}
\label{sec:non-increasing}
% \noindent
% \begin{minipage}{0.7\textwidth}
%     % \setlength{\parindent}{1em}
%     Let \( \rho = (L \overset{l}{\leftarrowtail} K \overset{r}{\rightarrowtail} R) \) be an injective DPO rewriting rule. 
%     A DPO diagram as illustrated on the right defines a rewriting step using $\rho$ with injective match \( m \).
% \end{minipage}%
% \hfill 
% \begin{minipage}{0.3\textwidth}
%     \hfill
%     % \begin{figure}[htbp] 
%     %     \center
%         \begin{tikzpicture}
%             % [node distance=15mm]
%             \node (I) {$K$};
%             \node (L) [left of=I] {$L$};
%             \node (R) [right of=I] {$R$};
%             \node (G) [below of=L] {$G$};
%             \node (C) [below of=I] {$C$};
%             \node (H) [below of=R] {$H$};
%             \draw [>->] (I) to  node [midway,below] {$l$} (L);
%             \draw [>->] (I) to  node [midway,below] {$r$} (R);
%             \draw [>->] (L) to node [midway,right] {$m$} (G);
%             \draw [>->] (I) to  node [midway,right] 
%             % {$u$}
%             {} (C);
%             \draw [>->] (R) to  node [midway,right] 
%             % {$m'$}
%             {} (H);
%             \draw [>->] (C) to node [midway,above] {$l'$} (G);
%             \draw [>->] (C) to node [midway,above] 
%             % {$r'$} 
%             {}
%             (H);
%             % \node [at=($(I)!.5!(G)$)] {\normalfont PO};
%             % \node [at=($(I)!.5!(H)$)] {\normalfont PO};
%         \end{tikzpicture}
%             % \caption{}
%     %         \label{diag:set_dpo_intuition_1}
%     % \end{figure}
% \end{minipage}

% \noindent
% \begin{minipage}{0.5\textwidth}
%     % \setlength{\parindent}{1em}
%      It follows that the DPO diagram depicted on the right, where all morphisms are inclusion functions and $H'$ is a graph isomorphic to $H$, holds.
% \end{minipage}%
% \hfill 
% \begin{minipage}{0.5\textwidth}  
%     \hfill
%     \begin{tikzpicture}
%         [node distance=20mm]
%         \node (I) at (0,0) {$\operatorname{Im}(l \star m)$};
%         \node (L) at (-2,0) {$\operatorname{Im}(m)$};
%         \node (R) at (2,0) {$\operatorname{Im}(m')$};
%         \node (G) at (-2,-1) {$G$};
%         \node (C) at (0,-1) {$\operatorname{Im}(l')$};
%         \node (H) at (2,-1) {$H'$};
%         \draw [>->] (I) to  node [midway,above] {} (L);
%         \draw [>->] (I) to  node [midway,above] {} (R);
%         \draw [>->] (L) to node [midway,right] {} (G);
%         \draw [>->] (I) to  node [midway,right] {} (C);
%         \draw [>->] (R) to  node [midway,right] {} (H);
%         \draw [>->] (C) to node [midway,above] {} (G);
%         \draw [>->] (C) to node [midway,above] {} (H);
%         % \node [at=($(I)!.5!(G)$)] {\normalfont PO};
%         % \node [at=($(I)!.5!(H)$)] {\normalfont PO};
%       \end{tikzpicture}
% \end{minipage}


%  \noindent
%     \begin{minipage}{0.6\textwidth}
%         By viewing a graph as a set o f nodes and edges and graph monomorphisms as set injections, the above DPO diagram can be considered as a DPO diagram in the category of sets, illustrated on the right.
%     \end{minipage}
%     \hfill
%     \begin{minipage}{0.3\textwidth}
%           \hfill
%           \resizebox{\textwidth}{!}{
%             \begin{tikzpicture}
%               \coordinate (k) at (0, 0);
%               \draw[fill=white] ($(k)+(0,0)$) rectangle ($(k)+(0.5,0.5)$);
%               \node () at ($(k)+(0.25,0.25)$) {\( \mathrm{K} \)};
          
%               \coordinate (c) at (0, -2.2);
%               \draw[fill=blue!20]
%               ($(c)+(0,-0.5)$)
%               -- ($(c)+(0,0.5)$) 
%               -- ($(c)+(1,0.5)$) 
%               arc[start angle=0, end angle=-90, radius=1]
%               -- cycle;
%               \node () at ($(c)+(0.75,0.25)$) {\( \mathrm{C'} \)};
%               \draw[fill=white] ($(c)+(0,0)$) rectangle ($(c)+(0.5,0.5)$);
%               \node () at ($(c)+(0.25,0.25)$) {\( \mathrm{K} \)};
          
%               \coordinate (l) at (-3, 0);
%               \draw[fill=orange!20] ($(l)+(-0.5,0)$) rectangle ($(l)+(0.5,1)$);
%               \node () at ($(l)+(-0.23,0.25)$) {\( \mathrm{L'} \)};
%               \draw[fill=white] ($(l)+(0,0)$) rectangle ($(l)+(0.5,0.5)$);
%               \node () at ($(l)+(0.25,0.25)$) {\( \mathrm{K} \)};
          
%               \coordinate (g) at (-3, -2.2);
%               \draw[fill=blue!20]
%               ($(g)+(0,-0.5)$)
%               -- ($(g)+(0,0.5)$)
%               -- ($(g)+(1,0.5)$) 
%               arc[start angle=0, end angle=-90, radius=1]
%               -- cycle;
%               \draw[fill=orange!20] ($(g)+(-0.5,0)$) rectangle ($(g)+(0.5,1)$);
%               \node () at ($(g)+(0.75,0.25)$) {\( \mathrm{C'} \)};
%               \node () at ($(g)+(-0.23,0.25)$) {\( \mathrm{L'} \)};
%               \draw[fill=white] ($(g)+(0,0)$) rectangle ($(g)+(0.5,0.5)$);
%               \node () at ($(g)+(0.25,0.25)$) {\( \mathrm{K} \)};
          
%               \coordinate (r) at (3,0);
%               \draw[fill=red!20] ($(r)+(-0.5,0)$)
%                 -- ($(r)+(-0.5,0.5)$)
%                 -- ($(r)+(0,1)$)
%                 --  ($(r)+(0.5,1)$)
%                 -- ($(r)+(0.5,0)$)
%                 -- cycle;
%               \node () at ($(r)+(-0.23,0.25)$) {\( \mathrm{R'} \)};
%               \draw[fill=white] ($(r)+(0,0)$) rectangle ($(r)+(0.5,0.5)$);
%               \node () at ($(r)+(0.25,0.25)$) {\( \mathrm{K} \)};
          
%               \coordinate (h) at (3, -2.2);
%               \draw[fill=blue!20]
%               ($(h)+(0,-0.5)$)
%               -- ($(h)+(0,0.5)$)
%               -- ($(h)+(1,0.5)$) 
%               arc[start angle=0, end angle=-90, radius=1]
%               -- cycle;
%               \draw[fill=red!20] ($(h)+(-0.5,0)$)
%               -- ($(h)+(-0.5,0.5)$)
%               -- ($(h)+(0,1)$)
%               --  ($(h)+(0.5,1)$)
%               -- ($(h)+(0.5,0)$)
%               -- cycle;
%              \node () at ($(h)+(0.75,0.25)$) {\( \mathrm{C'} \)};
%              \draw[fill=white] ($(h)+(0,0)$) rectangle ($(h)+(0.5,0.5)$);
%              \node () at ($(h)+(0.25,0.25)$) {\( \mathrm{K} \)};
%              \node () at ($(h)+(-0.23,0.25)$) {\( \mathrm{R'} \)};
          
%               \node[ font=\huge] (kl) at ($(k)!0.5!(l)+(0.25,0.25)$)
%              %    {\( \overset{l}{\leftarrowtail} \)}
%                 {\( \leftarrowtail \)}
%                ; 
%               \node[ font=\huge] (kr) at ($(k)!0.5!(r)+(0.25,0.25)$)
%                {\( \rightarrowtail \)}
%                ;  
%               \node[ font=\huge] (cg) at ($(c)!0.5!(g)+(0.25,0.25)$) 
%               {\( \leftarrowtail \)}
%              ;  
%               \node[ font=\huge] (ch) at ($(c)!0.5!(h)+(0.25,0.25)$)
%                {\( \rightarrowtail \)}
%              ; 
%               \node[ font=\huge] (kc) at ($(k)!0.5!(c)+(0.2,0.4)$) {\( \downarrowtail \)}; 
%             %   \node[ font=\LARGE] () at ($(l)!0.5!(g)+(0.5,0.4)$) {$m$}; 
%               \node[ font=\huge] (lg) at ($(l)!0.5!(g)+(0.1,0.4)$) {\( \downarrowtail \)}; 
%               \node[ font=\huge] (rh) at ($(r)!0.5!(h)+(0.1,0.4)$) {\( \downarrowtail \)}; 
%             %   \node[ font=\LARGE] () at ($(r)!0.5!(h)+(0.55,0.4)$) {$m'$}; 
%             %   \node[ font=\LARGE] () at ($(k)!0.5!(c)+(0.5,0.4)$) {$u$}; 
%             \end{tikzpicture}
%         }
%     \end{minipage}
\begin{minipage}{0.59\textwidth}
    % \setlength{\parindent}{1em}
    Let \( \rho = (L \overset{l}{\leftarrowtail} K \overset{r}{\rightarrowtail} R) \) be an injective DPO rewriting rule. 
    A rewriting step using $\rho$ with injective match \( m \) is defined by a DPO diagram as shown on the upper right.
     % \setlength{\parindent}{1em}
     It follows that both squares in the diagram depicted on the middle right, where all morphisms are inclusion functions and $H'$ is a graph isomorphic to $H$, are pushouts.
     By viewing a graph as a set of nodes and edges and graph monomorphisms as set injections, the DPO diagram on the middle right can be considered as a DPO diagram in the category of sets, illustrated on the lower right. In
\end{minipage}%
\hfill 
\begin{minipage}{0.39\textwidth}
    % \hfill 
    %     \resizebox{0.7\textwidth}{!}{
    %         \begin{tikzpicture}
    %             % [node distance=15mm]
    %             \node (I) {$K$};
    %             \node (L) [left of=I] {$L$};
    %             \node (R) [right of=I] {$R$};
    %             \node (G) [below of=L] {$G$};
    %             \node (C) [below of=I] {$C$};
    %             \node (H) [below of=R] {$H$};
    %             \draw [>->] (I) to  node [midway,below] {$l$} (L);
    %             \draw [>->] (I) to  node [midway,below] {$r$} (R);
    %             \draw [>->] (L) to node [midway,right] {$m$} (G);
    %             \draw [>->] (I) to  node [midway,right] 
    %             % {$u$}
    %             {} (C);
    %             \draw [>->] (R) to  node [midway,right] 
    %             % {$m'$}
    %             {} (H);
    %             \draw [>->] (C) to node [midway,above] {$l'$} (G);
    %             \draw [>->] (C) to node [midway,above] 
    %             % {$r'$} 
    %             {}
    %             (H);
    %             % \node [at=($(I)!.5!(G)$)] {\normalfont PO};
    %             % \node [at=($(I)!.5!(H)$)] {\normalfont PO};
    %         \end{tikzpicture}
    %     }
    \centering
    \resizebox{0.7\textwidth}{!}{
        \begin{tikzpicture}
            % [node distance=15mm]
            \node (I) at (0,0) {$K$};
            \node (L)  at (-2,0) {$L$};
            \node (R)  at (2,0) {$R$};
            \node (G)  at (-2,-2) {$G$};
            \node (C)  at (0,-2) {$C$};
            \node (H)  at (2,-2) {$H$};
            \draw [>->] (I) to  node [midway,below] {$l$} (L);
            \draw [>->] (I) to  node [midway,below] {$r$} (R);
            \draw [>->] (L) to node [midway,right] {$m$} (G);
            \draw [>->] (I) to  node [midway,right] 
            % {$u$}
            {} (C);
            \draw [>->] (R) to  node [midway,right] 
            {$m'$}
            (H);
            \draw [>->] (C) to node [midway,above] {$l'$} (G);
            \draw [>->] (C) to node [midway,above] 
            % {$r'$} 
            {}
            (H);
            % \node [at=($(I)!.5!(G)$)] {\normalfont PO};
            % \node [at=($(I)!.5!(H)$)] {\normalfont PO};
        \end{tikzpicture}
    }

        \resizebox{0.7\textwidth}{!}{
            \begin{tikzpicture}
                [node distance=20mm]
                \node (I) at (0,0) {$\operatorname{Im}(l \star m)$};
                \node (L) at (-2,0) {$\operatorname{Im}(m)$};
                \node (R) at (2,0) {$\operatorname{Im}(m')$};
                \node (G) at (-2,-2) {$G$};
                \node (C) at (0,-2) {$\operatorname{Im}(l')$};
                \node (H) at (2,-2) {$H'$};
                \draw [>->] (I) to  node [midway,above] {} (L);
                \draw [>->] (I) to  node [midway,above] {} (R);
                \draw [>->] (L) to node [midway,right] {} (G);
                \draw [>->] (I) to  node [midway,right] {} (C);
                \draw [>->] (R) to  node [midway,right] {} (H);
                \draw [>->] (C) to node [midway,above] {} (G);
                \draw [>->] (C) to node [midway,above] {} (H);
                % \node [at=($(I)!.5!(G)$)] {\normalfont PO};
                % \node [at=($(I)!.5!(H)$)] {\normalfont PO};
            \end{tikzpicture}
        }

          \resizebox{0.7\textwidth}{!}{
            \begin{tikzpicture}
                \coordinate (k) at (0, 0);
                \draw[fill=white] ($(k)+(0,0)$) rectangle ($(k)+(0.5,0.5)$);
                \node () at ($(k)+(0.25,0.25)$) {\( \mathrm{K} \)};
            
                \coordinate (c) at (0, -2.2);
                \draw[fill=blue!20]
                ($(c)+(0,-0.5)$)
                -- ($(c)+(0,0.5)$) 
                -- ($(c)+(1,0.5)$) 
                arc[start angle=0, end angle=-90, radius=1]
                -- cycle;
                \node () at ($(c)+(0.75,0.25)$) {\( \mathrm{C'} \)};
                \draw[fill=white] ($(c)+(0,0)$) rectangle ($(c)+(0.5,0.5)$);
                \node () at ($(c)+(0.25,0.25)$) {\( \mathrm{K} \)};
            
                \coordinate (l) at (-3, 0);
                \draw[fill=orange!20] ($(l)+(-0.5,0)$) rectangle ($(l)+(0.5,1)$);
                \node () at ($(l)+(-0.23,0.25)$) {\( \mathrm{L'} \)};
                \draw[fill=white] ($(l)+(0,0)$) rectangle ($(l)+(0.5,0.5)$);
                \node () at ($(l)+(0.25,0.25)$) {\( \mathrm{K} \)};
            
                \coordinate (g) at (-3, -2.2);
                \draw[fill=blue!20]
                ($(g)+(0,-0.5)$)
                -- ($(g)+(0,0.5)$)
                -- ($(g)+(1,0.5)$) 
                arc[start angle=0, end angle=-90, radius=1]
                -- cycle;
                \draw[fill=orange!20] ($(g)+(-0.5,0)$) rectangle ($(g)+(0.5,1)$);
                \node () at ($(g)+(0.75,0.25)$) {\( \mathrm{C'} \)};
                \node () at ($(g)+(-0.23,0.25)$) {\( \mathrm{L'} \)};
                \draw[fill=white] ($(g)+(0,0)$) rectangle ($(g)+(0.5,0.5)$);
                \node () at ($(g)+(0.25,0.25)$) {\( \mathrm{K} \)};
            
                \coordinate (r) at (3,0);
                \draw[fill=red!20] ($(r)+(-0.5,0)$)
                -- ($(r)+(-0.5,0.5)$)
                -- ($(r)+(0,1)$)
                --  ($(r)+(0.5,1)$)
                -- ($(r)+(0.5,0)$)
                -- cycle;
                \node () at ($(r)+(-0.23,0.25)$) {\( \mathrm{R'} \)};
                \draw[fill=white] ($(r)+(0,0)$) rectangle ($(r)+(0.5,0.5)$);
                \node () at ($(r)+(0.25,0.25)$) {\( \mathrm{K} \)};
            
                \coordinate (h) at (3, -2.2);
                \draw[fill=blue!20]
                ($(h)+(0,-0.5)$)
                -- ($(h)+(0,0.5)$)
                -- ($(h)+(1,0.5)$) 
                arc[start angle=0, end angle=-90, radius=1]
                -- cycle;
                \draw[fill=red!20] ($(h)+(-0.5,0)$)
                -- ($(h)+(-0.5,0.5)$)
                -- ($(h)+(0,1)$)
                --  ($(h)+(0.5,1)$)
                -- ($(h)+(0.5,0)$)
                -- cycle;
            \node () at ($(h)+(0.75,0.25)$) {\( \mathrm{C'} \)};
            \draw[fill=white] ($(h)+(0,0)$) rectangle ($(h)+(0.5,0.5)$);
            \node () at ($(h)+(0.25,0.25)$) {\( \mathrm{K} \)};
            \node () at ($(h)+(-0.23,0.25)$) {\( \mathrm{R'} \)};
            
                \node[ font=\huge] (kl) at ($(k)!0.5!(l)+(0.25,0.25)$)
            %    {\( \overset{l}{\leftarrowtail} \)}
                {\( \leftarrowtail \)}
                ; 
                \node[ font=\huge] (kr) at ($(k)!0.5!(r)+(0.25,0.25)$)
                {\( \rightarrowtail \)}
                ;  
                \node[ font=\huge] (cg) at ($(c)!0.5!(g)+(0.25,0.25)$) 
                {\( \leftarrowtail \)}
            ;  
                \node[ font=\huge] (ch) at ($(c)!0.5!(h)+(0.25,0.25)$)
                {\( \rightarrowtail \)}
            ; 
                \node[ font=\huge] (kc) at ($(k)!0.5!(c)+(0.2,0.4)$) {\( \downarrowtail \)}; 
            %   \node[ font=\LARGE] () at ($(l)!0.5!(g)+(0.5,0.4)$) {$m$}; 
                \node[ font=\huge] (lg) at ($(l)!0.5!(g)+(0.1,0.4)$) {\( \downarrowtail \)}; 
                \node[ font=\huge] (rh) at ($(r)!0.5!(h)+(0.1,0.4)$) {\( \downarrowtail \)}; 
            %   \node[ font=\LARGE] () at ($(r)!0.5!(h)+(0.55,0.4)$) {$m'$}; 
            %   \node[ font=\LARGE] () at ($(k)!0.5!(c)+(0.5,0.4)$) {$u$}; 
            \end{tikzpicture}
          }
\end{minipage}
the lower right diagram, $L', C', R', K$ are mutually disjoint sets such that $K$, 
         $K \uplus L'$, 
         $K \uplus R'$, 
         $K \uplus C'$, 
         $L' \uplus K \uplus C'$ and 
         $R' \uplus K \uplus C'$ 
            are corresponding sets of 
         $\operatorname{Im}(l \star m)$, 
         $\operatorname{Im}(m)$, 
         $\operatorname{Im}(m')$,
         $\operatorname{Im}(l')$,
         $G$ and
         $H'$ respectively,
         and all morphisms as inclusion functions.
    % Up to isomorphism, we can further consider all morphisms as inclusion functions, illustrated on the right. In this visualization, $L',R',C',K$ are mutually disjoint sets with $L' \uplus K = L$, $R' \uplus K= R$, $C' \uplus K = C$, $G = L \cup C$ and $H = R \cup C$.
    Thus, any subgraph of $G$ (resp. $H$) can be considered as a union of subgraphs of $C$ and $L$ (resp. $R$), up to isomorphism.
    % \color{red}\begin{definition}
%         Consider a rewriting step defined by the DPO diagram above. Let $X$ be a ruler-graph and $G$ a graph. An \textbf{$X$-occurrence} in $G$ is a subgraph of $G$ isomorphic to $X$. An \textbf{$X$-occurrence $x$ implicitly created by the rewriting step} is an \( X \)-occurrence in $H$ which is included neither in \( R \) nor in $C$. 
%         Let $R' = x \cap R$, $C' = x \cap C$ and $h_{R'L}:R' \rightarrowtail L$ a monomorphism preserving interface elements. The \textbf{corresponding $X$-occurrence relative to $h_{R'L}$} of $x$ is the $X$-occurrence in $G$ defined by $h_{R'L}(R') \cup C'$.
%         An \textbf{$X$-occurrence $x$ implicitly destroyed by the rewriting step} is an \( X \)-occurrence in $G$ which is included neither in \( L \) nor in $C$. 
%     \end{definition}
% \color{black}

With the above DPO diagrams in mind, we define:

\begin{definition}
    \label{def:corresponding_occurrence}
     An \textbf{$X$-occurrence} in $G$ is a subgraph of $G$ isomorphic to $X$. An \textbf{$X$-occurrence $x$ implicitly created by the rewriting step} is an \( X \)-occurrence in $H$ which is included neither in \( R \) nor in $C$. 
      The occurrence $x$ can be decomposed as the union of graphs $x = R_x \cup C_x$ where $R_x = x \cap R$ and $C_x = x \cap C$. For every graph monomorphism $h : R_x \rightarrowtail L$ that preserves interface elements (i.e., \( h(k) = k \) for all \( k \in K \)), we define the \textbf{corresponding $X$-occurrence of $x$ relative to $h$} as the $X$-occurrence in $G$ given by $h(R_x) \cup C_x$.
    An \textbf{$X$-occurrence implicitly destroyed by the rewriting step} is an \( X \)-occurrence in $G$ which is included neither in \( L \) nor in $C$.
\end{definition}
We define the set \(D(R,X)\) of all subgraphs of \( R \) which, when glued along some common interface elements with a subgraph $C'$ of a context graph \( C \), can form an \( X \)-occurrence implicitly created by the rewriting step.

\begin{definition}
    \label{def:rx}
    % Let \(X\) be a graph and 
    % \(
    %     L \overset{l}{\leftarrowtail} K \overset{r}{\rightarrowtail} R
    % \) an injective DPO rewriting rule.
    % %  such that $l$ is edge-injective and $r$ is node-injective.
    % The set \(D(R,X)\) consists of all subgraphs \( R' \subseteq R \) such that there exists a diagram (shown on the right) where
    % \(h_{R'X} \colon R' \rightarrowtail X \) and \(h_{C'X} \colon C' \rightarrowtail X \) are not isomorphisms, and \(h_{K'K} \colon K' \rightarrowtail K \) and \(h_{R'R} \colon R' \rightarrowtail R \) are inclusion functions. The union of all subgraphs in \( D(R,X) \) is denoted \( R_X \).
    % Let \(X\) be a graph and 
    % \(
    %     L \overset{l}{\leftarrowtail} K \overset{r}{\rightarrowtail} R
    % \) an injective DPO rewriting rule.
    % %  such that $l$ is edge-injective and $r$ is node-injective.
    % The set \(D(R,X)\) consists of all subgraphs \( R' \subseteq R \) such that 
    % we can construct a diagram as shown on the right where \(h_{K'K} \colon K' \rightarrowtail K \) and \(h_{R'R} \colon R' \rightarrowtail R \) are inclusion functions, and \(h_{R'X} \colon R' \rightarrowtail X \) and \(h_{C'X} \colon C' \rightarrowtail X \) are not isomorphisms. The union of all subgraphs in \( D(R,X) \) is denoted \( R_X \).
    Let \(X\) be a graph and 
    \(
        L \overset{l}{\leftarrowtail} K \overset{r}{\rightarrowtail} R
    \) an injective DPO rewriting rule.
    %  such that $l$ is edge-injective and $r$ is node-injective.
    The set \(D(R,X)\) consists of all subgraphs \( R' \subseteq R \)
    satisfying the following conditions:
    \newline
    \noindent
    \begin{minipage}{0.69\textwidth}\setlength{\parindent}{1em}
    \begin{itemize}
        \item a diagram as shown on the right where \(h_{K'K} \colon K' \rightarrowtail K \) and \(h_{R'R} \colon R' \rightarrowtail R \) are inclusion functions can be constructed,
        \item $R'$ is not a subgraph of $r(K)$,
        \item $R'$ is not isomorphic to $X$.
    \end{itemize}
    \end{minipage}
    \begin{minipage}{0.29\textwidth}
        \hfill
        \resizebox{0.9\textwidth}{!}{
        \begin{tikzpicture}[node distance=11mm]
            \node (k) {K};
            \node (r) [above=of k] {R};
            \node (r') [right=of r] {$R'$};
            \node (k') [right=of k] {K'};
            \node (x)  [right=of r'] {X};
            \node (c') [below=of x] {C'};
            \node () [at=($(r)!0.5!(k')$)] {PB};
            \node () [at=($(x)!0.5!(k')$)] {PO};
            \draw[>->] (k) -- (r) node[pos= .5, left] {$r$}; 
            \draw[>->] (r') -- (r) node[pos= .5, above] {$h_{R'R}$};
            \draw[>->] (k') -- (k) node[pos= .5, below] {$h_{K'K}$};
            \draw[>->] (k') -- (r');
            \draw[>->] (k') -- (c');
            \draw[>->] (c') -- (x);
            \draw[>->] (r') -- (x);
        \end{tikzpicture}
        }
    \end{minipage}
    % The union of all subgraphs in \( D(R,X) \) is denoted \( R_X \).
\end{definition}  
The first condition of the definition ensures that any $R' \in D(R,X)$ can be glued along some common interface elements with a subgraph $C'$ of a context graph \( C \) to form an \( X \)-occurrence. The second and third conditions ensure that such $X$-occurrences are included neither in \( C \) nor in \( R \), thereby guaranteeing that they are implicitly created as defined in \autoref{def:corresponding_occurrence}.

% In another word, the set $D(R,X)$ contains all subgraphs $R''$ of $R$ which is not a subgraph of $K$ such that there exists a rewriting step defined by the diagram above such that $R'' \cup C''$ for a subgraph $C''$ of $C$ is an $X$-occurrence created by the rewriting step. 
\begin{notation}
    We use the notation from~\cite[Notation 1]{overbeek2023apbpotutorial} to visualize edge-labeled graph homomorphisms. Labeled graphs are enclosed in boxes with their names displayed in the top-left corner. Nodes and edges are assigned subsets of \(\mathbb{N}\) as identifiers, and these identifiers are chosen such that: (i) Each node or edge \( y \) in the codomain graph is assigned the union of the identifiers of all nodes or edges in the domain graph that are mapped to \( y \); (ii) The graph homomorphism is uniquely determined by this assignment.
    
    \noindent To further improve readability, we represent sets by listing their elements. Additionally, we omit identifiers when doing so does not cause confusion. This is illustrated in the following representation of a homomorphism \( h: G \to H \).
    
    \begin{center}
        \resizebox{0.45\textwidth}{!}{
        \begin{tikzpicture}
            \graphbox{\( G \)}{00mm}{-20mm}{45mm}{20mm}{2mm}{-5mm}{
                \coordinate (o) at (-5mm,-8mm); 
                \node[draw,circle] (l1) at ($(o)+(-10mm,0mm)$) {1};
                \node[draw,circle] (l2) at ($(l1)+(3,0)$) {2};
                \node[draw,circle] (l3) at ($(l1)+(1,0)$) {3};
                \node[draw,circle] (l4) at ($(l1)+(2,0)$) {4};
                \draw[->] (l1) -- (l3) node[midway,above] {a};
                \draw[->] (l3) -- (l4) node[midway,above] {b};
                \draw[->] (l4) -- (l2) node[midway,above] {a};
            }  
            \graphbox{\( H \)}{50mm}{-20mm}{34mm}{20mm}{2mm}{-5mm}{
                \coordinate (o) at (0mm,-8mm); 
                \node[draw,circle] (l1) at ($(o)+(-10mm,0mm)$) {1};
                \node[draw,circle] (l2) at ($(l1)+(2,0)$) {2};
                \node[draw,circle] (l3) at ($(l1)+(1,0)$) {3\ 4};
                \draw[->] (l1) -- (l3) node[midway,above] {a};
                \draw[->] (l3) edge[loop above] (l3) node[midway,above] {b};
                \draw[->] (l3) -- (l2) node[midway,above] {a};
            }      
            % \node () at (53mm,-30mm) {$\rightarrow$};
        \end{tikzpicture}
    }
    \end{center} 
    In this example, the sets \(\{1\}\), \(\{2\}\), \(\{3\}\), \(\{4\}\), and \(\{3,4\}\) are represented as \(1\), \(2\), \(3\), \(4\), and \(3\ 4\), respectively. Edge identifiers are omitted.
\end{notation} 

\begin{example}
    \label{ex:grsaa_rx}
    Consider the injective DPO rule from \cite[Example 6]{bruggink2014termination}:  
    \begin{center} 
        \resizebox{0.7\textwidth}{!}{
        \begin{tikzpicture}
            \graphbox{$L$}{0mm}{0mm}{34mm}{15mm}{2mm}{-5mm}{
                \coordinate (o) at (0mm,-3mm); 
                \node[draw,circle] (l1) at ($(o)+(-10mm,0mm)$) {1};
                \node[draw,circle] (l2) at ($(l1)+(2,0)$) {2};
                \node[draw,circle] (l3) at ($(l1) + (1,0)$) {3};
                \draw[->] (l1) -- (l3) node[midway,above] {a};
                \draw[->] (l3) -- (l2) node[midway,above] {a};
            }     
            \graphbox{$K$}{40mm}{0mm}{24mm}{15mm}{2mm}{-5mm}{
                \coordinate (o) at (5mm,-3mm); 
                \node[draw,circle] (l1) at ($(o)+(-10mm,0mm)$) {1};
                \node[draw,circle] (l2) at ($(l1)+(1,0)$) {2};
                % \node[draw,circle] (l3) at ($(l1) + (1,0)$) {$\ $};
                % \draw[->] (l1) -- (l3) node[midway,above] {a};
                % \draw[->] (l3) -- (l2) node[midway,above] {a};
            }    
            \graphbox{$R$}{70mm}{0mm}{45mm}{15mm}{2mm}{-5mm}{
                \coordinate (o) at (-5mm,-3mm); 
                \node[draw,circle] (l1) at ($(o)+(-10mm,0mm)$) {1};
                \node[draw,circle] (l2) at ($(l1)+(3,0)$) {2};
                \node[draw,circle] (l3) at ($(l1) + (1,0)$) {4};
                \node[draw,circle] (l4) at ($(l1) + (2,0)$) {5};
                \draw[->] (l1) -- (l3) node[midway,above] {a};
                \draw[->] (l3) -- (l4) node[midway,above] {b};
                \draw[->] (l4) -- (l2) node[midway,above] {a};
            }    

            \node () at (37mm,-8mm) {$\leftarrowtail$};
            \node () at (67mm,-8mm) {$\rightarrowtail$};

            % \draw[>->] (51mm,2mm) -- (52mm,3mm);
        \end{tikzpicture}
        }
    \end{center}
    Let $X$ be \tikz[baseline=-0.5ex]{
        \node (x) at (0,0) {$\bullet$};
        \node (y) at (1,0) {$\bullet$ };
        \node (z) at (2,0) { $\bullet$};
        \draw[->] (x) -- (y) node[midway, above] {$a$};
        \draw[->] (y) -- (z) node[midway, above] {$a$};
    } and $R' \in D(R,X)$. 
    $R'$ is neither the empty graph nor isomorphic to $X$, by Conditions 2 and 3 of \autoref{def:rx}. 
    Therefore, $R'$ must contain an interface node (either node 1 or node 2), because if it contains no interface node, Condition 1 of \autoref{def:rx} would be violated due to $X$'s connectivity.
    Additionally, $R'$ must contain node 4 or node 5 to avoid being a subgraph of $r(K)$ (Condition 2). 
    Therefore, $R'$ must contain either $1 \overset{a}{\to} 4$ or $5 \overset{a}{\to} 2$, because if node 4 or node 5 is isolated, Condition 1 would be violated due to $X$'s connectivity.
    Furthermore, $R'$ must not contain $4 \overset{b}{\to} 5$ by Condition 1.
    The following graph is not in $D(R,X)$ because it violates Condition 1.
    \begin{center}
        \resizebox{0.24\textwidth}{!}{
            \begin{tikzpicture}
                \graphbox{}{70mm}{0mm}{45mm}{15mm}{2mm}{-5mm}{
                    \coordinate (o) at (-5mm,-3mm); 
                    \node[draw,circle] (l1) at ($(o)+(-10mm,0mm)$) {1};
                    \node[draw,circle] (l3) at ($(l1) + (1,0)$) {4};
                    \node[draw,circle] (l2) at ($(l1)+(3,0)$) {2};
                    \node[draw,circle] (l4) at ($(l1) + (2,0)$) {5};
                    \draw[->] (l4) -- (l2) node[midway,above] {a};
                    \draw[->] (l1) -- (l3) node[midway,above] {a};
                }    
            \end{tikzpicture}
        } 
    \end{center}
    Thus, $D(R,X)$ consists of the four graphs $R'_1$, $R'_2$, $R'_3$ and $R'_4$ shown below.
    For each graph, the construction of the pullback square in Condition 1 is demonstrated in \autoref{example:grs_aa:has_more_left} and the construction of the pushout square in Condition 1 is straightforward. 
    \begin{center}
        \resizebox{0.24\textwidth}{!}{
            \begin{tikzpicture}
                \graphbox{$R'_1$}{70mm}{0mm}{45mm}{15mm}{2mm}{-5mm}{
                    \coordinate (o) at (-5mm,-3mm); 
                    \node[circle] (l1) at ($(o)+(-10mm,0mm)$) {};
                    \node[draw,circle] (l2) at ($(l1)+(3,0)$) {2};
                    \node[draw,circle] (l4) at ($(l1) + (2,0)$) {5};
                    \draw[->] (l4) -- (l2) node[midway,above] {a};
                }    
            \end{tikzpicture}
        }
        \resizebox{0.24\textwidth}{!}{
            \begin{tikzpicture}
                \graphbox{$R'_2$}{70mm}{0mm}{45mm}{15mm}{2mm}{-5mm}{
                    \coordinate (o) at (-5mm,-3mm); 
                    \node[draw,circle] (l1) at ($(o)+(-10mm,0mm)$) {1};
                    \node[draw,circle] (l2) at ($(l1)+(3,0)$) {2};
                    \node[draw,circle] (l4) at ($(l1) + (2,0)$) {5};
                    \draw[->] (l4) -- (l2) node[midway,above] {a};
                }    
            \end{tikzpicture}
        } 
        \resizebox{0.24\textwidth}{!}{
            \begin{tikzpicture}
                \graphbox{$R'_3$}{70mm}{0mm}{45mm}{15mm}{2mm}{-5mm}{
                    \coordinate (o) at (-5mm,-3mm); 
                    \node[draw,circle] (l1) at ($(o)+(-10mm,0mm)$) {1};
                    \node[draw,circle] (l3) at ($(l1) + (1,0)$) {4};
                    \draw[->] (l1) -- (l3) node[midway,above] {a};
                }   
        \end{tikzpicture}
        } 
        \resizebox{0.24\textwidth}{!}{
            \begin{tikzpicture}
                \graphbox{$R'_4$}{70mm}{0mm}{45mm}{15mm}{2mm}{-5mm}{
                    \coordinate (o) at (-5mm,-3mm); 
                    \node[draw,circle] (l1) at ($(o)+(-10mm,0mm)$) {1};
                    \node[draw,circle] (l2) at ($(l1)+(3,0)$) {2};
                    \node[draw,circle] (l3) at ($(l1) + (1,0)$) {4};
                    \draw[->] (l1) -- (l3) node[midway,above] {a};
                }    
            \end{tikzpicture}
        }
    \end{center}
\end{example} 
 
\begin{example}
    \label{ex:rx_counter_ex}
    Let $X$ be the graph 
    \tikz[baseline=-0.5ex]{ 
            \node (x) at (0,0) {$\bullet$}; 
            \node (y) at (1,0) {$\bullet$};
            \node (z) at (2,0) {$\bullet$};
            \draw[->] (x) -- (y)   {};
            \draw[->] (y) -- (z)   {};
    }. Consider the following rule.
    \begin{center}
        \resizebox{0.6\textwidth}{!}{
            \begin{tikzpicture}
                \graphbox{$L$}{0mm}{0mm}{34mm}{15mm}{2mm}{-5mm}{
                    \coordinate (o) at (0mm,-3mm); 
                    \node[draw,circle] (l1) at ($(o)+(-10mm,0mm)$) {1};
                    % \node[draw,circle] (l2) at ($(l1)+(2,0)$) {2};
                    \node[draw,circle] (l3) at ($(l1) + (1,0)$) {2};
                    \draw[->] (l3) -- (l1) node[midway,above] {};
                    % \draw[->] (l3) -- (l2) node[midway,above] {a};
                }     
                \graphbox{$K$}{40mm}{0mm}{24mm}{15mm}{2mm}{-5mm}{
                    \coordinate (o) at (5mm,-3mm); 
                    \node[draw,circle] (l1) at ($(o)+(-10mm,0mm)$) {1};
                    \node[draw,circle] (l2) at ($(l1)+(1,0)$) {2};
                    % \node[draw,circle] (l3) at ($(l1) + (1,0)$) {$\ $};
                    \draw[->] (l2) -- (l1) node[midway,above] {};
                    % \draw[->] (l3) -- (l2) node[midway,above] {a};
                }    
                \graphbox{$R$}{70mm}{0mm}{40mm}{15mm}{2mm}{-5mm}{
                    \coordinate (o) at (-5mm,-3mm); 
                    \node[draw,circle] (l1) at ($(o)+(-10mm,0mm)$) {1};
                    % \node[draw,circle] (l2) at ($(l1)+(3,0)$) {2};
                    \node[draw,circle] (l3) at ($(l1) + (1,0)$) {2};
                    \node[draw,circle] (l4) at ($(l1) + (2,0)$) {3};
                    \draw[->] (l3) -- (l1) node[midway,above] {};
                    \draw[->] (l3) -- (l4) node[midway,above] {};
                    % \draw[->] (l4) -- (l2) node[midway,above] {a};
                }    
                \node () at (37mm,-8mm) {$\leftarrowtail$};
                \node () at (67mm,-8mm) {$\rightarrowtail$};
            \end{tikzpicture}
        }
    \end{center}
    Let $R' \in D(R,X)$. $R'$ must contain node 3, otherwise it would be a subgraph of $r(K)$, violating Condition 2 of \autoref{def:rx}. Node 3 must be connected to an interface node, otherwise Condition 1 of \autoref{def:rx} would be violated. Therefore, $R'$ must include \raisebox{2pt}{\scalebox{0.6}{\tikz[baseline=-0.5ex]{
        \node [draw,circle] (x) at (0,0) {2};
        \node[draw,circle] (y) at (1,0) {3};
        \draw[->] (x) -- (y) {};
    }}}.
    The graph $R$ is not in $D(R,X)$ because it violates Condition 1. Thus, the set \( D(R,X) \) consists of $R'_1$:
    \raisebox{2pt}{\scalebox{0.6}{\tikz[baseline=-0.5ex]{
        \node [draw,circle] (x) at (0,0) {2};
        \node[draw,circle] (y) at (1,0) {3};
        \draw[->] (x) -- (y) {};
    }}} and $R'_2$:
    \raisebox{2pt}{\scalebox{0.6}{\tikz[baseline=-0.5ex]{
        \node [draw,circle] (node1) at (-1,0) {1};
        \node [draw,circle] (x) at (0,0) {2};
        \node[draw,circle] (y) at (1,0) {3};
        \draw[->] (x) -- (y) {};
    }}}.
\end{example}
We define the notion of \( X \)-non-increasing rules. Intuitively, if a rule $\rho$ is $X$-non-increasing, then for any rewriting step using $\rho$, there exists an injective mapping from the set of \( X \)-occurrences implicitly created by the step to the set of \( X \)-occurrences implicitly destroyed by the step. This will be proved in \autoref{lem:w_u_l_not_geq_r_not}.
\begin{definition}[Non-increasing rule]
    \label{def:creates_more_x_on_the_left}
    Let \(\rho = (L \overset{l}{\leftarrowtail} K \overset{r}{\rightarrowtail} R)\) be a rule and \(X\) a ruler-graph. 
    Rule \( \rho \) is \emph{\( X \)-non-increasing} if the following four conditions hold:
    \newline
    \vspace{1mm}
    \noindent
    \begin{minipage}{0.64\textwidth} 
    \begin{enumerate}
        \item For all $R' \in D(R,X)$, there exists a morphism \(h_{R'L}: R' \rightarrowtail L \) such that both squares in the diagram shown on the right are pullbacks, where all morphisms other than $h_{R'L}$, $l$ and $r$ are inclusion functions, 
        \item \label{def:non_increasing:non_clapse} For all $R' \in D(R,X)$, for all node or edge $r'$ in $R'$, if $ r' \notin \operatorname{Im}(r)$ then $h_{R'L}(r') \notin \operatorname{Im}(l)$,
        \item For all $R',R'' \in D(R,X)$, for all edges $x \in R'$ and $y \in R''$, if $x \neq y$, then $h_{R'L}(x) \neq h_{R''L}(y)$,
        \item If $X$ has isolated nodes, then for all $R',R'' \in D(R,X)$ and nodes $x \in R', y \in R''$, if $x \neq y$, then $h_{R'L}(x) \neq h_{R''L}(y)$.
    \end{enumerate} 
    \end{minipage}
    \hfill
    \begin{minipage}{0.4\textwidth}
        % \hfill
        \begin{center}
            % \resizebox{0.9\textwidth}{!}{ 
                % \begin{tikzpicture}[rotate=90]
                %     \node (k) {K}; 
                %     \node (k') [above=of k] {$K'$};
                %     \node (l) [left=of k] {$L$};
                %     \node (rb) [above=of l] {$R'$};
                %     \node (r) [right=of k] {$R$};
                %     \node (rb') [above=of r] {$R'$};
                %     \draw[<-<]  (l) -- (k) node [midway,below] {l};
                %     \draw[<-<]  (r) -- (k) node [midway,below] {r};
                %     \draw[>->]  (k') -- (rb');
                %     \draw[>->]  (rb') -- (r);
                %     \draw[>->]  (k') -- (rb);
                %     \draw[>->]  (k') -- (k);
                %     \draw[>->]  (rb) -- (l) node [midway,above,sloped] {$h_{R'L}$};
                %     \node () [at=($(k)!0.5!(rb)$)] {PB};
                %     \node () [at=($(k)!0.5!(rb')$)] {PB};
                % \end{tikzpicture}
                \begin{tikzpicture} 
                    \node (k) {K}; 
                    \node (k') [left=of k] {$K'$};
                    \node (l) [above=of k] {$L$};
                    \node (rb) [left=of l] {$R'$};
                    \node (r) [below=of k] {$R$};
                    \node (rb') [left=of r] {$R'$};
                    \draw[<-<]  (l) -- (k) node [midway,right] {l};
                    \draw[<-<]  (r) -- (k) node [midway,right] {r};
                    \draw[>->]  (k') -- (rb');
                    \draw[>->]  (rb') -- (r);
                    \draw[>->]  (k') -- (rb);
                    \draw[>->]  (k') -- (k);
                    \draw[>->]  (rb) -- (l) node [midway,above] {$h_{R'L}$};
                    \node () [at=($(k)!0.5!(rb)$)] {PB};
                    \node () [at=($(k)!0.5!(rb')$)] {PB};
                \end{tikzpicture}
            % }
        \end{center}
    \end{minipage}

    %  \begin{enumerate}
    %     % \item[(2)] \label{def:non_increasing:non_clapse} For all node or arrow $r'$ in $R'$, if $ r' \notin \operatorname{Im}(r)$ then $h_{R'L}(r') \notin \operatorname{Im}(l)$,
    %     % \item there exists a morphism $h_{R_XL}$ such that
    %     %     \( h_{R'R_X} \star h_{R_XL} = h_{R'L} \) for all $R' \in D(R,X)$;
    %     % \item[(3)] \label{def:non_increasing:edge_injective} For all $R',R'' \in D(R,X)$, for all edges $x \in R'$ and $y \in R''$, if $x \neq y$, then $h_{R'L}(x) \neq h_{R''L}(y)$,
    %     %  \item $h$  is edge-injective 
    %     % \item $\bigcup_{R' \in D(R,X)} h_{R'L}$ is an edge-injective morphism from $R_X$ to $L$;
    %     % \item[(4)] \label{def:non_increasing:isolated_nodes} If $X$ has isolated nodes, then for all $R',R'' \in D(R,X)$ and nodes $x \in R', y \in R''$, if $x \neq y$, then $h_{R'L}(x) \neq h_{R''L}(y)$.
    %     % \item $h$ is node-injective if $X$ has isolated nodes;
    %  \end{enumerate}
    %  \todo{est il util d'avoir "or arrow?" c'est une partie de edge injective, non ?}
\end{definition}
The first condition ensures that whenever $R' \in D(R,X)$ forms an $X$-occurrence in the result graph $H$, this occurrence has a corresponding $X$-occurrence relative to the morphism $h_{R'L}$ (as defined in \autoref{def:corresponding_occurrence}) in the host graph $G$.
The second condition guarantees that any $X$-occurrence implicitly created by the rewriting step has its corresponding $X$-occurrence not included in the context (Consequently, its corresponding $X$-occurrence is implicitly destroyed by the rewriting step). Counterexamples are provided in \autoref{ex:cond_2_necessaire} and \autoref{ex:cond_2_necessaire2}.
The third and fourth conditions ensure that distinct $X$-occurrences created implicitly have distinct corresponding $X$-occurrences.
Counterexamples for these conditions are given in \autoref{ex:cond3_necessaire} and \autoref{ex:cond4_necessaire}.

The following example shows an $X$-non-increasing rule. 
\begin{example} 
    \label{example:grs_aa:has_more_left}
    Let $X$ be \tikz[baseline=-0.5ex]{
        \node (x) at (0,0) {$\bullet$};
        \node (y) at (1,0) {$\bullet$};
        \node (z) at (2,0) {$\bullet$};
        \draw[->] (x) -- (y) node[midway, above] {$a$};
        \draw[<-] (z) -- (y) node[midway, above] {$a$};
    }. Consider the rule presented
    in \autoref{ex:grsaa_rx}. By \autoref{ex:grsaa_rx}, the set \( D(R,X) \) consists of exactly four graphs: $R'_1$, $R'_2$, $R'_3$ and $R'_4$.
    The rule is $X$-non-increasing because the following diagrams can be constructed, and $h_{R'_1L}, h_{R'_2L}, h_{R'_3L}$ and $h_{R'_4L}$ satisfy all conditions in \autoref{def:creates_more_x_on_the_left}. 
%1
\begin{center}
    \resizebox{0.7\textwidth}{!}{
        \begin{tikzpicture}
            \graphbox{$R'_1$}{0mm}{0mm}{34mm}{15mm}{2mm}{-5mm}{
                \coordinate (o) at (0mm,-3mm); 
                \node (l1) at ($(o)+(-10mm,0mm)$) {};
                \node[draw,circle] (l2) at ($(l1)+(2,0)$) {2};
                \node[draw,circle] (l3) at ($(l1) + (1,0)$) {5};
                % \draw[->] (l1) -- (l3) node[midway,above] {a};
                \draw[->] (l3) -- (l2) node[midway,above] {a};
            }     
            \graphbox{$K'$}{40mm}{0mm}{24mm}{15mm}{2mm}{-5mm}{
                \coordinate (o) at (5mm,-3mm); 
                \node (l1) at ($(o)+(-10mm,0mm)$) {};
                \node[draw,circle] (l2) at ($(l1)+(1,0)$) {2};
                % \node[draw,circle] (l3) at ($(l1) + (1,0)$) {$\ $};
                % \draw[->] (l1) -- (l3) node[midway,above] {a};
                % \draw[->] (l3) -- (l2) node[midway,above] {a};
            }    
            \graphbox{$R'_1$}{70mm}{0mm}{45mm}{15mm}{2mm}{-5mm}{
                \coordinate (o) at (-5mm,-3mm); 
                \node (l1) at ($(o)+(-10mm,0mm)$) {};
                \node[draw,circle] (l2) at ($(l1)+(3,0)$) {2};
                % \node[draw,circle] (l3) at ($(l1) + (1,0)$) {4};
                \node[draw,circle] (l4) at ($(l1) + (2,0)$) {5};
                % \draw[->] (l1) -- (l3) node[midway,above] {a};
                % \draw[->] (l3) -- (l4) node[midway,above] {b};
                \draw[->] (l4) -- (l2) node[midway,above] {a};
            }    
            \graphbox{$L$}{0mm}{-20mm}{34mm}{15mm}{2mm}{-5mm}{
                \coordinate (o) at (0mm,-3mm); 
                \node[draw,circle] (l1) at ($(o)+(-10mm,0mm)$) {1};
                \node[draw,circle] (l2) at ($(l1)+(2,0)$) {2};
                \node[draw,circle] (l3) at ($(l1) + (1,0)$) {5};
                \draw[->] (l1) -- (l3) node[midway,above] {a};
                \draw[->] (l3) -- (l2) node[midway,above] {a};
            }     
            \graphbox{$K$}{40mm}{-20mm}{24mm}{15mm}{2mm}{-5mm}{
                \coordinate (o) at (5mm,-3mm); 
                \node[draw,circle] (l1) at ($(o)+(-10mm,0mm)$) {1};
                \node[draw,circle] (l2) at ($(l1)+(1,0)$) {2};
                % \node[draw,circle] (l3) at ($(l1) + (1,0)$) {$\ $};
                % \draw[->] (l1) -- (l3) node[midway,above] {a};
                % \draw[->] (l3) -- (l2) node[midway,above] {a};
            }    
            \graphbox{$R$}{70mm}{-20mm}{45mm}{15mm}{2mm}{-5mm}{
                \coordinate (o) at (-5mm,-3mm); 
                \node[draw,circle] (l1) at ($(o)+(-10mm,0mm)$) {1};
                \node[draw,circle] (l2) at ($(l1)+(3,0)$) {2};
                \node[draw,circle] (l3) at ($(l1) + (1,0)$) {4};
                \node[draw,circle] (l4) at ($(l1) + (2,0)$) {5};
                \draw[->] (l1) -- (l3) node[midway,above] {a};
                \draw[->] (l3) -- (l4) node[midway,above] {b};
                \draw[->] (l4) -- (l2) node[midway,above] {a};
            }    
            \node () at (37mm,-8mm) {$\leftarrowtail$};
            \node () at (17mm,-17mm) {$\downarrowtail$};
            \node () at (11.5mm,-17.5mm) {$h_{R'_1L}$};
            \node () at (52mm,-17mm) {$\downarrowtail$};
            \node () at (92mm,-17mm) {$\downarrowtail$};
            \node () at (37mm,-18mm) {\text{PB}};
            \node () at (68mm,-18mm) {\text{PB}};
            \node () at (67mm,-8mm) {$\rightarrowtail$};
            \node () at (37mm,-28mm) {$\leftarrowtail$};
            \node () at (67mm,-28mm) {$\rightarrowtail$};
        \end{tikzpicture}
    } 
\end{center}   
% 2
\begin{center}
    \resizebox{0.7\textwidth}{!}{
        \begin{tikzpicture}
            \graphbox{$R'_2$}{0mm}{0mm}{34mm}{15mm}{2mm}{-5mm}{
                \coordinate (o) at (0mm,-3mm); 
                \node[draw,circle] (l1) at ($(o)+(-10mm,0mm)$) {1};
                \node[draw,circle] (l2) at ($(l1)+(2,0)$) {2};
                \node[draw,circle] (l3) at ($(l1) + (1,0)$) {5};
                % \draw[->] (l1) -- (l3) node[midway,above] {a};
                \draw[->] (l3) -- (l2) node[midway,above] {a};
            }     
            \graphbox{$K'$}{40mm}{0mm}{24mm}{15mm}{2mm}{-5mm}{
                \coordinate (o) at (5mm,-3mm); 
                \node[draw,circle] (l1) at ($(o)+(-10mm,0mm)$) {1};
                \node[draw,circle] (l2) at ($(l1)+(1,0)$) {2};
                % \node[draw,circle] (l3) at ($(l1) + (1,0)$) {$\ $};
                % \draw[->] (l1) -- (l3) node[midway,above] {a};
                % \draw[->] (l3) -- (l2) node[midway,above] {a};
            }    
            \graphbox{$R'_2$}{70mm}{0mm}{45mm}{15mm}{2mm}{-5mm}{
                \coordinate (o) at (-5mm,-3mm); 
                \node[draw,circle] (l1) at ($(o)+(-10mm,0mm)$) {1};
                \node[draw,circle] (l2) at ($(l1)+(3,0)$) {2};
                % \node[draw,circle] (l3) at ($(l1) + (1,0)$) {4};
                \node[draw,circle] (l4) at ($(l1) + (2,0)$) {5};
                % \draw[->] (l1) -- (l3) node[midway,above] {a};
                % \draw[->] (l3) -- (l4) node[midway,above] {b};
                \draw[->] (l4) -- (l2) node[midway,above] {a};
            }    
            \graphbox{$L$}{0mm}{-20mm}{34mm}{15mm}{2mm}{-5mm}{
                \coordinate (o) at (0mm,-3mm); 
                \node[draw,circle] (l1) at ($(o)+(-10mm,0mm)$) {1};
                \node[draw,circle] (l2) at ($(l1)+(2,0)$) {2};
                \node[draw,circle] (l3) at ($(l1) + (1,0)$) {5};
                \draw[->] (l1) -- (l3) node[midway,above] {a};
                \draw[->] (l3) -- (l2) node[midway,above] {a};
            }      
            \graphbox{$K$}{40mm}{-20mm}{24mm}{15mm}{2mm}{-5mm}{
                \coordinate (o) at (5mm,-3mm); 
                \node[draw,circle] (l1) at ($(o)+(-10mm,0mm)$) {1};
                \node[draw,circle] (l2) at ($(l1)+(1,0)$) {2};
                % \node[draw,circle] (l3) at ($(l1) + (1,0)$) {$\ $};
                % \draw[->] (l1) -- (l3) node[midway,above] {a};
                % \draw[->] (l3) -- (l2) node[midway,above] {a};
            }     
            \graphbox{$R$}{70mm}{-20mm}{45mm}{15mm}{2mm}{-5mm}{
                \coordinate (o) at (-5mm,-3mm); 
                \node[draw,circle] (l1) at ($(o)+(-10mm,0mm)$) {1};
                \node[draw,circle] (l2) at ($(l1)+(3,0)$) {2};
                \node[draw,circle] (l3) at ($(l1) + (1,0)$) {4};
                \node[draw,circle] (l4) at ($(l1) + (2,0)$) {5};
                \draw[->] (l1) -- (l3) node[midway,above] {a};
                \draw[->] (l3) -- (l4) node[midway,above] {b};
                \draw[->] (l4) -- (l2) node[midway,above] {a};
            }    
            \node () at (37mm,-8mm) {$\leftarrowtail$};
            \node () at (17mm,-17mm) {$\downarrowtail$};
            \node () at (11.5mm,-17.5mm) {$h_{R'_2L}$};
            \node () at (52mm,-17mm) {$\downarrowtail$};
            \node () at (92mm,-17mm) {$\downarrowtail$};
            \node () at (37mm,-18mm) {\text{PB}};
            \node () at (68mm,-18mm) {\text{PB}};
            \node () at (67mm,-8mm) {$\rightarrowtail$};
            \node () at (37mm,-28mm) {$\leftarrowtail$};
            \node () at (67mm,-28mm) {$\rightarrowtail$};
        \end{tikzpicture}
    } 
\end{center}
%3
\begin{center}
    \resizebox{0.7\textwidth}{!}{
        \begin{tikzpicture}
            \graphbox{$R'_3$}{0mm}{0mm}{34mm}{15mm}{2mm}{-5mm}{
                \coordinate (o) at (0mm,-3mm); 
                \node[draw,circle] (l1) at ($(o)+(-10mm,0mm)$) {1};
                % \node[draw,circle] (l2) at ($(l1)+(2,0)$) {2};
                \node[draw,circle] (l3) at ($(l1) + (1,0)$) {4};
                \draw[->] (l1) -- (l3) node[midway,above] {a};
                % \draw[->] (l3) -- (l2) node[midway,above] {a};
            }     
            \graphbox{$K'$}{40mm}{0mm}{24mm}{15mm}{2mm}{-5mm}{
                \coordinate (o) at (5mm,-3mm); 
                \node[draw,circle] (l1) at ($(o)+(-10mm,0mm)$) {1};
                % \node[draw,circle] (l2) at ($(l1)+(1,0)$) {2};
                % \node[draw,circle] (l3) at ($(l1) + (1,0)$) {$\ $};
                % \draw[->] (l1) -- (l3) node[midway,above] {a};
                % \draw[->] (l3) -- (l2) node[midway,above] {a};
            }     
            \graphbox{$R'_3$}{70mm}{0mm}{45mm}{15mm}{2mm}{-5mm}{
                \coordinate (o) at (-5mm,-3mm); 
                \node[draw,circle] (l1) at ($(o)+(-10mm,0mm)$) {1};
                % \node[draw,circle] (l2) at ($(l1)+(3,0)$) {4};
                \node[draw,circle] (l3) at ($(l1) + (1,0)$) {4};
                % \node[draw,circle] (l4) at ($(l1) + (2,0)$) {5};
                \draw[->] (l1) -- (l3) node[midway,above] {a};
                % \draw[->] (l3) -- (l4) node[midway,above] {b};
                % \draw[->] (l4) -- (l2) node[midway,above] {a};
            }    
            \graphbox{$L$}{0mm}{-20mm}{34mm}{15mm}{2mm}{-5mm}{
                \coordinate (o) at (0mm,-3mm); 
                \node[draw,circle] (l1) at ($(o)+(-10mm,0mm)$) {1};
                \node[draw,circle] (l2) at ($(l1)+(2,0)$) {2};
                \node[draw,circle] (l3) at ($(l1) + (1,0)$) {4};
                \draw[->] (l1) -- (l3) node[midway,above] {a};
                \draw[->] (l3) -- (l2) node[midway,above] {a};
            }     
            \graphbox{$K$}{40mm}{-20mm}{24mm}{15mm}{2mm}{-5mm}{
                \coordinate (o) at (5mm,-3mm); 
                \node[draw,circle] (l1) at ($(o)+(-10mm,0mm)$) {1};
                \node[draw,circle] (l2) at ($(l1)+(1,0)$) {2};
                % \node[draw,circle] (l3) at ($(l1) + (1,0)$) {$\ $};
                % \draw[->] (l1) -- (l3) node[midway,above] {a};
                % \draw[->] (l3) -- (l2) node[midway,above] {a};
            }    
            \graphbox{$R$}{70mm}{-20mm}{45mm}{15mm}{2mm}{-5mm}{
                \coordinate (o) at (-5mm,-3mm); 
                \node[draw,circle] (l1) at ($(o)+(-10mm,0mm)$) {1};
                \node[draw,circle] (l2) at ($(l1)+(3,0)$) {2};
                \node[draw,circle] (l3) at ($(l1) + (1,0)$) {4};
                \node[draw,circle] (l4) at ($(l1) + (2,0)$) {5};
                \draw[->] (l1) -- (l3) node[midway,above] {a};
                \draw[->] (l3) -- (l4) node[midway,above] {b};
                \draw[->] (l4) -- (l2) node[midway,above] {a};
            }    
            \node () at (37mm,-8mm) {$\leftarrowtail$};
            \node () at (17mm,-17mm) {$\downarrowtail$};
            \node () at (11.5mm,-17.5mm) {$h_{R'_3L}$};
            \node () at (52mm,-17mm) {$\downarrowtail$};
            \node () at (92mm,-17mm) {$\downarrowtail$};
            \node () at (37mm,-18mm) {\text{PB}};
            \node () at (68mm,-18mm) {\text{PB}};
            \node () at (67mm,-8mm) {$\rightarrowtail$};
            \node () at (37mm,-28mm) {$\leftarrowtail$};
            \node () at (67mm,-28mm) {$\rightarrowtail$};
        \end{tikzpicture}
    } 
\end{center}
%4
\begin{center}
    \resizebox{0.7\textwidth}{!}{
        \begin{tikzpicture}
            \graphbox{$R'_4$}{0mm}{0mm}{34mm}{15mm}{2mm}{-5mm}{
                \coordinate (o) at (0mm,-3mm); 
                \node[draw,circle] (l1) at ($(o)+(-10mm,0mm)$) {1};
                \node[draw,circle] (l2) at ($(l1)+(2,0)$) {2};
                \node[draw,circle] (l3) at ($(l1) + (1,0)$) {4};
                \draw[->] (l1) -- (l3) node[midway,above] {a};
                % \draw[->] (l3) -- (l2) node[midway,above] {a};
            }     
            \graphbox{$K'$}{40mm}{0mm}{24mm}{15mm}{2mm}{-5mm}{
                \coordinate (o) at (5mm,-3mm); 
                \node[draw,circle] (l1) at ($(o)+(-10mm,0mm)$) {1};
                \node[draw,circle] (l2) at ($(l1)+(1,0)$) {2};
                % \node[draw,circle] (l3) at ($(l1) + (1,0)$) {$\ $};
                % \draw[->] (l1) -- (l3) node[midway,above] {a};
                % \draw[->] (l3) -- (l2) node[midway,above] {a};
            }     
            \graphbox{$R'_4$}{70mm}{0mm}{45mm}{15mm}{2mm}{-5mm}{
                \coordinate (o) at (-5mm,-3mm); 
                \node[draw,circle] (l1) at ($(o)+(-10mm,0mm)$) {1};
                \node[draw,circle] (l2) at ($(l1)+(3,0)$) {2};
                \node[draw,circle] (l3) at ($(l1) + (1,0)$) {4};
                % \node[draw,circle] (l4) at ($(l1) + (2,0)$) {5};
                \draw[->] (l1) -- (l3) node[midway,above] {a};
                % \draw[->] (l3) -- (l4) node[midway,above] {b};
                % \draw[->] (l4) -- (l2) node[midway,above] {a};
            }    
            \graphbox{$L$}{0mm}{-20mm}{34mm}{15mm}{2mm}{-5mm}{
                \coordinate (o) at (0mm,-3mm); 
                \node[draw,circle] (l1) at ($(o)+(-10mm,0mm)$) {1};
                \node[draw,circle] (l2) at ($(l1)+(2,0)$) {2};
                \node[draw,circle] (l3) at ($(l1) + (1,0)$) {4};
                \draw[->] (l1) -- (l3) node[midway,above] {a};
                \draw[->] (l3) -- (l2) node[midway,above] {a};
            }     
            \graphbox{$K$}{40mm}{-20mm}{24mm}{15mm}{2mm}{-5mm}{
                \coordinate (o) at (5mm,-3mm); 
                \node[draw,circle] (l1) at ($(o)+(-10mm,0mm)$) {1};
                \node[draw,circle] (l2) at ($(l1)+(1,0)$) {2};
                % \node[draw,circle] (l3) at ($(l1) + (1,0)$) {$\ $};
                % \draw[->] (l1) -- (l3) node[midway,above] {a};
                % \draw[->] (l3) -- (l2) node[midway,above] {a};
            }    
            \graphbox{$R$}{70mm}{-20mm}{45mm}{15mm}{2mm}{-5mm}{
                \coordinate (o) at (-5mm,-3mm); 
                \node[draw,circle] (l1) at ($(o)+(-10mm,0mm)$) {1};
                \node[draw,circle] (l2) at ($(l1)+(3,0)$) {2};
                \node[draw,circle] (l3) at ($(l1) + (1,0)$) {4};
                \node[draw,circle] (l4) at ($(l1) + (2,0)$) {5};
                \draw[->] (l1) -- (l3) node[midway,above] {a};
                \draw[->] (l3) -- (l4) node[midway,above] {b};
                \draw[->] (l4) -- (l2) node[midway,above] {a};
            }    
            \node () at (37mm,-8mm) {$\leftarrowtail$};
            \node () at (17mm,-17mm) {$\downarrowtail$};
            \node () at (11.5mm,-17.5mm) {$h_{R'_4L}$};
            \node () at (52mm,-17mm) {$\downarrowtail$};
            \node () at (92mm,-17mm) {$\downarrowtail$}; 
            \node () at (37mm,-18mm) {\text{PB}};
            \node () at (68mm,-18mm) {\text{PB}};
            \node () at (67mm,-8mm) {$\rightarrowtail$};
            \node () at (37mm,-28mm) {$\leftarrowtail$};
            \node () at (67mm,-28mm) {$\rightarrowtail$};
        \end{tikzpicture}
    } 
\end{center}
To illustrate the $X$-non-increasing property, consider the rewriting step defined by the DPO diagram below. Whenever \( R'_i \) (for \( 1 \leq i \leq 4 \)) forms an \( X \)-occurrence $x$ that is included neither in \( R \) nor $C$ with a subgraph $C' \subseteq C \cap x \subseteq C$, the graph \( h_{R'_iL}(R'_i) \) forms a \( X \)-occurrence in $G$ with $C'$. Furthermore, distinct $X$-occurrences created in $H$ have distinct corresponding $X$-occurrences in $G$.
    \begin{center} 
        \resizebox{0.7\textwidth}{!}{
        \begin{tikzpicture}
            \graphbox{\( L \)}{0mm}{-3mm}{34mm}{12mm}{2mm}{2mm}{
                \coordinate (o) at (0mm,-8mm); 
                \node[draw,circle] (l1) at ($(o)+(-10mm,0mm)$) {1};
                \node[draw,circle] (l2) at ($(l1)+(2,0)$) {2};
                \node[draw,circle] (l3) at ($(l1) + (1,0)$) {3};
                \draw[->] (l1) -- (l3) node[midway,above] {a};
                \draw[->] (l3) -- (l2) node[midway,above] {a};
            } 

            \graphbox{\( K \)}{40mm}{-3mm}{34mm}{12mm}{2mm}{2mm}{
                \coordinate (o) at (0mm,-8mm); 
                \node[draw,circle] (l1) at ($(o)+(-10mm,0mm)$) {1};
                \node[draw,circle] (l2) at ($(l1)+(2,0)$) {2};
            }  

            \graphbox{\( R \)}{80mm}{-3mm}{45mm}{12mm}{2mm}{2mm}{
                \coordinate (o) at (-5mm,-8mm); 
                \node[draw,circle] (l1) at ($(o)+(-10mm,0mm)$) {1};
                \node[draw,circle] (l2) at ($(l1)+(3,0)$) {2};
                \node[draw,circle] (l3) at ($(l1) + (1,0)$) {4};
                \node[draw,circle] (l4) at ($(l1) + (2,0)$) {5};
                \draw[->] (l1) -- (l3) node[midway,above] {a};
                \draw[->] (l3) -- (l4) node[midway,above] {b};
                \draw[->] (l4) -- (l2) node[midway,above] {a};
            }    

            \graphbox{\( G \)}{0mm}{-22mm}{34mm}{22mm}{2mm}{-3mm}{
                \coordinate (o) at (0mm,-3mm); 
                \node[draw,circle] (l1) at ($(o)+(-10mm,0mm)$) {1};
                \node[draw,circle] (l2) at ($(l1)+(2,0)$) {2};
                \node[draw,circle] (l3) at ($(l1) + (1,0)$) {3};
                \node[draw,circle] (l4) at ($(l2) + (0,-1)$) {6};
                \draw[->] (l1) -- (l3) node[midway,above] {a};
                \draw[->] (l3) -- (l2) node[midway,above] {a};
                \draw[->] (l2) -- (l4) node[midway,right] {a};
                \node[draw,circle] (l6) at ($(l1) + (0,-1)$) {7};
                \draw[<-] (l1) -- (l6) node[midway,left] {a};
                \draw[->] (l2) edge[out=-135,in=-45]node[midway,below] {a} (l1) ;
            }    

            \graphbox{\( C  \)}{40mm}{-22mm}{34mm}{22mm}{2mm}{-3mm}{
                \coordinate (o) at (0mm,-3mm); 
                \node[draw,circle] (l1) at ($(o)+(-10mm,0mm)$) {1};
                \node[draw,circle] (l2) at ($(l1)+(2,0)$) {2};
                \node[draw,circle] (l4) at ($(l2) + (0,-1)$) {6};
                \draw[->] (l2) -- (l4) node[midway,right] {a};
                \draw[->] (l2) edge[out=-135,in=-45]node[midway,below] {a} (l1) ;
                \node[ draw,circle] (l6) at ($(l1) + (0,-1)$) {7};
                \draw[<-] (l1) -- (l6) node[midway,left] {a};
            }    

            \graphbox{\( H \)}{80mm}{-22mm}{45mm}{22mm}{2mm}{-3mm}{
                \coordinate (o) at (-5mm,-3mm); 
                \node[draw,circle] (l1) at ($(o)+(-10mm,0mm)$) {1};
                \node[draw,circle] (l2) at ($(l1)+(3,0)$) {2};
                \node[draw,circle] (l3) at ($(l1) + (1,0)$) {4};
                \node[draw,circle] (l4) at ($(l1) + (2,0)$) {5};
                \node[ draw,circle] (l5) at ($(l2) + (0,-1)$) {6};
                \node[ draw,circle] (l6) at ($(l1) + (0,-1)$) {7};
                \draw[<-] (l1) -- (l6) node[midway,left] {a};
                \draw[->] (l1) -- (l3) node[midway,above] {a};
                \draw[->] (l3) -- (l4) node[midway,above] {b};
                \draw[->] (l4) -- (l2) node[midway,above] {a};
                \draw[->] (l2) -- (l5) node[midway,right] {a};
                \draw[->] (l2) edge[out=-135,in=-45]node[midway,below] {a} (l1) ;
            }    

            \node () at (37mm,-8mm) {\( \leftarrowtail \)}; % K -> L
            \node () at (77mm,-8mm) {\( \rightarrowtail \)}; % K -> R
            \node () at (15mm,-18mm) {\( m\ \downarrowtail \)};
            \node () at (37mm,-33mm) {\( \leftarrowtail \)};
            \node () at (58mm,-18mm) {\( \downarrowtail \)};
            \node () at (102mm,-18mm) {\( \downarrowtail \)};
            \node () at (77mm,-33mm) {\( \rightarrowtail \)}; % C -> H
        \end{tikzpicture}
        }
    \end{center}
\end{example}

% cond 2 necessaire  1
\begin{example} 
    \label{ex:cond_2_necessaire}
    Let $X$ be the graph 
    \tikz[baseline=-0.5ex]{ 
            \node (x) at (0,0) {$\bullet$}; 
            \node (y) at (1,0) {$\bullet$};
            \node (z) at (2,0) {$\bullet$};
            \draw[->] (x) -- (y)   {};
            \draw[->] (y) -- (z)   {};
    }. The following rewriting rule is potentially $X$-increasing.
    \begin{center}
        \resizebox{0.7\textwidth}{!}{
            \begin{tikzpicture}
                \graphbox{$L$}{0mm}{0mm}{34mm}{15mm}{2mm}{-5mm}{
                    \coordinate (o) at (0mm,-3mm); 
                    \node[draw,circle] (l1) at ($(o)+(-10mm,0mm)$) {1};
                    % \node[draw,circle] (l2) at ($(l1)+(2,0)$) {2};
                    \node[draw,circle] (l3) at ($(l1) + (1,0)$) {2};
                    \draw[->] (l3) -- (l1) node[midway,above] {};
                    % \draw[->] (l3) -- (l2) node[midway,above] {a};
                }     
                \graphbox{$K$}{40mm}{0mm}{24mm}{15mm}{2mm}{-5mm}{
                    \coordinate (o) at (5mm,-3mm); 
                    \node[draw,circle] (l1) at ($(o)+(-10mm,0mm)$) {1};
                    \node[draw,circle] (l2) at ($(l1)+(1,0)$) {2};
                    % \node[draw,circle] (l3) at ($(l1) + (1,0)$) {$\ $};
                    \draw[->] (l2) -- (l1) node[midway,above] {};
                    % \draw[->] (l3) -- (l2) node[midway,above] {a};
                }    
                \graphbox{$R$}{70mm}{0mm}{45mm}{15mm}{2mm}{-5mm}{
                    \coordinate (o) at (-5mm,-3mm); 
                    \node[draw,circle] (l1) at ($(o)+(-10mm,0mm)$) {1};
                    % \node[draw,circle] (l2) at ($(l1)+(3,0)$) {2};
                    \node[draw,circle] (l3) at ($(l1) + (1,0)$) {2};
                    \node[draw,circle] (l4) at ($(l1) + (2,0)$) {3};
                    \draw[->] (l3) -- (l1) node[midway,above] {};
                    \draw[->] (l3) -- (l4) node[midway,above] {};
                    % \draw[->] (l4) -- (l2) node[midway,above] {a};
                }    
                \node () at (37mm,-8mm) {$\leftarrowtail$};
                \node () at (67mm,-8mm) {$\rightarrowtail$};
            \end{tikzpicture}
        }
    \end{center}
    As shown in \autoref{ex:rx_counter_ex}, \( D(R,X) \) contains the graph $R'$:
    \raisebox{2pt}{
        \scalebox{0.6}{\tikz[baseline=-0.5ex]{
        \node [draw,circle] (x) at (0,0) {2};
        \node[draw,circle] (y) at (1,0) {3};
        \draw[->] (x) -- (y) {};
    }}}. The unique monomorphism $h_{R'L}:R' \rightarrowtail L$ fails the second condition of \autoref{def:creates_more_x_on_the_left}. 
    %  as illustrated below 
    % \begin{center}
    %     \resizebox{0.6\textwidth}{!}{
    %     \begin{tikzpicture}
    %             \graphbox{$R'$}{40mm}{-20mm}{24mm}{15mm}{2mm}{-5mm}{
    %                 \coordinate (o) at (5mm,-3mm); 
    %                 \node[draw,circle] (l1) at ($(o)+(-10mm,0mm)$) {2};
    %                 \node[draw,circle] (l2) at ($(l1)+(1,0)$) {3};
    %                 % \node[draw,circle] (l3) at ($(l1) + (1,0)$) {$\ $};
    %                 \draw[->] (l1) -- (l2) node[midway,above] {};
    %                 % \draw[->] (l3) -- (l2) node[midway,above] {a};
    %             }    
    %             \graphbox{$L$}{70mm}{-20mm}{45mm}{15mm}{2mm}{-5mm}{
    %                 \coordinate (o) at (-5mm,-3mm); 
    %                 \node[draw,circle] (l1) at ($(o)+(-10mm,0mm)$) {3};
    %                 % \node[draw,circle] (l2) at ($(l1)+(3,0)$) {2};
    %                 \node[draw,circle] (l2) at ($(l1) + (1,0)$) {2};
    %                 \draw[->] (l2) -- (l1) node[midway,above] {};
    %                 % \draw[->] (l4) -- (l2) node[midway,above] {a};
    %             }    
    %             \node () at (67mm,-28mm) {$\rightarrowtail$};
    %         \end{tikzpicture}
    %     }
    % \end{center}
    For any rewriting step using this rule, the corresponding $X$-occurrence relative to $h_{R'L}$ of any implicitly created $X$-occurrence is included in the context.
\end{example}
% cond 2 necessaire  2
\begin{example}
    \label{ex:cond_2_necessaire2}
    Let $X$ be the graph 
    \tikz[baseline=-0.5ex]{ 
            \node (x) at (0,0) {$\bullet$}; 
            \node (y) at (1,0) {$\bullet$};
            \node (z) at (2,0) {$\bullet$};
            \draw[->] (x) -- (y)   {};
            % \draw[->] (y) -- (z)   {};
    }. The following rewriting rule is potentially $X$-increasing.
    \begin{center}
        \resizebox{0.7\textwidth}{!}{
            \begin{tikzpicture}
                \graphbox{$L$}{0mm}{0mm}{34mm}{15mm}{2mm}{-5mm}{
                    \coordinate (o) at (0mm,-3mm); 
                    \node[draw,circle] (l1) at ($(o)+(-10mm,0mm)$) {1};
                    \node[draw,circle] (l3) at ($(l1) + (1,0)$) {2};
                }     
                \graphbox{$K$}{40mm}{0mm}{24mm}{15mm}{2mm}{-5mm}{
                    \coordinate (o) at (5mm,-3mm); 
                    \node[draw,circle] (l1) at ($(o)+(-10mm,0mm)$) {1};
                    \node[draw,circle] (l2) at ($(l1)+(1,0)$) {2};
                }    
                \graphbox{$R$}{70mm}{0mm}{45mm}{15mm}{2mm}{-5mm}{
                    \coordinate (o) at (-5mm,-3mm); 
                    \node[draw,circle] (l1) at ($(o)+(-10mm,0mm)$) {1};
                    \node[draw,circle] (l3) at ($(l1) + (1,0)$) {2};
                    \node[draw,circle] (l4) at ($(l1) + (2,0)$) {3};
                }    
                \node () at (37mm,-8mm) {$\leftarrowtail$};
                \node () at (67mm,-8mm) {$\rightarrowtail$};
            \end{tikzpicture}
        }
    \end{center}
    The set \( D(R,X) \) contains $R'$:
    \raisebox{2pt}{ \scalebox{0.6}{\tikz[baseline=-0.5ex]{
        \node [draw,circle] (x) at (0,0) {1};
        \node[draw,circle] (y) at (1,0) {3};
    }}}.  
    Despite the existence of a unique monomorphism $h_{R'L}: R' \rightarrowtail L$ preserving interface elements, this rule fails the second condition of \autoref{def:creates_more_x_on_the_left}.
    For any rewriting step using this rule, the corresponding $X$-occurrence relative to $h_{R'L}$ of any implicitly created $X$-occurrence is included in the context.
\end{example}
\begin{example}
    \label{ex:cond3_necessaire}
    Let $X$ be the graph 
    \tikz[baseline=-0.5ex]{ 
            \node (x) at (0,0) {$\bullet$}; 
            \node (y) at (1,0) {$\bullet$};
            \node (z) at (2,0) {$\bullet$};
            \draw[->] (x) -- (y)   {};
            \draw[->] (y) -- (z) {};
    }. The following rewriting rule is potentially $X$-increasing.
    \begin{center}
        \resizebox{0.6\textwidth}{!}{
            \begin{tikzpicture}
                \graphbox{$L$}{0mm}{0mm}{30mm}{20mm}{0}{0}{
                    \node[draw,circle]  (x) at (-6mm,-12mm) {x};
                    \node[draw,circle] (y) at (6mm,-12mm) {};
                    % \node[draw,circle]  (z) at (6mm,0mm) {};
                    \draw[->]  (x) to (y);
                    % \draw[->] (y) to[bend right=20] (z);
                    % \draw[->]  (z) to[bend right=20] (y);
                }
                \graphbox{$K$}{40mm}{0mm}{30mm}{20mm}{0}{0}{
                    \node[draw,circle]  (x) at (-6mm,-12mm) {x};
                    % \node[draw,circle]  (y) at (6mm,-12mm) {y};
                }
                \graphbox{$R$}{80mm}{0mm}{30mm}{20mm}{0}{0}{
                    \node[draw,circle]  (x) at (-6mm,-12mm) {x};
                        \node[draw,circle]  (y) at (6mm,-8mm) {y};
                        \node[draw,circle]  (z) at (6mm,-16mm) {z};
                        \draw[->]  (x) to (y);
                        \draw[->]  (x) to (z);
                }
                \node () at (35mm,-12mm) {$\leftarrowtail$};
                \node () at (75mm,-12mm) {$\rightarrowtail$};
            \end{tikzpicture}
        }
    \end{center}
    The set \( D(R,X) \) contains exactly two elements $R'$:
    \raisebox{2pt}{ 
        \scalebox{0.6}{\tikz[baseline=-0.5ex]{
        \node [draw,circle] (x) at (0,0) {x};
        \node[draw,circle] (y) at (1,0) {y};
        \draw[->] (x) -- (y) {};
    }}} and $R''$:\raisebox{2pt}{ 
        \scalebox{0.6}{\tikz[baseline=-0.5ex]{
        \node [draw,circle] (x) at (0,0) {x};
        \node[draw,circle] (y) at (1,0) {z};
        \draw[->] (x) -- (y) {};
    }}}. 
    Despite unique monomorphisms \( h_{R'L}: R' \rightarrowtail L \) and \( h_{R''L}: R'' \rightarrowtail L \) preserving interface elements, they fail the third condition of \autoref{def:creates_more_x_on_the_left}. 
    For any rewriting step using this rule, implicitly created $X$-occurrences with the same subgraph of the context have the same corresponding $X$-occurrence. 
\end{example}
 
\begin{example}
    \label{ex:cond4_necessaire}
    Let $X$ be the graph 
    \tikz[baseline=-0.5ex]{ 
            \node (x) at (0,0) {$\bullet$}; 
            \node (y) at (1,0) {$\bullet$};
            \node (z) at (2,0) {$\bullet$};
            \draw[->] (x) -- (y)   {};
    }, which has an isolated node. The following rewriting rule is potentially $X$-increasing.
    \begin{center}
        \resizebox{0.6\textwidth}{!}{
            \begin{tikzpicture}
                \graphbox{$L$}{0mm}{0mm}{30mm}{20mm}{0}{0}{
                    \node[draw,circle]  (x) at (-6mm,-10mm) {x};
                    \node[draw,circle] (y) at (6mm,-10mm) {y};
                }
                \graphbox{$K$}{40mm}{0mm}{30mm}{20mm}{0}{0}{
                    \node[draw,circle]  (x) at (-6mm,-10mm) {x};
                }
                \graphbox{$R$}{80mm}{0mm}{30mm}{20mm}{0}{0}{
                    \node[draw,circle]  (x) at (-6mm,-10mm) {x};
                        \node[draw,circle]  (y) at (6mm,-6mm) {z};
                        \node[draw,circle]  (z) at (6mm,-14mm) {w};
                }
                \node () at (35mm,-10mm) {$\leftarrowtail$};
                \node () at (75mm,-10mm) {$\rightarrowtail$};
            \end{tikzpicture}
        }
    \end{center}
    The set \( D(R,X) \) contains exactly two elements $R'$:
    \raisebox{2pt}{\scalebox{0.6}{\tikz[baseline=-0.5ex]{
        \node [draw,circle] (x) at (0,0) {x};
        \node[draw,circle] (y) at (1,0) {z};
        % \draw[->] (x) -- (y) {};
    }}} and $R''$:\raisebox{2pt}{\scalebox{0.6}{\tikz[baseline=-0.5ex]{
        \node [draw,circle] (x) at (0,0) {x};
        \node[draw,circle] (y) at (1,0) {w};
        % \draw[->] (x) -- (y) {};
    }}}. 
    There are unique monomorphisms $h_{R'L}:R' \rightarrowtail L$ and $h_{R''L}:R'' \rightarrowtail L$ preserving interface elements, but they fail the fourth condition of \autoref{def:creates_more_x_on_the_left}.
    For any rewriting step using this rule, implicitly created $X$-occurrences with the same subgraph of the context have the same corresponding $X$-occurrence.
\end{example}


\section{Solution to the key challenge}
\label{sec:solution_to_the_key_challenge}
The following lemma (see \iflongversion
Appendix \hyperref[proof:lem:w_u_l_not_geq_r_not]{Section C}
\else
\cite[Lemma 40]{qiu2025termination}
\fi for a proof) ensures that if $\rho$ is $X$-non-increasing, then, for any rewriting step using $\rho$, 
more X-occurrences are implicitly destroyed than implicitly created by the rewriting step.
% the cardinality of the set of \( X \)-occurrences implicitly destroyed by the rewriting step is at least as large as the cardinality of the set of \( X \)-occurrences implicitly created by the rewriting step.
% There are more implicitly destroyed X-occurrences than implicitly created X-occurrences.

\begin{lemma} 
    \label{lem:w_u_l_not_geq_r_not}
    \ \newline 
    \noindent 
    \begin{minipage}{0.7\textwidth}
        Let $X$ be a ruler-graph and $\rho = (L \overset{l}{\leftarrowtail} K \overset{r}{\rightarrowtail} R)$ an injective DPO rewriting rule.
        Suppose that $\rho$ is $X$-non-increasing. For every rewriting step induced by the DPO diagram (shown on the right), we have 
        \[
            |\operatorname{Mono}(X, G, \lnot m, \lnot l')| \geq |\operatorname{Mono}(X, H, \lnot m', \lnot r')|
        \]
    \end{minipage}
    \hfill
    \begin{minipage}{0.29\textwidth}
        \hfill
        \resizebox{0.9\textwidth}{!}{
                \begin{tikzpicture}
            \node (k) at (0,1) {K};
            \node (l) at (-2,1) {L};
            \node (r) at (2,1) {R};
            \node (c) at (0,-1) {C};
            \node (g) at (-2,-1) {G};
            \node (h) at (2,-1) {H};
            \draw[<-<]  (l) -- (k) node [midway,below] {l};
            \draw[>->]  (k) -- (r) node [midway,below] {r};
            \draw[>->] (c) -- (g) node [midway, above] {l'};
            \draw[>->] (c) -- (h) node [midway,above] {r'};
            \draw[>->] (l) -- (g) node[midway, right] {m};
            \draw[>->] (r) -- (h) node[midway, left] {m'};
            \draw[>->] (k) -- (c) node[midway, left] {};
            \node () [at=($(l)!0.5!(c)$)] {PO};
            \node () [at=($(r)!0.5!(c)$)] {PO};
        \end{tikzpicture}
        }
    \end{minipage}
\end{lemma}
 
 
\subsection{Termination criterion}
\label{sec:termination} 
Finally, we present a lemma that allows us to overapproximate the change in weight of a graph when a rule is applied, followed by our termination criterion. 

\begin{lemma}[Decreasing step]
    \label{subgraph_counting:lem:w_g_geq_w_h_leq}
    Let $\rho = (L \overset{l}{\leftarrowtail} K \overset{r}{\rightarrowtail} R)$ be an injective DPO rewriting rule,
    \( \mathbb{X} \) a set of ruler-graphs,
    \( s_{\mathbb{X}} \colon \mathbb{X} \to \mathbb{N} \) a weight function,
    and \( G \Rightarrow_{\rho,\mathfrak{M}} H \) a rewriting step. 
    If $\rho$ is \( X \)-non-increasing for every ruler-graph \( X \in \mathbb{X} \), then $
        w_{s_\mathbb{X}}(G) - w_{s_\mathbb{X}}(H) 
        \geq 
        w_{s_\mathbb{X}}(L) - w_{s_\mathbb{X}}(R)
    $.
\end{lemma}
\begin{proof}
    \iflongversion
        See~\textsection~\ref{subgraph_counting:proof:lem:w_g_geq_w_h_leq}.
    \else
        See \cite[Lemma 41]{qiu2025termination}.
    \fi 
\end{proof} 
\begin{theorem}[Termination] 
    \label{subgraph_counting:thm:termination_grs}
    Let \(\mathcal{A}\) and \(\mathcal{B}\) be sets of injective DPO rewriting rules, $\mathbb{X}$ a set of ruler-graphs, and $s_\mathbb{X}$ a weight function. If the following hold:
    \begin{enumerate}
        \item  for every $\rho \in \mathcal{A} \cup \mathcal{B}$ and for every $X \in \mathbb{X}$, the rule $\rho$ is $X$-non-increasing,
        \item for every \(\rho \in \mathcal{A}\), we have \( w_{s_\mathbb{X}}(lhs(\rho)) > w_{s_\mathbb{X}}(rhs(\rho)) \),
        \item for every \(\rho \in \mathcal{B}\), we have \( w_{s_\mathbb{X}}(lhs(\rho)) \geq w_{s_\mathbb{X}}(rhs(\rho)) \).
    \end{enumerate}
    Then \(\Rightarrow_{\mathcal{A},\mathcal{M}}\) terminates relative to \(\Rightarrow_{\mathcal{B},\mathcal{M}}\).
\end{theorem}
\begin{proof}
    \iflongversion
        See~\textsection~\ref{subgraph_counting:proof:lem:w_g_geq_w_h_leq}.
    \else
        See \cite[Lemma 41]{qiu2025termination}.
    \fi 
\end{proof} 
\noindent Our technique is described for DPO rewriting with injective matches due to their expressiveness \cite{habel2001double}, but it also extends to non-injective matches.

\begin{definition}[Quotient Rule \cite{habel2001double}] 
    \ \newline 
    \noindent
    \begin{minipage}{0.7\textwidth}
        Given a rule $\rho = L \leftarrow K \rightarrow R$, a rule $\rho' = (L' \leftarrow K' \rightarrow R')$ is called a \textbf{quotient rule} of $\rho$ if there exists a DPO diagram (shown on the right)
    where all vertical graph morphisms are surjective. The set of quotient rules of $\rho$ is denoted by $Q(\rho)$, and 
    for a rule set $\mathcal{A}$, we write $Q(\mathcal{A}) = \bigcup_{\rho\in\mathcal{A}} Q(\rho)$.
    \end{minipage}
    \hfill
    \begin{minipage}{0.29\textwidth}
        \hfill
        \begin{tikzpicture}[scale=0.7]
            \node (k) at (0,0) {K};
            \node (l) at (-2,0) {L};
            \node (r) at (2,0) {R};
            \node (k') at (0,-2) {K'};
            \node (l') at (-2,-2) {L'};
            \node (r') at (2,-2)  {R'};
            \draw[->] (k) -> (l);
            \draw[->] (k) -> (r); 
            \draw[->] (k') -> (l'); 
            \draw[->] (k') -> (r'); 
            \draw[->] (k) -> (k'); 
            \draw[->] (l) -> (l'); 
            \draw[->] (r) -> (r'); 
            \node () [at=($(l)!0.5!(k')$)] {PO};
            \node () [at=($(r)!0.5!(k')$)] {PO};
        \end{tikzpicture}
    \end{minipage}
\end{definition}

\begin{lemma}[Quotient Lemma \cite{habel2001double}]
    Let $\rho$ be a DPO rewriting rule. For any graphs $G$ and $H$,
    $G \Rightarrow_{\rho,\mathfrak{F}} H$ if and only if $G \Rightarrow_{Q(\rho),\mathfrak{M}} H$.
\end{lemma}
\begin{corollary}
    \label{cor:termination}
    Let \(\mathcal{A}\) and \(\mathcal{B}\) be sets of injective DPO rewriting rules. 
    The rewriting relation $\Rightarrow_{\mathcal{A},\mathfrak{F}}$ is terminating relative to $\Rightarrow_{\mathcal{B},\mathfrak{F}}$ 
    if 
    $\Rightarrow_{Q(\mathcal{A}),\mathfrak{M}}$ is terminating relative to $\Rightarrow_{Q(\mathcal{B}),\mathfrak{M}}$.
\end{corollary} 

\section{Examples}
\label{sec:examples}
While the type graph method~\cite{bruggink2014termination, bruggink2015proving,endrullis2024generalized_arxiv_v2} can prove termination for this example, it requires constructing an appropriate weighted type graph, a task that is difficult in general~\cite[\textsection 6]{bruggink2015proving}. 
In contrast, our criterion requires checking simpler conditions.
\begin{example}
    \label{subgraph_counting:ex:termination:grsaa}
    Consider the rewriting rule $\rho$ in Example~\ref{subgraph_counting:ex:grsaa_rx}, let $X$ be the graph \tikz[baseline=-0.5ex]{
        \node[draw,circle] (x) at (0,0) { };
        \node[draw,circle] (y) at (1,0) { }; 
        \node[draw,circle] (z) at (2,0) { };
        \draw[->] (x) -- (y) node[midway, above] {$a$};
        \draw[<-] (z) -- (y) node[midway, above] {$a$};
    }, $\mathbb{X} = \{X\}$ and $s_\mathbb{X}(X) = 1$. The rule is $X$-non-increasing as shown in Example~\ref{example:grs_aa:has_more_left}. 
    Since \(w_{s_\mathbb{X}}(\operatorname{lhs}(\rho)) = 1 > 0 = w_{s_\mathbb{X}}(\operatorname{rhs}(\rho)) \),
    it terminates by Theorem~\ref{subgraph_counting:thm:termination_grs}.
\end{example}
We consider an example for which the techniques from \cite{bruggink2014termination,bruggink2015proving,endrullis2024generalized_arxiv_v2,plump2018modular,overbeek2024termination_lmcs} fail.
\begin{example} 
    \label{subgraph_counting:ex_contrib_variant}
    The rewriting rule shown below 
    % in Figure~\ref{fig:subgraph_counting:ex_confdkjfakljl} 
    is a variant of a rule presented in \cite[Example 6]{plump2018modular}, obtained by removing all edges in the interface:
     
    % \begin{figure}[H]
    %     \centering
    \begin{center}
        \resizebox{0.6\textwidth}{!}{
            \begin{tikzpicture}
                \graphbox{$L$}{0mm}{0mm}{35mm}{35mm}{2mm}{-5mm}{
                    \coordinate (delta) at (0,-18mm);
                    \node[draw,circle] (l1) at ($(delta) + (-1,1.5)$) {$1$};
                    \node[draw,circle] (l2) at ($(delta) + (1,1.5)$) {$2$};
                    \node[draw,circle] (l3) at ($(delta) + (0,0)$) {3};
                    \draw[->] (l1) -- (l3) node[midway,left] {$s$};
                    \draw[->] (l2) -- (l3) node[midway,right] {$s$};
                    \draw[->] (l3) edge [loop below] node {$0$} (l3);
                }
                \graphbox{$K$}{40mm}{0mm}{35mm}{35mm}{2mm}{-5mm}{
                    \coordinate (delta) at (0,-18mm);
                    \coordinate (interfaceorigin) at ($(delta) +(5,0)$);
                    \node[draw,circle] (r1) at ($(delta) +(-1,1.5)$) {$1$};
                    \node[draw,circle] (r2) at ($(delta) +(0.5,1.5)$) {$2$};
                    \node[draw,circle] (r3) at ($(delta) + (0,0)$) {3};
                    % \draw[->] (r1) -- (r3) node[midway,left] {$s$};
                    % \draw[->] (r3) edge [loop below] node {$0$} (r3);
                }
                \graphbox{$R$}{80mm}{0mm}{50mm}{35mm}{2mm}{-5mm}{
                    \coordinate (delta) at (-10mm,-18mm);
                    \node[draw,circle] (r1) at ($(delta) + (-1,1.5)$) {$1$};
                    \node[draw,circle] (r2) at ($(delta) + (0.5,1.5)$) {$2$};
                    \node[draw,circle] (r3) at ($(delta) + (0,0)$) {3};
                    \node[draw,circle] (r4) at ($(delta) + (1,0)$) {4};
                    \draw[->] (r1) -- (r3) node[midway,left] {$s$};
                    \draw[->] (r2) -- (r4) node[midway,right] {$s$};
                    \draw[->] (r4) edge [loop below] node {$0$} (r4);
                    \draw[->] (r3) edge [loop below] node {$0$} (r3);
                    \node[draw,circle] (r5) at ($(r2) + (1.5,0)$) {};
                    \draw[->] (r5) edge [loop below] node {$0$} (r5);
                    \draw[->] (r5) edge [loop right] node {$0$} (r5);
                    \draw[->] (r5) edge [loop left] node {$0$} (r5);
                }
                % \graphbox{$R_x$}{40mm}{40mm}{35mm}{35mm}{2mm}{-5mm}{
                %     \coordinate (delta) at (0,-18mm);
                %     \coordinate (rxorigin) at ($(interfaceorigin)+(0,6)$);
                %     \node[draw,circle] (r1) at ($(delta) + (-1,1.5)$) {$1$};
                %     \node[draw,circle] (r2) at ($(delta) +  (0.5,1.5)$) {$2$};
                %     \node[draw,circle] (r3) at ($(delta) +  (0,0)$) {3};
                %     \draw[->] (r1) -- (r3) node[midway,left] {$s$};
                %     % \draw[->] (r3) edge [loop below] node {$0$} (r3);
                % }
                \node () at (38mm,-18mm) {$\leftarrowtail$};
                \node () at (77mm,-18mm) {$\rightarrowtail$};
                % \node () at (57mm,2mm) {$\uparrowtail$};
                % \node () at (38mm,2mm) {$\swarrowtail$};
                % \node () at (79mm,2mm) {$\searrowtail$};
            \end{tikzpicture}
            }
          \end{center}
    %         \caption{}
    %         \label{fig:subgraph_counting:ex_confdkjfakljl}
    % \end{figure}
        %     \caption{Diagram of Example~\ref{ex_contrib_variant}}
        %     \label{fig:contrib_variant}
        % \end{figure}
        Let $X$ be the ruler-graph 
        \tikz[baseline=-0.5ex]{ 
                \node[draw,circle] (x) at (0,0) { }; 
                \node[draw,circle] (y) at (1,0) { };
                \node[draw,circle] (z) at (2,0) { };
                \draw[->] (x) -- (y) node[midway, above] {$s$};
                \draw[->] (z) -- (y) node[midway, above] {$s$};
        }, $\mathbb{X} = \{X\}$ and $s_\mathbb{X}(X)=1$.
        The set \( D(R,X) \) consists of two elements $R'_1$:
        \raisebox{2pt}{
            \scalebox{0.7}{\tikz[baseline=-0.5ex]{
            \node [draw,circle] (x) at (0,0) {$1$};
            \node[draw,circle] (y) at (1,0) {3};
            \draw[->] (x) -- (y) node[midway, above] {$s$};
        }}} and $R'_2$:
        \raisebox{2pt}{
            \scalebox{0.7}{\tikz[baseline=-0.5ex]{
            \node [draw,circle] (z) at (-1,0) {$2$};
            \node [draw,circle] (x) at (0,0) {$1$};
            \node[draw,circle] (y) at (1,0) {3};
            \draw[->] (x) -- (y) node[midway, above] {$s$};
        }}}. 
        Each \( R'_i \) admits a unique monomorphism \( h_{R'_i L} \colon R'_i \rightarrowtail L \) preserving interface elements. 
        Therefore, the rule is $X$-non-increasing under $h_{R'_1L}$ and $h_{R'_2L}$. This is because all conditions in Definition~\ref{subgraph_counting:def:creates_more_x_on_the_left} are satisfied.
        Since \(w_{s_\mathbb{X}}(L) = 1 > 0 = w_{s_\mathbb{X}}(R)\).
        It follows that the system terminates by Theorem~\ref{subgraph_counting:thm:termination_grs}.
\end{example}

The following example demonstrates the usefulness of a weight function that assigns distinct weight to measurements from different ruler-graphs.
\begin{example} 
    \label{ex:overbeek_5d6}
    Consider the rewriting rules presented in \cite[Example 5.6]{overbeek2024termination_lmcs} illustrated 
    below.
    % in Figure~\ref{fig:subgraph_counting:ex_overbdfjkahsdfkjfhew}.
    % \begin{figure}[H]
    %       \centering
    \begin{center}
        \begin{center} 
          $\rho = $\scalebox{0.7}{{
\begin{tikzpicture}[baseline=-10mm]
    \graphbox{$L$}{0mm}{0mm}{31mm}{20mm}{2mm}{-13mm}{
      \node [draw, circle] (x) at (-7mm,0mm) {1};
      \node [draw, circle] (y) at (0mm,0mm) {2};
      \draw[->] (x) edge[loop above] node  {$a$} (x);
      \draw[->] (y) edge [loop above] node {$a$} (y);
    }
    \graphbox{$K$}{32mm}{-0mm}{18mm}{20mm}{0mm}{-10mm}{
    }
    \begin{scope}[opacity=1]        
    \graphbox{$R$}{51mm}{-0mm}{38mm}{20mm}{2mm}{-13mm}{
      \node [draw, circle] (x) at (-7mm,0mm) {3};
      \node [draw, circle] (y) at (0mm,0mm) {4};
      \node [draw, circle] (z) at (7mm,0mm) {5};
      \draw[->] (x) edge[loop above] node  {$b$} (x);
      \draw[->] (y) edge[loop above] node  {$b$} (y);
      \draw[->] (z) edge[loop above] node  {$b$} (z);
    }
    \end{scope}
\end{tikzpicture}
}}
        \end{center}
        \begin{center}
        $\tau = $\scalebox{0.7}{{\begin{tikzpicture}[baseline=-10mm]
    \graphbox{L}{0mm}{0}{31mm}{20mm}{2mm}{-13mm}{
        \node[draw,circle] (x) at (-7mm,0mm) {1};
        \node[draw,circle] (y) at (0mm,0mm) {2};
        \draw[->] (x) edge[loop above] node {$b$} (x);
        \draw[->] (y) edge [loop above] node {$b$} (y);
      }
      \graphbox{K}{32mm}{0}{18mm}{20mm}{2mm}{-13mm}{
      }
      \begin{scope}[opacity=1]        
      \graphbox{R}{51mm}{0}{25mm}{20mm}{1mm}{-13mm}{
        \node[draw,circle] (x) at (-3.5mm,0mm) {3};
        \draw[->] (x) edge[loop above] node {$a$} (x);
      }
      \end{scope}
  \end{tikzpicture}}
        }
        \end{center}
  %       \caption{}
  %       \label{fig:subgraph_counting:ex_overbdfjkahsdfkjfhew}
  % \end{figure}
      \end{center}
     Let $X$ be the ruler-graph
    \tikz[baseline=-0.5ex]{
        \node[draw,circle] (x) at (0,0) { };
        \draw[->] (x) edge [loop right] node {$a$} (x);
    }, $Y$ the ruler-graph 
    \tikz[baseline=-0.5ex]{  
        \node[draw,circle] (x) at (0,0) { };
        \draw[->] (x) edge [loop right] node {$b$} (x);
    } and $\mathbb{X} = \{X, Y\}$.
    The sets $D(R_\rho, X)$, $D(R_\rho, Y)$, $D(R_\tau, X)$ and $D(R_\tau, Y)$ are all empty. Therefore, both rules $\rho$ and $\tau$ are $X$- and $Y$-non-increasing as all conditions in Definition~\ref{subgraph_counting:def:creates_more_x_on_the_left} are met.
    
    For $s_\mathbb{X}(X) = 5$ and $s_\mathbb{X}(Y) = 3$, we have $
    w_{s_\mathbb{X}}(L_\rho) = 10 > 9 = w_{s_\mathbb{X}}(R_\rho)
    $ and $
    w_{s_\mathbb{X}}(L_\tau) = 6 > 5 = w_{s_\mathbb{X}}(R_\tau)$.
    Therefore, the system terminates by Theorem~\ref{subgraph_counting:thm:termination_grs}.

  However, for $s_\mathbb{X}'(X) = 1$ and $s_\mathbb{X}'(Y) = 1$, we have
    \(
        w_{s_\mathbb{X}'}(L_\rho) = 2 \not > 3 = w_{s_\mathbb{X}'}(R_\rho)
    \) which prevents us from applying Theorem~\ref{subgraph_counting:thm:termination_grs} to prove termination of the rewriting system.
\end{example}

% \begin{example}[Limitation] 
%     \label{ex:plump95_4d1}
%    Consider a rule presented in \cite[Example 4.1]{Plump1995}:
 
%     \begin{center}
%         \resizebox{0.7\textwidth}{!}{
%       $\rho = $  \begin{tikzpicture}[baseline=-20mm]
%         \graphbox{$L$}{0mm}{0}{32mm}{40mm}{0}{0}{
%             \node[draw,circle]  (n1) at (0,-6mm) {$1$};
%             \node[draw,circle]   (n2) at (0mm,-26mm) {$2$};
%             \node[draw,circle]   (n3) at (-10mm,-26mm) {3};
%             \node[draw,circle]   (n4) at (10mm,-26mm) {4};
%             \draw[->]  (n2) edge [loop below] node  {$a$} (n2);
%             \draw[->]  (n3) edge [loop below] node  {$b$} (n3);
%             \draw[->]  (n1) to node [right] {$f$} (n2) ;
%             \draw[->]  (n1) to node [right] {$f$}  (n3);
%             \draw[->]  (n1) to node [right] {$f$}  (n4);
%           }
%           \graphbox{$K$}{33mm}{0}{32mm}{40mm}{0}{-0}{
%             \node[draw,circle]  (n1) at (0,-6mm) {$1$};
%             \node[draw,circle]   (n2) at (0mm,-26mm) {$2$};
%             \node[draw,circle]   (n3) at (-10mm,-26mm) {3};
%             \node[draw,circle]   (n4) at (10mm,-26mm) {4};
%             \draw[->]  (n2) edge [loop below] node  {$a$} (n2);
%             \draw[->]  (n3) edge [loop below] node  {$b$} (n3);
%           }
%           \graphbox{$R$}{66mm}{0}{32mm}{40mm}{0}{-0}{
%             \node[draw,circle]  (n1) at (0,-6mm) {$1$};
%             \node[draw,circle]   (n2) at (0mm,-26mm) {$2$};
%             \node[draw,circle]   (n3) at (-10mm,-26mm) {3};
%             \node[draw,circle]   (n4) at (10mm,-26mm) {4};
%             \draw[->]  (n2) edge [loop below] node  {$a$} (n2);
%             \draw[->]  (n3) edge [loop below] node  {$b$} (n3);
%             \draw[->]  (n1) edge [bend left] node [right] {$f$} (n4) ;
%             \draw[->]  (n1) edge [bend right] node [left] {$f$}  (n4);
%             \draw[->]  (n1) to node [right] {$f$}  (n4);
%           }   
%     \end{tikzpicture}
%         }
%   \end{center}
%   \noindent
%   \begin{minipage}{0.7\textwidth}
%     To prove its termination with our method, a ruler-graph $X$ containing an edge labeled by $f$ must be used, because the number of occurrences of every graph containing only edges labeled by $a$ and $b$ does not change by any rewriting step using $\rho$. But in this case, condition \eqref{def:non_increasing:edge_injective} of Definition~\ref{def:creates_more_x_on_the_left} cannot be satisfied.
%   \end{minipage}
%   \begin{minipage}{0.3\textwidth}
%     \hfill
%   \begin{center}
%     \begin{tikzpicture} 
%         \graphbox{}{0}{0}{30mm}{20mm}{-10mm}{-10mm}{
%             \node[draw,circle]  (n1) at (0,0mm) {$1$};
%             % \node[draw,circle]   (n2) at (0mm,-26mm) {$2$};
%             % \node[draw,circle]   (n3) at (-10mm,-26mm) {3};
%             \node[draw,circle]   (n4) at (20mm,0mm) {4};
%             % \draw[->]  (n2) edge [loop below] node  {$a$} (n2);
%             % \draw[->]  (n3) edge [loop below] node  {$b$} (n3);
%             \draw[->]  (n1) edge [bend left] node [above] {$f$} (n4) ;
%             \draw[->]  (n1) edge [bend right] node [below] {$f$}  (n4);
%             \draw[->]  (n1) to node {$f$}  (n4);
%           }  
%     \end{tikzpicture}
%   \end{center}
% \end{minipage}
  
%   \end{example}
% Finally, we present a limitation of our approach.
 

\begin{example}
  \label{ex_endrullis_6d3_endrullis_5d8}
  Consider the following rule from~\cite[Example 6.3]{endrullis2024generalized_arxiv_v2}: 
  %  in Figure~\ref{fig:subgraph_counting:ex_endrullis_6d3_endrullis_5d8}.
  % \begin{figure}[H]
  %     \centering
  \begin{center}
      \resizebox{0.6\textwidth}{!}{
          \begin{tikzpicture}
              \graphbox{$L$}{0mm}{0mm}{35mm}{25mm}{2mm}{-8mm}{
                  \coordinate (delta) at (0,-18mm);
                  \node[draw,circle] (x) at (-6mm,-10mm) {x};
                  \node[draw,circle] (y) at (6mm,-10mm) {y};
                  \node[draw,circle]  (z) at (6mm,0mm) {z};
                  \draw[->]  (x) to (y);
                  \draw[->] (y) to[bend right=20] (z);
                  \draw[->]  (z) to[bend right=20] (y);
              }
              \graphbox{$K$}{40mm}{0mm}{35mm}{25mm}{2mm}{-8mm}{
                  \node[draw,circle]  (x) at (-6mm,-10mm) {x};
                  \node[draw,circle]  (y) at (6mm,-10mm) {y};
              }
              \graphbox{$R$}{80mm}{0mm}{35mm}{25mm}{2mm}{-8mm}{
                  \node[draw,circle]  (x) at (-6mm,-10mm) {x};
                      \node[draw,circle]  (y) at (6mm,-10mm) {y};
                      \node[draw,circle]  (z) at (0mm,0mm) {w};
                      \draw[->]  (x) to (y);
                      \draw[->]  (y) to [bend left=0] (z);
                      \draw[->]  (z) to [bend left=0] (x);
              }
              % \graphbox{$R_x$}{40mm}{20mm}{35mm}{15mm}{2mm}{2mm}{
              %     \node[draw,circle]  (x) at (-6mm,-10mm) {x};
              %     \node[draw,circle]  (y) at (6mm,-10mm) {y};
              %     \draw[->]  (x) to (y);
              % }
              \node () at (37mm,-12mm) {$\leftarrowtail$};
              \node () at (78mm,-12mm) {$\rightarrowtail$};
              % \node () at (57mm,2mm) {$\uparrowtail$};
              % \node () at (38mm,2mm) {$\swarrowtail$};
              % \node () at (79mm,2mm) {$\searrowtail$};
          \end{tikzpicture}
      }
  %     \caption{}
  %     \label{fig:subgraph_counting:ex_endrullis_6d3_endrullis_5d8}
  % \end{figure}
  \end{center}

  Let $X$ be the graph 
  \tikz[baseline=-0.5ex]{
      \node[draw,circle] (x) at (0,0) { };
      \node[draw,circle] (y) at (1,0) { };
      \draw[->] (x) to[bend right=20] (y);
      \draw[->] (y) to[bend right=20] (x); 
  } and $\mathbb{X} = \{X\}$. Let $s_\mathbb{X}$ be the weight function associating the weight of $1$ to $X$. There is only one element 
  \scalebox{0.7}{\tikz[baseline=-0.5ex]{
          \node [draw,circle] (x) at (0,0) {x};
          \node[draw,circle] (y) at (1,0) {y};
          \draw[->] (x) -- (y) node[midway, above] {};
      }}
  in $D(X, R)$. 
   Since $w_{s_\mathbb{X}}(L) = 1 > 0 = w_{s_\mathbb{X}}(R)
      $, the rule terminates by Theorem~\ref{subgraph_counting:thm:termination_grs}.
\end{example}


\begin{example} 
  \label{ex:overbeek_5d5}
  Consider the following rewriting rul from~\cite[Example 5.5]{overbeek2024termination_lmcs}:
  \begin{center}
  % \begin{figure}[H]
  %   \centering
      \resizebox{0.6\textwidth}{!}{
          \begin{tikzpicture}
              \graphbox{$L$}{0mm}{0mm}{35mm}{25mm}{2mm}{-8mm}{
                  \coordinate (delta) at (0,-18mm);
                  \node[draw,circle] (x) at (-6mm,-10mm) {x};
                  \node[draw,circle] (y) at (6mm,-10mm) {y};
                  \node[draw,circle]  (z) at (6mm,0mm) {z};
                  \draw[->]  (x) to (y);
                  \draw[->] (y) to[bend right=20] (z);
                  \draw[->]  (z) to[bend right=20] (y);
              }
              \graphbox{$K$}{40mm}{0mm}{35mm}{25mm}{2mm}{-8mm}{
                  \node[draw,circle]  (x) at (-6mm,-10mm) {x};
                  \node[draw,circle]  (y) at (6mm,-10mm) {y};
                  \draw[->] (x) to (y);
              }
              \graphbox{$R$}{80mm}{0mm}{35mm}{25mm}{2mm}{-8mm}{
                  \node[draw,circle]  (x) at (-6mm,-10mm) {x};
                      \node[draw,circle]  (y) at (6mm,-10mm) {y};
                      \node[draw,circle]  (z) at (0mm,0mm) {w};
                      \draw[->]  (x) to (y);
                      \draw[->]  (y) to [bend left=0] (z);
                      \draw[->]  (z) to [bend left=0] (x);
              }
              % \graphbox{$R_x$}{40mm}{20mm}{35mm}{15mm}{2mm}{2mm}{
              %     \node[draw,circle]  (x) at (-6mm,-10mm) {x};
              %     \node[draw,circle]  (y) at (6mm,-10mm) {y};
              %     \draw[->]  (x) to (y);
              % }
              \node () at (37mm,-12mm) {$\leftarrowtail$};
              \node () at (78mm,-12mm) {$\rightarrowtail$};
              % \node () at (57mm,2mm) {$\uparrowtail$};
              % \node () at (38mm,2mm) {$\swarrowtail$};
              % \node () at (79mm,2mm) {$\searrowtail$};
          \end{tikzpicture}
      }
  %     \caption{}
  %     \label{fig:subgraph_counting:overbeek_5d5}
  % \end{figure}
  \end{center}
    The termination of the system in the DPO rewriting framework with monic matches can be established using our approach; the argument is identical to that of Example~\ref{ex_endrullis_6d3_endrullis_5d8}.
\end{example}

\begin{example}
  \label{ex:plump_ex4}
  Our method cannot prove the termination of the DPO rewriting system with monic matches from~\cite[Example 4]{plump2018modular}:
  %  in Figure~\ref{fig:subgraph_counting:ex_plump_ex4}.
  % \begin{figure}[H]
\begin{center}
  $r_1 = ${ 
   \begin{tikzpicture}[baseline=-20mm]
         \graphbox{$L$}{0mm}{-11mm}{35mm}{15mm}{0.5mm}{-9mm}{
           \node [draw, circle] (x) at (-10mm,0mm) {$1$};
           \node [draw, circle] (y) at (0mm,0mm) {$2$};
           \node [draw, circle] (z) at (10mm,0mm) {3};
           \draw[->] (x) to node [above] {$0$} (y);
           \draw[->] (y) to node [above] {$L$} (z);
         }
         \graphbox{$K$}{36mm}{-11mm}{35mm}{15mm}{0.5mm}{-9mm}{
           \node [draw, circle] (x) at (-10mm,0mm) {$1$};
           %\node [draw, circle] (y) at (0mm,0mm) {$2$};
           \node [draw, circle] (z) at (10mm,0mm) {3};
         }
         \begin{scope}[opacity=1]        
         \graphbox{$R$}{72mm}{-11mm}{40mm}{15mm}{-4mm}{-9mm}{
           \node [draw, circle] (x) at (-10mm,0mm) {$1$};
           \node [draw, circle] (y) at (0mm,0mm) {$2$};
           \node [draw, circle] (z) at (10mm,0mm) {3};
           \node [draw, circle] (w) at (20mm,0mm) {3};
           \draw[->] (x) to node [above] {$L$} (y);
           \draw[->] (y) to node [above] {$1$} (z);
           \draw[->] (z) to node [above] {$1$} (w);
         }
         \end{scope}
       \end{tikzpicture}
  }
 \end{center}
 
 \begin{center}
   $r_2 = ${ 
    \begin{tikzpicture}[baseline=-20mm]
       \graphbox{$L$}{0mm}{-11mm}{35mm}{15mm}{0.5mm}{-9mm}{
           \node [draw, circle] (x) at (-10mm,0mm) {$1$};
           \node [draw, circle] (y) at (0mm,0mm) {$2$};
           \node [draw, circle] (z) at (10mm,0mm) {3};
           \draw[->] (x) to node [above] {$R$} (y);
           \draw[->] (y) to node [above] {$1$} (z);
       }
       \graphbox{$K$}{36mm}{-11mm}{35mm}{15mm}{0.5mm}{-9mm}{
           \node [draw, circle] (x) at (-10mm,0mm) {$1$};
           %\node [draw, circle] (y) at (0mm,0mm) {$2$};
           \node [draw, circle] (z) at (10mm,0mm) {3};
       }
       \begin{scope}[opacity=1]        
       \graphbox{$R$}{72mm}{-11mm}{35mm}{15mm}{0.5mm}{-9mm}{
           \node [draw, circle] (x) at (-10mm,0mm) {$1$};
           \node [draw, circle] (y) at (0mm,0mm) {$2$};
           \node [draw, circle] (z) at (10mm,0mm) {3};
           \draw[->] (x) to node [above] {$0$} (y);
           \draw[->] (y) to node [above] {$R$} (z);
       }
       \end{scope}
       \end{tikzpicture}
   }
   \end{center}

\begin{center}
  $r_3 = ${ 
   \begin{tikzpicture}[baseline=-20mm]
         \graphbox{$L$}{0mm}{-11mm}{35mm}{15mm}{0.5mm}{-9mm}{
           \node [draw, circle] (x) at (-10mm,0mm) {$1$};
           \node [draw, circle] (y) at (0mm,0mm) {$2$};
           \node [draw, circle] (z) at (10mm,0mm) {3};
           \draw[->] (x) to node [above] {$B$} (y);
           \draw[->] (y) to node [above] {$L$} (z);
         }
         \graphbox{$K$}{36mm}{-11mm}{35mm}{15mm}{0.5mm}{-9mm}{
           \node [draw, circle] (x) at (-10mm,0mm) {$1$};
           %\node [draw, circle] (y) at (0mm,0mm) {$2$};
           \node [draw, circle] (z) at (10mm,0mm) {3};
         }
         \begin{scope}[opacity=1]        
         \graphbox{$R$}{72mm}{-11mm}{35mm}{15mm}{0.5mm}{-9mm}{
           \node [draw, circle] (x) at (-10mm,0mm) {$1$};
           \node [draw, circle] (y) at (0mm,0mm) {$2$};
           \draw[->] (x) to node [above] {$R$} (y);
         }
         \end{scope}
       \end{tikzpicture}
  }
 \end{center}
 
 \begin{center}
   $r_4 = ${ 
    \begin{tikzpicture}[baseline=-20mm]
       \graphbox{$L$}{0mm}{-11mm}{35mm}{15mm}{0.5mm}{-9mm}{
           \node [draw, circle] (x) at (-10mm,0mm) {$1$};
           \node [draw, circle] (y) at (0mm,0mm) {$2$};
           \node [draw, circle] (z) at (10mm,0mm) {3};
           \draw[->] (x) to node [above] {$R$} (y);
           \draw[->] (y) to node [above] {$B$} (z);
       }
       \graphbox{$K$}{36mm}{-11mm}{35mm}{15mm}{0.5mm}{-9mm}{
           \node [draw, circle] (x) at (-10mm,0mm) {$1$};
           %\node [draw, circle] (y) at (0mm,0mm) {$2$};
           \node [draw, circle] (z) at (10mm,0mm) {3};
       }
       \begin{scope}[opacity=1]        
       \graphbox{$R$}{72mm}{-11mm}{35mm}{15mm}{0.5mm}{-9mm}{
           \node [draw, circle] (x) at (-10mm,0mm) {$1$};
           \node [draw, circle] (y) at (0mm,0mm) {$2$};
           \node [draw, circle] (z) at (10mm,0mm) {3};
           \draw[->] (x) to node [above] {$L$} (y);
           \draw[->] (y) to node [above] {$B$} (z);
       }
       \end{scope}
       \end{tikzpicture}
   }
   \end{center}
  %  \caption{}
  %  \label{fig:subgraph_counting:ex_plump_ex4}
  % \end{figure}
  The rule $r_3$ can be eliminated by considering occurrences of \tikz[baseline=-0.5ex]{
    \node[draw,circle] (x) at (0,0) { };
    \node[draw,circle] (y) at (1,0) { };
    \draw[->] (x) -- (y) node[midway, above] {$B$};
  }, and the rule $r_4$ can be eliminated by considering occurrences of \tikz[baseline=-0.5ex]{
  \node[draw,circle] (x) at (0,0) { };
  \node[draw,circle] (y) at (1,0) { };
  \draw[->] (x) -- (y) node[midway, above] {$R$};
}.

However, our method cannot eliminate $r_1$ and $r_2$. Yet, combined with the technique proposed by Plump~\cite{plump2018modular}, we can prove the termination of this rewriting system: $\{r_1\}$ and $\{r_2\}$ are both terminating, and their union is terminating~\cite{plump2018modular}.
\end{example}


\begin{example}
  \label{ex:overbeek_5d8_plump1995_3d8_plump2018_3_overbeek_5d8}
  Consider the rewriting rules from~\cite[Example~3.8]{plump1995ontermination}:
  % in Figure~\ref{fig:subgraph_counting:overbeek_5d8_plump1995_3d8_plump2018_3_overbeek_5d8}.

  % \begin{figure}[H]
  %   \centering
\begin{center}
 $\rho = ${ 
  \begin{tikzpicture}[baseline=-20mm]
        \graphbox{$L$}{0mm}{-11mm}{35mm}{15mm}{0.5mm}{-9mm}{
          \node [draw, circle] (x) at (-10mm,0mm) {$1$};
          \node [draw, circle] (y) at (0mm,0mm) {$2$};
          \node [draw, circle] (z) at (10mm,0mm) {3};
          \draw[->] (x) to node [above] {$a$} (y);
          \draw[->] (y) to node [above] {$b$} (z);
        }
        \graphbox{$K$}{36mm}{-11mm}{35mm}{15mm}{0.5mm}{-9mm}{
          \node [draw, circle] (x) at (-10mm,0mm) {$1$};
          %\node [draw, circle] (y) at (0mm,0mm) {$2$};
          \node [draw, circle] (z) at (10mm,0mm) {3};
        }
        \begin{scope}[opacity=1]        
        \graphbox{$R$}{72mm}{-11mm}{35mm}{15mm}{0.5mm}{-9mm}{
          \node [draw, circle] (x) at (-10mm,0mm) {$1$};
          \node [draw, circle] (y) at (0mm,0mm) {$2$};
          \node [draw, circle] (z) at (10mm,0mm) {3};
          \draw[->] (x) to node [above] {$a$} (y);
          \draw[->] (y) to node [above] {$c$} (z);
        }
        \end{scope}
      \end{tikzpicture}
 }
\end{center}

\begin{center}
  $\tau = ${ 
   \begin{tikzpicture}[baseline=-20mm]
      \graphbox{$L$}{0mm}{-11mm}{35mm}{15mm}{0.5mm}{-9mm}{
          \node [draw, circle] (x) at (-10mm,0mm) {$1$};
          \node [draw, circle] (y) at (0mm,0mm) {$2$};
          \node [draw, circle] (z) at (10mm,0mm) {3};
          \draw[->] (x) to node [above] {$c$} (y);
          \draw[->] (y) to node [above] {$d$} (z);
      }
      \graphbox{$K$}{36mm}{-11mm}{35mm}{15mm}{0.5mm}{-9mm}{
          \node [draw, circle] (x) at (-10mm,0mm) {$1$};
          %\node [draw, circle] (y) at (0mm,0mm) {$2$};
          \node [draw, circle] (z) at (10mm,0mm) {3};
      }
      \begin{scope}[opacity=1]        
      \graphbox{$R$}{72mm}{-11mm}{35mm}{15mm}{0.5mm}{-9mm}{
          \node [draw, circle] (x) at (-10mm,0mm) {$1$};
          \node [draw, circle] (y) at (0mm,0mm) {$2$};
          \node [draw, circle] (z) at (10mm,0mm) {3};
          \draw[->] (x) to node [above] {$d$} (y);
          \draw[->] (y) to node [above] {$b$} (z);
      }
      \end{scope}
      \end{tikzpicture}
  }
  \end{center}
%   \caption{}
%   \label{fig:subgraph_counting:overbeek_5d8_plump1995_3d8_plump2018_3_overbeek_5d8}
% \end{figure}
  Let $X$ be  \tikz[baseline=-0.5ex]{
      \node[draw,circle] (x) at (0,0) { };
      \node[draw,circle] (y) at (1,0) { };
      \node[draw,circle] (z) at (2,0) { };
      \draw[->] (x) -- (y) node[midway, above] {$a$};
      \draw[<-] (z) -- (y) node[midway, above] {$b$};
  } 
  We have $D(R_\rho,X) = \emptyset$ and $D(R_\tau,X) = \emptyset$. Therefore, both rules are $X$-non-increasing. 
  We have 
  $|\operatorname{Mono}(X,L)| = 1 > 0 = |\operatorname{Mono}(X,R)|$ for $\rho$, and 
  $|\operatorname{Mono}(X,L)| = 0 = 0 = |\operatorname{Mono}(X,R)|$ for $\tau$, therefore $\rho$ can be eliminated.
  
  Let $Y$ be  \tikz[baseline=-0.5ex]{
      \node[draw,circle] (x) at (0,0) {};
      \node[draw,circle] (y) at (1,0) {};
      \draw[->] (x) -- (y) node[midway, above] {$c$};
  }. We have $D(R_\tau,Y) = \emptyset$. Therefore, $\tau$ is $Y$-non-increasing. Moreover, we have $|\operatorname{Mono}(Y,L)| = 1 > 0 = |\operatorname{Mono}(Y,R)|$.
  Thus, this rewriting system terminates by Theorem~\ref{subgraph_counting:thm:termination_grs}.
\end{example} 


\begin{example}
  \label{ex:overbeek_5d3}
  Consider the following rewriting rules from~\cite[Example 5.3]{overbeek2024termination_lmcs}:
  % in Figure~\ref{fig:subgraph_counting:ex_overbeek_5d3}.
  % \begin{figure}[H]
  %   \centering
\begin{center}
  $\rho = ${ 
  \begin{tikzpicture} 
      \graphbox{$L$}{0mm}{0mm}{23mm}{12mm}{2mm}{-5mm}{
          \node[draw,circle] (x) at (0mm,0mm) {$1$};
          \draw[->] (x) edge [loop right] node {} (x);
      }
      \graphbox{$K$}{24mm}{0mm}{23mm}{12mm}{2mm}{-5mm}{
        \node[draw,circle] (x) at (0mm,0mm) {$1$};
      }
      \begin{scope}[opacity=1]        
      \graphbox{$R$}{48mm}{0mm}{30mm}{12mm}{-1mm}{-5mm}{
        \node[draw,circle] (x) at (0mm,0mm) {$1$};
        \node[draw,circle] (y) at (10mm,0mm) {$2$};
      }
      \end{scope}
    \end{tikzpicture}
  }
  \end{center}

  \begin{center}
  $\tau = $
  \begin{tikzpicture}
      \graphbox{$L$}{0mm}{-11mm}{32mm}{12mm}{2mm}{-8mm}{
      \node[draw,circle]  (x) at (-6mm,0mm) {$1$};
      \node[draw,circle] (y) at (6mm,0mm) {$2$};
      \draw[->]  (x) to (y);
      }
      \graphbox{$K$}{33mm}{-11mm}{32mm}{12mm}{2.5mm}{-8mm}{
          \node[draw,circle]  (x) at (-6mm,0mm) {$1$};
          \node[draw,circle]  (y) at (6mm,0mm) {$2$};
      }
      \begin{scope}[opacity=1]        
      \graphbox{$R$}{66mm}{-11mm}{40mm}{12mm}{-2mm}{-8mm}{
          \node[draw,circle]  (x) at (-6mm,0mm) {$1$};
          \node[draw,circle]  (y) at (6mm,0mm) {$2$};
          \node[draw,circle]  (z) at (17mm,0mm) {3};
      }
      \end{scope}
    \end{tikzpicture}
  \end{center}
%   \caption{}
%   \label{fig:subgraph_counting:ex_overbeek_5d3}
% \end{figure}
  Let $X$ be 
  \tikz[baseline=-0.5ex]{
      \node[draw,circle] (x) at (0,0) { };
      \node[draw,circle] (y) at (1,0) { };
      \draw[->] (x) to (y);
  }. Both rules are $X$-non-increasing, and $|\operatorname{Mono}(X,L)| = 1 > 0 = |\operatorname{Mono}(X,R)|$ holds for both rules. Thus, this rewriting system terminates by Theorem~\ref{subgraph_counting:thm:termination_grs}.
\end{example}

\begin{example} 
  \label{ex:endrullis2024_6d2}  
  Consider the rewriting rule from~\cite[Example 6.2]{endrullis2024generalized_arxiv_v2}:
  %  in Figure~\ref{fig:subgraph_counting:ex_endrullis2024_6d2}.

  % \begin{figure}[H]
  %   \centering
  \begin{center}
  \begin{tikzpicture}  
     \graphbox{$L$}{0mm}{-11mm}{32mm}{15mm}{2mm}{-8mm}{  
         \node[draw,circle]  (x) at (-6mm,0mm) {$1$};  
        \draw[->] (x) edge[loop above] node  {} (x);

         \node[draw,circle]  (y) at (6mm,0mm) {$2$};  
       }  
       \graphbox{$K$}{33mm}{-11mm}{32mm}{15mm}{2mm}{-8mm}{  
        \node[draw,circle]  (x) at (-6mm,0mm) {$1$};  
        \node[draw,circle]  (y) at (6mm,0mm) {$2$};  
       }  
       \graphbox{$R$}{66mm}{-11mm}{32mm}{15mm}{1mm}{-8mm}{  
        \node[draw,circle]  (x) at (-6mm,0mm) {$1$};  
        \node[draw,circle]  (y) at (6mm,0mm) {$2$};  
        \draw[->]  (x) to (y);  
       }    
 \end{tikzpicture}  
\end{center}
%  \caption{}
%  \label{fig:subgraph_counting:ex_endrullis2024_6d2}
% \end{figure}

 Our method can prove its termination with the ruler-graph \tikz[baseline=-0.5ex]{
  \node[draw,circle] (x) at (0,0) {};
  \draw[->] (x) edge[loop above] node  {} (x);
} of weight $1$.
\end{example} 

\begin{example}[Limitation]
  \label{ex:plump95_4d1}
 Consider the rewriting rule from~\cite[Example 4.1]{plump1995ontermination}:
%  in Figure~\ref{fig:subgraph_counting:ex_plump95_4d1}.
  % \begin{figure}[H]
  %   \centering
\begin{center}
      \resizebox{0.7\textwidth}{!}{
    $\rho = $  \begin{tikzpicture}[baseline=-20mm]
      \graphbox{$L$}{0mm}{0}{32mm}{40mm}{0}{0}{
          \node[draw,circle]  (n1) at (0,-6mm) {$1$};
          \node[draw,circle]   (n2) at (0mm,-26mm) {$2$};
          \node[draw,circle]   (n3) at (-10mm,-26mm) {3};
          \node[draw,circle]   (n4) at (10mm,-26mm) {4};
          \draw[->]  (n2) edge [loop below] node  {$a$} (n2);
          \draw[->]  (n3) edge [loop below] node  {$b$} (n3);
          \draw[->]  (n1) to node [right] {$f$} (n2) ;
          \draw[->]  (n1) to node [right] {$f$}  (n3);
          \draw[->]  (n1) to node [right] {$f$}  (n4);
        }
        \graphbox{$K$}{33mm}{0}{32mm}{40mm}{0}{-0}{
          \node[draw,circle]  (n1) at (0,-6mm) {$1$};
          \node[draw,circle]   (n2) at (0mm,-26mm) {$2$};
          \node[draw,circle]   (n3) at (-10mm,-26mm) {3};
          \node[draw,circle]   (n4) at (10mm,-26mm) {4};
          \draw[->]  (n2) edge [loop below] node  {$a$} (n2);
          \draw[->]  (n3) edge [loop below] node  {$b$} (n3);
        }
        \graphbox{$R$}{66mm}{0}{32mm}{40mm}{0}{-0}{
          \node[draw,circle]  (n1) at (0,-6mm) {$1$};
          \node[draw,circle]   (n2) at (0mm,-26mm) {$2$};
          \node[draw,circle]   (n3) at (-10mm,-26mm) {3};
          \node[draw,circle]   (n4) at (10mm,-26mm) {4};
          \draw[->]  (n2) edge [loop below] node  {$a$} (n2);
          \draw[->]  (n3) edge [loop below] node  {$b$} (n3);
          \draw[->]  (n1) edge [bend left] node [right] {$f$} (n4) ;
          \draw[->]  (n1) edge [bend right] node [left] {$f$}  (n4);
          \draw[->]  (n1) to node [right] {$f$}  (n4);
        }   
  \end{tikzpicture}
      }
  %   \caption{}
  %   \label{fig:subgraph_counting:ex_plump95_4d1}
  % \end{figure}
\end{center}
  To prove its termination with our method, a ruler-graph $X$ containing an edge labeled by $f$ must be used, because the number of occurrences of every graph containing only edges labeled by $a$ and $b$ does not change by any rewriting step using $\rho$. But in this case, condition~\eqref{def:non_increasing_rule_img_edges_distinct} of Definition~\ref{subgraph_counting:def:creates_more_x_on_the_left} cannot be satisfied.
\todo{to do todo: more details}
%   \begin{center}
%   \begin{tikzpicture} 
%       \graphbox{}{0}{0}{30mm}{20mm}{-10mm}{-10mm}{
%           \node[draw,circle]  (n1) at (0,0mm) {$1$};
%           % \node[draw,circle]   (n2) at (0mm,-26mm) {$2$};
%           % \node[draw,circle]   (n3) at (-10mm,-26mm) {3};
%           \node[draw,circle]   (n4) at (20mm,0mm) {4};
%           % \draw[->]  (n2) edge [loop below] node  {$a$} (n2);
%           % \draw[->]  (n3) edge [loop below] node  {$b$} (n3);
%           \draw[->]  (n1) edge [bend left] node [above] {$f$} (n4) ;
%           \draw[->]  (n1) edge [bend right] node [below] {$f$}  (n4);
%           \draw[->]  (n1) to node {$f$}  (n4);
%         }  
%   \end{tikzpicture}
% \end{center}
\end{example}

\begin{example}[Limitation]
    \label{ex:endrullis:d3:limitation}  
    Our method fails to prove termination of the following rewriting rule from \cite[Example D.3]{endrullis2024generalized_arxiv_v2}:
    % \begin{figure}[H]
    %   \centering
    \begin{center}
        \resizebox{0.7\textwidth}{!}{
    \begin{tikzpicture}  
       \graphbox{$L$}{0mm}{-11mm}{32mm}{15mm}{2mm}{-8mm}{  
           \node[draw,circle]  (x) at (-6mm,0mm) {$1$};  
           \node[draw,circle]  (y) at (6mm,0mm) {};  
         }  
         \graphbox{$K$}{33mm}{-11mm}{32mm}{15mm}{2mm}{-8mm}{  
           \node[draw,circle]  (x) at (-6mm,0mm) {$1$};  
         }  
         \graphbox{$R$}{66mm}{-11mm}{32mm}{15mm}{1mm}{-8mm}{  
          \node[draw,circle]  (x) at (-6mm,0mm) {$1$};  
          \node[draw,circle]  (y) at (6mm,0mm) {};  
          \draw[->]  (x) to (y);  
         }    
   \end{tikzpicture}
        }  
    %     \caption{}
    %     \label{fig:subgraph_counting:ex_endrullis_d3_limitation}
    % \end{figure}
  \end{center}
    %  in Figure~\ref{fig:subgraph_counting:ex_endrullis_d3_limitation}.
    % because our approach cannot take into account antipatterns(see~\cite[Remark 6.2]{overbeek2024termination_lmcs}).
    This happens because, for any system $\mathcal{R}$, rewriting step $G \Rightarrow_{\mathcal{R},\mathfrak{M}} H$, set of ruler-graphs $\mathbb{X}$, and weight function $s_\mathbb{X}$, we have $w_{s_\mathbb{X}}(G) \geq w_{s_\mathbb{X}}(H)$ because there is an injective graph homomorphism from $G$ to $H$.
  \end{example} 

\begin{example}[Limitation]
    \label{ex:overbeek:5d2:limitation}
    Consider a rewriting rule given by Overbeek and Endrullis in~\cite[Example 5.2]{overbeek2024termination_lmcs} shown below.
    % in Figure~\ref{fig:subgraph_counting:ex_overbeek_5d2}.
  % \begin{figure}[H]
  %   \centering
  \begin{center} 
    \begin{tikzpicture}
        \graphbox{$L$}{0mm}{-11mm}{32mm}{15mm}{2mm}{-8mm}{
            \node[draw,circle]  (x) at (-6mm,0mm) {$1$};
            \node[draw,circle]   (y) at (6mm,0mm) {$2$};
            \draw[->]  (x) to (y);
          }
          \graphbox{$K$}{33mm}{-11mm}{32mm}{15mm}{2mm}{-8mm}{
            \node[draw,circle]  (x) at (-6mm,0mm) {$1$};
            \node[draw,circle]  (y) at (6mm,0mm) {$2$};
            \draw[->]  (x) to (y);
          }
          \graphbox{$R$}{66mm}{-11mm}{32mm}{15mm}{1mm}{-8mm}{
            \node[draw,circle]  (x) at (0mm,0mm) {1 2};
            \draw[->]  (x) edge [loop right] (x);
          }  
    \end{tikzpicture} 
  %   \caption{}
  %   \label{fig:subgraph_counting:ex_overbeek_5d2}
  % \end{figure}
  \end{center}
    The termination of this rule is not in the scope of our method due to its non-node-injective right-hand side morphism. It can be proved using the method proposed by Overbeek et al.~\cite{overbeek2024termination_lmcs}.
\end{example}
 
\begin{example}[Limitation]
  \label{ex:overbeek:5d6_bis:limitation} 
  Consider the rewriting rule with monic matches from~\cite[Example 5.6]{overbeek2024termination_lmcs}:
  \begin{center}
    $\rho' = $\scalebox{0.9} { {
\begin{tikzpicture}[baseline=-10mm]
    \graphbox{$L$}{0mm}{0mm}{38mm}{20mm}{2mm}{-13mm}{
      \node [draw, circle] (x) at (-7mm,0mm) {1};
      \node [draw, circle] (y) at (3mm,0mm) {2 3};
      \draw[->] (x) edge[loop above] node  {$a$} (x);
      \draw[->] (y) edge [loop above] node {$a$} (y);
    }
    \graphbox{$K$}{42mm}{-0mm}{38mm}{20mm}{0mm}{-10mm}{
        \node [draw, circle] (x) at (-7mm,0mm) {1};
        \node [draw, circle] (y) at (0mm,0mm) {2};
        \node [draw, circle] (z) at (7mm,0mm) {3};    
    }
    \begin{scope}[opacity=1]        
    \graphbox{$R$}{85mm}{-0mm}{38mm}{20mm}{2mm}{-13mm}{
      \node [draw, circle] (x) at (-7mm,0mm) {1};
      \node [draw, circle] (y) at (0mm,0mm) {2};
      \node [draw, circle] (z) at (7mm,0mm) {3};
      \draw[->] (x) edge[loop above] node  {$b$} (x);
      \draw[->] (y) edge[loop above] node  {$b$} (y);
      \draw[->] (z) edge[loop above] node  {$b$} (z);
    }
    \end{scope}
    \node () at (40mm,-10mm) {$\leftarrow$};
    \node () at (83mm,-10mm) {$\rightarrow$};
\end{tikzpicture}
}

}
    % \caption{}
    % \label{fig:subgraph_counting:ex_overbeek_5d6_bis}
  \end{center} 
  Its termination is not in the scope of our method due to its non-node-injective left-hand side morphism. 
\end{example}

 
\section{Related work}
\label{sec:related_work} 
A comparative analysis of termination techniques for DPO graph rewriting systems, drawn from prior work~\cite{plump1995ontermination,plump2018modular,bruggink2014termination,bruggink2015proving,endrullis2024generalized_arxiv_v2,overbeek2024termination_lmcs}, is summarized in~\autoref{tab:comparison:subgraph_counting}. Our approach successfully proves termination for 14 of these systems. In the remainder of this section, we compare our method with some existing methods in more detail.


{% local group
  \setlength{\tabcolsep}{3pt}
  \renewcommand{\arraystretch}{1}
\begin{table}[H]
   \centering
  \small % Reduce font size
  \caption{Applicability of termination techniques to DPO rewriting examples.
  The symbol \ding{51} indicates termination can be proved by the technique,
  \ding{55} indicates it cannot be proved, and 
  $-$ denotes irrelevance or out-of-scope cases.
        }
 \label{tab:comparison:subgraph_counting}
% \setlength{\tabcolsep}{4pt} % Reduce horizontal padding
   \begin{NiceTabular}{ccccccccc}[vlines] % <-- 9 columns now (was 7)
    \Hline
                % \diagbox{\enskip \textbf{Examples}}{\textbf{Techniques}} 
    \Block{1-2}{\diagbox{\enskip \textbf{Examples}}{\textbf{Techniques}}} & 
    &
    \RowStyle{\rotate}
     \makecell{Forward closure \cite{plump1995ontermination}} % NEW column #1
    & \RowStyle{\rotate}
     \makecell{Modular criterion \cite{plump2018modular}} % NEW column #2
    & \RowStyle{\rotate}
     \makecell{Type graph \cite{bruggink2014termination}}  
    & \RowStyle{\rotate}
     \makecell{Type graph \cite{bruggink2015proving}} 
    & \RowStyle{\rotate}
     \makecell{Type graph \cite{endrullis2024generalized_arxiv_v2}} 
    & \RowStyle{\rotate}
     \makecell{Subgraph 
            counting \cite{overbeek2024termination_lmcs}} 
    & \RowStyle{\rotate}
     \makecell{Morphism
                Counting 
                \\(Chapter~\ref{chap:subgraph_counting})} \\
    \Hline
    \Hline
% ok examples
Chapter~\ref{chap:subgraph_counting} 
&~\autoref{subgraph_counting:ex_contrib_variant}
  & \ding{51} & \ding{55} & \ding{55} & \ding{55} & \ding{55} & \ding{55} & \ding{51} \\ \Hline
  
\cite{plump1995ontermination} &
\hyperref[ex:overbeek_5d8_plump1995_3d8_plump2018_3_overbeek_5d8]{Example 3.8}
             & \ding{51} & -- & -- & -- & -- &
          --
             & \ding{51}\\ 
\hline

\Block{2-1}{\cite{plump2018modular}} & \hyperref[ex:overbeek_5d8_plump1995_3d8_plump2018_3_overbeek_5d8]{Example 3} 
          & -- & \ding{51} &  -- & -- & -- & 
          --
          & \ding{51}\\ 

\Hline

%~\cite{plump2018modular} 
& Example 5 &  -- &  \ding{51} &   -- & -- & -- &  
            --
          & \ding{51}\\  
\Hline

\cite{bruggink2014termination} & \hyperref[ex:termination:grsaa]{Example 4 and 6}  
  & -- & -- & \ding{51} & -- & -- & 
            --
          & \ding{51} \\ \Hline

\Block{2-1}{\cite{bruggink2015proving}} & \hyperref[ex:termination:grsaa]{Example 2}  
  & -- & -- & -- & \ding{51} & -- & 
  -- &  \ding{51}\\ \Hline
  
%~\cite{bruggink2015proving}
 & Example 4 
  & -- & -- & -- & \ding{51} & -- & 
  --& \ding{51} \\ \Hline


% ----- from endrullis2024generalized_arxiv_v2 -----
\Block{3-1}{\cite{endrullis2024generalized_arxiv_v2}} & \hyperref[ex:endrullis2024_6d2]{Example 6.2}  
  & -- & -- & -- & -- & \ding{51} & -- & \ding{51}\\ \Hline

%~\cite{endrullis2024generalized_arxiv_v2}
 & \hyperref[ex_endrullis_6d3_endrullis_5d8]{Example 6.3}
  & -- & -- & -- & -- & \ding{51} &% \ding{55} 
  \ding{55} & \ding{51}\\ \Hline

%~\cite{endrullis2024generalized_arxiv_v2} 
& \hyperref[ex:overbeek_5d8_plump1995_3d8_plump2018_3_overbeek_5d8]{Example D.1}
  & -- & -- & -- & -- & \ding{51} & -- & \ding{51}\\ \Hline

  \Block{4-1}{\cite{overbeek2024termination_lmcs}} & \hyperref[ex:overbeek_5d3]{Example 5.3}
  & -- & -- & -- & -- & -- & \ding{51} & \ding{51}\\ \Hline

%~\cite{overbeek2024termination_lmcs} \hyperref[ex:overbeek_5d3]{Example 5.3 monic matches}
%     & -- & -- & -- & -- & -- & \ding{51} & \ding{51}\\ \Hline
%~\cite{overbeek2024termination_lmcs} 
& \hyperref[ex:overbeek_5d5]{Example 5.5} 
  & -- & -- & -- & -- & -- & \ding{51} & \ding{51}\\ \Hline

%~\cite{overbeek2024termination_lmcs} 
& \hyperref[ex:overbeek_5d6]{Example 5.6}
  & -- & -- & -- & -- & -- & \ding{51} & \ding{51} \\ \Hline

%~\cite{overbeek2024termination_lmcs} \hyperref[ex:overbeek_5d6]{Example 5.6 bis}
%     & -- & -- & -- & -- & -- & \ding{51} & -- \\ \Hline


%~\cite{overbeek2024termination_lmcs}  
& \hyperref[ex:overbeek_5d8_plump1995_3d8_plump2018_3_overbeek_5d8]{Example 5.8}
  & -- & -- & -- & -- & -- & \ding{51} & \ding{51}\\ \Hline

      % not supported examples  
     ~\cite{plump2018modular} &  \hyperref[ex:plump2018_ex6_endrullis_d4]{Example 6} &  -- & \ding{51} & -- & -- & -- & 
      --
          & -- \\
      \Hline

\Block{3-1}{\cite{endrullis2024generalized_arxiv_v2}}
 & Example 6.4  
      & -- & -- & -- & -- & \ding{51} & -- & -- \\ \Hline

%  ~\cite{endrullis2024generalized_arxiv_v2}
  &  Example 6.5  
      & -- & -- & -- & -- &  \ding{51} & -- & -- \\ \Hline

    %  ~\cite{endrullis2024generalized_arxiv_v2}
       &\hyperref[ex:plump2018_ex6_endrullis_d4]{Example D.4} 
      & -- & -- & -- & -- & \ding{51} & -- & --\\ \Hline

   % ----- from overbeek2024termination_lmcs -----
   \Block{3-1}{\cite{overbeek2024termination_lmcs}} & Example 5.2
      & -- & -- & -- & -- & -- & \ding{51} & -- \\ \Hline

    %  ~\cite{overbeek2024termination_lmcs} 
      & Example 5.7 
      & -- & -- & -- & -- & -- & \ding{51} & -- \\ \Hline
      
%  ~\cite{overbeek2024termination_lmcs} 
  & Example 5.9 
      & -- & -- & -- & -- & -- & \ding{51} & --\\ \Hline
 

    % not ok examples
   ~\cite{plump1995ontermination} & \hyperref[ex:plump95_4d1]{Example 4.1} & \ding{51} & -- & -- & -- & -- & 
              \ding{51}
              
              & \ding{55}\\ 
   \Hline
  ~\cite{plump2018modular} & \hyperref[ex:plump_ex4]{Example 4} &  -- &  \ding{51} &  -- & -- & -- & 
               --
               & \ding{55}\\ 
   \Hline

   \Block{3-1}{\cite{bruggink2014termination}} & Example 1 
   & -- & -- & \ding{51} & -- & -- & 
                 --
               &  \ding{55}\\ 
   \Hline

%   ~\cite{bruggink2014termination} 
   & Routing Protocol
       & -- & -- & \ding{51} & -- & -- & 
           --
           &  \ding{55}\\  
           \Hline
% \Block{2-1}{\cite{bruggink2014termination}}
 & \hyperref[ex:plump_ex4]{Example 5}
   & -- & -- & \ding{51} & -- & -- & -- &  \ding{55}\\ 
\Hline

\Block{2-1}{\cite{bruggink2015proving}} & \hyperref[ex:bruggink2015_ex5]{Example 5}
   & -- & -- & -- & \ding{51} & -- &  
   -- &  \ding{55}\\
   \Hline

%   ~\cite{bruggink2015proving} 
   & \hyperref[ex:bruggink2015_ex6_endrullis2024_d2]{Example 6} 
   & -- & -- & -- & \ding{51} & -- &  
   --&  \ding{55}\\  
   \Hline

   \Block{2-1}{\cite{endrullis2024generalized_arxiv_v2}} & \hyperref[ex:bruggink2015_ex6_endrullis2024_d2]{Example D.2} 
   & -- & -- & -- & -- & \ding{51} & -- & \ding{55}\\ 
   \Hline

%   ~\cite{endrullis2024generalized_arxiv_v2}
   & \hyperref[rem:d3_limitation]{Example D.3}
   & -- & -- & -- & -- & \ding{51} & \ding{55} & \ding{55}\\ \Hline

  \end{NiceTabular}
  % \vspace{2mm}
  % \caption{Applicability of termination techniques to DPO rewriting examples.
  %  The symbol \ding{51} indicates termination can be proved by the technique,
  %  \ding{55} indicates it cannot be proved, and 
  %  $-$ denotes irrelevance or out-of-scope cases.
  %        }
  % \label{tab:comparison}
  \end{table}
} 


The subgraph-counting method in~\cite{overbeek2024termination_lmcs} by Overbeek and Endrullis is designed for the PBPO+~\cite{overbeek2023graph, overbeek2023apbpotutorial}—a rewriting formalism capable of simulating left-injective DPO rewriting. It can prove termination for systems like~\cite[Examples 5.2, 5.7, 5.9]{overbeek2024termination_lmcs} and~\cite[Example 6]{plump2018modular}, which lie beyond the scope of our technique. Additionally, this method can be applied for many categories. 

It is very closely related to our work in the setting where they are both defined. Both methods weigh objects by summing weighted morphisms targeting them, and, for both methods, the key challenge lies in estimating weights for morphisms whose images partially overlap with the rewriting context, because morphisms fully embedded in the context are shared between host and resulting graphs, while those entirely within left- or right-hand-side graphs are trivial to quantify.  

The two methods employ distinct strategies to overcome this challenge. The subgraph-counting method relies on type-morphisms (see~\cite[page 9, remark 4.11, Lemma 4.23]{overbeek2024termination_lmcs}), whereas our approach requires injective mappings from (i) subgraph occurrences partially overlapping the rewriting context in the result graph to (ii) those in the host graph.

Both approaches have limitations and strengths. Their method suffers from discrete interface, as mentioned in~\cite[Example 5.5]{overbeek2024termination_lmcs}. For instance, it proves termination for~\cite[Example 5.5]{overbeek2024termination_lmcs} but fails for~\cite[Example 6.3]{endrullis2024generalized_arxiv_v2}, where the rules differ only by a discrete interface. 
Similarly, it fails to prove termination for~\autoref{subgraph_counting:ex_contrib_variant}, but if an edge labeled by ``s'' from node $1$ to node $3$ is added, then it succeeds. Our approach, however, can handle all the aforementioned cases. Nevertheless, their method can address~\cite[Example 4.1]{plump1995ontermination}, which remains beyond the scope of our technique. Finally, both approaches cannot handle~\cite[Example D.3]{endrullis2024generalized_arxiv_v2} (see~\cite[Remark 6.2]{overbeek2024termination_lmcs} and~\autoref{ex:endrullis:d3:limitation}).

The type graph method, which weighs an object by summing the weights of morphisms from the object to a type graph, was initially introduced by Zantema, K{\"o}nig and Bruggink~\cite{zantema2014termination} for cycle-rewriting systems. 
This method has since been generalized to edge-labeled multigraphs by Bruggink et al.~\cite{bruggink2014termination} for DPO rewriting with monic matches and injective rules, later extended to DPO rewriting in general by Bruggink et al.~\cite{bruggink2015proving}, and further adapted to more categories and different DPO variants by Endrullis et al.~\cite{endrullis2024generalized_arxiv_v2}. 
These methods are not directly comparable with our technique in general.

On the one hand, the termination of~\autoref{subgraph_counting:ex_contrib_variant} can be proved by our method but not by the type graph methods due to the exsitence of a surjection from the output graph to the input graph as explained in~\cite[Example D.4]{endrullis2024generalized_arxiv_v2}. On the other hand, our method cannot prove the termination of~\cite[Example 1, 5 and Ad-hoc Routing Protocol]{bruggink2014termination}, nor~\cite[Example 5, 6]{bruggink2015proving}, nor~\cite[Examples D2 and D3]{endrullis2024generalized_arxiv_v2}.

Plump~\cite{plump1995ontermination} introduced a necessary and sufficient termination condition for left-injective DPO hypergraph rewriting via forward closure, though verifying this condition is undecidable. 
While this method proves termination of~\autoref{subgraph_counting:ex_contrib_variant}, our approach succeeds on~\cite[Example 3.8]{plump1995ontermination} but fails in proving termination of~\cite[Example 4.1]{plump1995ontermination}. 

Plump~\cite{plump2018modular} later proposed a modular critical pair-based strategy for left-injective DPO hypergraph rewriting with monic matches. 
Our method complements this: while modularity reduces global complexity, each subsystem requires individual termination proofs. For example, the measure based on the indegree proposed in~\cite{plump2018modular} cannot prove the termination of~\autoref{subgraph_counting:ex_contrib_variant} due to the additional loops. Specifically, if the context graph is identical to the interface graph, then we have $3^k < 2^k+2^k+3^k$ for all $k \in \mathbb{N}$. In contrast, our method succeeds.
Additionally, ur method can proe termination for~\cite[Examples 1 and 5]{plump2018modular} but not~\cite[Example 4]{plump2018modular}, and non-injective rules (e.g.,~\cite[Example 6]{plump2018modular}) fall outside our scope.

In~\cite{levendovszky2007termination}, a terminaiton criterion for DPO rewriting with monic matches, injective rules and negative application conditions on finite typed attributed graphs is proposed. It is based on the fact that if the application of every infinite sequence of rules requires the initial graph to be infinite, then the system terminates. This technique is theoretically very interesting, but it is hard to automatically check the termination condition.

~\cite{bottoni2005termination} presents a termination criterion for DPO and SPO rewriting on high-level replacement units. The high-level replacement units that they consider are rewriting systems with very restrictive external control mechanisms. The method relies on a mesearing function satisfying a very strong constrainte, and the instance of such a mesearing function proposed is node and edge counting, which is subsumed by our method in the setting of DPO rewriting with injective rules.

~\cite{bottoni2010atermination} presents a criterion for termination of DPO rewriting with monic matches, injective rules and negative application conditions, based on the construction of a labeled transition system.
whose states represent overlaps between the negative application condition and the right hand side that can give rise to cycles.

Levendovszky et al.~\cite{levendovszky2007termination} propose a termination criterion for DPO rewriting (monic matches, injective rules, negative application condition), though automated verification is hard as explained in~\cite[\textsection 6]{levendovszky2007termination}. 

Bottoni et al.~\cite{bottoni2005termination} propose a termination criterion for DPO/SPO rewriting on high-level replacement units. Their method imposes a strongly constrained measuring function and the only concretes measuring function proposed are node-counting and edge-counting.

Bottoni et al.~\cite{bottoni2010atermination} presents a criterion for termination of DPO rewriting with monic matches, injective rules and negative application conditions, based on the construction of a labeled transition system. 



\section{Conclusion and future work}
\label{sec:conclusion} 
We have presented a machine-checkable sufficient condition for relative termination of DPO systems with injective rules on edge-labeled multigraphs.  
This method resolves cases where prior interpretation-based techniques \cite{zantema2014termination,bruggink2014termination,bruggink2015proving,endrullis2024generalized_arxiv_v2,
overbeek2024termination_lmcs} fail (e.g.,Example~\ref{subgraph_counting:ex_contrib_variant}) and, unlike the subgraph-counting approach in \cite{overbeek2024termination_lmcs}, avoids limitations arising from interfaces with no edges.
The method has been implemented in an OCaml tool. 

An interesting direction for future work is to extend our method to DPO systems with negative application conditions and to address the limitation highlighted in Example~\ref{ex:endrullis:d3:limitation}.

\part{Termination Tool: LyonParallel}

\part{Conclusion and Futur Works}
\section{Conclusion}
\begin{itemize}
    \item[bruggink2014] the paper "high-level replacement units and their termination properties 2005" considers high-level replacement units (HLRU), which are transformation systems with externam control expressions. The paper introduces a general framework for proving termination of such HLRUs, but the only concrete termination criteria considered are node and edge counting, which are subsumed by the weighted type graph method.
    \item[bruggink2014] in "termination criteria for model transformation 2005" layered graph transformation systems are considered, which are graph transformation systems where interleaving creation and deletion of edges with the same label is prohibited and creation of nodes is bounded. The paper shows such graph transformation systems are terminating.
    \item[bruggink2014] The paper "Termination analysis of model transformations by petri nets. 2006" simulates a graph transformation system by a Petri-net. Thus, the presence of edges with certain labels and the causal relationships between them are modeled, but no other structural properties of the graph. The paper uses typed graph transformation systems; thus, a type graph is used but, unlike in our weighted type graph method, it is fixed by the graph transformation system. Finally, [3] was one of the inspirations for this paper. 
    \item[bruggink2014] "Towards a systematic method for proving termination of graph transformation systems" is one of the inspirations for the paper of bruggink2014. its termination argument is subsumed by the weighted type graph technique.
\end{itemize}

\section{Future Works}

\begin{itemize}
    \item[bruggink2014]Although all theorems have been stated and proved for (binary) multigraphs, a generalization to hypergraphs would be trivial. On the other hand, transferring the results to other graph transformation formalisms is harder. For example, in the single pushout approach, the graph transformation system corresponding to the one of Ex. 1 is non-terminating, so the result of Ex. 2 (in which it is proved that this system is terminating) shows that the weighted type graph technique cannot be transferred one-to-one to single pushout graph transformation. It is left as future research to find similar arguments for the single pushout approach and other formalisms.
    
    Another direction for further research is to allow for graph transformation systems with negative application conditions or more general application conditions. Note, however, that the implicit negative application condition of double pushout graph transformation, the dangling edge condition, can in some cases already be handled (see Ex. 2). 
    
    Finally, for interesting real-world applications, it would be interesting to generalize the technique to more expressive methods of specifying the initial graph languages, so that we can, for example, restrict to trees or rings of arbitrary size (both graph languages cannot be expressed by a type graph). 
\end{itemize}

\begin{idea}
    example majority A or B\\
    implementation of the algorithms proposed by plump\\
    propose a reinforced version of the algorithms\\
\end{idea}


\part{Appendix}
\chapter{Appendix for Subgraph Counting}
\section*{Pullback, Adhesive Category and VK-square}
\begin{proposition}[\text{\cite[Lemma 13]{lack2004adhesive}}]
    \label{prop:pb_eq_po}
    In the category \textbf{Graph}, pushouts along monomorphisms are also pullbacks. 
\end{proposition}

\begin{definition}[Van Kampen Square \text{\cite[Definition 2]{lack2004adhesive}}]
    \label{def:vk_square}
    A \textit{van Kampen (VK) square} (i) is a pushout square which satisfies the following condition: given a commutative cube (ii) of which (i) forms the bottom face and the back faces are pullbacks, the front faces are pullbacks if and only if the top face is a pushout.

    
    \begin{enumerate}
        \item[(i)]
        \begin{tikzcd}
            & C \arrow[ld, "m"'] \arrow[rd, "f"]  &\\
            A \arrow[rd,"g"'] & & B \arrow[ld,"n"]\\
            &D &  
        \end{tikzcd}
    
        \item[(ii)]
        \begin{tikzcd}
            A' \arrow[dd, "a" description, near start] \arrow[dr, "g'" description] & & C' \arrow[ll, "m'" description, near start]  \arrow[dd, "c" description, near start] \arrow[dr, "f'" description] & \\
            & D'  \arrow[dd, "d" description, near start] & & B' \arrow[ll, "n'" description, near start] \arrow[dd, "b" description, near start] \\
            A  \arrow[dr, "g" description] & & C \arrow[dr, "f" description] \arrow[ll, "m" description, near start] & \\
            & D  & & B \arrow[ll, "n" description]
        \end{tikzcd}
    \end{enumerate}
\end{definition} 

\begin{definition}[Adhesive category\text{\cite[Definition 5]{lack2004adhesive}}] 
    \label{def:adhesive_cat}
    A category $\mathcal{C}$ is said to be \textbf{adhesive} if 
    \begin{itemize}
        \item $\mathcal{C}$ has pushouts along  monomorphisms, and
        \item $\mathcal{C}$ has pullbacks, and
        \item pushouts along monomorphisms are VK squares.
    \end{itemize}
\end{definition}

\begin{proposition}[\text{\cite[Proposition 8]{lack2004adhesive}}]
    \label{prop:graph_adhesive}
    The category \textbf{Graph} of directed labeled multigraphs is adhesive.
\end{proposition}
% \newpage
\section*{Auxiliary lemmas}
\begin{lemma}
    \label{lem:b_c_same_img_exist_a}
    Consider the pushout square in $\mathbf{Graph}$ of form
    \begin{center} 
        \begin{tikzpicture}[node distance=12mm]
          \node (A) at (0,0) {$A$};
          \node (B) at (2,0) {$B$}; 
          \node (C) at (0,-2) {$C$}; 
          \node (D) at (2,-2) {$D$}; 
          \node () [at=($(A)!0.5!(D)$)] {\normalfont PO};
          \draw [>->] (A) to (B);
          \draw [>->] (A) to (C);
          \draw [>->] (B) to (D); 
          \draw [>->] (C) to (D);
        \end{tikzpicture}
    \end{center} 
    The following statements hold:
    \begin{enumerate}
        \item  If $ b\in B $ and $ c \mathop{\in} C $ such that $ h_{BD}(b) \mathop{=} h_{CD}(c) $ then there is $ a \mathop{\in} A $ such that $ h_{AB}(a) \mathop{=} b $ and $ h_{AC}(a)= c $;
        \item  For all \( b \mathop{\in} B \mathop{\setminus} h_{AB}(A) \) and \( c \mathop{\in} C \), we have \( \beta'(b) \mathop{\neq} \alpha'(c) \).
    \end{enumerate}
\end{lemma}
\begin{proof}
     Let \( b \mathop{\in} B \) and \( c \mathop{\in} C \) such that \( h_{BD}(b) \mathop{=} h_{CD}(c) \). 
     
     There are two cases: if $b$ and $c$ are two nodes then let \( S \) be the singleton graph with a unique node \( u \), otherwise\textemdash$b$ and $c$ must be two edges with the same label $l$\textemdash let \( S \) be the graph with two distinct nodes and a unique edge \( u \) labeled by $l$.
     
     Define morphisms \( h_{SB}: S \mathop{\to} B \) and \( h_{SC}: S \mathop{\to} C \) such that \( h_{SB}(u) \mathop{=} b \) and \( h_{SC}(u) \mathop{=} c \). We have \( h_{SB} \mathop{\star} h_{BD} \mathop{=} h_{SC} \mathop{\star} h_{CD} \). The pushout square \( ACDB \) along monomorphisms is also a pullback square by Proposition~\ref{prop:pb_eq_po}, and the universal property of the pullback guarantees the existence of a unique morphism \( h_{SA}: S \mathop{\to} A \) such that \( h_{SA} \mathop{\star} h_{AB} \mathop{=} h_{SB} \) and \( h_{SA} \mathop{\star} h_{AC} \mathop{=} h_{SC} \). Let \( a \mathop{=} h_{SA}(u) \). We have \( a \mathop{\in} A \), and by composition we obtain: 
     \begin{itemize}
        \item \(h_{AB}(a) \mathop{=} (h_{SA} \mathop{\star} h_{AB})(u) \mathop{=} h_{SB}(u) \mathop{=} b \), and
        \item \(h_{AC}(a) \mathop{=} (h_{SA} \mathop{\star} h_{AC})(u) \mathop{=} h_{SC}(u) \mathop{=} c\).
     \end{itemize} 

    The second assertion is a direct consequentce of the first. Suppose, for contradiction, there are \( b \mathop{\in} B \mathop{\setminus} h_{AB}(A) \) and \( c \mathop{\in} C \) such that \( h_{BD}(b) \mathop{=} h_{CD}(c) \). Since the first assertion holds, there must exist \( a \mathop{\in} A \) with \( h_{AB}(a) \mathop{=} b \) and \( h_{AC}(a) \mathop{=} c \). However, this contradicts the assumption that \( b \notin h_{AB}(A) \).
    \qed
\end{proof}

\begin{lemma}
    \label{kpcpck_pullback}
    Consider the following commutative diagram in \textbf{Graph}.
    \begin{center}
        % \resizebox{6cm}{!}{ 
            \begin{tikzpicture}[scale=1.2]
                \node (k) at (0,0) {K};
                \node (r) at (4,0) {R};
                \node (c) at (0,-3) {C};
                \node (h) at (4,-3) {H};
                \draw[<-<]  (r) -- (k) node [midway,above] {};
                \draw[>->] (c) -- (h) node [midway, below] {};
                \draw[>->] (r) -- (h) node[midway, left] {};
                \draw[>->] (k) -- (c) node[midway, left] {};
                \node (k') at (1,-1) {K'};
                \node (r') at (2,-1) {R'};
                \node (c') at (1,-2) {C'};
                \node () at (1.5,-1.5) {$PB$};
                \node () at (3,-1.5) {$PB$};
                \node () at (1.5,-2.5) {$PB$};
                \node (x) at (2,-2) {X};
                \draw [>->] (c') -- (x);
                \draw [>->] (r') -- (x);
                \draw [>->] (k') -- (r');
                \draw [>->] (k') -- (c');
                \draw [>->] (c') -- (c);
                \draw[>->] (r') -- (r);
                \draw[>->] (x) -- (h);
                \draw[->] (k') -- (k);
            \end{tikzpicture}
        % }
        \end{center}  
        The squares $K'KCC'$ and $K'KRR'$ are pulllback square.
\end{lemma}
\begin{proof}
    Let $K \mathop{\leftarrow} A \mathop{\rightarrow} C'$ be a span such that 
    \begin{flalign}
        h_{AK} \mathop{\star} h_{KC} \mathop{=} h_{AC'} \mathop{\star} h_{C'C}. \label{hakhkchacp}
    \end{flalign}
    We have the following commutative diagram 

    \begin{center}
        % \resizebox{0.3\textwidth}{!}{
            \begin{tikzpicture}[scale=1.2]
                \node (k) at (0,0) {K};
                \node (r) at (4,0) {R};
                \node (c) at (0,-3) {C};
                \node (h) at (4,-3) {H};
                \node (h) at (4,-3) {H};
                \node (a) at (-1,0) {A};
                \draw[<-<]  (r) -- (k) node [midway,above] {};
                \draw[>->] (c) -- (h) node [midway, below] {};
                \draw[>->] (r) -- (h) node[midway, left] {};
                \draw[>->] (k) -- (c) node[midway, left] {};
                \node (k') at (1,-1) {K'};
                \node (r') at (2,-1) {R'};
                \node (c') at (1,-2) {C'};
                \node () at (1.5,-1.5) {$PB$};
                \node () at (3,-1.5) {$PB$};
                \node () at (1.5,-2.5) {$PB$};
                \node (x) at (2,-2) {X};
                \draw [>->] (c') -- (x);
                \draw [>->] (r') -- (x);
                \draw [>->] (k') -- (r');
                \draw [>->] (k') -- (c');
                \draw [>->] (c') -- (c);
                \draw[>->] (r') -- (r);
                \draw[>->] (x) -- (h);
                \draw[->] (k') -- (k);
                \draw[->,red] (a) -- (c');
                \draw[->,red] (a) -- (k);
            \end{tikzpicture}
        % }
        \end{center} 
    We show first the existence of a morphism $h_{AK'} \mathop{\colon} A \mathop{\to} K'$ such that the following statements hold:
    \begin{flalign*}
        h_{AK'} \mathop{\star} h_{K'C'} \mathop{=} h_{AC'},  
        \\
        h_{AK'} \mathop{\star} h_{K'K} \mathop{=} h_{AK}. 
    \end{flalign*}
    From Equation~\eqref{hakhkchacp}, we obtain 
    \begin{flalign}
        h_{AK} \mathop{\star} h_{KC} \mathop{\star} h_{CH}= h_{AC'} \mathop{\star} h_{C'C} \mathop{\star} h_{CH}. \label{hakhkchchhacp}
    \end{flalign}
    From Equation~\eqref{hakhkchchhacp} and commutativity of $C'CHX$ and $KCHR$, we have $h_{KC} \mathop{\star} h_{CH} \mathop{=} h_{KR} \mathop{\star} h_{RH}$ and $h_{C'C} \mathop{\star} h_{CH} \mathop{=} h_{C'X} \mathop{\star} h_{XH}$ and deduce  
     \begin{flalign}
        h_{AK} \mathop{\star} h_{KR} \mathop{\star} h_{RH}= h_{AC'} \mathop{\star} h_{C'X} \mathop{\star} h_{XH}. \label{hakhkrhrhhacp}
     \end{flalign}
    Therefore, the following commutative diagram holds 

    \begin{center}
        % \resizebox{0.3\textwidth}{!}{
            \begin{tikzpicture}[scale=1.2]
                \node (k) at (0,0) {K};
                \node (r) at (4,0) {R};
                \node (h) at (4,-3) {H};
                \node (h) at (4,-3) {H};
                \node (a) at (-1,0) {A};
                \draw[<-<]  (r) -- (k) node [midway,above] {};
                \draw[>->] (r) -- (h) node[midway, left] {};
                \node (r') at (2,-1) {R'};
                \node (c') at (1,-2) {C'};
                \node () at (3,-1.5) {$PB$};
                \node (x) at (2,-2) {X};
                \draw [>->] (c') -- (x);
                \draw [>->] (r') -- (x);
                \draw[>->] (r') -- (r);
                \draw[>->] (x) -- (h);
                \draw[->] (a) -- (c');
                \draw[->] (a) -- (k);
            \end{tikzpicture}
        % }
        \end{center} 
    Since $R'XHR$ is a pullback square, by the universal property Definition~\ref{def:cat:pb}, there is a unique morphism $h_{AR'} \mathop{\colon} A \mathop{\to} R'$ such that 
    \begin{flalign}
        h_{AR'} \mathop{\star} h_{R'R} \mathop{=} h_{AK} \mathop{\star} h_{KR},  \label{harprprakkr}
        \\
        h_{AR'} \mathop{\star} h_{R'X} \mathop{=} h_{AC'} \mathop{\star} h_{C'X}. \label{harprpxacpcpx}
    \end{flalign}
    We have thus the following commutative diagram 

        \begin{center}
            % \resizebox{0.2\textwidth}{!}{
                \begin{tikzpicture}[scale=1.2]
                    \node (a) at (-1,0) {A};
                    \node (k') at (1,-1) {K'};
                    \node (r') at (2,-1) {R'};
                    \node (c') at (1,-2) {C'};
                    \node () at (1.5,-1.5) {$PB$};
                    \node (x) at (2,-2) {X};
                    \draw [>->] (c') -- (x);
                    \draw [>->] (r') -- (x);
                    \draw [>->] (k') -- (r');
                    \draw [>->] (k') -- (c');
                    \draw[->] (a) -- (c');
                    \draw[->] (a) -- (r');
                \end{tikzpicture}
            % }
        \end{center} 
    Since $K'C'XR'$ is a pullback square, there is a unique morphism $h_{AK'} \mathop{\colon} A \mathop{\to} K'$ such that 
    \begin{flalign}
        h_{AK'} \mathop{\star} h_{K'C'} \mathop{=} h_{AC'},  \label{hakphkpcphacp}
        \\
        h_{AK'} \mathop{\star} h_{K'R'} \mathop{=} h_{AR'}. \label{hakpkprparp}
    \end{flalign}
    We have $h_{AK'} \mathop{\star} h_{K'C'} \mathop{=} h_{AC'}$ by Equation~\eqref{hakphkpcphacp}, and $h_{AK'} \mathop{\star} h_{K'K} \mathop{=} h_{AK}$ because $h_{KR}$ is injective and the following holds:
    \begin{flalign*}
        (h_{AK'} \mathop{\star} h_{K'K}) \mathop{\star} h_{KR} &= h_{AK} \mathop{\star} h_{K'R'} \mathop{\star} h_{R'R} \hspace{0.5cm} &\text{by commutativity of $K'KRR'$}\\
        &= h_{AR'} \mathop{\star} h_{R'R} &\text{by Equation~\eqref{hakpkprparp}} \\
        &= h_{AK} \mathop{\star} h_{KR}. &\text{by Equation~\eqref{harprprakkr}}
    \end{flalign*}
    Now, we show the unicity of such a morphism from $A$ to $K'$. 
    Let $h_{AK'}' \mathop{\colon} A \mathop{\to} K'$ be a morphism  such that 
    \begin{flalign*}
        h_{AK'}' \mathop{\star} h_{K'C'} \mathop{=} h_{AC'},  
        \\
        h_{AK'}' \mathop{\star} h_{K'K} \mathop{=} h_{AK}.
    \end{flalign*}
    From $h_{AK'} \mathop{\star} h_{K'C'} \mathop{=} h_{AC'}$, $h_{AK'}' \mathop{\star} h_{K'C'} \mathop{=} h_{AC'}$ and injectivity of $h_{K'C'}$, we deduce $h_{AK'} \mathop{=} h_{AK'}'$.

    \qed
\end{proof}


\begin{lemma}
    \label{kpcpxrp_po}
    Consider the following commutative diagram in \textbf{Graph} where $KCHR$ is a pushout.
    \begin{center}
        % \resizebox{6cm}{!}{
            \begin{tikzpicture}
                \node (k) at (0,0) {K};
                \node (r) at (4,0) {R};
                \node (c) at (0,-3) {C};
                \node (h) at (4,-3) {H};
                % \node (rb) at ($\scl*(1.5,-0.5)$) {$R_X$};
                % \node (h') at ($\scl*(1.5,-1.5)$) {$H'$};
                % \draw[>->]  (rb) -- (h') node [midway,above] {};
                % \draw[>->]  (c) -- (h') node [midway,above] {};
                \draw[<-<]  (r) -- (k) node [midway,above] {};
                \draw[>->] (c) -- (h) node [midway, below] {};
                \draw[>->] (r) -- (h) node[midway, left] {};
                \draw[>->] (k) -- (c) node[midway, left] {};
                % \draw[->] (rb) to (l);
                % \draw[<-<] (rb) to (k);
                \node (k') at (1,-1) {K'};
                \node (r') at (2,-1) {R'};
                \node (c') at (1,-2) {C'};
                % \node () at (1.5,-1.5) {$PB$};
                \node () at (3,-1.5) {$PB$};
                \node () at (1.5,-2.5) {$PB$};
                \node () at (1.5,-0.5) {$PB$};
                \node () at (0.5,-1.5) {$PB$};
                \node (x) at (2,-2) {X};
                % \draw [->] (x) -- (h') node[midway] {!};
                \draw [>->] (c') -- (x);
                \draw [>->] (r') -- (x);
                \draw [>->] (k') -- (r');
                \draw [>->] (k') -- (c');
                \draw [>->] (c') -- (c);
                % \draw [->] (r') -- (rb);
                % \draw [->] (r') -- (rb);
                % \draw [->] (h') -- (h);
                \draw[>->] (r') -- (r);
                \draw[>->] (x) -- (h) node[midway,right] {$h$};
                % \node (rb) at ($\scl*(\sclx*1.5,-0.5)$) {$R_X$};
                % \node (h') at ($\scl*(\sclx*1.5,-1.2)$) {$H'$};
                % \draw[>->]  (rb) -- (h') node [midway,above] {};
                % \draw[>->]  (c) -- (h') node [midway,above] {};
                % \draw[>->]  (rb) -- (h') node [midway,above] {};
                % \draw[>->] (rb) to (r);
                % \draw[<-<] (rb) to (k);
                \draw[>->] (k') -- (k) ;
            \end{tikzpicture}
        % }
        \end{center} 
        The square $K'C'XR'$ is a pushout.
\end{lemma}
\begin{proof}
    \textbf{Graph} is adhesive by Proposition~\ref{prop:graph_adhesive}. By Definition~\ref{def:adhesive_cat}, pushouts along monomorphisms are VK squares in \textbf{Graph}. The square $KCHR$ is a pushout along monomorphisms in \textbf{Graph}. Hence $KCHR$ is VK square. By Definition~\ref{def:vk_square}, the square $K'C'XR'$ is pushout if the squares $C'CHX$ and $R'XHR$ are pullbacks. By assumption, the squares $C'CHX$ and $R'XHR$ are pullbacks. Therefore, $K'C'XR'$ is pushout.
\end{proof} 

\begin{lemma}
    \label{lem:xinXcpinCrpinR}
        Consider the following pushout square in \textbf{Graph}.
    \begin{center}
        \begin{tikzpicture}
            \node (k') at (0,0) {K'};
            \node (r') at (-2,0) {R'};
            \node (c') at (0,-2) {C'};
            \node (x) at (-2,-2) {X};
            \node () at (-1,-1) {$PO$};
            \draw [>->] (c') -- (x);
            \draw [>->] (r') -- (x);
            \draw [>->] (k') -- (r');
            \draw [>->] (k') -- (c');
        \end{tikzpicture}
    \end{center} 
     Let $x$ be a node (resp. an edge) in $X$. There is either a node (resp. an edge) $r'$ in $R'$ such that $h_{R'X}(r') \mathop{=} x$ or a node (resp. an edge) $c'$ in $C'$ such that $h_{C'X}(c') \mathop{=} x$.
\end{lemma}
\begin{proof}
    Suppose, by contradiction, that there is neither a node (resp. an edge) $r'$ in $R'$ such that $h_{R'X}(r') \mathop{=} x$ nor a node (resp. an edge) $c'$ in $C'$ such that $h_{C'X}(c') \mathop{=} x$.

    Let $G$ be a graph with two distinct nodes (resp. edges) $y$ and $y'$.
    By the universal property of pushout square Definition~\ref{def:cat:po}, there is a unique morphism $h_{XG}:X \rightarrowtail G$ such that $h_{C'G} \mathop{=} h_{C'X} \mathop{\star} h_{XG}$ and  $h_{R'G} \mathop{=} h_{R'X} \mathop{\star} h_{XG}$.
    \begin{center}
        \begin{tikzpicture}
            \node (g) at (-4,-4) {$G$};
            \node (k') at (-1,-1) {$K'$};
            \node (r') at (-2,-1) {$R'$};
            \node (c') at (-1,-2) {$C'$};
            \node (x) at (-2,-2) {$X$};
            \node () at (-1.5,-1.5) {$PO$};
            \draw [>->] (c') -- (x);
            \draw [>->] (r') -- (x);
            \draw [>->] (k') -- (r');
            \draw [>->] (k') -- (c');
            \draw [>->] (c') -- (g);
            \draw [>->] (r') -- (g);
            \draw [->] (x) -- (g) node[midway] {!};
        \end{tikzpicture}
    \end{center}
    However, if we define two morphisms $f$ and $g$ by 
    % \begin{flalign*}
    %     f(z) &= h_{XG}(z) \text{if $z \mathop{\neq} x$}\\
    %     f(x) &= y \\
    %     g(z) &= h_{XG}(z) \text{if $z \mathop{\neq} x$}\\
    %     g(x) &= y' \\
    % \end{flalign*}
    \[
    f(z)=
    \begin{cases}
    h_{XG}(z), & \text{if $z \mathop{\neq} x$},\\
    y,         & \text{if $z \mathop{=} x$},
    \end{cases}
    \]

    \[
    g(z)=
    \begin{cases}
    h_{XG}(z), & \text{if $z \mathop{\neq} x$},\\
    y',        & \text{if $z \mathop{=} x$}.
    \end{cases}
    \]
    we have $f \mathop{\neq} g$ and 
    \begin{flalign*}
        h_{C'G} \mathop{=} h_{C'X} \mathop{\star} f,\\
        h_{R'G} \mathop{=} h_{R'X} \mathop{\star} f,\\
        h_{C'G} \mathop{=} h_{C'X} \mathop{\star} g,\\
        h_{R'G} \mathop{=} h_{R'X} \mathop{\star} g.
    \end{flalign*}
    which is a contradiction.
\end{proof}

\begin{lemma}
    \label{lem:g_monic}
    Consider the following commutative diagram in \textbf{Graph}.
    \begin{center}
        \begin{tikzpicture}
            \node (k) at (0,0) {K};
            \node (c) at (0,-4) {C};
            \node (l) at (-4,0) {$L$};
            \node (g) at (-4,-4) {$G$};
            \draw[>->] (l) -- (g) node [midway,above] {};
            \draw[>->] (c) -- (g) node [midway,above] {};
            \draw[>->] (k) -- (c) node[midway, left] {};
            \draw[<-<] (l) to (k);
            \node (k') at (-1,-1) {$K'$};
            \node (r') at (-2,-1) {$R'$};
            \node (c') at (-1,-2) {$C'$};
            \node (x) at (-2,-2) {$X$};
            \node () at (-1.5,-1.5) {$PO$};
            \node () at (-0.5,-1.5) {$PB$};
            \draw [>->] (c') -- (x);
            \draw [>->] (r') -- (x);
            \draw [>->] (k') -- (r');
            \draw [>->] (k') -- (c');
            \draw [>->] (c') -- (c);
            \draw [>->] (r') -- (l);
            \draw [>->] (k') -- (k);
        \end{tikzpicture}
    \end{center}
    there is a unique monomorphism $h_{XG}:X \rightarrowtail G$ such that
    $C'XGC$ and $R'XGL$ are commutative squares.
    %  $h_{C'C} \mathop{\star} h_{CG} \mathop{=} h_{C'X} \mathop{\star} h_{XG}$ and  $h_{R'L} \mathop{\star} h_{LG} \mathop{=} h_{R'X} \mathop{\star} h_{XG}$.
\end{lemma}
\begin{proof}
    there is a unique morphism $h_{XG}:X \mathop{\to} G$, by Definition~\ref{def:cat:po}, because $K'C'XR'$ is a pushout\todo{$K'C'XR'$ PO} and 
    $h_{K'C'} \mathop{\star} h_{C'C} \mathop{\star} h_{CG}
    =h_{K'R'} \mathop{\star} h_{R'L} \mathop{\star} h_{LG}$.
    % the following hold
    % \begin{flalign*}
    %     &h_{K'C'} \mathop{\star} h_{C'C} \mathop{\star} h_{CG} &\\ 
    %    =&h_{K'K} \mathop{\star} h_{KC} \mathop{\star} h_{CG} &\hspace{2cm}\text{by commutativity of $K'C'CK$}\\
    %    =&h_{K'K} \mathop{\star} h_{KL} \mathop{\star} h_{LG} &\text{by the commutativity of $KCGL$} \\
    %    =&h_{K'R'} \mathop{\star} h_{R'L} \mathop{\star} h_{LG} &\text{by the commutativity of $K'KLR'$}
    % \end{flalign*}
    Therefore, the following commutative diagram holds
    \begin{center}
        \begin{tikzpicture}
            \node (k) at (0,0) {$K$};
            \node (c) at (0,-4) {$C$};
            \node (l) at (-4,0) {$L$};
            \node (g) at (-4,-4) {$G$};
            \draw[>->] (l) -- (g) node [midway,above] {};
            \draw[>->] (c) -- (g) node [midway,above] {};
            \draw[>->] (k) -- (c) node[midway, left] {};
            \draw[<-<] (l) to (k);
            \node (k') at (-1,-1) {$K'$};
            \node (r') at (-2,-1) {$R'$};
            \node (c') at (-1,-2) {$C'$};
            \node (x) at (-2,-2) {$X$};
            \node () at (-1.5,-1.5) {$PO$};
            \node () at (-0.5,-1.5) {$PB$};
            \draw [>->] (c') -- (x);
            \draw [>->] (r') -- (x);
            \draw [>->] (k') -- (r');
            \draw [>->] (k') -- (c');
            \draw [>->] (c') -- (c);
            \draw [>->] (r') -- (l);
            \draw [>->] (k') -- (k);
            \draw [->] (x) -- (g) node[midway,above] {!};
        \end{tikzpicture}
    \end{center}
    It remains to show that $h_{XG}$ is injective.
    
    Let $x_1,x_2$ be two elements in the graph $X$ such that $x_1 \mathop{\neq} x_2$. Since $K'R'XC'$ is a pushout\todo{$K'R'XC'$ P.O.} square, by Lemma~\ref{lem:xinXcpinCrpinR}, there are 4 possible cases:
    \begin{enumerate}
        \item there are $c_1, c_2 \mathop{\in} C'$ such that $h_{C'X}(c_1) \mathop{=} x_1$ and $h_{C'X}(c_2) \mathop{=} x_2$;
        \item there are $r_1, r_2 \mathop{\in} R'$ such that $h_{R'X}(r_1) \mathop{=} x_1$ and $h_{R'X}(r_2) \mathop{=} x_2$;
        \item there are $r' \mathop{\in} R'$ and $c' \mathop{\in} C'$ such that $h_{C'X}(c') \mathop{=} x_1$ and $h_{R'X}(r') \mathop{=} x_2$;
        \item there are $r' \mathop{\in} R'$ and $c' \mathop{\in} C'$ such that $h_{C'X}(c') \mathop{=} x_2$ and $h_{R'X}(r') \mathop{=} x_1$.
    \end{enumerate}
    Since the second and fourth cases are symmetric to the first and third case, respectively, it suffices to show that in the first and third cases, the inequality $h_{XG}(x_1) \mathop{\neq} h_{XG}(x_2)$ holds. 
    \begin{itemize}
        \item[Case (1)] Suppose that there are $c_1, c_2 \mathop{\in} C'$ such that $h_{C'X}(c_1) \mathop{=} x_1$ and $h_{C'X}(c_2) \mathop{=} x_2$. 
        
        We have $c_1 \mathop{\neq} c_2$ since $x_1 \mathop{\neq} x_2$ and $h_{C'X}$ is injective.
        
        Therefore we have 
        \begin{flalign*}
            h_{XG}(x_1)&= h_{XG}(h_{C'X}(c_1)) & \text{by definition of $c_1$}\\
                        &= (h_{C'X} \mathop{\star} h_{XG})(c_1)  \\
                        &= (h_{C'C} \mathop{\star} h_{CG})(c_1) & \text{by commutativity of $C'CGX$} \\
                        &\mathop{\neq} (h_{C'C} \mathop{\star} h_{CG})(c_2) & \text{by injectivity of $h_{C'C} \mathop{\star} h_{CG}$ and $c_1 \mathop{\neq} c_2$} \\
                        &= (h_{C'X} \mathop{\star} h_{XG})(c_2) & \text{by commutativity of $C'CGX$} \\
                        &= h_{XG}(h_{C'X}(c_2)) \\
                        &= h_{XG}(x_2). & \text{by definition of $c_2$}
        \end{flalign*}
        
        % \begin{center}
        %     \begin{tikzpicture}
        %         \node (k) at (0,0) {K};
        %         \node (c) at (0,-4) {C};
        %         \node (l) at (-4,0) {$L$};
        %         \node (g) at (-4,-4) {$G$};
        %         \draw[>->] (l) -- (g) node [midway,above] {};
        %         \draw[>->] (c) -- (g) node [midway,above] {};
        %         \draw[>->] (k) -- (c) node[midway, left] {};
        %         \draw[<-<] (l) to (k);
        %         \node (k') at (-1,-1) {K'};
        %         \node (r') at (-2,-1) {R'};
        %         \node (c') at (-1,-2) {C'};
        %         \node (x) at (-2,-2) {X};
        %         \node () at (-1.5,-1.5) {$PO$};
        %         \node () at (-0.5,-1.5) {$PB$};
        %         \draw [>->] (c') -- (x);
        %         \draw [>->] (r') -- (x);
        %         \draw [>->] (k') -- (r');
        %         \draw [>->] (k') -- (c');
        %         \draw [>->] (c') -- (c);
        %         \draw [>->] (r') -- (l);
        %         \draw [>->] (k') -- (k);
        %         \draw [->] (x) -- (g) node[midway,above] {!};
        %     \end{tikzpicture}
        % \end{center}
        % \item[(2)] Suppose that there are $r_1, r_2 \mathop{\in} R'$ such that $h_{R'X}(r_1) \mathop{=} x_1$ and $h_{R'X}(r_2) \mathop{=} x_2$. The proof of this case is analogous to the previous case.

        % We have $r_1 \mathop{\neq} r_2$ because $x_1 \mathop{\neq} x_2$ and $h_{R'X}$ is injective, and
        % \begin{flalign*}
        %     h_{XG}(x_1) &= h_{XG}(h_{R'X}(r_1)) & \text{by definition of $r_1$}\\
        %                  &= (h_{R'X} \mathop{\star} h_{XG})(r_1)  \\
        %                  &= (h_{R'L} \mathop{\star} h_{LG})(r_1) & \text{by commutativity of $R'XGL$} \\
        %                  &\mathop{\neq} (h_{R'L} \mathop{\star} h_{LG})(r_2) & \text{by injectivity of $h_{R'L} \mathop{\star} h_{LG}$ and $r_1 \mathop{\neq} r_2$} \\
        %                  &= (h_{R'X} \mathop{\star} h_{XG})(r_2) & \text{by commutativity of $R'XGL$} \\
        %                  &= h_{XG}(h_{R'X}(r_2)) \\
        %                  &= h_{XG}(x_2) & \text{by definition of $r_2$}
        % \end{flalign*}
        
        % \begin{center}
        %     \begin{tikzpicture}
        %         \node (k) at (0,0) {K};
        %         \node (c) at (0,-4) {C};
        %         \node (l) at (-4,0) {$L$};
        %         \node (g) at (-4,-4) {$G$};
        %         \draw[>->] (l) -- (g) node [midway,above] {};
        %         \draw[>->] (c) -- (g) node [midway,above] {};
        %         \draw[>->] (k) -- (c) node[midway, left] {};
        %         \draw[<-<] (l) to (k);
        %         \node (k') at (-1,-1) {K'};
        %         \node (r') at (-2,-1) {R'};
        %         \node (c') at (-1,-2) {C'};
        %         \node (x) at (-2,-2) {X};
        %         \node () at (-1.5,-1.5) {$PO$};
        %         \node () at (-0.5,-1.5) {$PB$};
        %         \draw [>->] (c') -- (x);
        %         \draw [>->] (r') -- (x);
        %         \draw [>->] (k') -- (r');
        %         \draw [>->] (k') -- (c');
        %         \draw [>->] (c') -- (c);
        %         \draw [>->] (r') -- (l);
        %         \draw [>->] (k') -- (k);
        %         \draw [->] (x) -- (g) node[midway,above] {!};
        %     \end{tikzpicture}
        % \end{center}
        \item[Case (3)] Suppose that there are $r' \mathop{\in} R'$ and $c' \mathop{\in} C'$ such that $h_{C'X}(c') \mathop{=} x_1$ and $h_{R'X}(r') \mathop{=} x_2$.
        % , or $h_{C'X}(c') \mathop{=} x_2$ and $h_{R'X}(r') \mathop{=} x_1$. By symmetry, we suppose that the first case holds. Let $c \mathop{=} h_{C'C}(c')$ and $l \mathop{=} h_{R'L}(r')$
        
        We have three cases:
            \begin{itemize}
                \item[(3.1)] Suppose that $h_{C'C}(c') \notin \operatorname{Im}(h_{KC})$. We have
                    \begin{flalign*}
                        h_{XG}(x_1) &= h_{XG}(h_{C'X}(c')) & \text{by definition of $c'$}\\
                                     &= (h_{C'X} \mathop{\star} h_{XG})(c') \\
                                     &= (h_{C'C} \mathop{\star} h_{CG})(c') & \text{by commutativity of $C'CGX$} \\
                                     &= h_{CG}(h_{C'C}(c')) & \\
                                    %  &= h_{CG}(c) & \\
                                     &\mathop{\neq} 
                                    %  h_{LG}(l) & \text{by Lemma~\ref{lem:b_c_same_img_exist_a} and $c\notin \operatorname{Im}(h_{KC})$}  
                                    h_{LG}(h_{R'L}(r')) & \text{by Lemma~\ref{lem:b_c_same_img_exist_a} and $h_{C'C}(c') \notin \operatorname{Im}(h_{KC})$}   
                                    \\
                                     &= (h_{R'L} \mathop{\star} h_{LG})(r') 
                                     \\
                                     &= (h_{R'X} \mathop{\star} h_{XG})(r') & \text{by commutativity of $R'XGL$} 
                                     \\
                                     &= h_{XG}(h_{R'X}(r'))
                                     \\
                                     &= h_{XG}(x_2). & \text{by definition of $r'$}
                    \end{flalign*}
                    
                    % \begin{center}
                    %     \begin{tikzpicture}
                    %         \node (k) at (0,0) {K};
                    %         \node (c) at (0,-4) {C};
                    %         \node (l) at (-4,0) {$L$};
                    %         \node (g) at (-4,-4) {$G$};
                    %         \draw[>->] (l) -- (g) node [midway,above] {};
                    %         \draw[>->] (c) -- (g) node [midway,above] {};
                    %         \draw[>->] (k) -- (c) node[midway, left] {};
                    %         \draw[<-<] (l) to (k);
                    %         \node (k') at (-1,-1) {K'};
                    %         \node (r') at (-2,-1) {R'};
                    %         \node (c') at (-1,-2) {C'};
                    %         \node (x) at (-2,-2) {X};
                    %         \node () at (-1.5,-1.5) {$PO$};
                    %         \node () at (-0.5,-1.5) {$PB$};
                    %         \draw [>->] (c') -- (x);
                    %         \draw [>->] (r') -- (x);
                    %         \draw [>->] (k') -- (r');
                    %         \draw [>->] (k') -- (c');
                    %         \draw [>->] (c') -- (c);
                    %         \draw [>->] (r') -- (l);
                    %         \draw [>->] (k') -- (k);
                    %         \draw [->] (x) -- (g) node[midway,above] {!};
                    %     \end{tikzpicture}
                    % \end{center}
                \item[(3.2)] Suppose $h_{R'L}(r') \notin \operatorname{Im}(h_{KL})$. This case is symmetric to the previous case.
                % where $c \notin \operatorname{Im}(h_{KC})$.
                \item[(3.3)] Suppose that there are $k_1, k_2 \mathop{\in} K$ such that $h_{KC}(k_1) \mathop{=} h_{C'C}(c')$ and $h_{KL}(k_2) \mathop{=} h_{R'L}(r')$. 
                
                We cannot have $k_1=k_2$. 
                
                Suppose, by contradiction, that we have 
                \begin{flalign}
                    k_1=k_2. \label{k1eqk2}
                \end{flalign} 
                
                % \begin{center}
                %     \begin{tikzpicture}
                %         \node (k) at (0,0) {K};
                %         \node (c) at (0,-4) {C};
                %         \node (l) at (-4,0) {$L$};
                %         \node (g) at (-4,-4) {$G$};
                %         \draw[>->] (l) -- (g) node [midway,above] {};
                %         \draw[>->] (c) -- (g) node [midway,above] {};
                %         \draw[>->] (k) -- (c) node[midway, left] {};
                %         \draw[<-<] (l) to (k);
                %         \node (k') at (-1,-1) {K'};
                %         \node (r') at (-2,-1) {R'};
                %         \node (c') at (-1,-2) {C'};
                %         \node (x) at (-2,-2) {X};
                %         \node () at (-1.5,-1.5) {$PO$};
                %         \node () at (-0.5,-1.5) {$PB$};
                %         \draw [>->] (c') -- (x);
                %         \draw [>->] (r') -- (x);
                %         \draw [>->] (k') -- (r');
                %         \draw [>->] (k') -- (c');
                %         \draw [>->] (c') -- (c);
                %         \draw [>->] (r') -- (l);
                %         \draw [>->] (k') -- (k);
                %         \draw [->] (x) -- (g) node[midway,above] {!};
                %     \end{tikzpicture}
                % \end{center}
                Since $K'C'CK$ is also a pullback square by Proposition~\ref{prop:pb_eq_po} and $h_{KC}(k_1) \mathop{=} h_{C'C}(c')$ and $h_{C'C}(c') \mathop{=}  h_{C'C}(c')$, there is $k' \mathop{\in} K'$ such that $h_{K'K}(k') \mathop{=} k_1$ and $h_{K'C'}(k') \mathop{=} c'$.
 
                We have $h_{K'R'}(k') \mathop{=} r'$ because $h_{R'L}$ is injective, $h_{R'L}(r') \mathop{=} h_{R'L}(r')$ and 
                \begin{flalign*}
                   &h_{R'L}(h_{K'R'}(k')) \hspace{2cm}&\\
                  =&(h_{K'R'} \mathop{\star} h_{R'L})(k') &\\
                  =&(h_{K'K} \mathop{\star} h_{KL})(k') & \text{by commutativity of $K'R'LK$}\\
                  =&(h_{KL})(h_{K'K}(k'))\\
                  =&(h_{KL})(k_1)\\
                  =&(h_{KL})(k_2). & \text{by assumption~\eqref{k1eqk2}}
                %  \\
                %   =& l
                \end{flalign*}
 
             Therefore, we have the contradiction:
                     \begin{flalign*}
                         x_1 &= h_{C'X}(c')  & \text{by definition of $c'$}\\
                             &= h_{C'X}(h_{K'C'}(k')) & \text{by definition of $k'$}\\
                             &= (h_{K'C'} \mathop{\star} h_{C'X})(k') &  \\
                             &= (h_{K'R'} \mathop{\star} h_{R'X})(k') & \text{by the commutativity of $K'C'XR'$} \\
                             &= h_{R'X}(h_{K'R'}(k')) &  \\
                             &= h_{R'X}(r') &  \text{by definition of $r'$}\\
                             &= x_2.& \text{by definition of $x_2$}
                     \end{flalign*}

                Thus, the following inequality holds 
                \begin{flalign}
                    k_1 \mathop{\neq} k_2. \label{k1neqk2}
                \end{flalign}

                \begin{flalign*}
                    h_{XG}(x_1) &= h_{XG}(h_{C'X}(c')) & \text{by definition of $c'$} \\
                            &= (h_{C'X} \mathop{\star} h_{XG})(c') &  \\
                            &= (h_{C'C} \mathop{\star} h_{CG})(c') & \text{by commutativity of $C'CGX$} \\
                            % &= (h_{CG})(c) & \text{by definition of $c$} \\
                            &= (h_{CG})(h_{KC}(k_1)) & \text{by definition of $k_1$} 
                            \\
                            &= (h_{KC} \mathop{\star} h_{CG})(k_1) & 
                            \\
                            &\mathop{\neq} (h_{KC} \mathop{\star} h_{CG})(k_2)  & \text{by injectivity of $h_{KC}\star h_{CG}$ and~\eqref{k1neqk2}} 
                            \\
                            &= (h_{KL} \mathop{\star} h_{LG})(k_2) & \text{by commutativity of $KLGC$}
                            \\
                            &= h_{LG}(h_{KL}(k_2)) & 
                            \\
                            % &= h_{LG}(l) & \text{by the definition of $k_2$}
                            % \\
                            &= h_{LG}(h_{R'L}(r')) & \text{by the definition of $k_2$}
                            \\
                            &= (h_{R'L} \mathop{\star} h_{LG})(r') & 
                            \\
                            &= (h_{R'X} \mathop{\star} h_{XG})(r') & \text{by commutativity of $R'XLG$} \\
                            &= h_{XG}(h_{R'X}(r'))\\
                                &= h_{XG}(x_2). & \text{by definition of $r'$}
                \end{flalign*}
                
                % \begin{center}
                %     \begin{tikzpicture}
                %         \node (k) at (0,0) {K};
                %         \node (c) at (0,-4) {C};
                %         \node (l) at (-4,0) {$L$};
                %         \node (g) at (-4,-4) {$G$};
                %         \draw[>->] (l) -- (g) node [midway,above] {};
                %         \draw[>->] (c) -- (g) node [midway,above] {};
                %         \draw[>->] (k) -- (c) node[midway, left] {};
                %         \draw[<-<] (l) to (k);
                %         \node (k') at (-1,-1) {K'};
                %         \node (r') at (-2,-1) {R'};
                %         \node (c') at (-1,-2) {C'};
                %         \node (x) at (-2,-2) {X};
                %         \node () at (-1.5,-1.5) {$PO$};
                %         \node () at (-0.5,-1.5) {$PB$};
                %         \draw [>->] (c') -- (x);
                %         \draw [>->] (r') -- (x);
                %         \draw [>->] (k') -- (r');
                %         \draw [>->] (k') -- (c');
                %         \draw [>->] (c') -- (c);
                %         \draw [>->] (r') -- (l);
                %         \draw [>->] (k') -- (k);
                %         \draw [->] (x) -- (g) node[midway,above] {!};
                %     \end{tikzpicture}
                % \end{center}
            \end{itemize}            
    \end{itemize}
\end{proof}



\begin{lemma}
    \label{lem:h_hp_diff_g_gp_diff}
    
    Let $X$ be a graph. Let $\rho=(L \overset{l}{\leftarrowtail} K \overset{r}{\rightarrowtail} R)$ be an injective DPO rewriting rule which is \(X\)-non-increasing under $\Psi$. Consider the following DPO diagram:

    \begin{center}
        \begin{tikzpicture}
            \node (k) at (0,1) {$K$};
            \node (l) at (-2,1) {$L$};
            \node (r) at (2,1) {$R$};
            \node (c) at (0,-1) {$C$};
            \node (g) at (-2,-1) {$G$};
            \node (h) at (2,-1) {$H$};
            % \node (rb) at (1.5,-0.5) {$R_X$};
            % \node (h') at (1.5,-1.5) {$H'$};
            % \draw[>->]  (rb) -- (h') node [midway,above] {};
            % \draw[>->]  (c) -- (h') node [midway,above] {};
            \draw[<-<]  (l) -- (k) node [midway,below] {$l$};
            \draw[>->]  (k) -- (r) node [midway,below] {$r$};
            \draw[>->] (c) -- (g) node [midway, above] {$l'$};
            \draw[>->] (c) -- (h) node [midway,above] {$r'$};
            \draw[>->] (l) -- (g) node[midway, right] {$m$};
            \draw[>->] (r) -- (h) node[midway, left] {$m'$};
            \draw[>->] (k) -- (c) node[midway, left] {};
            \node () [at=($(l)!0.5!(c)$)] {$PO$};
            \node () [at=($(r)!0.5!(c)$)] {$PO$};
            % \draw[>->] (rb) to (r);
            % \draw[->] (rb) to (l);
            % \draw[<-<] (rb) to (k);
        \end{tikzpicture}
    \end{center}
   
    Let $h_{XH}',h_{XH}'':X \rightarrowtail H$ be monomorphisms.
    In the category \textbf{Graph}, we can construct the following commutative diagram
    \begin{center}
        \begin{tikzpicture} 
            \node (k) at (0,0) {K};
            \node (c) at (0,-4) {C};
            \node (l) at (-4,0) {$R$};
            \node (g) at (-4,-4) {$H$};
            \node (r) at (4,0) {$L$};
            \node (h) at (4,-4) {$G$};
            \draw [>->] (k) -- (r); 
            \draw [>->] (c) -- (h);
            \draw [>->] (r) -- (h);
            \draw[>->] (l) -- (g) node [midway,above] {};
            \draw[>->] (c) -- (g) node [midway,above] {};
            \draw[>->] (k) -- (c) node[midway, left] {};
            \draw[<-<] (l) to (k);
            \node (k') at (-2,-1) {K'};
            \node (r') at (-3.5,-1.5) {R'};
            \node (c') at (-1.5,-2.2) {C'};
            \node (x) at (-3,-3) {X};
            \draw [>->] (c') -- (x);
            \draw [>->] (r') -- (x); 
            \draw [>->] (k') -- (r');
            \draw [>->] (k') -- (c');
            \draw [>->] (c') -- (c);
            \draw [>->] (r') -- (l);
            \draw [>->] (k') -- (k);

            \node (k'') at (-0.5,-1.5) {K''};
            \node (r'') at (-2.5,-2) {R''};
            \node (c'') at (-0.5,-2.5) {C''};
            \draw [>->] (c'') -- (x);
            \draw [>->] (r'') -- (x);
            \draw [>->] (k'') -- (r'');
            \draw [>->] (k'') -- (c'');
            \draw [>->] (k'') -- (k);
            \draw [>->] (c'') -- (c);
            \draw [>->] (r'') -- (l);

            \draw [>->] (x) edge[bend left] node[pos=0.1,below]{$h_{XH}'$} (g);
            \draw [>->] (x) edge[bend right] node[midway,above]{$h_{XH}''$} (g);

        \end{tikzpicture}
    \end{center}
    where $K'R'XC'$, $K''R''XC''$, $K'KCC'$, $K''KCC''$, $R'RHX$ and $R''RHX$ are pullbacks.

    By Lemma~\ref{kpcpck_pullback}, $K'KRR'$ and $K''KCC''$ are pullbacks. 

    By Lemma~\ref{kpcpxrp_po}, $K'R'XC'$ and $K''R''XC''$ are pushouts.

    Since $\rho$ is $X$-non-increasing under $\Psi$, there are $\Psi(R'): R' \rightarrowtail L$ and $\Psi(R''):R'' \rightarrowtail L$ such that $K'R'LK$ and $K''R''LK$ are pullbacks.

    By Lemma~\ref{lem:g_monic}, we have the following commutative diagram
    \begin{center}
        \begin{tikzpicture} 
            \node (k) at (0,0) {K};
            \node (c) at (0,-4) {C};
            \node (l) at (-4,0) {L};
            \node (g) at (-4,-4) {G};
            \node (r) at (4,0) {L};
            \node (h) at (4,-4) {G};
            \draw [>->] (k) -- (r); 
            \draw [>->] (c) -- (h);
            \draw [>->] (r) -- (h);
            \draw[>->] (l) -- (g) node [midway,above] {};
            \draw[>->] (c) -- (g) node [midway,above] {};
            \draw[>->] (k) -- (c) node[midway, left] {};
            \draw[<-<] (l) to (k);
            \node (k') at (-2,-1) {K'};
            \node (r') at (-3.5,-1.5) {R'};
            \node (c') at (-1.5,-2.2) {C'};
            \node (x) at (-3,-3) {X};
            \draw [>->] (c') -- (x);
            \draw [>->] (r') -- (x); 
            \draw [>->] (k') -- (r');
            \draw [>->] (k') -- (c');
            \draw [>->] (c') -- (c);
            \draw [>->] (r') -- (l);
            \draw [>->] (k') -- (k);

            \node (k'') at (-0.5,-1.5) {K''};
            \node (r'') at (-2.5,-2) {R''};
            \node (c'') at (-0.5,-2.5) {C''};
            \draw [>->] (c'') -- (x);
            \draw [>->] (r'') -- (x);
            \draw [>->] (k'') -- (r'');
            \draw [>->] (k'') -- (c'');
            \draw [>->] (k'') -- (k);
            \draw [>->] (c'') -- (c);
            \draw [>->] (r'') -- (l);

            \draw [>->] (x) edge[bend left] node[pos=0.1,below]{$h_{XG}'$} (g);
            \draw [>->] (x) edge[bend right] node[midway,above]{$h_{XG}''$} (g);

        \end{tikzpicture}
    \end{center}
    where  

    \begin{itemize}
        \item $K'R'XC'$ is pushout and pullback;
        \item $K''R''XC''$ is pushout and pullback;
        \item $K'KCC'$ and $K''KCC''$ are pullback;
        \item $h_{XG}':X \mathop{\to} G$ is the unique morphism such that $h_{C'C} \mathop{\star} h_{CG} \mathop{=} h_{C'X} \mathop{\star} h_{XG}'$ and  $h_{R'L} \mathop{\star} h_{LG} \mathop{=} h_{R'X} \mathop{\star} h_{XG}'$ where $h_{R'L} \mathop{=} \Psi(R')$;
        \item $h_{XG}'':X \mathop{\to} G$ is the unique morphism such that $h_{C''C} \mathop{\star} h_{CG} \mathop{=} h_{C''X} \mathop{\star} h_{XG}''$ and  $h_{R''L} \mathop{\star} h_{LG} \mathop{=} h_{R''X} \mathop{\star} h_{XG}''$ where $h_{R''L} \mathop{=} \Psi(R'')$.
    \end{itemize}
    If $h_{XH}' \mathop{\neq} h_{XH}''$ then $h_{XG}' \mathop{\neq} h_{XG}''$.
\end{lemma}
\begin{proof}
    Suppose that $h_{XH}' \mathop{\neq} h_{XH}''$. We are going to show $h_{XG}' \mathop{\neq} h_{XG}''$.

    From $h_{XH}' \mathop{\neq} h_{XH}''$, we deduce there is an element $x \mathop{\in} X$ such that the inequality $h_{XH}(x) \mathop{\neq} h_{XH}''(x)$ holds. We can suppose that 
        \begin{flalign}
            \text{$x$ is either an isolated nodes or an edge,} \label{x_isolated_or_edge}
        \end{flalign} 
         \todo{$x$ is either an isolated node or an edge}
    because otherwise $x$ would be a non-isolated node and we can take an incident edge of $x$.

    \noindent
    Since $K'R'XC'$ and $K''R''XC''$ are pushouts, by Lemma~\ref{lem:xinXcpinCrpinR}, there are 4 possible cases:
    \begin{enumerate}
        \item there are $c' \mathop{\in} C', c'' \mathop{\in} C''$ such that $h_{C'X}(c') \mathop{=} h_{C'X}(c'') \mathop{=} x$;
        \item  there are $r' \mathop{\in} R', r'' \mathop{\in} R''$ such that $h_{R'X}(r') \mathop{=} h_{R'X}(r'') \mathop{=} x$;
        \item there are arrows $r' \mathop{\in} R'$ and $c'' \mathop{\in} C''$ such that 
        \begin{flalign*}
            h_{C''X}(c'') \mathop{=} x
            \\
            h_{R'X}(r') \mathop{=} x
            \\
            \nexists c' \mathop{\in} C'. h_{C'X}(c') \mathop{=} x 
            \\
            \nexists r'' \mathop{\in} R''. h_{R''X}(r'') \mathop{=} x;
        \end{flalign*}
        \item there are arrows $r'' \mathop{\in} R''$ and $c' \mathop{\in} C'$ such that 
        \begin{flalign*}
            h_{C'X}(c') \mathop{=} x
            \\
            h_{R''X}(r'') \mathop{=} x
            \\
            \nexists c'' \mathop{\in} C''. h_{C''X}(c'') \mathop{=} x 
            \\
            \nexists r' \mathop{\in} R'. h_{R'X}(r') \mathop{=} x.  
        \end{flalign*}
    \end{enumerate}

    We are going to show that in each case we have $h_{XG}'(x) \mathop{\neq} h_{XG}''(x)$. Since the fourth case is symmetric to the third case, it suffices to analyse the first three cases.

    \begin{itemize}
        \item[(1)] Suppose that there are $c' \mathop{\in} C', c'' \mathop{\in} C''$ such that $h_{C'X}(c') \mathop{=} h_{C'X}(c'') \mathop{=} x$. 
        Suppose, by contradiction, that the following equality holds
        \begin{flalign}
            h_{C'C}(c') \mathop{=} h_{C''C}(c''). \label{cpccp_eq_cppccpp}
        \end{flalign}
         We have the following contradiction:
        \begin{flalign*}
            &h_{XH}'(x) \\
           =&h_{XH}'(h_{C'X}(c')) \\
           =&(h_{C'X} \mathop{\star} h_{XH}')(c') \\
           =&(h_{C'C} \mathop{\star} h_{CH})(c') &\text{By commutativity of $C'XHC$}\\
           =&h_{CH}(h_{C'C}(c')) \\
           =&h_{CH}(h_{C''C}(c'')) & \text{ by assumption~\eqref{cpccp_eq_cppccpp}}\\
           =&(h_{C''C} \mathop{\star} h_{CH})(c'') \\
           =&(h_{C''X} \mathop{\star} h_{XH}'')(c'') &\text{By commutativity of $C''CHX$}\\
           =&h_{XH}''(h_{C''X}(c'')) \\
           \mathop{=} &h_{XH}''(x). \\
        \end{flalign*}
        
        % \begin{center}
        %     \begin{tikzpicture} 
        %         \node (k) at (0,0) {K};
        %         \node (c) at (0,-4) {C};
        %         \node (l) at (-4,0) {$L$};
        %         \node (g) at (-4,-4) {$G$};
        %         \draw[>->] (l) -- (g) node [midway,above] {};
        %         \draw[>->] (c) -- (g) node [midway,above] {};
        %         \draw[>->] (k) -- (c) node[midway, left] {};
        %         \draw[<-<] (l) to (k);
        %         \node (k') at (-2,-1) {K'};
        %         \node (r') at (-3.5,-1.5) {R'};
        %         \node (c') at (-1.5,-2.2) {C'};
        %         \node (x) at (-3,-3) {X};
        %         \draw [>->] (c') -- (x);
        %         \draw [>->] (r') -- (x); 
        %         \draw [>->] (k') -- (r');
        %         \draw [>->] (k') -- (c');
        %         \draw [>->] (c') -- (c);
        %         \draw [>->] (r') -- (l);
        %         \draw [>->] (k') -- (k);
    
        %         \node (k'') at (-0.5,-1.5) {K''};
        %         \node (r'') at (-2.5,-2) {R''};
        %         \node (c'') at (-0.5,-2.5) {C''};
        %         \draw [>->] (c'') -- (x);
        %         \draw [>->] (r'') -- (x);
        %         \draw [>->] (k'') -- (r'');
        %         \draw [>->] (k'') -- (c'');
        %         \draw [>->] (k'') -- (k);
        %         \draw [>->] (c'') -- (c);
        %         \draw [>->] (r'') -- (l);
    
        %         \draw [>->] (x) edge[bend left] node[pos=0.3,below]{!g''} (g);
        %         \draw [>->] (x) edge[bend right] node[midway,above]{!g'} (g);
    
        %     \end{tikzpicture}
        % \end{center}
        Therefore, we have 
        \begin{flalign}
            h_{C'C}(c') \mathop{\neq} h_{C''C}(c''). \label{cpccp_neq_cppccpp}
        \end{flalign} 
        \begin{flalign*}
            h_{XG}'(x)&= h_{XG}'(h_{C'X}(c')) & \text{by definition of $c'$}\\
                        &= (h_{C'X} \mathop{\star} h_{XG}')(c')  \\
                        &= (h_{C'C} \mathop{\star} h_{CG})(c') & \text{by commutativity of $C'CGX$} \\
                        &= h_{CG}(h_{C'C}(c')) &  \\
                        &\mathop{\neq} h_{CG}(h_{C''C}(c'')) & \text{by~\eqref{cpccp_neq_cppccpp} and the injectivity of $h_{CG}$} \\
                        &= (h_{C''C} \mathop{\star} h_{CG})(c'') &  \\
                        &= (h_{C''X} \mathop{\star} h_{XG}'')(c'') & \text{by commutativity of $C''CGX$} \\
                        &= h_{XG}''(h_{C''X}(c'')) \\
                        &= h_{XG}''(x). & \text{by definition of $c''$}
        \end{flalign*}
  
        \item[(2)] Suppose that there are $r' \mathop{\in} R', r'' \mathop{\in} R''$ such that $h_{R'X}(r') \mathop{=} h_{R'X}(r'') \mathop{=} x$. We have 
        \begin{flalign}
            h_{R'R}(r') \mathop{\neq} h_{R''R}(r'') \label{rprrppr}
        \end{flalign}  
        by an argument analogous to the argument for $h_{C'C}(c') \mathop{\neq} h_{C''C}(c'')$ in the previous case. Since $\rho$ is $X$-non-increasing, by~\eqref{rprrppr} and~\eqref{x_isolated_or_edge}, we have \todo{use : x is either an isolated node or an edge}       
        \begin{flalign}
            h_{R'L}(r') \mathop{\neq} h_{R''L}(r''). \label{rplrprpplrpp}
        \end{flalign}  

        % Therefore, we have  
        % \begin{flalign*}
        %     h_{XG}'(x)&= h_{XG}'(h_{R'X}(r')) & \text{by definition of $r'$}\\
        %         &= (h_{R'X} \mathop{\star} h_{XG}')(r')  \\
        %         &= (h_{R'L} \mathop{\star} h_{LG})(r') & \text{by commutativity of $R'LGX$} \\
        %         &= h_{LG}(h_{R'L}(r')) &  \\
        %         &\mathop{\neq} h_{LG}(h_{R''L}(r'')) & \text{by $h_{R'R}(r') \mathop{\neq} h_{R''R}(r'')$ and the injectivity of $h_{RG}$} \\
        %         &= (h_{R''L} \mathop{\star} h_{LG})(r'') &  \\
        %         &= (h_{R''X} \mathop{\star} h_{XG}'')(r'') & \text{by commutativity of $R''LGX$} \\
        %         &= h_{XG}''(h_{R''X}(r'')) \\
        %         &= h_{XG}''(x) & \text{by definition of $r''$}
        % \end{flalign*}
         
        % \todo{$h_{R_XL}$ edge-injective} Since $h_{R_XL}$ is edge-injective by assumption, because it is $X$-non-increasing by assumption. Therefore, we have $h_{R'L}(r') \mathop{\neq} h_{R''L}(r'')$ because \( h_{R'R}(r') \mathop{\neq} h_{R''R}(r'') \). Hence, from~\ref{rprneqrppr}, we have 
        % \begin{flalign}
        %     h_{R'L}(r') \mathop{\neq} h_{R''L}(r'') \label{rplrprpplrpp}
        % \end{flalign}
        Hence, the following inequality holds
        
        % \begin{center}
        %     \begin{tikzpicture} 
        %         \node (k) at (0,0) {K};
        %         \node (c) at (0,-4) {C};
        %         \node (l) at (-4,0) {$L$};
        %         \node (g) at (-4,-4) {$G$};
        %         \draw[>->] (l) -- (g) node [midway,above] {};
        %         \draw[>->] (c) -- (g) node [midway,above] {};
        %         \draw[>->] (k) -- (c) node[midway, left] {};
        %         \draw[<-<] (l) to (k);
        %         \node (k') at (-2,-1) {K'};
        %         \node (r') at (-3.5,-1.5) {R'};
        %         \node (c') at (-1.5,-2.2) {C'};
        %         \node (x) at (-3,-3) {X};
        %         \draw [>->] (c') -- (x);
        %         \draw [>->] (r') -- (x); 
        %         \draw [>->] (k') -- (r');
        %         \draw [>->] (k') -- (c');
        %         \draw [>->] (c') -- (c);
        %         \draw [>->] (r') -- (l);
        %         \draw [>->] (k') -- (k);
    
        %         \node (k'') at (-0.5,-1.5) {K''};
        %         \node (r'') at (-2.5,-2) {R''};
        %         \node (c'') at (-0.5,-2.5) {C''};
        %         \draw [>->] (c'') -- (x);
        %         \draw [>->] (r'') -- (x);
        %         \draw [>->] (k'') -- (r'');
        %         \draw [>->] (k'') -- (c'');
        %         \draw [>->] (k'') -- (k);
        %         \draw [>->] (c'') -- (c);
        %         \draw [>->] (r'') -- (l);
    
        %         \draw [>->] (x) edge[bend left] node[pos=0.3,below]{!g''} (g);
        %         \draw [>->] (x) edge[bend right] node[midway,above]{!g'} (g);
    
        %     \end{tikzpicture}
        % \end{center}
        \begin{flalign*}
            h_{XG}'(x) &= h_{XG}'(h_{R'X}(r')) & \text{by definition of $r'$}\\
                         &= (h_{R'X} \mathop{\star} h_{XG}')(r')  \\
                         &= (h_{R'L} \mathop{\star} h_{LG})(r') & \text{by commutativity of $R'XGL$} \\
                         &= h_{LG}(h_{R'L}(r')) &  \\
                         &\mathop{\neq} h_{LG}(h_{R''L}(r'')) &  \text{By injectivity of $h_{LG}$ and~\eqref{rplrprpplrpp}}\\
                         &= (h_{R''L} \mathop{\star} h_{LG})(r'') &  \\
                         &= (h_{R''X} \mathop{\star} h_{XG}'')(r'') & \text{by commutativity of $R''XGL$} \\
                         &= h_{XG}''(h_{R''X}(r'')) \\
                         &= h_{XG}''(x). & \text{by definition of $r''$}
        \end{flalign*}  
        
        % \begin{center}
        %     \begin{tikzpicture} 
        %         \node (k) at (0,0) {K};
        %         \node (c) at (0,-4) {C};
        %         \node (l) at (-4,0) {$L$};
        %         \node (g) at (-4,-4) {$G$};
        %         \draw[>->] (l) -- (g) node [midway,above] {};
        %         \draw[>->] (c) -- (g) node [midway,above] {};
        %         \draw[>->] (k) -- (c) node[midway, left] {};
        %         \draw[<-<] (l) to (k);
        %         \node (k') at (-2,-1) {K'};
        %         \node (r') at (-3.5,-1.5) {R'};
        %         \node (c') at (-1.5,-2.2) {C'};
        %         \node (x) at (-3,-3) {X};
        %         \draw [>->] (c') -- (x);
        %         \draw [>->] (r') -- (x); 
        %         \draw [>->] (k') -- (r');
        %         \draw [>->] (k') -- (c');
        %         \draw [>->] (c') -- (c);
        %         \draw [>->] (r') -- (l);
        %         \draw [>->] (k') -- (k);
    
        %         \node (k'') at (-0.5,-1.5) {K''};
        %         \node (r'') at (-2.5,-2) {R''};
        %         \node (c'') at (-0.5,-2.5) {C''};
        %         \draw [>->] (c'') -- (x);
        %         \draw [>->] (r'') -- (x);
        %         \draw [>->] (k'') -- (r'');
        %         \draw [>->] (k'') -- (c'');
        %         \draw [>->] (k'') -- (k);
        %         \draw [>->] (c'') -- (c);
        %         \draw [>->] (r'') -- (l);
    
        %         \draw [>->] (x) edge[bend left] node[pos=0.3,below]{!g''} (g);
        %         \draw [>->] (x) edge[bend right] node[midway,above]{!g'} (g);
    
        %     \end{tikzpicture}
        % \end{center}
        \item[(3)] Suppose, that there are arrows $r' \mathop{\in} R'$ and $c'' \mathop{\in} C''$ such that the following statements hold:
        \begin{flalign}
            h_{C''X}(c'') \mathop{=} x, \nonumber
            \\
            h_{R'X}(r') \mathop{=} x, \nonumber
            \\
            \nexists c' \mathop{\in} C'. h_{C'X}(c') \mathop{=} x, \nonumber
            \\
            \nexists r'' \mathop{\in} R''. h_{R''X}(r'') \mathop{=} x.  \label{assumption_norpp}
        \end{flalign}

        Suppose, by contradiction, that we have \begin{flalign}
            h_{C''C}(c'') \mathop{\in} \operatorname{Im}(h_{KC}).  \label{assump_c_in_imhkc}
        \end{flalign} 
        Then,
        there are $k_2 \mathop{\in} K$ such that $h_{KC}(k_2) \mathop{=} h_{C''C}(c'')$.

        There is $k'' \mathop{\in} K''$ such that $h_{K''K}(k'') \mathop{=} k_2$ and $h_{K''C''}(k'') \mathop{=} c''$ becase $K''C''CK$ is a pullback square and $h_{KC}(k_2) \mathop{=} h_{C''C}(c'')$.

        We have $h_{K''R''}(k'') \mathop{\in} R''$ and 
        \begin{flalign*}
          &h_{R''X}(h_{K''R''}(k''))\\
         =&(h_{K''R''} \mathop{\star} h_{R''X})(k'') \\
         =&(h_{K''C''} \mathop{\star} h_{C''X})(k'') & \text{by commutativity of $K''R''XC''$}\\
         =&h_{C''X}(h_{K''C''}(k''))\\
         =&h_{C''X}(c'')\\
         =&x.
        \end{flalign*}
        which contradicts the assumption~\eqref{assumption_norpp}.
        
        Thus, the following holds
        \begin{flalign}
            h_{C''C}(c'') \notin \operatorname{Im}(h_{KC}).  \label{assump_c_notin_imhkc}
        \end{flalign} 

        We deduce
        \begin{flalign*}
            h_{XG}'(x) &= h_{XG}'(h_{C''X}(c'')) & \text{by definition of $c''$}\\
                         &= (h_{C''X} \mathop{\star} h_{XG}')(c'') \\
                         &= (h_{C''C} \mathop{\star} h_{CG})(c'') & \text{by commutativity of $C''CGX$} \\
                         &= h_{CG}(h_{C''C}(c'')) &  \\
                         &\mathop{\neq}  h_{LG}(h_{R'L}(r')) & \text{by Lemma~\ref{lem:b_c_same_img_exist_a} and the assumption~\eqref{assump_c_notin_imhkc}} \\
                         &= (h_{R'L} \mathop{\star} h_{LG})(r') \\
                         &= (h_{R'X} \mathop{\star} h_{XG}'')(r') & \text{by commutativity of $R'XGL$} \\
                         &= h_{XG}''(h_{R'X}(r'))\\
                         &= h_{XG}''(x). & \text{by definition of $r'$}
        \end{flalign*} 

    %     We analyse the two cases where $h_{C''C}(c'') \notin \operatorname{Im}(h_{KC}) $ and 
    %         \begin{itemize}
    %             \item[(3.1)] Suppose, that the following holds
    %                 \begin{flalign}
    %                     h_{C''C}(c'') \notin \operatorname{Im}(h_{KC})  \label{assump_c_notin_imhkc}
    %                 \end{flalign} then, we have 
    %                 \begin{flalign*}
    %                     h_{XG}'(x) &= h_{XG}'(h_{C''X}(c'')) & \text{by definition of $c''$}\\
    %                                  &= (h_{C''X} \mathop{\star} h_{XG}')(c'') \\
    %                                  &= (h_{C''C} \mathop{\star} h_{CG})(c'') & \text{by commutativity of $C''CGX$} \\
    %                                  &= h_{CG}(h_{C''C}(c'')) &  \\
    %                                  &\mathop{\neq}  h_{LG}(h_{R'L}(r')) & \text{by Lemma~\ref{lem:b_c_same_img_exist_a} and the assumption Equation~\eqref{assump_c_notin_imhkc}} \\
    %                                  &= (h_{R'L} \mathop{\star} h_{LG})(r') \\
    %                                  &= (h_{R'X} \mathop{\star} h_{XG}'')(r') & \text{by commutativity of $R'XGL$} \\
    %                                  &= h_{XG}''(h_{R'X}(r'))\\
    %                                  &= h_{XG}''(x) & \text{by definition of $r'$}
    %                 \end{flalign*} 
    %             \item[(3.2)] Suppose, that the following holds
    %             \begin{flalign}
    %                 h_{R'R}(r') \notin \operatorname{Im}(h_{KR}) \label{rprrp_im_kr}
    %             \end{flalign} 
    %             Since $\rho$ is $X$-non-increasing, by the assumption Equation~\eqref{rprrp_im_kr}, (Definition~\ref{def:creates_more_x_on_the_left})$h_{R'L}$ does not identify non-interface elements with interface elements. Therefore, the following holds 
    %             \begin{flalign}
    %                 h_{R'L}(r') \notin \operatorname{Im}(h_{KL}) \label{lnotinhkl}
    %             \end{flalign} 
    %             \todo{assumption: does not identify non-interface elements with interface elements} 
    %             Therefore, 
    %     
    %     % \begin{center}
    %     %     \begin{tikzpicture} 
    %     %         \node (k) at (0,0) {K};
    %     %         \node (c) at (0,-4) {C};
    %     %         \node (l) at (-4,0) {$L$};
    %     %         \node (g) at (-4,-4) {$G$};
    %     %         \draw[>->] (l) -- (g) node [midway,above] {};
    %     %         \draw[>->] (c) -- (g) node [midway,above] {};
    %     %         \draw[>->] (k) -- (c) node[midway, left] {};
    %     %         \draw[<-<] (l) to (k);
    %     %         \node (k') at (-2,-1) {K'};
    %     %         \node (r') at (-3.5,-1.5) {R'};
    %     %         \node (c') at (-1.5,-2.2) {C'};
    %     %         \node (x) at (-3,-3) {X};
    %     %         \draw [>->] (c') -- (x);
    %     %         \draw [>->] (r') -- (x); 
    %     %         \draw [>->] (k') -- (r');
    %     %         \draw [>->] (k') -- (c');
    %     %         \draw [>->] (c') -- (c);
    %     %         \draw [>->] (r') -- (l);
    %     %         \draw [>->] (k') -- (k);
    
    %     %         \node (k'') at (-0.5,-1.5) {K''};
    %     %         \node (r'') at (-2.5,-2) {R''};
    %     %         \node (c'') at (-0.5,-2.5) {C''};
    %     %         \draw [>->] (c'') -- (x);
    %     %         \draw [>->] (r'') -- (x);
    %     %         \draw [>->] (k'') -- (r'');
    %     %         \draw [>->] (k'') -- (c'');
    %     %         \draw [>->] (k'') -- (k);
    %     %         \draw [>->] (c'') -- (c);
    %     %         \draw [>->] (r'') -- (l);
    
    %     %         \draw [>->] (x) edge[bend left] node[pos=0.3,below]{!g''} (g);
    %     %         \draw [>->] (x) edge[bend right] node[midway,above]{!g'} (g);
    
    %     %     \end{tikzpicture}
    %     % \end{center}
    %                 \begin{flalign*}
    %                     h_{XG}'(x) &= h_{XG}'(h_{R'X}(r')) & \text{by definition of $r'$}\\
    %                                  &= (h_{R'X} \mathop{\star} h_{XG}')(r') \\
    %                                  &= (h_{R'L} \mathop{\star} h_{LG})(r') & \text{by commutativity of $R'LGX$} \\
    %                                  &= h_{LG}(h_{R'L}(r')) &  \\
    %                                  &= h_{CG}(h_{C''C}(c'')) & \text{by Lemma~\ref{lem:b_c_same_img_exist_a} and the assumption Equation~\eqref{lnotinhkl}}  \\
    %                                  &= (h_{C''C} \mathop{\star} h_{CG})(c'') \\
    %                                  &= (h_{C''X} \mathop{\star} h_{XG}'')(c'') & \text{by commutativity of $C''XGC$} \\
    %                                  &= h_{XG}''(h_{C''X}(c''))\\
    %                                  &= h_{XG}''(x) & \text{by definition of $c''$}
    %                 \end{flalign*}
    %                  \newpage  
    %          
    % \begin{center}
    %         \begin{tikzpicture} 
    %             \node (k) at (0,0) {K};
    %             \node (c) at (0,-4) {C};
    %             \node (l) at (-4,0) {$L$};
    %             \node (g) at (-4,-4) {$G$};
    %             \draw[>->] (l) -- (g) node [midway,above] {};
    %             \draw[>->] (c) -- (g) node [midway,above] {};
    %             \draw[>->] (k) -- (c) node[midway, left] {};
    %             \draw[<-<] (l) to (k);
    %             \node (k') at (-2,-1) {K'};
    %             \node (r') at (-3.5,-1.5) {R'};
    %             \node (c') at (-1.5,-2.2) {C'};
    %             \node (x) at (-3,-3) {X};
    %             \draw [>->] (c') -- (x);
    %             \draw [>->] (r') -- (x); 
    %             \draw [>->] (k') -- (r');
    %             \draw [>->] (k') -- (c');
    %             \draw [>->] (c') -- (c);
    %             \draw [>->] (r') -- (l);
    %             \draw [>->] (k') -- (k);
    
    %             \node (k'') at (-0.5,-1.5) {K''};
    %             \node (r'') at (-2.5,-2) {R''};
    %             \node (c'') at (-0.5,-2.5) {C''};
    %             \draw [>->] (c'') -- (x);
    %             \draw [>->] (r'') -- (x);
    %             \draw [>->] (k'') -- (r'');
    %             \draw [>->] (k'') -- (c'');
    %             \draw [>->] (k'') -- (k);
    %             \draw [>->] (c'') -- (c);
    %             \draw [>->] (r'') -- (l);
     
    %             \draw [>->] (x) edge[bend left] node[pos=0.3,below]{!g''} (g);
    %             \draw [>->] (x) edge[bend right] node[midway,above]{!g'} (g);
    
    %         \end{tikzpicture}
    %     \end{center}

    %             \item[(3.3)]  
    %             Suppose, 
    %             \color{red}
    %             that \begin{flalign}
    %                 h_{R'R}(r') \mathop{\in} h_{KR}(K) 
    %                 \\
    %                 h_{C''C}(c'') \mathop{\in} h_{KC}(K)
    %             \end{flalign}
    %             Thus,
    %             there are $k_1, k_2 \mathop{\in} K$ such that $h_{KR}(k_1) \mathop{=} h_{R'R}(r')$ and $h_{KC}(k_2) \mathop{=} h_{C''C}(c'')$.
    %             \color{black}

    %             There is $k'' \mathop{\in} K''$ such that $h_{K''K}(k'') \mathop{=} k_2$ and $h_{K''C''}(k'') \mathop{=} c''$ becase $K''C''CK$ is a pullback square and $h_{KC}(k_2) \mathop{=} c \mathop{=} h_{C''C}(c'')$.

    %             We cannot have $k_1=k_2$. Suppose, by contradiction, that we have $k_1=k_2$. 
    %             \color{red} 
    %             Let $r'' \mathop{=} h_{K''R''}(k'')$. We have 
    %             \begin{flalign}
    %                 &h_{R''R}(r'')\\
    %                =&h_{R''R}(h_{K''R''}(k''))\\
    %                =&(h_{K''R''} \mathop{\star} h_{R''R})(k'') & \text{by definition of $r''$}\\
    %                =&(h_{K''K} \mathop{\star} h_{KR})(k'') &\text{by commutativity of $K''R''RK$}\\
    %                =&h_{KR}(h_{K''K}(k''))  \\
    %                =&h_{KR}(k_1) & \text{by definition of $k_1$} \\
    %                =&r
    %             \end{flalign}
    %             \color{black}
    %             % We have therefore $h_{K''C''}(k'') \mathop{=} c''$ because $h_{C''C}$ is injective, $h_{C''C}(c'') \mathop{=} c$ and 
    %             % \begin{flalign*}
    %             %    &h_{C''C}(h_{K''C''}(k'')) \\
    %             %   =&(h_{K''C''} \mathop{\star} h_{C''C})(k'')\\
    %             %   =&(h_{K''K} \mathop{\star} h_{KC})(k'')\\
    %             %   =&h_{KC}(h_{K''K}(k''))\\
    %             %   =&h_{KC}(k_2)\\
    %             %   =&c
    %             % \end{flalign*}
    %             and 
    %             \begin{flalign*}
    %               &h_{R''X}(r'')\\
    %              =&h_{R''X}(h_{K''R''}(k'')) \\
    %              =&(h_{K''R''} \mathop{\star} h_{R''X})(k'') \\
    %              =&(h_{K''C''} \mathop{\star} h_{C''X})(k'') & \text{by commutativity of $K''R''XC''$}\\
    %              =&h_{C''X}(h_{K''C''}(k''))\\
    %              =&h_{C''X}(c'')\\
    %              =&x
    %             \end{flalign*}


    %             \begin{flalign*}
    %                =&h_{R''X}(h_{K''R''}(k'')) \\
    %                =&(h_{K''R''} \mathop{\star} h_{R''X})(k'') \\
    %                =&(h_{K''C''} \mathop{\star} h_{C''X})(k'') & \text{by commutativity of $K''R''XC''$}\\
    %                =&h_{C''X}(h_{K''C''}(k''))\\
    %                =&h_{C''X}(c'')\\
    %                =&x
    %               \end{flalign*}

    %                     \begin{flalign*}
    %                         x &= h_{C''X}(c'')  & \text{by definition of $r''$}\\
    %                             &= h_{C''X}(h_{K'C''}(k')) & \text{by definition of $k'$}\\
    %                             &= (h_{K'C''} \mathop{\star} h_{C''X})(k') &  \\
    %                             &= (h_{K'R'} \mathop{\star} h_{R'X})(k') & \text{by the commutativity of $K'C'XR'$} \\
    %                             &= h_{R'X}(h_{K'R'}(k')) &  \\
    %                             &= h_{R'X}(r') &  \text{by definition of $k'$}\\
    %                             &= x& \text{by definition of $r'$}
    %                     \end{flalign*}
    %             which contradic the assumption Equation~\eqref{assumption_norpp}.

    %             Hence \begin{flalign}
    %                 k_1 \mathop{\neq} k_2 \label{k1_neq_k2}
    %             \end{flalign}

    %             From Equation~\eqref{k1_neq_k2}, we have 
    %             \begin{flalign}
    %                 h_{KR}(k_1) \mathop{\neq} h_{KR}(k_2) \label{hkrk1_neq_hkrk2}
    %             \end{flalign} 

    %             Since Equation~\eqref{hkrk1_neq_hkrk2} holds and $\rho$ is $X$-non-increasing, we have 
    %             \begin{flalign}
    %                 h_{KL}(k_1) \mathop{\neq} h_{KR}(k_2) \label{hkrk1_neq_hkrk2}
    %             \end{flalign} 

    %             By assumption on $\rho$ \todo{assumption: non-increasing}, we have $h_{KL}(k_2) \mathop{=} h_{R'L}(r') \mathop{=} l$ for some $k_2 \mathop{\in} K$.
                
    %             \begin{flalign*}
    %                 h_{XG}(x) &= h_{XG}(h_{C''X}(c'')) & \text{by definition of $c''$} \\
    %                         &= (h_{C''X} \mathop{\star} h_{XG})(c') &  \\
    %                         &= (h_{C''C} \mathop{\star} h_{CG})(c'') & \text{by commutativity of $C''CGX$} \\
    %                         &= h_{CG}(h_{C''C}(c'')) &  \\
    %                         &= h_{CG}(c) & \text{by definition of $c$} \\
    %                         &= h_{CG}(h_{KC}(k_1)) & \text{by definition of $k_1$} 
    %                         \\
    %                         &= (h_{KC} \mathop{\star} h_{CG})(k_1) & 
    %                         \\
    %                         &\mathop{\neq} (h_{KC} \mathop{\star} h_{CG})(k_2) & \text{by Equation~\eqref{k1_neq_k2} and injectivity of $h_{KC} \mathop{\star} h_{CG}$} 
    %                         \\
    %                         &= (h_{KL} \mathop{\star} h_{LG})(k_2) & 
    %                         \\
    %                         &= h_{LG}(h_{KL}(k_2)) & 
    %                         \\
    %                         &= h_{LG}(h_{R'L}(r')) & 
    %                         \\
    %                         &= (h_{R'L} \mathop{\star} h_{LG})(r') & 
    %                         \\
    %                         &= (h_{R'X} \mathop{\star} h_{XG})(r') & \text{by commutativity of $R'LGX$} \\
    %                         &= h_{XG}(h_{R'X}(r'))\\
    %                             &= h_{XG}(x) & \text{by definition of $r'$}
    %             \end{flalign*}
              
    %         \end{itemize}            
    \end{itemize}
\end{proof}
% \newpage
\section*{Proofs of~\autoref{lem:decomp_w_u},~\autoref{lem:xlnlmxrnr},~\autoref{lem:w_u_l_not_geq_r_not} and~\autoref{lem:w_g_geq_w_h_leq}}
\label{sec:appendix:a}
\begin*{\textbf{\autoref{subgraph_counting:lem:decomp_w_u}}}
\ \newline
\noindent
\begin{minipage}{0.7\textwidth}
    Let $X$ be a graph. For a pushout square as shown on the right, we have 
    % $\card{\operatorname{Mono}(X, B)} = \card{\operatorname{Mono}(X, D, \beta')}$ and $\card{\operatorname{Mono}(X, C, \lnot \beta)} = \card{\operatorname{Mono}(X, D, \lnot \beta', \alpha')}$.
    \begin{flalign*}
        \card{\operatorname{Mono}(X, B)} &= \card{\operatorname{Mono}(X, D, \beta')}
        \\
        \card{\operatorname{Mono}(X, C, \lnot \beta)} &= \card{\operatorname{Mono}(X, D, \lnot \beta', \alpha')}
    \end{flalign*}
    
\end{minipage}
\hfill
\begin{minipage}{0.3\textwidth}
    \hfill
    \begin{tikzpicture}
        \node (A) {$A$};
        \node [below of=A] (B) {$B$}; 
        \node [left of=A] (C) {$C$}; 
        \node [left of=B] (D) {$D$}; 
        \begin{scope}[nodes=rectangle]          
        \draw [>->] (A) to node [right,label,pos=0.5] {$\alpha$} (B);
        \draw [>->] (A) to node [above,label,pos=0.5] {$\beta$} (C);
        \draw [>->] (B) to node [below,label,pos=0.45] {$\beta'$} (D); 
        \draw [>->] (C) to node [right,label,pos=0.45] {$\alpha'$} (D);
        \end{scope}
    \end{tikzpicture}
\end{minipage}
\end*{}
\begin{proof}
    \label{subgraph_counting:proof:dcomp_w_u}
%    Since $\operatorname{Mono}(X, B) = \operatorname{Mono}(X, D, \beta')$ by definition of $\operatorname{Mono}(X, B)$ and $\operatorname{Mono}(X, D, \beta')$, the equality $\card{\operatorname{Mono}(X, B)} = \card{\operatorname{Mono}(X, D, \beta')}$ holds.
\begin{claim}
    $\card{\operatorname{Mono}(X, B)} = \card{\operatorname{Mono}(X, D, \beta')}$
\end{claim}
\begin{itemize} 
    \item Inclusion \(\operatorname{Mono}(X, B) \star \beta' \subseteq \operatorname{Mono}(X, D, \beta')\) holds by the definition of $\operatorname{Mono}(X, B)$ and $\operatorname{Mono}(X, D, \beta')$. 
    In fact, if \(\iota  \in \operatorname{Mono}(X, B) \star \beta' \), then there is a monomorphism \(\zeta : X \rightarrowtail B\) satisfying \(\iota = \zeta \star \beta'\). Since $\zeta \star \beta'$ is also a monomorphism\todo{monicity of $\beta'$}, \(\iota\) is also an element of \(\operatorname{Mono}(X, D, \beta')\).
    \item Inclusion \(\operatorname{Mono}(X, B) \star \beta' \supseteq \operatorname{Mono}(X, D, \beta')\) holds by the definition of $\operatorname{Mono}(X, B)$ and $\operatorname{Mono}(X, D, \beta')$.
    Suppose \(\iota \in \operatorname{Mono}(X, D, \beta')\). There is a monomorphism \(\zeta : X \rightarrowtail B\) satisfying \(\iota = \zeta \star \beta'\). It follows that \(\iota = \zeta \star \beta' \in \operatorname{Mono}(X, B) \star \beta'\).
    \item Hence, we have \(\operatorname{Mono}(X, B) \star \beta' = \operatorname{Mono}(X, D, \beta')\).
    \item $\card{\operatorname{Mono}(X, B)} \leq \card{\operatorname{Mono}(X, D, \beta')}$ follows by the injectivity of $\beta'$\todo{monicity of $\beta'$}.
\end{itemize}
\begin{claim}
     $\card{\operatorname{Mono}(X, C, \lnot \beta)} = \card{\operatorname{Mono}(X, D, \lnot \beta', \alpha')}$
   \end{claim}
    \begin{itemize}
        % \item The inclusion \(\operatorname{Mono}(X, C, \lnot \beta) \star \alpha'  \subseteq \operatorname{Mono}(X, D, \lnot \beta', \alpha')\) can be justified as follows. Suppose that \(\iota : X \to D\) is an element of \(\operatorname{Mono}(X, C, \lnot \beta) \star \alpha'\). According to the definition of \(\operatorname{Mono}(X, C, \lnot \beta) \star \alpha'\), there exists \(\eta : X \to C\) in \(\operatorname{Mono}(X, C, \lnot \beta)\) such that \(\iota = \eta \star \alpha'\). We need to show that no \(\zeta : X \to B\) exists such that \(\iota = \zeta \star \beta'\). Assuming the contrary, that such a \(\zeta\) exists. Under this assumption, the following commutative diagram holds:
        \item The inclusion \(\operatorname{Mono}(X, C, \lnot \beta) \star \alpha'  \subseteq \operatorname{Mono}(X, D, \lnot \beta', \alpha')\) can be justified as follows: Suppose \(
            \eta \star \alpha' \in \operatorname{Mono}(X, C, \lnot \beta) \star \alpha'\).
        %  By the definition of \(\operatorname{Mono}(X, C, \lnot \beta) \star \alpha'\), there exists a morphism \(\eta : X \to C\) in \(\operatorname{Mono}(X, C, \lnot \beta)\) such that \(\iota = \). 
         Since $\eta \star \alpha'$ is a monomorphism\todo{monicity of $\alpha'$}, it suffices to show that there is no \(\zeta : X \rightarrowtail B\) such that \(\eta \star \alpha' = \zeta \star \beta'\). Suppose, by contradiction, that such a \(\zeta\) exists, then the following commutative diagram holds:
        \begin{center}
            \begin{tikzpicture}[node distance=10mm]
                \node (A) {A}; 
                \node (B) [above right=of A] {B};
                \node (X) [left = of A] {X};
                \node (D) [below right=of B] {D};
                \node (C) [below right=of A] {C};
                \draw[>->] (A) to node[pos=.7, below] {$\alpha$} (B) ;
                \draw[>->] (C) to  node[pos=0.7, below] {$\alpha'$} (D);
                \draw[>->] (A) -- (C) node[pos=.4, right] {$\beta$};
                \draw[>->] (B) -- (D) node[pos=.4, right] {$\beta'$};
                \draw[>->] (X) -- node[above] {$\zeta$} (B);
                \draw[>->] (X) -- node[below] {$\eta$} (C);
            \end{tikzpicture}
        \end{center} 
        The pushout square \(\square ABDC\) is also a pullback square, by~\autoref{prop:pb_eq_po}. The universal property of the pullback provides a morphism \(\gamma : X \rightarrow A\) such that \(\eta = \gamma \star \beta\). \(\gamma\) is a monomorphism, because \(\eta = \gamma \star \beta\) and $\eta$ is a monomorphism. Therefore, the existence of $\gamma$ contradicts the assumption that \(\eta \in \operatorname{Mono}(X, C, \lnot \beta)\). Thus, \(\iota\) is also an element of \(\operatorname{Mono}(X, D, \lnot \beta', \alpha')\). 
        \item The inclusion \(\operatorname{Mono}(X, C, \lnot \beta) \star \alpha'  \supseteq \operatorname{Mono}(X, D, \lnot \beta', \alpha')\) can be justified as follows. Suppose that \(\iota : X \rightarrowtail D\) is an element of \(\operatorname{Mono}(X, D, \lnot \beta', \alpha')\). According to the definition of \(\operatorname{Mono}(X, D, \lnot \beta', \alpha')\), there exists \(\eta : X \rightarrowtail C\) making 
            \begin{flalign}
                \iota = \eta \star \alpha' \label{eq:etastaralphap}
            \end{flalign}
        We show that \(\eta\) is an element of 
        \(\operatorname{Mono}(X, C, \lnot \beta)\) by contradiction.
        
        Suppose that there is a monomorphism \(\zeta : X \rightarrowtail A\) such that 
        \begin{flalign}
            \eta = \zeta \star \beta \label{eq:zetastarbeta}
        \end{flalign}
        and observe:
        \begin{flalign*}
           \hspace{2cm} \iota &\overset{\operatorname{def}}{=} \eta \star \alpha' &\text{by \eqref{eq:etastaralphap}} 
           \\
                  &\overset{\operatorname{def}}{=} (\zeta \star \beta) \star \alpha' & \text{by \eqref{eq:zetastarbeta}}
                  \\
                  &= \zeta \star (\beta \star \alpha') & \text{by associativity} 
                  \\
                  &= \zeta \star (\alpha \star \beta') &\text{by commutativity of $ABDC$}\\
                  &= (\zeta \star \alpha) \star \beta'
        \end{flalign*}
          There is a contradiction because no such factorization of \(\iota\) should exist as \(\iota \in \operatorname{Mono}(X, D, \lnot \beta', \alpha')\).
        %    Therefore, there is no \(\zeta : X \rightarrow A\) such that \(\eta = \zeta \star \beta\), confirming \(\eta\) as an element of \(\operatorname{Mono}(X, C, \lnot \beta)\). Consequently, \(\iota = \eta \star \alpha'\) is an element of \(\operatorname{Mono}(X, C, \lnot \beta) \star \alpha'\).
        \item We conclude by the monicity of \(\alpha'\). \todo{monicity of $\alpha'$}
        % From \(\operatorname{Mono}(X, C, \lnot \beta) \star \alpha'  \subseteq \operatorname{Mono}(X, D, \lnot \beta', \alpha')\) and \(\operatorname{Mono}(X, C, \lnot \beta) \star \alpha'  \supseteq \operatorname{Mono}(X, D, \lnot \beta', \alpha')\), we conclude \(\operatorname{Mono}(X, C, \lnot \beta) \star \alpha' = \operatorname{Mono}(X, D, \lnot \beta', \alpha')\) and \(\card{\operatorname{Mono}(X, C, \lnot \beta) \star \alpha'} = \card{\operatorname{Mono}(X, D, \lnot \beta', \alpha')}\).
        % Since \(\alpha'\) is monic, \(\card{\operatorname{Mono}(X, C, \lnot \beta)} = \card{\operatorname{Mono}(X, C, \lnot \beta) \star \alpha'}\). Therefore, \(\card{\operatorname{Mono}(X, C, \lnot \beta)} = \card{\operatorname{Mono}(X, D, \lnot \beta', \alpha')}\).
    \end{itemize}
\end{proof} 



\noindent
\begin*{\textbf{\autoref{subgraph_counting:lem:xlnlmxrnr}}}
Let $X$ be a graph. Let $L \overset{l}{\leftarrow} K \overset{r}{\rightarrow} R$ be an injective DPO graph rewriting rule. We have 
\[
   \card{\operatorname{Mono}(X, L, \lnot l)}  - \card{\operatorname{Mono}(X, R, \lnot r)} 
   = 
   \card{\operatorname{Mono}(X, L)}  - \card{\operatorname{Mono}(X, R)} 
    \]
\end*{}
\begin{proof}
    \label{subgraph_counting:proof:lem:xlnlmxrnr}
    \begin{claim}
       $\card{\operatorname{Mono}(X, L, l)} = \card{\operatorname{Mono}(X, R, r)}$
    \end{claim}
    \begin{itemize}
        \item We prove $\card{\operatorname{Mono}(X, L, l)} \leq \card{\operatorname{Mono}(X, R, r)}$ by constructing an injection from $\operatorname{Mono}(X, L, l)$ to $\operatorname{Mono}(X, R, r)$. Let $h \in \operatorname{Mono}(X, L, l)$. By definition, we have $h = g \star l$ for some $g: X \rightarrowtail K$. Thus, $g \star r \in \operatorname{Mono}(X, L, r)$. Let $h' \in \operatorname{Mono}(X, L, l)$ such that $h \neq h'$. We have $h' = g' \star l$ for some $g':X \rightarrowtail K$. We have $g \neq g'$ because otherwise we would have $h = g \star l = g' \star l = h'$. From the monicity of $r$, we conclude $g \star r \neq g' \star r$.
        \item Analoguously, we can prove $\card{\operatorname{Mono}(X, L, l)} \geq \card{\operatorname{Mono}(X, R, r)}$ by constructing an injection from $\operatorname{Mono}(X, R, r)$ to $\operatorname{Mono}(X, L, l)$.
        %  Let $h \in \operatorname{Mono}(X, R, r)$. By definition, we have $h = g \star r$ for some $g: X \rightarrowtail K$. Thus, $g \star l \in \operatorname{Mono}(X, L, l)$.
        % Let $h' \in \operatorname{Mono}(X, R, r)$ such that $h \neq h'$. We have $h' = g' \star l$ for some $g':X \rightarrowtail K$. Since $r$ is monic, we deduce $g \neq g'$. From the monicity of $l$, we conclude $g \star r \neq g' \star r$.
    \end{itemize}


    % we have \[\card{\operatorname{Mono}(X, L, l)} - \card{\operatorname{Mono}(X, R, r)} = 0 \] becase 
    
    Hence, the following equality holds
    \begin{flalign*}
         &\card{\operatorname{Mono}(X, L, \lnot l)}  - \card{\operatorname{Mono}(X, R, \lnot r)}\\
        =&\card{\operatorname{Mono}(X, L) \setminus \operatorname{Mono}(X, L, l)}  - 
            \card{\operatorname{Mono}(X, R) \setminus \operatorname{Mono}(X, R, r)}\\
        =&(\card{\operatorname{Mono}(X, L)} - \card{\operatorname{Mono}(X, L, l)}) - (\card{\operatorname{Mono}(X, R)} - \card{\operatorname{Mono}(X, R, r)}) \\
        =&(\card{\operatorname{Mono}(X, L)} - \card{\operatorname{Mono}(X, R)})
         - (\card{\operatorname{Mono}(X, L, l)} - \card{\operatorname{Mono}(X, R, r)})\\
        = &\card{\operatorname{Mono}(X, L)} - \card{\operatorname{Mono}(X, R)}
    \end{flalign*}
\end{proof}

\noindent
\begin*{\textbf{\autoref{subgraph_counting:lem:w_u_l_not_geq_r_not}}}
\ \newline 
    \noindent 
    \begin{minipage}{0.7\textwidth}
        Let $X$ be a ruler-graph and $\rho = (L \overset{l}{\leftarrowtail} K \overset{r}{\rightarrowtail} R)$ an injective DPO rewriting rule.
        Suppose that $\rho$ is $X$-non-increasing. For every rewriting step induced by the DPO diagram (shown on the right), we have 
        \[
            |\operatorname{Mono}(X, G, \lnot m, \lnot l')| \geq |\operatorname{Mono}(X, H, \lnot m', \lnot r')|
        \]
    \end{minipage}
    \hfill
    \begin{minipage}{0.29\textwidth}
        \hfill
        \resizebox{0.9\textwidth}{!}{
                \begin{tikzpicture}
            \node (k) at (0,1) {K};
            \node (l) at (-2,1) {L};
            \node (r) at (2,1) {R};
            \node (c) at (0,-1) {C};
            \node (g) at (-2,-1) {G};
            \node (h) at (2,-1) {H};
            % \node (rb) at (1.5,-0.5) {$R_X$};
            % \node (h') at (1.5,-1.5) {$H'$};
            % \draw[>->]  (rb) -- (h') node [midway,above] {};
            % \draw[>->]  (c) -- (h') node [midway,above] {};
            \draw[<-<]  (l) -- (k) node [midway,below] {$l$};
            \draw[>->]  (k) -- (r) node [midway,below] {$r$};
            \draw[>->] (c) -- (g) node [midway, above] {$l'$};
            \draw[>->] (c) -- (h) node [midway,above] {$r'$};
            \draw[>->] (l) -- (g) node[midway, right] {$m$};
            \draw[>->] (r) -- (h) node[midway, left] {$m'$};
            \draw[>->] (k) -- (c) node[midway, left] {};
            \node () [at=($(l)!0.5!(c)$)] {$PO$};
            \node () [at=($(r)!0.5!(c)$)] {$PO$};
            % \draw[>->] (rb) to (r);
            % \draw[->] (rb) to (l);
            % \draw[<-<] (rb) to (k);
        \end{tikzpicture}
        }
    \end{minipage}
\end*{}
\begin{proof}
    \label{proof:lem:w_u_l_not_geq_r_not}
    We are going to show that for two arbitrary distinct monomorphisms $h_{XH}', h_{XH}'' \in \operatorname{Mono}(X, H, \lnot m', \lnot r')$, there are two distinct monomorphisms $h_{XG}', h_{XG}'' \in \operatorname{Mono}(X, G, \lnot m, \lnot l')$.

    Let $h_{XH}'\in \operatorname{Mono}(X, H, \lnot m', \lnot r')$ be a monomorphism. In the category \textbf{Graph}, we can construct the following commutative diagram
    \begin{center}
        \resizebox{6cm}{!}{
            \begin{tikzpicture}
                \node (k) at (0,0) {K};
                \node (r) at (4,0) {R};
                \node (c) at (0,-3) {C};
                \node (h) at (4,-3) {H};
                % \node (rb) at ($\scl*(1.5,-0.5)$) {$R_X$};
                
                % \node (h') at ($\scl*(1.5,-1.5)$) {$H'$};
                % \draw[>->]  (rb) -- (h') node [midway,above] {};
                % \draw[>->]  (c) -- (h') node [midway,above] {};
    
                \draw[<-<]  (r) -- (k) node [midway,above] {};
                \draw[>->] (c) -- (h) node [midway, below] {};
                \draw[>->] (r) -- (h) node[midway, left] {};
                \draw[>->] (k) -- (c) node[midway, left] {};
    
                % \draw[->] (rb) to (l);
                % \draw[<-<] (rb) to (k);
            
                \node (k') at (1,-1) {K'};
                \node (r') at (2,-1) {R'};
                \node (c') at (1,-2) {C'};
                \node () at (1.5,-1.5) {$PB$};
                \node () at (3,-1.5) {$PB$};
                \node () at (1.5,-2.5) {$PB$};
                \node (x) at (2,-2) {X};
                % \draw [->] (x) -- (h') node[midway] {!};
                \draw [>->] (c') -- (x);
                \draw [>->] (r') -- (x);
                \draw [>->] (k') -- (r');
                \draw [>->] (k') -- (c');
                \draw [>->] (c') -- (c);
                % \draw [->] (r') -- (rb);
                % \draw [->] (r') -- (rb);
                % \draw [->] (h') -- (h);
    
                \draw[>->] (r') -- (r);
                \draw[>->] (x) -- (h) node[midway,right] {$h$};
                % \node (rb) at ($\scl*(\sclx*1.5,-0.5)$) {$R_X$};
                % \node (h') at ($\scl*(\sclx*1.5,-1.2)$) {$H'$};
                % \draw[>->]  (rb) -- (h') node [midway,above] {};
                % \draw[>->]  (c) -- (h') node [midway,above] {};
                % \draw[>->]  (rb) -- (h') node [midway,above] {};
                % \draw[>->] (rb) to (r);
                % \draw[<-<] (rb) to (k);
            \end{tikzpicture}
        }
        \end{center} 
    
    $KCHR$ is also pullback, by~\autoref{prop:pb_eq_po}, because it is pushout. 
    Therefore, since $KCHR$ is also pullback and \(
        h_{K'C'} \star h_{C'C} \star h_{CH} =h_{K'R'} \star h_{R'R} \star h_{R'R}
    \) holds, there exists $h_{K'K}:K' \rightarrowtail K$ such that 
    \begin{flalign}
        h_{K'C'} \star h_{C'C} &= h_{K'K} \star h_{KC} \label{kpcpcpckpkkc}
        \\
        h_{K'R'} \star h_{R'R} &= h_{K'K} \star h_{KR} \label{kprprprkpk}
    \end{flalign} 
    
    The square $K'KCC'$ and $K'R'RK$ are pullbacks, by~\autoref{kpcpck_pullback}. Hence, the following commutative diagram holds
    \begin{center}
        \resizebox{6cm}{!}{
            \begin{tikzpicture}
                \node (k) at (0,0) {K};
                \node (r) at (4,0) {R};
                \node (c) at (0,-3) {C};
                \node (h) at (4,-3) {H};
                % \node (rb) at ($\scl*(1.5,-0.5)$) {$R_X$};
                % \node (h') at ($\scl*(1.5,-1.5)$) {$H'$};
                % \draw[>->]  (rb) -- (h') node [midway,above] {};
                % \draw[>->]  (c) -- (h') node [midway,above] {};
                \draw[<-<]  (r) -- (k) node [midway,above] {};
                \draw[>->] (c) -- (h) node [midway, below] {};
                \draw[>->] (r) -- (h) node[midway, left] {};
                \draw[>->] (k) -- (c) node[midway, left] {};
                % \draw[->] (rb) to (l);
                % \draw[<-<] (rb) to (k);
                \node (k') at (1,-1) {K'};
                \node (r') at (2,-1) {R'};
                \node (c') at (1,-2) {C'};
                \node () at (1.5,-1.5) {$PB$};
                \node () at (3,-1.5) {$PB$};
                \node () at (1.5,-2.5) {$PB$};
                \node () at (1.5,-0.5) {$PB$};
                \node () at (0.5,-1.5) {$PB$};
                \node (x) at (2,-2) {X};
                % \draw [->] (x) -- (h') node[midway] {!};
                \draw [>->] (c') -- (x);
                \draw [>->] (r') -- (x);
                \draw [>->] (k') -- (r');
                \draw [>->] (k') -- (c');
                \draw [>->] (c') -- (c);
                % \draw [->] (r') -- (rb);
                % \draw [->] (r') -- (rb);
                % \draw [->] (h') -- (h);
                \draw[>->] (r') -- (r);
                \draw[>->] (x) -- (h) node[midway,right] {$h$};
                % \node (rb) at ($\scl*(\sclx*1.5,-0.5)$) {$R_X$};
                % \node (h') at ($\scl*(\sclx*1.5,-1.2)$) {$H'$};
                % \draw[>->]  (rb) -- (h') node [midway,above] {};
                % \draw[>->]  (c) -- (h') node [midway,above] {};
                % \draw[>->]  (rb) -- (h') node [midway,above] {};
                % \draw[>->] (rb) to (r);
                % \draw[<-<] (rb) to (k);
                \draw[>->] (k') -- (k) ;
            \end{tikzpicture}
        }
        \end{center} 

    \todo{use: assumption: non increasing}
    Since the rule is by assumption $X$-non-increasing, there is a morphism $h_{R'L}$ such that $K'KLR'$ is commutative.
    
    $K'C'XR'$ is a pushout by~\autoref{kpcpxrp_po}. 
    
    Thus, the following commutative diagram holds
    \begin{center}
        \resizebox{6cm}{!}{
            \begin{tikzpicture}
                \node (k) at (0,0) {K};
                \node (r) at (4,0) {L};
                \node (c) at (0,-3) {C};
                \node (h) at (4,-3) {G};
                % \node (rb) at ($\scl*(1.5,-0.5)$) {$R_X$};
                
                % \node (h') at ($\scl*(1.5,-1.5)$) {$H'$};
                % \draw[>->]  (rb) -- (h') node [midway,above] {};
                % \draw[>->]  (c) -- (h') node [midway,above] {};
    
                \draw[<-<]  (r) -- (k) node [midway,above] {};
                \draw[>->] (c) -- (h) node [midway, below] {};
                \draw[>->] (r) -- (h) node[midway, left] {};
                \draw[>->] (k) -- (c) node[midway, left] {};
    
                % \draw[->] (rb) to (l);
                % \draw[<-<] (rb) to (k);
            
                \node (k') at (1,-1) {K'};
                \node (r') at (2,-1) {R'};
                \node (c') at (1,-2) {C'};
                \node () at (1.5,-1.5) {$PO$};
                % \node () at (1.5,-0.5) {$PB$};
                \node () at (0.5,-1.5) {$PB$};
                % \node () at (3,-1.5) {$PB$};
                % \node () at (1.5,-2.5) {$PB$};
                \node (x) at (2,-2) {X};
                % \draw [->] (x) -- (h') node[midway] {!};
                \draw [>->] (c') -- (x);
                \draw [>->] (r') -- (x);
                \draw [>->] (k') -- (r');
                \draw [>->] (k') -- (c');
                \draw [>->] (c') -- (c);
                \draw [->] (k') -- (k);
                % \draw [->] (r') -- (rb);
                % \draw [->] (r') -- (rb);
                % \draw [->] (h') -- (h);
    
                \draw[>->] (r') -- (r);
                % \draw[>->] (x) -- (h) node[midway,right] {$h$};
                % \node (rb) at ($\scl*(\sclx*1.5,-0.5)$) {$R_X$};
                % \node (h') at ($\scl*(\sclx*1.5,-1.2)$) {$H'$};
                % \draw[>->]  (rb) -- (h') node [midway,above] {};
                % \draw[>->]  (c) -- (h') node [midway,above] {};
                % \draw[>->]  (rb) -- (h') node [midway,above] {};
                % \draw[>->] (rb) to (r);
                % \draw[<-<] (rb) to (k);
            \end{tikzpicture}
        }
        \end{center} 

    By~\autoref{lem:g_monic}, there is a unique monomorphism $h_{XG}':X \rightarrowtail G$ such that $h_{C'X} \star h_{XG}' = h_{C'C} \star h_{CG}$ and $h_{R'X} \star h_{XG}' = h_{R'L} \star h_{LG}$.

    The monomorphism $h_{XG}'$ is in the set $\operatorname{Mono}(X, G, \lnot m, \lnot l')$, by $h_{XH}'\in \operatorname{Mono}(X, H, \lnot m', \lnot r')$ and~\autoref{subgraph_counting:def:creates_more_x_on_the_left}. 
    Intuitively, $h_{XH}'\in \operatorname{Mono}(X, H, \lnot m', \lnot r')$ implies that some elements in $X$ are mapped onto $h_{RH}(R) \setminus (r \star h_{RH})(K)$. By the first assumption of~\autoref{subgraph_counting:def:creates_more_x_on_the_left}, these elements are mapped onto $h_{LG}(L) \setminus (l \star h_{LG})(K)$.

    
    The monomorphism \( h_{XG}' \) belongs to the set \( \operatorname{Mono}(X, G, \lnot m, \lnot l') \), due to \( h_{XH}' \in \operatorname{Mono}(X, H, \lnot m', \lnot r') \) and~\autoref{subgraph_counting:def:creates_more_x_on_the_left}. Intuitively, \( h_{XH}' \in \operatorname{Mono}(X, H, \lnot m', \lnot r') \) implies that certain elements of \( X \) are mapped into \( h_{RH}(R) \setminus (r \star h_{RH})(K) \). By the second assumption of~\autoref{subgraph_counting:def:creates_more_x_on_the_left}, these elements are also mapped into \( h_{LG}(L) \setminus (l \star h_{LG})(K) \).

     \todo{use; def non-increase; condition 1}

    Let $h_{XH}'':X \rightarrowtail H$ be a monomorphism such that 
    \begin{flalign}
        h_{XH}' \neq h_{XH}'' \label{hpneqh}
    \end{flalign}
    
    In the category \textbf{Graph}, we can construct the following commutative diagram:
    \begin{center}
        \resizebox{6cm}{!}{
            \begin{tikzpicture}
                \node (k) at (0,0) {K};
                \node (r) at (4,0) {L};
                \node (c) at (0,-3) {C};
                \node (h) at (4,-3) {G};
                % \node (rb) at ($\scl*(1.5,-0.5)$) {$R_X$};
                
                % \node (h') at ($\scl*(1.5,-1.5)$) {$H'$};
                % \draw[>->]  (rb) -- (h') node [midway,above] {};
                % \draw[>->]  (c) -- (h') node [midway,above] {};
    
                \draw[<-<]  (r) -- (k) node [midway,above] {};
                \draw[>->] (c) -- (h) node [midway, below] {};
                \draw[>->] (r) -- (h) node[midway, left] {};
                \draw[>->] (k) -- (c) node[midway, left] {};
    
                % \draw[->] (rb) to (l);
                % \draw[<-<] (rb) to (k);
            
                \node (k') at (1,-1) {K'};
                \node (r') at (2,-1) {R'};
                \node (c') at (1,-2) {C'};
                \node () at (1.5,-1.5) {$PO$};
                % \node () at (1.5,-0.5) {$PB$};
                \node () at (0.5,-1.5) {$PB$};
                % \node () at (3,-1.5) {$PB$};
                % \node () at (1.5,-2.5) {$PB$};
                \node (x) at (2,-2) {X};
                % \draw [->] (x) -- (h') node[midway] {!};
                \draw [>->] (c') -- (x);
                \draw [>->] (r') -- (x);
                \draw [>->] (k') -- (r');
                \draw [>->] (k') -- (c');
                \draw [>->] (c') -- (c);
                \draw [->] (k') -- (k);
                % \draw [->] (r') -- (rb);
                % \draw [->] (r') -- (rb);
                % \draw [->] (h') -- (h);
    
                \draw[>->] (r') -- (r);
                % \draw[>->] (x) -- (h) node[midway,right] {$h$};
                % \node (rb) at ($\scl*(\sclx*1.5,-0.5)$) {$R_X$};
                % \node (h') at ($\scl*(\sclx*1.5,-1.2)$) {$H'$};
                % \draw[>->]  (rb) -- (h') node [midway,above] {};
                % \draw[>->]  (c) -- (h') node [midway,above] {};
                % \draw[>->]  (rb) -- (h') node [midway,above] {};
                % \draw[>->] (rb) to (r);
                % \draw[<-<] (rb) to (k);
            \end{tikzpicture}
        }
        \end{center} 

    Analoguously, there exists a unique monomorphism $h_{XG}'':X \rightarrowtail G$ such that $h_{C''X} \star h_{XG} ''= h_{C''C} \star h_{CG}$ and $h_{R''X} \star h_{XG}'' = h_{R''L} \star h_{LG}$.
    
    By \eqref{hpneqh} and~\autoref{lem:h_hp_diff_g_gp_diff}, we have $h_{XG} \neq h_{XG}'$.
\end{proof}



\noindent
\begin*{\textbf{\autoref{subgraph_counting:lem:w_g_geq_w_h_leq}}}[Decreasing step]
Let $\rho = (L \overset{l}{\leftarrowtail} K \overset{r}{\rightarrowtail} R)$ be an injective DPO rewriting rule,
\( \mathbb{X} \) a set of ruler-graphs,
\( s_{\mathbb{X}} \colon \mathbb{X} \to \mathbb{N} \) a weight function,
and \( G \Rightarrow_{\rho,\mathfrak{M}} H \) a rewriting step. 
If $\rho$ is \( X \)-non-increasing for every ruler-graph \( X \in \mathbb{X} \), then:\[
    w_{s_\mathbb{X}}(G) - w_{s_\mathbb{X}}(H) 
    \geq 
    w_{s_\mathbb{X}}(L) - w_{s_\mathbb{X}}(R)
\]
\end*{}
\begin{proof}
    \label{subgraph_counting:proof:lem:w_g_geq_w_h_leq}
    Let the following diagram be the witness diagram of the rewriting step
    \begin{center}
        \begin{tikzpicture}
            \node (k) at (0,1) {K};
            \node (l) at (-2,1) {L};
            \node (r) at (2,1) {R};
            \node (c) at (0,-1) {C};
            \node (g) at (-2,-1) {G};
            \node (h) at (2,-1) {H};
            \draw[<-<]  (l) -- (k) node [midway,below] {$l$};
            \draw[>->]  (k) -- (r) node [midway,below] {$r$};
            \draw[>->] (c) -- (g) node [midway, above] {$l'$};
            \draw[>->] (c) -- (h) node [midway,above] {$r'$};
            \draw[>->] (l) -- (g) node[midway, right] {$m$};
            \draw[>->] (r) -- (h) node[midway, left] {$m'$};
            \draw[>->] (k) -- (c) node[midway, left] {};
            \node () [at=($(l)!0.5!(c)$)] {$PO$};
            \node () [at=($(r)!0.5!(c)$)] {$PO$};
        \end{tikzpicture}
    \end{center}
    Let $X \in \mathbb{X}$.

    We have 
    \begin{flalign*}
        & \card{\operatorname{Mono}(X,G)}
            % \overset{\operatorname{def}}{=} 
            % w_{T_\Sigma}(t_G) 
        \\
        % =& \Psi_{G}\\
        =& 
        \card{\operatorname{Mono}(X,G,l')
        \uplus
         \operatorname{Mono}(X,G,\lnot l',m)
        \uplus
         \operatorname{Mono}(X,G,\lnot l',\lnot m)}\\
        =& 
        \card{\operatorname{Mono}(X,G,l')}
        +
         \card{\operatorname{Mono}(X,G,\lnot l',m)}
        +
        \card{\operatorname{Mono}(X,G,\lnot l',\lnot m)}
        \\
    =& \card{\operatorname{Mono}(X,C)}
            +
            \card{\operatorname{Mono}(X, L, \lnot l)} 
            +
            \card{\operatorname{Mono}(X, G,\lnot l', \lnot m)} 
        &  \text{by~\autoref{subgraph_counting:lem:decomp_w_u}}
    \end{flalign*}
    and
    \begin{flalign*}
        % ligne 1
        &\card{\operatorname{Mono}(X,H)}
        % \overset{\operatorname{def}}{=}  
        %     w_{T_\Sigma}(t_H )
        \\
        =& 
        \card{\operatorname{Mono}(X,H,r')
        \uplus
            \operatorname{Mono}(X,H,\lnot r',m')
        \uplus
            \operatorname{Mono}(X,H,\lnot r',\lnot m')}
        \\
        =& 
        \card{\operatorname{Mono}(X,H,r')}
        +
         \card{\operatorname{Mono}(X,H,\lnot r',m')}
        +
        \card{\operatorname{Mono}(X,H,\lnot r',\lnot m')}
        \\
        % = & \Psi_{H} \\
        % ligne 2
        = &
            \card{\operatorname{Mono}(X,C)}
            +
            \card{\operatorname{Mono}(X, R, \lnot r)} 
            +
            \card{\operatorname{Mono}(X, H, \lnot r', \lnot m')}
        % ( 
        %             w_{T_\Sigma} (h_{RH} \star t_H ) + w_{T_\Sigma} (h_{CH} \star t_H  -h_{KC}) 
        % )
        %         +  w_{T_\Sigma}(t_H  - \{h_{RH},h_{CH} \})  
        &  \text{by~\autoref{subgraph_counting:lem:decomp_w_u}}
    \end{flalign*}
    Therefore, the following inequality holds
     \begin{flalign*}
           &\card{\operatorname{Mono}(X,G)} - \card{\operatorname{Mono}(X,H)} &\\
           =&(\card{\operatorname{Mono}(X, L, \lnot l)} - \card{\operatorname{Mono}(X, R, \lnot r)}) + 
            (\card{\operatorname{Mono}(X, G, \lnot l', \lnot m)} - \card{\operatorname{Mono}(X, H, \lnot r', \lnot m')}) \\
       \geq& \card{\operatorname{Mono}(X, L, \lnot l)} - \card{\operatorname{Mono}(X, R, \lnot r)} \hspace{3cm} 
       \text{by~\autoref{subgraph_counting:lem:w_u_l_not_geq_r_not}} \\
          =& \card{\operatorname{Mono}(X, L)} - \card{\operatorname{Mono}(X, R)} \hspace{3cm} 
          \text{by~\autoref{subgraph_counting:lem:xlnlmxrnr}}
     \end{flalign*}
     Thus, the following inequality holds
     \begin{flalign*}
          &w_{s_\mathbb{X}}(G) - w_{s_\mathbb{X}}(H)
          \\
         =&\sum_{X \in \mathbb{X}}^{}w(X) * |\operatorname{Mono}(X,G)| - \sum_{X \in \mathbb{X}}^{}w(X) * |\operatorname{Mono}(X,H)|
         \\
         =&\sum_{X \in \mathbb{X}}^{}w(X) * \left( \card{\operatorname{Mono}(X,G)} -  \card{\operatorname{Mono}(X,H)} \right)
         \\
         \geq&\sum_{X \in \mathbb{X}}^{}w(X) * \left(  \card{\operatorname{Mono}(X,L)}- \card{\operatorname{Mono}(X,R)} \right)
         \\
         = &\sum_{X \in \mathbb{X}}^{}w(X) * \card{\operatorname{Mono}(X,L)} -  \sum_{X \in \mathbb{X}}^{}w(X) * \card{\operatorname{Mono}(X,R)} 
         \\
         = & w_{s_\mathbb{X}}(L) - w_{s_\mathbb{X}}(R)
     \end{flalign*}
    \qed
\end{proof}

\noindent
\begin*{\textbf{\autoref{subgraph_counting:thm:termination_grs}}}[Termination]
    Let \(\mathcal{A}\) and \(\mathcal{B}\) be sets of injective DPO rewriting rules, $\mathbb{X}$ a set of ruler-graphs and $s_\mathbb{X}$ a weight function. If the following hold
    \begin{enumerate}
        \item\label{subgraph_counting:thm:termination_grs:assump:1} $\rho$ is $X$-non-increasing for every rule $\rho \in \mathcal{A} \cup \mathcal{B}$ and for every ruler-graph $X \in \mathbb{X}$,
        \item\label{thm:termination_grs:assump:2} \( w_{s_\mathbb{X}}(lhs(\rho)) > w_{s_\mathbb{X}}(rhs(\rho)) \) for every rule \(\rho \in \mathcal{A}\),
        \item\label{thm:termination_grs:assump:3} \( w_{s_\mathbb{X}}(lhs(\rho)) \geq w_{s_\mathbb{X}}(rhs(\rho)) \) for every rule \(\rho \in \mathcal{B}\).
    \end{enumerate}
    Then \(\Rightarrow_{\mathcal{A},\mathcal{M}}\) terminates relative to \(\Rightarrow_{\mathcal{B},\mathcal{M}}\).
\end*{}
\begin{proof}
    \label{subgraph_counting:proof:thm:termination_grs}
    
    By assumption \eqref{subgraph_counting:thm:termination_grs:assump:1} and ~\autoref{subgraph_counting:lem:w_g_geq_w_h_leq}, the inequality \(
        w_{s_\mathbb{X}}(G) - w_{s_\mathbb{X}}(H) 
        \geq 
        w_{s_\mathbb{X}}(\operatorname{lhs}(\rho)) - w_{s_\mathbb{X}}(\operatorname{rhs}(\rho))
    \) holds for all rules $\rho \in \mathcal{A} \cup \mathcal{B}$ and for all rewriting steps $G \Rightarrow_{\rho, \mathcal{M}} H$.
    
    \noindent From assumptions \eqref{thm:termination_grs:assump:2} and \eqref{thm:termination_grs:assump:3}, we deduce 
    \begin{itemize}
        \item \( w_{s_\mathbb{X}}(G) > w_{s_\mathbb{X}}(H) \) for every rule \(\rho \in \mathcal{A}\);
        \item  \( w_{s_\mathbb{X}}(G) \geq w_{s_\mathbb{X}}(H) \) for every rule \(\rho \in \mathcal{B}\).
    \end{itemize}
    Since $w_{s_\mathbb{X}}(G) \in \mathbb{N}$ for all graph $G$. Every rewriting chain can only have a finite number of rewriting steps using rules in $\mathcal{A}$.
\end{proof} 
\section*{Additional Examples}
% % \begin{table}[H]
% %   \begin{NiceTabular}{ccccccccc}[vlines] % <-- 9 columns now (was 7)
% %    \Hline
% %    \Block{1-2}{\diagbox{\enskip \textbf{Examples}}{\textbf{Techniques}}} & &
% %    \RowStyle{\rotate}
% %     \makecell{Forward\\ closure \\~\cite{Plump1995}} % NEW column #1
% %    & \RowStyle{\rotate}
% %     \makecell{Modular \\ criterion \\~\cite{plump2018modular}} % NEW column #2
% %    & \RowStyle{\rotate}
% %     \makecell{Type graph \\~\cite{bruggink2014termination}}  
% %    & \RowStyle{\rotate}
% %     \makecell{Type graph \\~\cite{bruggink2015proving}} 
% %    & \RowStyle{\rotate}
% %     \makecell{Type graph \\~\cite{endrullis2024generalized_arxiv_v2}} 
% %    & \RowStyle{\rotate}
% %     \makecell{Subgraph \\counting \\~\cite{overbeek2024termination_lmcs}}
% %    & \RowStyle{\rotate}
% %     \makecell{Our proposal} \\
% %    \Hline
% %    \Hline
% %    % ----- from plump 1995 -----
% %    \Block{2-1}{\cite{Plump1995}} 
% %     & \hyperref[ex:overbeek_5d8_plump1995_3d8_plump2018_3_overbeek_5d8]{Example 3.8} 
% %                  & \ding{51} & -- & -- & -- & -- &\ding{51} & \ding{51}\\ 
% %     \Hline
% %     & Example 4.1 & \ding{51} & -- & -- & -- & -- & \ding{51}& \ding{55}\\ 
% %     \hline
% %     % ----- from plump 2018 -----
% %    \Block{4-1}{\cite{plump2018modular}} 
% %    & \hyperref[ex:overbeek_5d8_plump1995_3d8_plump2018_3_overbeek_5d8]{Example 3} 
% %               & -- & \ding{51} &  -- & -- & -- & \ding{51} & \ding{51}\\ 
% %    \Hline
% %    & \hyperref[ex:plump_ex4]{Example 4} &  -- &  \ding{51} &  -- & -- & -- & \ding{55} &\ding{55}\\ 
% %    \Hline
% %    & Example 5 &  -- &  \ding{51} &   -- & -- & -- &  \ding{51} & \ding{51}\\ 
% %    \Hline
% %    & Example 6 s->s&  -- & \ding{51} & -- & -- & -- & \ding{51} & \ding{55}\\ 
% %    \Hline
% %    % ----- from bruggink2014 -----
% %    \Block{4-1}{\cite{bruggink2014termination}} 
% %     & Example 1 
% %       & -- & -- & \ding{51} & -- & -- & \ding{55} & \ding{55}\\ \Hline
% %    & \hyperref[ex:plump_ex4]{Example 5}
% %       & -- & -- & \ding{51} & -- & -- & -- &  \ding{55}\\ \Hline
% %    & \hyperref[ex:termination:grsaa]{Example 4 and 6}  grsaa
% %       & -- & -- & \ding{51} & -- & -- & ? & \ding{51} \\ \Hline
% %    & Routing Protocol
% %       & -- & -- & \ding{51} & -- & -- & \ding{55} &  \ding{55}\\ \Hline
% %    % Example~\ref{ex_contrib_variant} 
% %    %   & \ding{55} & -- & -- & -- & \ding{51}\\ \Hline
   
% %    % ----- from bruggink2015 -----
% %    \Block{4-1}{\cite{bruggink2015proving}}
% %     &\hyperref[ex:termination:grsaa]{Example 2} grsaa
% %       & -- & -- & -- & \ding{51} & -- & ? &  \ding{51}\\ \Hline
% %    & Example 4 
% %       & -- & -- & -- & \ding{51} & -- & \ding{51} & \ding{51} \\ \Hline
% %    & Example 5 
% %       & -- & -- & -- & \ding{51} & -- & \ding{55} & \ding{55}\\ \Hline
% %    & Example 6 
% %       & -- & -- & -- & \ding{51} & -- & \ding{55} & \ding{55}\\ \Hline
% %    % Example~\ref{ex_contrib_variant} 
% %    %   & -- & \ding{55} & -- & -- & \ding{51}\\ \Hline
   
% %    % ----- from endrullis2024generalized_arxiv_v2 -----
% %    \Block{8-1}{\cite{endrullis2024generalized_arxiv_v2}}
% %     & Example 6.2  
% %       & -- & -- & -- & -- & \ding{51} & -- & \ding{51}\\ \Hline
% %    & \hyperref[ex_endrullis_5d8_overbeek]{Example 6.3}
% %       & -- & -- & -- & -- & \ding{51} & \ding{55} & \ding{51}\\ \Hline
% %    & Example 6.4 SGraph
% %       & -- & -- & -- & -- & -- & -- & -- \\ \Hline
% %    & Example 6.5  SGraph
% %       & -- & -- & -- & -- &  -- & -- & -- \\ \Hline
% %    & \hyperref[ex:overbeek_5d8_plump1995_3d8_plump2018_3_overbeek_5d8]{Example D.1}
% %       & -- & -- & -- & -- & \ding{51} & -- & \ding{51}\\ \Hline
% %    & Example D.2 
% %       & -- & -- & -- & -- & \ding{51} & -- & \ding{55}\\ \Hline
% %    & Example D.3 
% %       & -- & -- & -- & -- & \ding{51} & -- & \ding{55}\\ \Hline
% %    & Example D.4 
% %       & -- & -- & -- & -- & \ding{51} & -- & \ding{55}\\ \Hline
% %    % & Example~\ref{ex_contrib_variant}
% %    %   & -- & -- & \ding{55} & -- & \ding{51}\\ \Hline
 
% %    % ----- from overbeek2024termination_lmcs -----
% %    \Block{7-1}{\cite{overbeek2024termination_lmcs}}
% %    & Example 5.2
% %       & -- & -- & -- & -- & -- & \ding{51} & \ding{55}\\ \Hline
% %    & \hyperref[ex:overbeek_5d3]{Example 5.3}
% %       & -- & -- & -- & -- & -- & \ding{51} & \ding{51}\\ \Hline
% %  %   & \hyperref[ex:overbeek_5d3]{Example 5.3 monic matches}
% %  %      & -- & -- & -- & -- & -- & \ding{51} & \ding{51}\\ \Hline
% %    & \hyperref[ex:overbeek_5d5]{Example 5.5} 
% %       & -- & -- & -- & -- & -- & \ding{51} & \ding{51}\\ \Hline
% %    & \hyperref[ex:overbeek_5d6]{Example 5.6}
% %       & -- & -- & -- & -- & -- & \ding{51} & \ding{51} \\ \Hline
% %  %   & \hyperref[ex:overbeek_5d6]{Example 5.6 bis}
% %  %      & -- & -- & -- & -- & -- & \ding{51} & -- \\ \Hline
% %    & Example 5.7 
% %       & -- & -- & -- & -- & -- & -- & -- \\ \Hline
% %    & Example 5.9 
% %       & -- & -- & -- & -- & -- & \ding{51} & \ding{55}\\ \Hline
% %    & \hyperref[ex:overbeek_5d8_plump1995_3d8_plump2018_3_overbeek_5d8]{Example 5.8}
% %       & -- & -- & \ding{55} & \ding{55} & -- & \ding{51} & \ding{51}\\ \Hline
% %    Current paper 
% %     & Example~\ref{ex_contrib_variant}
% %       & \ding{55} & \ding{55} & \ding{55} & \ding{55} & \ding{55} & \ding{55} & \ding{51} \\ \Hline
% %  \end{NiceTabular}
% %  \caption{Applicability of the termination techniques to the examples with DPO rewriting with monic matches.
% %    \ding{51} indicates that the method can be applied to the example,
% %          \ding{55} indicates that the method cannot be applied,
% %          and the symbol $-$ of a symbol signifies that its applicability
% %          is not relevant to this discussion.
% %         }
% %  \label{tab:comparaison}
% %  \end{table}
 

% \begin{example} 
%   \label{ex:endrullis2024_6d2}  
%   Consider the following rewriting rule presented in~\cite[Example 6.2]{endrullis2024generalized_arxiv_v2}. 

%   \begin{tikzpicture}  
%      \graphbox{$L$}{0mm}{-11mm}{32mm}{15mm}{2mm}{-8mm}{  
%          \node[draw,circle]  (x) at (-6mm,0mm) {1};  
%         \draw[->] (x) edge[loop above] node  {} (x);

%          \node[draw,circle]  (y) at (6mm,0mm) {2};  
%        }  
%        \graphbox{$K$}{33mm}{-11mm}{32mm}{15mm}{2mm}{-8mm}{  
%         \node[draw,circle]  (x) at (-6mm,0mm) {1};  
%         \node[draw,circle]  (y) at (6mm,0mm) {2};  
%        }  
%        \graphbox{$R$}{66mm}{-11mm}{32mm}{15mm}{1mm}{-8mm}{  
%         \node[draw,circle]  (x) at (-6mm,0mm) {1};  
%         \node[draw,circle]  (y) at (6mm,0mm) {2};  
%         \draw[->]  (x) to (y);  
%        }    
%  \end{tikzpicture}  
%  Our method can prove its termination with the ruler-graph  \tikz[baseline=-0.5ex]{
%   \node (x) at (0,0) {$\bullet$};
%   \draw[->] (x) edge[loop above] node  {} (x);
% } of weight $1$.
% \end{example} 

% \begin{example}
%     Consider a rewriting rule given by Overbeek and Endrullis in~\cite[Example 5.2]{overbeek2024termination_lmcs} depicted in the bottom span of the following diagram

%     \begin{tikzpicture}
%         \graphbox{$L$}{0mm}{-11mm}{32mm}{15mm}{2mm}{-8mm}{
%             \node[draw,circle]  (x) at (-6mm,0mm) {1};
%             \node[draw,circle]   (y) at (6mm,0mm) {2};
%             \draw[->]  (x) to (y);
%           }
%           \graphbox{$K$}{33mm}{-11mm}{32mm}{15mm}{2mm}{-8mm}{
%             \node[draw,circle]  (x) at (-6mm,0mm) {1};
%             \node[draw,circle]  (y) at (6mm,0mm) {2};
%             \draw[->]  (x) to (y);
%           }
%           \graphbox{$R$}{66mm}{-11mm}{32mm}{15mm}{1mm}{-8mm}{
%             \node[draw,circle]  (x) at (0mm,0mm) {1 2};
%             \draw[->]  (x) edge [loop right] (x);
%           }  
%     \end{tikzpicture}  

%     The termination of this rule is not in the scope of our method due to its non-node-injective right-hand side morphism. It can be proved using the method proposed by Overbeek et al. in~\cite{overbeek2024termination_lmcs}.
% \end{example}


% \begin{example}
%   \label{ex:overbeek_5d3}
%   Consider the rewriting rules presented in~\cite[Example 5.3]{overbeek2024termination_lmcs} that can be depicted as the bottom spans of

%   \begin{tikzpicture} 
%       \graphbox{$L$}{0mm}{0mm}{23mm}{12mm}{2mm}{-5mm}{
%           \node[draw,circle] (x) at (0mm,0mm) {1};
%           \draw[->] (x) edge [loop right] node {} (x);
%       }
%       \graphbox{$K$}{24mm}{0mm}{23mm}{12mm}{2mm}{-5mm}{
%         \node[draw,circle] (x) at (0mm,0mm) {1};
%       }
%       \graphbox{$R_X$}{24mm}{13mm}{23mm}{12mm}{2mm}{-5mm}{
%           % \node[draw,circle] (x) at (0mm,0mm) {1};
%         }
%       \begin{scope}[opacity=1]        
%       \graphbox{$R$}{48mm}{0mm}{30mm}{12mm}{-1mm}{-5mm}{
%         \node[draw,circle] (x) at (0mm,0mm) {1};
%         \node[draw,circle] (y) at (10mm,0mm) {2};
%       }
%       \end{scope}
%     \end{tikzpicture}

%   \begin{tikzpicture}
%       \graphbox{$L$}{0mm}{-11mm}{32mm}{12mm}{2mm}{-8mm}{
%       \node[draw,circle]  (x) at (-6mm,0mm) {1};
%       \node[draw,circle] (y) at (6mm,0mm) {2};
%       \draw[->]  (x) to (y);
%       }
%       \graphbox{$K$}{33mm}{-11mm}{32mm}{12mm}{2.5mm}{-8mm}{
%           \node[draw,circle]  (x) at (-6mm,0mm) {1};
%           \node[draw,circle]  (y) at (6mm,0mm) {2};
%       }
%       \graphbox{$R_X$}{33mm}{2mm}{32mm}{12mm}{2.5mm}{-8mm}{
%           % \node[draw,circle]  (x) at (-6mm,0mm) {1};
%           % \node[draw,circle]  (y) at (6mm,0mm) {2};
%       }
%       \begin{scope}[opacity=1]        
%       \graphbox{$R$}{66mm}{-11mm}{40mm}{12mm}{-2mm}{-8mm}{
%           \node[draw,circle]  (x) at (-6mm,0mm) {1};
%           \node[draw,circle]  (y) at (6mm,0mm) {2};
%           \node[draw,circle]  (z) at (17mm,0mm) {3};
%       }
%       \end{scope}
%     \end{tikzpicture}

%   Let $X$ be 
%   \tikz[baseline=-0.5ex]{
%       \node (x) at (0,0) {$\bullet$};
%       \node (y) at (1,0) {$\bullet$};
%       \draw[->] (x) to (y);
%   }.
%   We have $|\operatorname{Mono}(X,L)| \mathop{=} 1 \mathop{>} 0 \mathop{=} |\operatorname{Mono}(X,R)|$ for both rules. Thus, this rewriting systems terminates by Theorem~\ref{thm:termination_grs}.
% \end{example}



% \begin{example}
%   \label{ex:overbeek_5d6_bis} 
%   Consider the DPO rewriting system with monic matches presented in~\cite[Example 5.6]{overbeek2024termination_lmcs} with $\rho$ replaced by $\rho'$ as depicted by the span 
%   \begin{center}
%     $\rho' \mathop{=} $\scalebox{0.9} { {
\begin{tikzpicture}[baseline=-10mm]
    \graphbox{$L$}{0mm}{0mm}{38mm}{20mm}{2mm}{-13mm}{
      \node [draw, circle] (x) at (-7mm,0mm) {1};
      \node [draw, circle] (y) at (3mm,0mm) {2 3};
      \draw[->] (x) edge[loop above] node  {$a$} (x);
      \draw[->] (y) edge [loop above] node {$a$} (y);
    }
    \graphbox{$K$}{42mm}{-0mm}{38mm}{20mm}{0mm}{-10mm}{
        \node [draw, circle] (x) at (-7mm,0mm) {1};
        \node [draw, circle] (y) at (0mm,0mm) {2};
        \node [draw, circle] (z) at (7mm,0mm) {3};    
    }
    \begin{scope}[opacity=1]        
    \graphbox{$R$}{85mm}{-0mm}{38mm}{20mm}{2mm}{-13mm}{
      \node [draw, circle] (x) at (-7mm,0mm) {1};
      \node [draw, circle] (y) at (0mm,0mm) {2};
      \node [draw, circle] (z) at (7mm,0mm) {3};
      \draw[->] (x) edge[loop above] node  {$b$} (x);
      \draw[->] (y) edge[loop above] node  {$b$} (y);
      \draw[->] (z) edge[loop above] node  {$b$} (z);
    }
    \end{scope}
    \node () at (40mm,-10mm) {$\leftarrow$};
    \node () at (83mm,-10mm) {$\rightarrow$};
\end{tikzpicture}
}

}
%   \end{center} 

%   Its termination is not in the scope of our method due to its non-node-injective left-hand side morphism. 
% \end{example}

% \begin{example}
%   \cite[Example 5.7]{overbeek2024termination_lmcs} has conditions on the host graph, which cannot be modeled by DPO rewriting rules. 
% \end{example}
% \begin{example}
%   The POPO+ rewriting system presented in~\cite[Example 5.7]{overbeek2024termination_lmcs} cannot be modeled by DPO rewriting rules due to the application condition on the context.
%   % Consider Example 5.7  without application condition.

%   % \begin{tikzpicture}
%   %     \graphbox{L}{0mm}{-9mm}{35mm}{8mm}{1mm}{-4mm}{
%   %         \node [draw,circle] (y) at (0mm,0mm) {1};
%   %         \draw [->] (y) edge [loop right] node {} (y);
%   %     }
%   %     \graphbox{K}{36mm}{-9mm}{35mm}{8mm}{1mm}{-4mm}{
%   %         \node [draw,circle] (y) at (0mm,-0mm) {1};
%   %     }
%   %     \begin{scope}[opacity=1]        
%   %         \graphbox{R}{72mm}{-9mm}{35mm}{8mm}{1mm}{-4mm}{
%   %         \node [draw,circle] (y) at (0mm,-0mm) {1};
%   %     }
%   %     \end{scope}
%   %   \end{tikzpicture}

%   % \begin{tikzpicture}
%   %     \graphbox{L}{0mm}{-9mm}{35mm}{8mm}{1mm}{-4mm}{
%   %       \node [draw,circle] (y) at (0mm,0mm) {1};
%   %     }
%   %     \graphbox{K}{36mm}{-9mm}{35mm}{8mm}{1mm}{-4mm}{
%   %     }
%   %     \begin{scope}[opacity=1]        
%   %     \graphbox{R}{72mm}{-9mm}{35mm}{8mm}{1mm}{-4mm}{
%   %     }
%   %     \end{scope}
%   %   \end{tikzpicture}
% \end{example}



% \begin{example}
%   \label{ex:overbeek_5d8_plump1995_3d8_plump2018_3_overbeek_5d8}
%   Consider the DPO rewriting system presented in~\cite[Example 5.8]{overbeek2024termination_lmcs},~\cite[Example 3]{plump2018modular} and~\cite[Example 3.8]{plump1995ontermination}. It rewriting rules are depicted by the bottom spans of the following diagram

% \begin{center}
%  $\rho \mathop{=} ${ 
%   \begin{tikzpicture}[baseline=-20mm]
%         \graphbox{$L$}{0mm}{-11mm}{35mm}{15mm}{0.5mm}{-9mm}{
%           \node [draw, circle] (x) at (-10mm,0mm) {1};
%           \node [draw, circle] (y) at (0mm,0mm) {2};
%           \node [draw, circle] (z) at (10mm,0mm) {3};
%           \draw[->] (x) to node [above] {$a$} (y);
%           \draw[->] (y) to node [above] {$b$} (z);
%         }
%         \graphbox{$K$}{36mm}{-11mm}{35mm}{15mm}{0.5mm}{-9mm}{
%           \node [draw, circle] (x) at (-10mm,0mm) {1};
%           %\node [draw, circle] (y) at (0mm,0mm) {2};
%           \node [draw, circle] (z) at (10mm,0mm) {3};
%         }
%         \begin{scope}[opacity=1]        
%         \graphbox{$R$}{72mm}{-11mm}{35mm}{15mm}{0.5mm}{-9mm}{
%           \node [draw, circle] (x) at (-10mm,0mm) {1};
%           \node [draw, circle] (y) at (0mm,0mm) {2};
%           \node [draw, circle] (z) at (10mm,0mm) {3};
%           \draw[->] (x) to node [above] {$a$} (y);
%           \draw[->] (y) to node [above] {$c$} (z);
%         }
%         \end{scope}
%       \end{tikzpicture}
%  }
% \end{center}

% \begin{center}
%   $\tau \mathop{=} ${ 
%    \begin{tikzpicture}[baseline=-20mm]
%       \graphbox{$L$}{0mm}{-11mm}{35mm}{15mm}{0.5mm}{-9mm}{
%           \node [draw, circle] (x) at (-10mm,0mm) {1};
%           \node [draw, circle] (y) at (0mm,0mm) {2};
%           \node [draw, circle] (z) at (10mm,0mm) {3};
%           \draw[->] (x) to node [above] {$c$} (y);
%           \draw[->] (y) to node [above] {$d$} (z);
%       }
%       \graphbox{$K$}{36mm}{-11mm}{35mm}{15mm}{0.5mm}{-9mm}{
%           \node [draw, circle] (x) at (-10mm,0mm) {1};
%           %\node [draw, circle] (y) at (0mm,0mm) {2};
%           \node [draw, circle] (z) at (10mm,0mm) {3};
%       }
%       \begin{scope}[opacity=1]        
%       \graphbox{$R$}{72mm}{-11mm}{35mm}{15mm}{0.5mm}{-9mm}{
%           \node [draw, circle] (x) at (-10mm,0mm) {1};
%           \node [draw, circle] (y) at (0mm,0mm) {2};
%           \node [draw, circle] (z) at (10mm,0mm) {3};
%           \draw[->] (x) to node [above] {$d$} (y);
%           \draw[->] (y) to node [above] {$b$} (z);
%       }
%       \end{scope}
%       \end{tikzpicture}
%   }
%   \end{center}

%   Let $X$ be  \tikz[baseline=-0.5ex]{
%       \node (x) at (0,0) {$\bullet$};
%       \node (y) at (1,0) {$\bullet$ };
%       \node (z) at (2,0) { $\bullet$};
%       \draw[->] (x) -- (y) node[midway, above] {$a$};
%       \draw[<-] (z) -- (y) node[midway, above] {$b$};
%   } 

%   For both rules, $R_X$ is 

%   \begin{tikzpicture}
%     \graphbox{$R_X$}{36mm}{6mm}{35mm}{15mm}{0.5mm}{-9mm}{
%       % \node [draw, circle] (x) at (-10mm,0mm) {1};
%       % %\node [draw, circle] (y) at (0mm,0mm) {2};
%       % \node [draw, circle] (z) at (10mm,0mm) {3};
%     }
%   \end{tikzpicture}

%   We have 
%   $|\operatorname{Mono}(X,L)| \mathop{=} 1 \mathop{>} 0 \mathop{=} |\operatorname{Mono}(X,R)|$ for $\rho$, and 
%   $|\operatorname{Mono}(X,L)| \mathop{=} 0 \mathop{=} 0 \mathop{=} |\operatorname{Mono}(X,R)|$ for $\tau$, therefore $\rho$ is eliminated.
  
%   Let $Y$ be  \tikz[baseline=-0.5ex]{
%       \node (x) at (0,0) {$\bullet$};
%       \node (y) at (1,0) {$\bullet$ };
%       \draw[->] (x) -- (y) node[midway, above] {$c$};
%   }. For $\tau$, we have

%   \begin{tikzpicture}
%     \graphbox{$R_X$}{36mm}{6mm}{35mm}{15mm}{0.5mm}{-9mm}{
%       % \node [draw, circle] (x) at (-10mm,0mm) {1};
%       % %\node [draw, circle] (y) at (0mm,0mm) {2};
%       % \node [draw, circle] (z) at (10mm,0mm) {3};
%     }
%   \end{tikzpicture} 
  
%   and $|\operatorname{Mono}(Y,L)| \mathop{=} 1 \mathop{>} 0 \mathop{=} |\operatorname{Mono}(Y,R)|$.

%   Thus, this rewriting systems terminates by Theorem~\ref{thm:termination_grs}.
% \end{example}

% \begin{example}
%   \label{ex:bruggink2015_ex5}
%   Our method cannot prove the termination of the DPO rewriting system with monic matches presented in~\cite[Example 5]{bruggink2015proving}. 
% \end{example}
% \begin{example}
%   \label{ex:bruggink2015_ex6_endrullis2024_d2}
%   Our method cannot prove the termination of the DPO rewriting system with monic matches presented in~\cite[Example 6]{bruggink2015proving}. 
% \end{example}

% \begin{example}
%   \label{ex:plump2018_ex6_endrullis_d4}
%   The DPO rewriting system with monic matches presented in~\cite[Example 6]{plump2018modular} is not in the scope of our method due to the non-injective rule.
% \end{example}


% \begin{example} 
%   \label{ex:termination:contrib} 
%   The rewriting rule illustrated below is from~\cite[Example 6]{plump2018modular}.
%   % \begin{figure}[H] 
%   \begin{center}
%     $\tau \mathop{=} ${ \resizebox{0.6\textwidth}{!}{
%       \begin{tikzpicture}[baseline=-17mm]
%           \graphbox{$L$}{0mm}{0mm}{35mm}{35mm}{2mm}{-5mm}{
%               \coordinate (delta) at (0,-18mm);
%               \node[draw,circle] (l1) at ($(delta)+(-1,1.5)$) {1};
%               \node[draw,circle] (l2) at ($(delta)+(1,1.5)$) {2};
%               \node[draw,circle] (l3) at ($(delta)+(0,0)$) {3};
%               \draw[->] (l1) -- (l3) node[midway,left] {$s$};
%               \draw[->] (l2) -- (l3) node[midway,right] {$s$};
%               \draw[->] (l3) edge [loop below] node {0} (l3);
%           }
%           \graphbox{$K$}{40mm}{0mm}{35mm}{35mm}{2mm}{-5mm}{
%               \coordinate (delta) at (0,-18mm);
%               \coordinate (interfaceorigin) at ($(delta) +(5,0)$);
%               \node[draw,circle] (r1) at ($(delta) +(-1,1.5)$) {1};
%               \node[draw,circle] (r2) at ($(delta) +(0.5,1.5)$) {2};
%               \node[draw,circle] (r3) at ($(delta)+(0,0)$) {3};
%               \draw[->] (r1) -- (r3) node[midway,left] {$s$};
%               \draw[->] (r3) edge [loop below] node {0} (r3);
%           }
%           \graphbox{$R$}{80mm}{0mm}{35mm}{35mm}{2mm}{-5mm}{
%               \coordinate (delta) at (0,-18mm);
%               \node[draw,circle] (r1) at ($(delta)+(-1,1.5)$) {1};
%               \node[draw,circle] (r2) at ($(delta)+(0.5,1.5)$) {2};
%               \node[draw,circle] (r3) at ($(delta)+(0,0)$) {3};
%               \node[draw,circle] (r4) at ($(delta)+(1,0)$) {};
%               \draw[->] (r1) -- (r3) node[midway,left] {$s$};
%               \draw[->] (r2) -- (r4) node[midway,right] {$s$};
%               \draw[->] (r4) edge [loop below] node {0} (r4);
%               \draw[->] (r3) edge [loop below] node {0} (r3);
%           }
%           % \graphbox{$R_x$}{40mm}{40mm}{35mm}{35mm}{2mm}{-5mm}{
%           %     \coordinate (delta) at (0,-18mm);
%           %     \coordinate (rxorigin) at ($(interfaceorigin)+(0,6)$);
%           %     \node[draw,circle] (r1) at ($(delta)+(-1,1.5)$) {1};
%           %     \node[draw,circle] (r2) at ($(delta)\mathop{+} (0.5,1.5)$) {2};
%           %     \node[draw,circle] (r3) at ($(delta)\mathop{+} (0,0)$) {3};
%           %     \draw[->] (r1) -- (r3) node[midway,left] {$s$};
%           %     \draw[->] (r3) edge [loop below] node {0} (r3);
%           % }
%           \node () at (37mm,-18mm) {$\leftarrowtail$};
%           \node () at (78mm,-18mm) {$\rightarrowtail$};
%           % \node () at (57mm,2mm) {$\uparrowtail$};
%           % \node () at (38mm,2mm) {$\swarrowtail$};
%           % \node () at (79mm,2mm) {$\searrowtail$};
%       \end{tikzpicture}
%       }
%     }
%   \end{center}
%   Its termination can be proved by our method.
% \end{example}  


% % \begin{example}
% %   \label{ex:termination:grsaa}
% %   Consider the rewriting rule in Example~\ref{ex:grsaa_rx}. Let $X$ be the graph \tikz[baseline=-0.5ex]{
% %       \node (x) at (0,0) {$\bullet$};
% %       \node (y) at (1,0) {$\bullet$ }; 
% %       \node (z) at (2,0) { $\bullet$};
% %       \draw[->] (x) -- (y) node[midway, above] {$a$};
% %       \draw[<-] (z) -- (y) node[midway, above] {$a$};
% %   } with weight $1$ and $\mathbb{X} \mathop{=} \{X\}$. The rule is $X$-non-increasing by Example~\ref{example:grs_aa:has_more_left}. 
% %   % Let $s_\mathbb{X}$ be the weight function which associates the weight of $1$ to $X$. 
% %   % Since \(w_{s_\mathbb{X}}(L) \mathop{=} 1 \mathop{>} 0 \mathop{=} w_{s_\mathbb{X}}(R)\), 
% %   Since \(w(L) \mathop{=} 1 \mathop{>} 0 \mathop{=} w(R)\),
% %   it terminates by Theorem~\ref{thm:termination_grs}.
% % \end{example}

% % We present an example where $\mathbb{X}$ is not a singleton.
% % \begin{example} 
% %   \label{ex:overbeek_5d6}
% %   Consider the rewriting rules presented in~\cite[Example 5.6]{overbeek2024termination_lmcs} illustrated below:
% %   \begin{center}
% %     $\rho \mathop{=} $\scalebox{0.7} { {
\begin{tikzpicture}[baseline=-10mm]
    \graphbox{$L$}{0mm}{0mm}{31mm}{20mm}{2mm}{-13mm}{
      \node [draw, circle] (x) at (-7mm,0mm) {1};
      \node [draw, circle] (y) at (0mm,0mm) {2};
      \draw[->] (x) edge[loop above] node  {$a$} (x);
      \draw[->] (y) edge [loop above] node {$a$} (y);
    }
    \graphbox{$K$}{32mm}{-0mm}{18mm}{20mm}{0mm}{-10mm}{
    }
    \begin{scope}[opacity=1]        
    \graphbox{$R$}{51mm}{-0mm}{38mm}{20mm}{2mm}{-13mm}{
      \node [draw, circle] (x) at (-7mm,0mm) {3};
      \node [draw, circle] (y) at (0mm,0mm) {4};
      \node [draw, circle] (z) at (7mm,0mm) {5};
      \draw[->] (x) edge[loop above] node  {$b$} (x);
      \draw[->] (y) edge[loop above] node  {$b$} (y);
      \draw[->] (z) edge[loop above] node  {$b$} (z);
    }
    \end{scope}
\end{tikzpicture}
}}
% %   \end{center}
% %   \begin{center}
% %   $\tau \mathop{=} $\scalebox{0.7}{ {\begin{tikzpicture}[baseline=-10mm]
    \graphbox{L}{0mm}{0}{31mm}{20mm}{2mm}{-13mm}{
        \node[draw,circle] (x) at (-7mm,0mm) {1};
        \node[draw,circle] (y) at (0mm,0mm) {2};
        \draw[->] (x) edge[loop above] node {$b$} (x);
        \draw[->] (y) edge [loop above] node {$b$} (y);
      }
      \graphbox{K}{32mm}{0}{18mm}{20mm}{2mm}{-13mm}{
      }
      \begin{scope}[opacity=1]        
      \graphbox{R}{51mm}{0}{25mm}{20mm}{1mm}{-13mm}{
        \node[draw,circle] (x) at (-3.5mm,0mm) {3};
        \draw[->] (x) edge[loop above] node {$a$} (x);
      }
      \end{scope}
  \end{tikzpicture}}
% %   }
% %   \end{center}
% %   Let $X$ be 
% %   \tikz[baseline=-0.5ex]{
% %       \node (x) at (0,0) {$\bullet$};
% %       \draw[->] (x) edge [loop right] node {$a$} (x);
% %   } with weight $5$, let $Y$ be
% %   \tikz[baseline=-0.5ex]{
% %       \node (x) at (0,0) {$\bullet$};
% %       \draw[->] (x) edge [loop right] node {$b$} (x);
% %   } with weight $3$ and $\mathbb{X} \mathop{=} \{X, Y\}$.
% %   The rules $\rho$ and $\tau$ are both $X$- and $Y$-non-increasing with $D(rhs(\rho), X) \mathop{=} D(rhs(\rho), Y) \mathop{=} D(rhs(\tau), X) \mathop{=} D(rhs(\tau), Y) \mathop{=} \emptyset$.   
% %   % Let $s_\mathbb{X}$ be the weight function associating the weight of $5$ to $X$ and the weight of $3$ to $Y$.
% %   % We have $
% %   %     w_{s_\mathbb{X}}(\operatorname{lhs}(\rho)) \mathop{=} 10 \mathop{>} 9 \mathop{=} w_{s_\mathbb{X}}(\operatorname{rhs}(\rho))
% %   %     $ and $
% %   %     w_{s_\mathbb{X}}(\operatorname{lhs}(\tau)) \mathop{=} 6 \mathop{>} 5 \mathop{=} w_{s_\mathbb{X}}(\operatorname{rhs}(\tau))$.
% %   We have $
% %   w(\operatorname{lhs}(\rho)) \mathop{=} 10 \mathop{>} 9 \mathop{=} w(\operatorname{rhs}(\rho))
% %   $ and $
% %   w(\operatorname{lhs}(\tau)) \mathop{=} 6 \mathop{>} 5 \mathop{=} w(\operatorname{rhs}(\tau))$.
% %   Thus, this rewriting systems terminates by Theorem~\ref{thm:termination_grs}.
% % \end{example}


\printbibliography
\end{document} 